\input{preamble}


% OK, start here.
%
\begin{document}

\title{Théorie des déformations formelles}


\maketitle

\phantomsection
\label{section-phantom}

\tableofcontents



\section{Introduction}
\label{section-introduction}

\noindent
Ce chapitre développe la théorie des déformations formelles sous une forme
applicable plus loin dans le projet Champs, en suivant de près Rim
\cite[Exposee VI]{SGA7-I} et Schlessinger \cite{Sch}. Nous recommandons vivement
au lecteur qui découvre ce sujet de lire d'abord l'article de Schlessinger :
il est suffisamment général pour la plupart des applications, et les résultats
de Schlessinger sont effectivement utilisés dans la plupart des articles
qui font appel à ce type de théorie des déformations formelles.

\medskip\noindent
Soit $\Lambda$ un anneau local noethérien complet de corps résiduel $k$,
et notons $\mathcal{C}_\Lambda$ la catégorie des $\Lambda$-algèbres
artiniennes locales de corps résiduel $k$. Étant donné un foncteur
$F : \mathcal{C}_\Lambda \to \textit{Ensembles}$ tel que $F(k)$
soit un ensemble à un élément, l'article de Schlessinger introduit des
conditions (H1)--(H4) telles que :
\begin{enumerate}
\item $F$ admet une « enveloppe » si et seulement si (H1)--(H3) sont satisfaites.
\item $F$ est pro-représentable si et seulement si (H1)--(H4) sont satisfaites.
\end{enumerate}
Le but de ce chapitre est de généraliser ces résultats de deux manières,
exactement comme dans l'article de Rim :
\begin{enumerate}
\item[(A)] Le foncteur $F$ est remplacé par une catégorie $\mathcal{F}$
cofibrée en groupoïdes sur $\mathcal{C}_\Lambda$ ; voir la
section \ref{section-CLambda}.
\item[(B)] On prend pour $\Lambda$ un anneau noethérien et pour
$\Lambda \to k$ un homomorphisme fini d'anneaux à valeurs dans un corps.
La catégorie $\mathcal{C}_\Lambda$ est celle des $\Lambda$-algèbres
artiniennes locales $A$ munies d'une identification donnée
$A/\mathfrak m_A = k$.
\end{enumerate}
L'analogue de la condition selon laquelle $F(k)$ est un ensemble à un élément
est que $\mathcal{F}(k)$ soit le groupoïde trivial. Si $\mathcal{F}$ satisfait
cette condition, on dit que c'est une {\it catégorie de prédéformation}, mais
on ne fait pas cette hypothèse en général. L'article de Rim
\cite[Exposee VI]{SGA7-I} est la source originale des résultats de ce document.
Mentionnons aussi l'utile article \cite{Talpo-Vistoli}, qui traite de la théorie
des déformations avec des groupoïdes, mais dans un cadre moins général qu'ici.

\medskip\noindent
Le « complété » $\widehat{\mathcal{C}}_\Lambda$ de la catégorie
$\mathcal{C}_\Lambda$ joue un rôle important. Un objet de
$\widehat{\mathcal{C}}_\Lambda$ est une $\Lambda$-algèbre locale
noethérienne complète $R$ dont le corps résiduel est identifié à $k$ ; voir la
section \ref{section-category-completion-CLambda}.
D'une part, $\mathcal{C}_\Lambda \subset \widehat{\mathcal{C}}_\Lambda$
est une sous-catégorie strictement pleine et, d'autre part,
$\widehat{\mathcal{C}}_\Lambda$ est une sous-catégorie pleine de la catégorie
des pro-objets de $\mathcal{C}_\Lambda$. Un foncteur
$\mathcal{C}_\Lambda \to \textit{Ensembles}$ est {\it pro-représentable}
s'il est isomorphe à la restriction d'un foncteur représentable
$\underline{R} = \Mor_{\widehat{\mathcal{C}}_\Lambda}(R, -)$
à $\mathcal{C}_\Lambda$, où
$R \in \Ob(\widehat{\mathcal{C}}_\Lambda)$.

\medskip\noindent
Les {\it catégories cofibrées en groupoïdes} sont duales des catégories
fibrées en groupoïdes ; nous les introduisons à la section
\ref{section-preliminary}. Un morphisme {\it lisse} de catégories cofibrées
en groupoïdes sur $\mathcal{C}_\Lambda$ est un morphisme qui satisfait le
critère infinitésimal de relèvement pour les objets ; voir la
section \ref{section-smooth-morphisms}.
Cette définition est analogue à celle d'un homomorphisme d'anneaux
formellement lisse ; voir Algèbre, définition
\ref{algebra-definition-formally-smooth}. Elle est exactement duale de la
notion de Critères de représentabilité, section
\ref{criteria-section-formally-smooth}. Cette notion est importante, car nous
voulons finalement démontrer que certaines catégories cofibrées en groupoïdes
admettent une présentation lisse pro-représentable, à l'instar de la
caractérisation des champs algébriques dans Champs algébriques, sections
\ref{algebraic-section-stack-to-presentation} et
\ref{algebraic-section-smooth-groupoid-gives-algebraic-stack}.
Un {\it objet formel versel} d'une catégorie $\mathcal{F}$ cofibrée en
groupoïdes sur $\mathcal{C}_\Lambda$ est un objet
$\xi \in \widehat{\mathcal{F}}(R)$ du complété tel que le morphisme
associé
$\underline{\xi} : \underline{R}|_{\mathcal{C}_\Lambda} \to \mathcal{F}$
soit lisse.

\medskip\noindent
Dans la section \ref{section-schlessinger-conditions}, nous définissons sur
$\mathcal{F}$ les conditions (S1) et (S2), qui généralisent les conditions
(H1) et (H2) de Schlessinger. L'analogue de la condition (H3) de
Schlessinger---à savoir que l'espace tangent de $\mathcal{F}$ est de dimension
finie---ne reçoit pas de nom. Une étape essentielle du développement de la
théorie consiste à établir l'existence d'objets formels versels pour les
catégories de prédéformation qui satisfont (S1), (S2) et (H3) ; voir le
lemme \ref{lemma-versal-object-existence}.
La notion d'{\it enveloppe} d'un foncteur, au sens de Schlessinger,
$F : \mathcal{C}_\Lambda \to \textit{Ensembles}$
est, dans notre terminologie, un objet formel versel
$\xi \in \widehat{F}(R)$ tel que l'application induite sur les espaces tangents
$d\underline{\xi} : T\underline{R}|_{\mathcal{C}_\Lambda} \to TF$
soit un isomorphisme. Dans la littérature, une enveloppe est souvent appelée
objet « miniversel ». Nous n'employons pas ce terme, pour la raison suivante.
Il peut arriver qu'un foncteur possède un objet formel versel sans posséder
d'enveloppe. De plus, nous montrons à la section
\ref{section-minimal-versal} que, si une catégorie de prédéformation possède
un objet formel versel, elle en possède toujours un qui est {\it minimal}
(au sens de la définition \ref{definition-minimal-versal}) et qui est unique
à isomorphisme près ; voir le lemme \ref{lemma-minimal-versal}.
Il peut cependant arriver que l'objet formel versel minimal n'induise pas
un isomorphisme sur les espaces tangents ! (Voir les exemples
\ref{example-do-not-get-S2} et
\ref{example-smooth-continued}.)

\medskip\noindent
Compte tenu des différences signalées ci-dessus, le théorème
\ref{theorem-miniversal-object-existence} est la généralisation directe de
(1) : il redonne le résultat de Schlessinger lorsque $\mathcal{F}$ est un
foncteur et, sous les conditions (S1) et (S2), il caractérise les objets
formels versels minimaux au moyen de l'application
$d\underline{\xi} : T\underline{R}|_{\mathcal{C}_\Lambda} \to TF$
entre espaces tangents.

\medskip\noindent
À la section \ref{section-RS-condition}, nous définissons sur $\mathcal{F}$
la condition (RS) de Rim, qui généralise la condition (H4) de Schlessinger.
Une {\it catégorie de déformation} est, par définition, une catégorie de
prédéformation satisfaisant (RS). Les analogues des foncteurs
pro-représentables sont les catégories cofibrées en groupoïdes sur
$\mathcal{C}_\Lambda$ qui admettent une {\it présentation par un groupoïde
en foncteurs lisse et pro-représentable} sur $\mathcal{C}_\Lambda$ ; voir les
définitions \ref{definition-groupoid-in-functors},
\ref{definition-prorepresentable-groupoid-in-functors}, et
\ref{definition-smooth-groupoid-in-functors}.
Cette notion de présentation tient compte de la structure de groupoïde des
fibres de $\mathcal{F}$. Dans le théorème
\ref{theorem-presentation-deformation-groupoid}, nous démontrons que
$\mathcal{F}$ admet une présentation par un groupoïde en foncteurs lisse et
pro-représentable si et seulement si $\mathcal{F}$ possède un espace tangent et un espace
d'automorphismes infinitésimaux sont de dimension finie. C'est la
généralisation de (2) ci-dessus : elle se réduit au résultat de Schlessinger
lorsque $\mathcal{F}$ est un foncteur. Enfin, la section
\ref{section-minimality} explique comment utiliser les objets formels versels
minimaux pour construire une présentation minimale (unique à isomorphisme
près) par un groupoïde en foncteurs lisse et pro-représentable.

\medskip\noindent
Nous obtenons également l'explication conceptuelle suivante des conditions
de Schlessinger. Si une catégorie de prédéformation $\mathcal{F}$ satisfait
(RS), alors le foncteur associé des classes d'isomorphie
$\overline{\mathcal{F}}: \mathcal{C}_\Lambda \to \textit{Ensembles}$
satisfait (H1) et (H2) (lemmes \ref{lemma-RS-implies-S1-S2} et
\ref{lemma-S1-S2-associated-functor}).
Réciproquement, si un foncteur $F : \mathcal{C}_\Lambda \to \textit{Ensembles}$
se présente naturellement comme le foncteur des classes d'isomorphie d'une
catégorie $\mathcal{F}$ cofibrée en groupoïdes, on constate en pratique qu'un
argument établissant que $F$ satisfait (H1) et (H2) établit également que
$\mathcal{F}$ satisfait (RS). Des exemples sont étudiés dans Problèmes de
déformation, section
\ref{examples-defos-section-introduction}.
De plus, si $\mathcal{F}$ satisfait (RS), la condition (H4) pour
$\overline{\mathcal{F}}$ admet une interprétation simple en termes de
prolongement des automorphismes des objets de $\mathcal{F}$
(lemme \ref{lemma-RS-associated-functor}). Ces observations suggèrent de
considérer (RS) comme la condition fondamentale de recollement de la théorie
des déformations.




\section{Notations et conventions}
\label{section-notations-conventions}

\noindent
Les anneaux sont commutatifs et possèdent un élément unité $1$. L'idéal maximal d'un anneau local
$A$ est noté $\mathfrak{m}_A$. L'ensemble des entiers strictement positifs est
noté $\mathbf{N} = \{1, 2, 3, \ldots\}$. Si $U$ est un objet d'une catégorie
$\mathcal{C}$, nous notons $\underline{U}$ le foncteur
$\Mor_\mathcal{C}(U, -): \mathcal{C} \to \textit{Ensembles}$ ; voir les remarques
\ref{remarks-cofibered-groupoids} (\ref{item-definition-yoneda}). Attention :
cette convention peut entrer en conflit avec la notation d'autres chapitres,
où nous employons parfois $\underline{U}$ pour désigner
$h_U(-) = \Mor_\mathcal{C}(-, U)$.

\medskip\noindent
Dans tout ce chapitre, $\Lambda$ est un anneau noethérien et
$\Lambda \to k$ est un homomorphisme fini de $\Lambda$ vers un corps.
Le noyau de cet homomorphisme est noté $\mathfrak m_\Lambda$ et son image
$k' \subset k$. On verra que $\mathfrak m_\Lambda$ est un idéal maximal,
que $k' = \Lambda/\mathfrak m_\Lambda$ est un corps et que l'extension
$k/k'$ est finie. Voir la discussion autour de
(\ref{equation-k-prime}).


\section{La catégorie de base}
\label{section-CLambda}

\noindent
Motivation. Une application importante de la théorie des déformations formelles
concerne les critères de représentabilité par des espaces algébriques.
Supposons donnés une base localement noethérienne $S$ et un foncteur
$F : (\Sch/S)_{fppf}^{opp} \to \textit{Ensembles}$.
Soit $k$ un corps de type fini sur $S$, c'est-à-dire que l'on se donne un
morphisme de type fini $\Spec(k) \to S$.
L'un des critères d'Artin exige que, pour tout élément
$x \in F(\Spec(k))$, le foncteur de prédéformation associé au triplet
$(S, k, x)$ soit pro-représentable. D'après Morphismes, lemme
\ref{morphisms-lemma-point-finite-type}, dire que $k$ est de type fini sur
$S$ signifie qu'il existe un ouvert affine $\Spec(\Lambda) \subset S$ tel que
$k$ soit une $\Lambda$-algèbre finie. Cela explique le choix, dans tout ce
chapitre, de la catégorie de base suivante.

\begin{definition}
\label{definition-CLambda}
Soit $\Lambda$ un anneau noethérien et soit $\Lambda \to k$ un homomorphisme
fini d'anneaux, où $k$ est un corps. On définit {\it $\mathcal{C}_\Lambda$}
comme la catégorie dont
\begin{enumerate}
\item les objets sont les couples $(A, \varphi)$, où $A$ est une
$\Lambda$-algèbre artinienne locale et où
$\varphi : A/\mathfrak m_A \to k$ est un isomorphisme de
$\Lambda$-algèbres ;
\item les morphismes $f : (B, \psi) \to (A, \varphi)$ sont les
homomorphismes locaux de $\Lambda$-algèbres tels que
$\varphi \circ (f \bmod \mathfrak m) = \psi$.
\end{enumerate}
On dit que l'on est dans le {\it cas classique} si $\Lambda$ est un anneau
local noethérien complet et si $k$ est son corps résiduel.
\end{definition}

\noindent
Remarquons que, si $\Lambda \to k$ est surjectif et si $A$ est une
$\Lambda$-algèbre artinienne locale, l'identification $\varphi$, lorsqu'elle
existe, est unique. De plus, dans ce cas, tout homomorphisme de
$\Lambda$-algèbres $A \to B$ est compatible avec les identifications. Ainsi,
$\mathcal{C}_\Lambda$ n'est alors que la catégorie des $\Lambda$-algèbres
artiniennes locales dont le corps résiduel « est » $k$. Par abus de notation,
dans le cas général, les objets de $\mathcal{C}_\Lambda$ seront aussi désignés
simplement par $A$. Nous écrirons en outre souvent
$A/\mathfrak m = k$, autrement dit nous feindrons que tous les anneaux de
$\mathcal{C}_\Lambda$ ont pour corps résiduel $k$ (comme tous les
homomorphismes d'anneaux de $\mathcal{C}_\Lambda$ sont compatibles avec les
identifications données, cela ne devrait jamais poser de difficulté).
Dans toute la suite du chapitre, l'anneau de base $\Lambda$ et le corps $k$
sont fixés. La catégorie $\mathcal{C}_\Lambda$ sera la catégorie de base des
catégories cofibrées considérées ci-dessous.

\begin{definition}
\label{definition-small-extension}
Soit $f: B \to A$ un homomorphisme d'anneaux dans $\mathcal{C}_\Lambda$.
On dit que $f$ est une {\it petite extension} s'il est surjectif et si
$\Ker(f)$ est un idéal principal non nul, annulé par $\mathfrak{m}_B$.
\end{definition}

\noindent
Le lemme suivant permet souvent de ramener les arguments portant sur des
homomorphismes d'anneaux surjectifs dans $\mathcal{C}_\Lambda$ au cas des
petites extensions.

\begin{lemma}
\label{lemma-factor-small-extension}
Soit $f: B \to A$ un homomorphisme d'anneaux surjectif dans
$\mathcal{C}_\Lambda$. Alors $f$ se décompose en un composé de petites
extensions.
\end{lemma}

\begin{proof}
Soit $I$ le noyau de $f$. L'idéal maximal $\mathfrak{m}_B$ est nilpotent,
car $B$ est artinien ; écrivons $\mathfrak{m}_B^n = 0$. On obtient donc une
décomposition
$$
B = B/I\mathfrak{m}_B^{n-1} \to B/I\mathfrak{m}_B^{n-2} \to
\ldots \to B/I \cong A
$$
de $f$ en un composé d'homomorphismes surjectifs dont les noyaux sont annulés
par l'idéal maximal. Il suffit donc de démontrer le lemme lorsque $f$ est lui-même
un tel homomorphisme, c'est-à-dire lorsque $I$ est annulé par
$\mathfrak{m}_B$. Dans ce cas, $I$ est un espace vectoriel sur $k$, de
dimension finie ; voir Algèbre, lemme
\ref{algebra-lemma-artinian-finite-length}. En choisissant une base
$x_1, \ldots, x_n$ de l'espace vectoriel $I$ sur $k$, on obtient une
décomposition
$$
B \to B/(x_1) \to \ldots \to  B/(x_1, \ldots, x_n) \cong  A
$$
de $f$ en un composé de petites extensions.
\end{proof}

\noindent
Le lemme suivant affirme que l'on peut calculer la longueur d'un module sur
une $\Lambda$-algèbre locale de corps résiduel $k$ en fonction de sa longueur
sur $\Lambda$. Pour expliquer la notation de l'énoncé, soit $k' \subset k$
l'image de l'homomorphisme fini fixé $\Lambda \to k$. Remarquons que
$k' \subset k$ est une extension finie d'anneaux. Par conséquent, $k'$ est un
corps et $k/k'$ une extension finie de corps ; voir Algèbre, lemme
\ref{algebra-lemma-integral-under-field}. De plus, puisque
$\Lambda \to k'$ est surjectif, son noyau est un idéal maximal
$\mathfrak m_\Lambda$. Ainsi
\begin{equation}
\label{equation-k-prime}
[k : k'] = [k : \Lambda/\mathfrak m_\Lambda] < \infty
\end{equation}
et, dans le cas classique, on a $k = k'$. La notation
$k' = \Lambda/\mathfrak m_\Lambda$ sera fixée dans tout ce chapitre.

\begin{lemma}
\label{lemma-length}
Soit $A$ une $\Lambda$-algèbre locale de corps résiduel $k$.
Soit $M$ un $A$-module. Alors
$[k : k'] \text{length}_A(M) = \text{length}_\Lambda(M)$.
Dans le cas classique, on a
$\text{length}_A(M) = \text{length}_\Lambda(M)$.
\end{lemma}

\begin{proof}
Si $M$ est un $A$-module simple, alors $M \cong k$ comme $A$-module ; voir
Algèbre, lemme \ref{algebra-lemma-characterize-length-1}.
Dans ce cas, $\text{length}_A(M) = 1$ et
$\text{length}_\Lambda(M) = [k' : k]$ ; voir
Algèbre, lemme \ref{algebra-lemma-dimension-is-length}.
Si $\text{length}_A(M)$ est finie, le résultat s'obtient en choisissant une
filtration de $M$ par des sous-$A$-modules à quotients simples, puis en
utilisant l'additivité ; voir Algèbre, lemme
\ref{algebra-lemma-length-additive}. Si $\text{length}_A(M)$ est infinie, le
résultat découle de l'inégalité évidente
$\text{length}_A(M) \leq \text{length}_\Lambda(M)$.
\end{proof}

\begin{lemma}
\label{lemma-surjective}
Soit $A \to B$ un homomorphisme d'anneaux dans $\mathcal{C}_\Lambda$.
Les conditions suivantes sont équivalentes :
\begin{enumerate}
\item $f$ est surjectif ;
\item $\mathfrak m_A/\mathfrak m_A^2 \to \mathfrak m_B/\mathfrak m_B^2$
est surjectif ;
\item $\mathfrak m_A/(\mathfrak m_\Lambda A + \mathfrak m_A^2)
\to \mathfrak m_B/(\mathfrak m_\Lambda B + \mathfrak m_B^2)$ est surjectif.
\end{enumerate}
\end{lemma}

\begin{proof}
Pour tout homomorphisme d'anneaux $f : A \to B$ dans
$\mathcal{C}_\Lambda$, on a $f(\mathfrak m_A) \subset \mathfrak m_B$ ; cela
résulte par exemple du fait que $\mathfrak m_A$ et $\mathfrak m_B$ sont les
ensembles des éléments nilpotents de $A$ et de $B$. Supposons $f$ surjectif.
Soit $y \in \mathfrak m_B$. Choisissons $x \in A$ tel que $f(x) = y$.
Comme $f$ induit un isomorphisme
$A/\mathfrak m_A \to B/\mathfrak m_B$, on a $x \in \mathfrak m_A$.
L'homomorphisme induit
$\mathfrak m_A/\mathfrak m_A^2 \to \mathfrak m_B/\mathfrak m_B^2$
est donc surjectif. Ainsi, (1) implique (2).

\medskip\noindent
Il est clair que (2) implique (3). L'homomorphisme $A \to B$ donne le
diagramme commutatif canonique
$$
\xymatrix{
\mathfrak m_\Lambda/\mathfrak m_\Lambda^2 \otimes_{k'} k \ar[r] \ar[d] &
\mathfrak m_A/\mathfrak m_A^2 \ar[r] \ar[d] &
\mathfrak m_A/(\mathfrak m_\Lambda A + \mathfrak m_A^2) \ar[r] \ar[d] & 0 \\
\mathfrak m_\Lambda/\mathfrak m_\Lambda^2 \otimes_{k'} k \ar[r] &
\mathfrak m_B/\mathfrak m_B^2 \ar[r] &
\mathfrak m_B/(\mathfrak m_\Lambda B + \mathfrak m_B^2) \ar[r] & 0
}
$$
dont les lignes sont exactes. Ainsi, si (3) est satisfaite, (2) l'est aussi.

\medskip\noindent
Supposons (2). Pour montrer que $A \to B$ est surjectif, il suffit, d'après
le lemme de Nakayama (Algèbre, lemme \ref{algebra-lemma-NAK}), de montrer que
$A/\mathfrak m_A \to B/\mathfrak m_AB$ est surjectif. (Remarquons que
$\mathfrak m_A$ est un idéal nilpotent.) Comme
$k = A/\mathfrak m_A = B/\mathfrak m_B$, il suffit de montrer que
$\mathfrak m_AB \to \mathfrak m_B$ est surjectif. En appliquant une nouvelle
fois le lemme de Nakayama, on voit qu'il suffit que
$\mathfrak m_AB/\mathfrak m_A\mathfrak m_B \to \mathfrak m_B/\mathfrak m_B^2$
soit surjectif, ce qui est précisément l'hypothèse.
\end{proof}

\noindent
Si $A \to B$ est un homomorphisme d'anneaux dans $\mathcal{C}_\Lambda$, alors
l'homomorphisme
$\mathfrak m_A/(\mathfrak m_\Lambda A + \mathfrak m_A^2)
\to \mathfrak m_B/(\mathfrak m_\Lambda B + \mathfrak m_B^2)$
est l'homomorphisme induit sur les espaces cotangents relatifs. En voici la
définition formelle.

\begin{definition}
\label{definition-tangent-space-ring}
Soit $R \to S$ un homomorphisme local d'anneaux locaux. L'{\it espace
cotangent relatif}\footnote{Attention : nous verrons plus loin que, dans
notre cadre général, l'espace tangent d'un objet
$A \in \mathcal{C}_\Lambda$ sur $\Lambda$ ne doit pas être défini simplement
comme le dual $k$-linéaire de l'espace cotangent relatif. En réalité, la bonne
définition de l'espace cotangent relatif est
$\Omega_{S/R} \otimes_S S/\mathfrak m_S$.} de $R$ sur $S$ est l'espace
vectoriel sur $S/\mathfrak m_S$
$\mathfrak m_S/(\mathfrak m_R S + \mathfrak m_S^2)$.
\end{definition}

\noindent
Si $f_1: A_1 \to A$ et $f_2: A_2 \to A$ sont deux homomorphismes d'anneaux,
le produit fibré $A_1 \times_A A_2$ est le sous-anneau de $A_1 \times A_2$
formé des éléments dont les deux projections dans $A$ sont égales. Dans tout
ce chapitre, nous considérerons des conditions faisant intervenir un tel
produit fibré lorsque $f_1$ et $f_2$ appartiennent à
$\mathcal{C}_\Lambda$. Le produit fibré n'est pas toujours un objet de
$\mathcal{C}_\Lambda$.

\begin{example}
\label{example-fibre-product}
Soient $p$ un nombre premier et $n \in \mathbf{N}$.
Posons $\Lambda = \mathbf{F}_p(t_1, t_2, \ldots, t_n)$ et
$k = \mathbf{F}_p(x_1, \ldots, x_n)$, l'homomorphisme $\Lambda \to k$ étant
donné par $t_i \mapsto x_i^p$. Posons $A = k[\epsilon] = k[x]/(x^2)$.
Alors $A$ est un objet de $\mathcal{C}_\Lambda$. Supposons que
$D : k \to k$ soit une dérivation de $k$ sur $\Lambda$, par exemple
$D = \partial/\partial x_i$. Alors l'homomorphisme
$$
f_D : k \longrightarrow k[\epsilon], \quad
a \mapsto a + D(a)\epsilon
$$
est un morphisme de $\mathcal{C}_\Lambda$. Posons $A_1 = A_2 = k$, puis
$f_1 = f_{\partial/\partial x_1}$ et $f_2(a) = a$. Alors
$A_1 \times_A A_2 = \{a \in k \mid \partial/\partial x_1(a) = 0\}$
n'a pas pour image tout $k$. Le produit fibré n'est donc pas un objet de
$\mathcal{C}_\Lambda$.
\end{example}

\noindent
On verra que ce problème ne peut se produire que si l'extension de corps
résiduels $k/k'$ (\ref{equation-k-prime}) est inséparable et si ni $f_1$ ni
$f_2$ n'est surjectif.

\begin{lemma}
\label{lemma-fiber-product-CLambda}
Soient $f_1 : A_1 \to A$ et $f_2 : A_2 \to A$ des homomorphismes d'anneaux
dans $\mathcal{C}_\Lambda$. Alors :
\begin{enumerate}
\item Si $f_1$ ou $f_2$ est surjectif, alors
$A_1 \times_A A_2$ appartient à $\mathcal{C}_\Lambda$.
\item Si $f_2$ est une petite extension, il en est de même de
$A_1 \times_A A_2 \to A_1$.
\item Si l'extension de corps $k/k'$ est séparable, alors
$A_1 \times_A A_2$ appartient à $\mathcal{C}_\Lambda$.
\end{enumerate}
\end{lemma}

\begin{proof}
L'anneau $A_1 \times_A A_2$ est une $\Lambda$-algèbre par l'homomorphisme
$\Lambda \to A_1 \times_A A_2$ induit par les homomorphismes
$\Lambda \to A_1$ et $\Lambda \to A_2$. C'est un anneau local dont l'unique
idéal maximal est
$$
\mathfrak m_{A_1} \times_{\mathfrak m_A} \mathfrak m_{A_2} =
\Ker(A_1 \times_A A_2 \longrightarrow k)
$$
Un anneau est artinien si et seulement s'il est de longueur finie comme module
sur lui-même ; voir Algèbre, lemme
\ref{algebra-lemma-artinian-finite-length}. Puisque $A_1$ et $A_2$ sont
artiniens, le lemme \ref{lemma-length} montre que
$\text{length}_\Lambda(A_1)$ et $\text{length}_\Lambda(A_2)$, donc aussi
$\text{length}_\Lambda(A_1 \times A_2)$, sont toutes finies. Comme
$A_1 \times_A A_2 \subset A_1 \times A_2$ est un sous-$\Lambda$-module, il
vient que
$\text{length}_{A_1 \times_A A_2}(A_1 \times_A A_2) \leq
\text{length}_\Lambda(A_1 \times_A A_2)$ est finie. Ainsi, $A_1
\times_A A_2$ est artinien. Le seul obstacle à ce que
$A_1 \times_A A_2$ soit un objet de $\mathcal{C}_\Lambda$ est donc la
possibilité que son corps résiduel s'envoie dans un sous-corps propre de $k$
par l'homomorphisme ci-dessus
$A_1 \times_A A_2 \to A \to A/\mathfrak m_A = k$.

\medskip\noindent
Démonstration de (1). Si $f_2$ est surjectif, la projection
$A_1 \times_A A_2 \to A_1$ est surjective. Par suite, le composé
$A_1 \times_A A_2 \to A_1 \to A_1/\mathfrak m_{A_1} = k$ est surjectif, et
l'on conclut que $A_1 \times_A A_2$ est un objet de $\mathcal{C}_\Lambda$.

\medskip\noindent
Démonstration de (2). Si $f_2$ est une petite extension, alors
$A_2 \to A$ et $A_1 \times_A A_2  \to A_1$ sont tous deux surjectifs et ont
le même noyau. Le noyau de $A_1 \times_A A_2  \to A_1$ est donc un espace
vectoriel de dimension $1$ sur $k$, et l'on voit que
$A_1 \times_A A_2  \to A_1$ est une petite extension.

\medskip\noindent
Démonstration de (3). Choisissons $\overline{x} \in k$ tel que
$k = k'(\overline{x})$ (voir Corps, lemme
\ref{fields-lemma-primitive-element}). Soit $P'(T) \in k'[T]$ le polynôme
minimal de $\overline{x}$ sur $k'$. Puisque $k/k'$ est séparable, on a
$\text{d}P/\text{d}T(\overline{x}) \not = 0$.
Choisissons un polynôme unitaire $P \in \Lambda[T]$ dont l'image est $P'$ par
l'homomorphisme surjectif $\Lambda[T] \to k'[T]$. Comme
$A, A_1, A_2$ sont henséliens, voir Algèbre, lemme
\ref{algebra-lemma-local-dimension-zero-henselian}, il existe
$x, x_1, x_2 \in A, A_1, A_2$ tels que $P(x) = 0, P(x_1) = 0,
P(x_2) = 0$ et tels que les images de $x, x_1, x_2$ dans $k$ soient toutes
$\overline{x}$.
Alors $(x_1, x_2) \in A_1 \times_A A_2$, car $x_1, x_2$ s'envoient sur
$x \in A$ par unicité ; voir Algèbre, lemme
\ref{algebra-lemma-uniqueness}. Ainsi, le corps résiduel de
$A_1 \times_A A_2$ contient un générateur de $k$ sur $k'$, ce qui conclut.
\end{proof}

\noindent
Définissons maintenant les surjections essentielles dans
$\mathcal{C}_\Lambda$. Une condition nécessaire et suffisante pour qu'une
surjection de $\mathcal{C}_\Lambda$ soit essentielle est donnée au lemme
\ref{lemma-essential-surjection}.

\begin{definition}
\label{definition-essential-surjection}
Soit $f: B \to A$ un homomorphisme d'anneaux dans $\mathcal{C}_\Lambda$.
On dit que $f$ est une {\it surjection essentielle} s'il possède les
propriétés suivantes :
\begin{enumerate}
\item $f$ est surjectif.
\item Si $g: C \to B$ est un homomorphisme d'anneaux dans
$\mathcal{C}_\Lambda$ tel que $f \circ g$ soit surjectif, alors $g$ est
surjectif.
\end{enumerate}
\end{definition}

\noindent
À l'aide du lemme \ref{lemma-surjective}, on peut caractériser comme suit les
surjections essentielles dans $\mathcal{C}_\Lambda$.

\begin{lemma}
\label{lemma-essential-surjection-mod-squares}
Soit $f: B \to A$ un homomorphisme d'anneaux dans $\mathcal{C}_\Lambda$.
Les conditions suivantes sont équivalentes :
\begin{enumerate}
\item $f$ est une surjection essentielle ;
\item l'homomorphisme $B/\mathfrak m_B^2 \to A/\mathfrak m_A^2$ est une
surjection essentielle ;
\item l'homomorphisme
$B/(\mathfrak m_\Lambda B + \mathfrak m_B^2) \to
A/(\mathfrak m_\Lambda A + \mathfrak m_A^2)$ est une surjection essentielle.
\end{enumerate}
\end{lemma}

\begin{proof}
Supposons (3). Soit $C \to B$ un homomorphisme d'anneaux dans
$\mathcal{C}_\Lambda$ tel que $C \to A$ soit surjectif. Alors
$C \to A/(\mathfrak m_\Lambda A + \mathfrak m_A^2)$ est également surjectif.
L'hypothèse entraîne que
$C \to B/(\mathfrak m_\Lambda B + \mathfrak m_B^2)$ est surjectif. Par deux
applications du lemme \ref{lemma-surjective}, $C \to B$ est donc surjectif.

\medskip\noindent
Supposons (1). Soit $C \to B/(\mathfrak m_\Lambda B + \mathfrak m_B^2)$ un
morphisme de $\mathcal{C}_\Lambda$ tel que
$C \to A/(\mathfrak m_\Lambda A + \mathfrak m_A^2)$ soit surjectif. Posons
$C' = C \times_{B/(\mathfrak m_\Lambda B + \mathfrak m_B^2)} B$
qui est un objet de $\mathcal{C}_\Lambda$ d'après le lemme
\ref{lemma-fiber-product-CLambda}. Remarquons que
$C' \to A/(\mathfrak m_\Lambda A + \mathfrak m_A^2)$ est encore surjectif ;
le lemme \ref{lemma-surjective} montre donc que $C' \to A$ est surjectif.
Notre hypothèse implique alors que $C' \to B$ est surjectif. Il en résulte que
$C' \to B/(\mathfrak m_\Lambda B + \mathfrak m_B^2)$ est surjectif, ce qui,
par construction de $C'$, implique que
$C \to B/(\mathfrak m_\Lambda B + \mathfrak m_B^2)$ est surjectif.

\medskip\noindent
Dans le premier paragraphe, nous avons démontré (3) $\Rightarrow$ (1), et dans
le second, (1) $\Rightarrow$ (3). L'équivalence de (2) et (3) est un cas
particulier de l'équivalence de (1) et (3), ce qui conclut.
\end{proof}

\noindent
Pour analyser plus avant les surjections essentielles dans
$\mathcal{C}_\Lambda$, introduisons quelques notations. Supposons que $A$ soit
un objet de $\mathcal{C}_\Lambda$, ou plus généralement une
$\Lambda$-algèbre quelconque munie d'une surjection de $\Lambda$-algèbres
$A \to k$. On dispose d'une suite exacte canonique
\begin{equation}
\label{equation-sequence}
\mathfrak m_A/\mathfrak m_A^2 \xrightarrow{\text{d}_A}
\Omega_{A/\Lambda} \otimes_A k \to
\Omega_{k/\Lambda} \to 0
\end{equation}
voir Algèbre, lemme \ref{algebra-lemma-differential-seq}.
Remarquons que $\Omega_{k/\Lambda} = \Omega_{k/k'}$, avec $k'$ comme dans
(\ref{equation-k-prime}). Soit $H_1(L_{k/\Lambda})$ le premier module
d'homologie du complexe cotangent naïf de $k$ sur $\Lambda$ ; voir Algèbre,
définition \ref{algebra-definition-naive-cotangent-complex}. On peut alors
prolonger (\ref{equation-sequence}) en la suite exacte
\begin{equation}
\label{equation-sequence-extended}
H_1(L_{k/\Lambda}) \to
\mathfrak m_A/\mathfrak m_A^2 \xrightarrow{\text{d}_A}
\Omega_{A/\Lambda} \otimes_A k \to
\Omega_{k/\Lambda} \to 0,
\end{equation}
voir Algèbre, lemme \ref{algebra-lemma-exact-sequence-NL}.
Si $B \to A$ est un homomorphisme d'anneaux dans $\mathcal{C}_\Lambda$, ou,
plus généralement, un homomorphisme de $\Lambda$-algèbres munies de
surjections de $\Lambda$-algèbres sur $k$, on obtient un diagramme commutatif
\begin{equation}
\label{equation-sequence-functorial}
\vcenter{
\xymatrix{
H_1(L_{k/\Lambda}) \ar[r] \ar@{=}[d] &
\mathfrak m_B/\mathfrak m_B^2 \ar[r]_{\text{d}_B} \ar[d] &
\Omega_{B/\Lambda} \otimes_B k \ar[r] \ar[d] &
\Omega_{k/\Lambda} \ar[r] \ar@{=}[d] & 0 \\
H_1(L_{k/\Lambda}) \ar[r] &
\mathfrak m_A/\mathfrak m_A^2 \ar[r]^{\text{d}_A} &
\Omega_{A/\Lambda} \otimes_A k \ar[r] &
\Omega_{k/\Lambda} \ar[r] & 0
}
}
\end{equation}
dont les lignes sont exactes.

\begin{lemma}
\label{lemma-H1-separable-case}
Il existe un homomorphisme canonique
$$
\mathfrak m_\Lambda/\mathfrak m_\Lambda^2 \longrightarrow H_1(L_{k/\Lambda}).
$$
Si $k' \subset k$ est séparable (par exemple si la caractéristique de $k$ est
nulle), cet homomorphisme induit un isomorphisme
$\mathfrak m_\Lambda/\mathfrak m_\Lambda^2 \otimes_{k'} k = H_1(L_{k/\Lambda})$.
Si $k = k'$ (par exemple dans le cas classique), alors
$\mathfrak m_\Lambda/\mathfrak m_\Lambda^2 = H_1(L_{k/\Lambda})$.
Le composé
$$
\mathfrak m_\Lambda/\mathfrak m_\Lambda^2 \longrightarrow
H_1(L_{k/\Lambda}) \longrightarrow \mathfrak m_A/\mathfrak m_A^2
$$
provient de l'homomorphisme canonique $\mathfrak m_\Lambda \to \mathfrak m_A$.
\end{lemma}

\begin{proof}
On a $H_1(L_{k'/\Lambda}) = \mathfrak m_\Lambda/\mathfrak m_\Lambda^2$, car
$\Lambda \to k'$ est surjectif de noyau $\mathfrak m_\Lambda$.
L'homomorphisme provient de la fonctorialité du complexe cotangent naïf.
Si $k' \subset k$ est séparable, alors $k' \to k$ est un homomorphisme
d'anneaux \'etale ; voir Algèbre, lemme
\ref{algebra-lemma-etale-over-field}. Son complexe cotangent naïf a donc des
groupes d'homologie triviaux ; voir Algèbre, définition
\ref{algebra-definition-etale}. Appliqué aux homomorphismes d'anneaux
$\Lambda \to k' \to k$, le lemme
\ref{algebra-lemma-exact-sequence-NL} d'Algèbre entraîne alors que
$\mathfrak m_\Lambda/\mathfrak m_\Lambda^2 \otimes_{k'} k = H_1(L_{k/\Lambda})$.
Nous omettons la démonstration de la dernière assertion.
\end{proof}

\begin{lemma}
\label{lemma-essential-surjection}
Soit $f: B \to A$ un homomorphisme d'anneaux dans $\mathcal{C}_\Lambda$.
On conserve les notations de (\ref{equation-sequence-functorial}).
\begin{enumerate}
\item Les conditions équivalentes du lemme \ref{lemma-surjective} qui
caractérisent la surjectivité de $f$ sont aussi équivalentes aux suivantes :
\begin{enumerate}
\item $\Im(\text{d}_B) \to \Im(\text{d}_A)$ est surjectif ;
\item l'homomorphisme $\Omega_{B/\Lambda} \otimes_B k \to
\Omega_{A/\Lambda} \otimes_A k$ est surjectif.
\end{enumerate}
\item Les conditions suivantes sont équivalentes :
\begin{enumerate}
\item $f$ est une surjection essentielle (voir le lemme
\ref{lemma-essential-surjection-mod-squares}) ;
\item l'homomorphisme $\Im(\text{d}_B) \to \Im(\text{d}_A)$ est un
isomorphisme ;
\item l'homomorphisme $\Omega_{B/\Lambda} \otimes_B k \to
\Omega_{A/\Lambda} \otimes_A k$ est un isomorphisme.
\end{enumerate}
\item Si $k/k'$ est séparable, alors $f$ est une surjection essentielle si et
seulement si l'homomorphisme
$\mathfrak m_B/(\mathfrak m_\Lambda B + \mathfrak m_B^2) \to
\mathfrak m_A/(\mathfrak m_\Lambda A + \mathfrak m_A^2)$
est un isomorphisme.
\item Si $f$ est une petite extension, alors $f$ n'est pas essentielle si et
seulement si $f$ admet une section $s: A \to B$ dans $\mathcal{C}_\Lambda$
telle que $f \circ s = \text{id}_A$.
\end{enumerate}
\end{lemma}

\begin{proof}
Démonstration de (1). Il résulte de (\ref{equation-sequence-functorial}) que
(1)(a) et (1)(b) sont équivalentes. En outre, si $A \to B$ est surjectif,
alors (1)(a) et (1)(b) sont satisfaites. Supposons (1)(a). Puisque le noyau de
$\text{d}_A$ est l'image de $H_1(L_{k/\Lambda})$, lequel s'envoie aussi dans
$\mathfrak m_B/\mathfrak m_B^2$, on en déduit que
$\mathfrak m_B/\mathfrak m_B^2 \to \mathfrak m_A/\mathfrak m_A^2$
est surjectif. Ainsi, $B \to A$ est surjectif d'après le lemme
\ref{lemma-surjective}, ce qui achève la démonstration de (1).

\medskip\noindent
Démonstration de (2). L'équivalence de (2)(b) et (2)(c) résulte immédiatement
de (\ref{equation-sequence-functorial}).

\medskip\noindent
Supposons (2)(b). Soit $g : C \to B$ un homomorphisme d'anneaux dans
$\mathcal{C}_\Lambda$ tel que $f \circ g$ soit surjectif. On en déduit que
$\mathfrak m_C/\mathfrak m_C^2 \to \mathfrak m_A/\mathfrak m_A^2$
est surjectif d'après le lemme \ref{lemma-surjective}. Par conséquent,
$\Im(\text{d}_C) \to \Im(\text{d}_A)$ est surjectif et, par hypothèse,
$\Im(\text{d}_C) \to \Im(\text{d}_B)$ l'est aussi. Il résulte de (1) que
$C \to B$ est surjectif.

\medskip\noindent
Supposons (2)(a). Alors $f$ est surjectif et l'on voit que
$\Omega_{B/\Lambda} \otimes_B k \to \Omega_{A/\Lambda} \otimes_A k$
est surjectif. Soit $K$ son noyau. Remarquons que
$K = \text{d}_B(\Ker(\mathfrak m_B/\mathfrak m_B^2 \to
\mathfrak m_A/\mathfrak m_A^2))$ d'après
(\ref{equation-sequence-functorial}). Choisissons une décomposition
$$
\Omega_{B/\Lambda} \otimes_B k =
\Omega_{A/\Lambda} \otimes_A k \oplus K
$$
d'espaces vectoriels sur $k$. L'homomorphisme
$\text{d} : B \to \Omega_{B/\Lambda}$ induit, par projection sur $K$, un
homomorphisme $D : B \to K$. Posons
$C = \{b \in B \mid D(b) = 0\}$. La règle de Leibniz montre que c'est une
sous-$\Lambda$-algèbre de $B$. Soit $\overline{x} \in k$. Choisissons
$x \in B$ d'image $\overline{x}$. Si $D(x) \not = 0$, il existe un élément
$y \in \mathfrak m_B$ tel que $D(y) = D(x)$. Ainsi, $x - y \in C$ est un
élément d'image $\overline{x}$. Par conséquent, $C \to k$ est surjectif et
$C$ est un objet de $\mathcal{C}_\Lambda$. De même, choisissons
$\omega \in \Im(\text{d}_A)$. D'après (1), il existe
$x \in \mathfrak m_B$ tel que $\text{d}_B(x)$ s'envoie sur $\omega$. Si
$D(x) \not = 0$, il existe un élément $y \in \mathfrak m_B$ dont l'image dans
$\mathfrak m_A/\mathfrak m_A^2$ est nulle et tel que $D(y) = D(x)$.
Alors $z = x - y$ est un élément de $\mathfrak m_C$ dont l'image
$\text{d}_C(z) \in \Omega_{C/k} \otimes_C k$ s'envoie sur $\omega$.
Ainsi, $\Im(\text{d}_C) \to \Im(\text{d}_A)$ est surjectif. On déduit de
(1) que $C \to A$ est surjectif, puis de l'hypothèse que $C \to B$ est
surjectif. Par suite, $D = 0$, c'est-à-dire $K = 0$, c'est-à-dire que (2)(c)
est satisfaite. Cela achève la démonstration de (2).

\medskip\noindent
Démonstration de (3). Si $k'/k$ est séparable, alors
$H_1(L_{k/\Lambda}) =
\mathfrak m_\Lambda/\mathfrak m_\Lambda^2 \otimes_{k'} k$ ; voir le lemme
\ref{lemma-H1-separable-case}. Par conséquent, $\Im(\text{d}_A) =
\mathfrak m_A/(\mathfrak m_\Lambda A + \mathfrak m_A^2)$
et il en va de même pour $B$. Ainsi, (3) résulte de (2).

\medskip\noindent
Démonstration de (4). Une section $s$ de $f$ n'est pas surjective (par
définition, une petite extension a un noyau non trivial) ; ainsi, $f$ n'est pas
une surjection essentielle. Réciproquement, supposons que $f$ soit une petite
extension, mais non une surjection essentielle. Choisissons un homomorphisme
d'anneaux $C \to B$ dans $\mathcal{C}_\Lambda$, non surjectif, tel que
$C \to A$ soit surjectif. Soit $C' \subset B$ l'image de $C \to B$. Alors
$C' \not = B$, mais l'application de $C'$ vers $A$ est surjective. Puisque
$f : B \to A$ est une petite extension,
$\text{length}_C(B) = \text{length}_C(A) + 1$. Ainsi,
$\text{length}_C(C') \leq \text{length}_C(A)$, car $C'$ est un sous-anneau
propre de $B$. Or $C' \to A$ est surjectif ; on a donc nécessairement
$\text{length}_C(C') = \text{length}_C(A)$, et $C' \to A$ est un
isomorphisme qui fournit la section cherchée.
\end{proof}

\begin{example}
\label{example-essential-surjection}
Soit $\Lambda = k[[x]]$ l'anneau des séries formelles en $1$ variable sur $k$.
Posons $A = k$ et $B = \Lambda/(x^2)$. Alors $B \to A$ est une surjection
essentielle d'après le lemme \ref{lemma-essential-surjection}, car c'est une
petite extension et l'homomorphisme $B \to A$ n'admet pas d'inverse à droite
(dans la catégorie $\mathcal{C}_\Lambda$). Mais l'homomorphisme
$$
k \cong \mathfrak m_B/\mathfrak m_B^2
\longrightarrow
\mathfrak m_A/\mathfrak m_A^2 = 0
$$
n'est pas un isomorphisme. Ainsi, dans le lemme
\ref{lemma-essential-surjection} (3), il est nécessaire de considérer
l'homomorphisme entre espaces cotangents relatifs
$\mathfrak m_B/(\mathfrak m_\Lambda B + \mathfrak m_B^2) \to
\mathfrak m_A/(\mathfrak m_\Lambda A + \mathfrak m_A^2)$.
\end{example}







\section{La catégorie de base complétée}
\label{section-category-completion-CLambda}

\noindent
Le « complété » suivant de la catégorie $\mathcal{C}_\Lambda$ servira de
catégorie de base au complété d'une catégorie cofibrée en groupoïdes sur
$\mathcal{C}_\Lambda$ (section \ref{section-formal-objects}).

\begin{definition}
\label{definition-completion-CLambda}
Soit $\Lambda$ un anneau noethérien et soit $\Lambda \to k$ un homomorphisme
fini d'anneaux, où $k$ est un corps. On définit
{\it $\widehat{\mathcal{C}}_\Lambda$} comme la catégorie dont
\begin{enumerate}
\item les objets sont les couples $(R, \varphi)$, où $R$ est une
$\Lambda$-algèbre locale noethérienne complète et où
$\varphi : R/\mathfrak m_R \to k$ est un isomorphisme de
$\Lambda$-algèbres ;
\item les morphismes $f : (S, \psi) \to (R, \varphi)$ sont les
homomorphismes locaux de $\Lambda$-algèbres tels que
$\varphi \circ (f \bmod \mathfrak m) = \psi$.
\end{enumerate}
\end{definition}

\noindent
Comme dans la discussion qui suit la définition \ref{definition-CLambda}, nous
désignerons en général un objet de $\widehat{\mathcal{C}}_\Lambda$ simplement
par $R$, l'identification
$R/\mathfrak m_R = k$ étant sous-entendue. Nous étudions dans cette section
quelques propriétés élémentaires des objets et des morphismes de la catégorie
$\widehat{\mathcal{C}}_\Lambda$, parallèlement à l'étude de la catégorie
$\mathcal{C}_\Lambda$ menée dans la section précédente.

\medskip\noindent
Observons d'abord que tout objet $A \in \mathcal{C}_\Lambda$ est un objet de
$\widehat{\mathcal{C}}_\Lambda$, puisqu'un anneau local artinien est toujours
noethérien et complet pour la topologie définie par son idéal maximal (qui est
un idéal nilpotent). De plus, il résulte clairement des définitions que
$\mathcal{C}_\Lambda \subset \widehat{\mathcal{C}}_\Lambda$
est la sous-catégorie strictement pleine formée de tous les anneaux artiniens.
Réciproquement, tout objet de $\widehat{\mathcal{C}}_\Lambda$ est en fait une
limite d'objets de $\mathcal{C}_\Lambda$.

\medskip\noindent
Supposons que $R$ soit un objet de $\widehat{\mathcal{C}}_\Lambda$.
Considérons les anneaux $R_n = R/\mathfrak m_R^n$, pour
$n \in \mathbf{N}$. Ce sont des anneaux locaux noethériens qui possèdent un
unique idéal premier, lequel est nilpotent ; ils sont donc artiniens, voir
Algèbre, proposition \ref{algebra-proposition-dimension-zero-ring}.
Les homomorphismes d'anneaux
$$
\ldots \to R_{n + 1} \to R_n \to \ldots \to R_2 \to R_1 = k
$$
sont tous surjectifs. Par définition, la complétude de $R$ signifie que
$R = \lim R_n$. Si $f : R \to S$ est un homomorphisme d'anneaux dans
$\widehat{\mathcal{C}}_\Lambda$, on obtient un système d'homomorphismes
$f_n : R_n \to S_n$ dont la limite est l'homomorphisme donné.

\begin{lemma}
\label{lemma-surjective-cotangent-space}
Soit $f: R \to S$ un homomorphisme d'anneaux dans
$\widehat{\mathcal{C}}_\Lambda$. Les conditions suivantes sont équivalentes :
\begin{enumerate}
\item $f$ est surjectif ;
\item l'homomorphisme
$\mathfrak m_R/\mathfrak m_R^2 \to \mathfrak m_S/\mathfrak m_S^2$
est surjectif ;
\item l'homomorphisme
$\mathfrak m_R/(\mathfrak m_\Lambda R + \mathfrak m_R^2) \to
\mathfrak m_S/(\mathfrak m_\Lambda S + \mathfrak m_S^2)$
est surjectif.
\end{enumerate}
\end{lemma}

\begin{proof}
Remarquons que, pour $n \geq 2$, on a l'égalité d'espaces cotangents relatifs
$$
\mathfrak m_R/(\mathfrak m_\Lambda R + \mathfrak m_R^2)
=
\mathfrak m_{R_n}/(\mathfrak m_\Lambda R_n + \mathfrak m_{R_n}^2)
$$
et de même pour $S$. Le lemme \ref{lemma-surjective} montre donc que
$R_n \to S_n$ est surjectif pour tout $n$.
Soit maintenant $K_n$ le noyau de $R_n \to S_n$. Les suites
$$
0 \to K_n \to R_n \to S_n \to 0
$$
forment une suite exacte de systèmes projectifs dirigés. Le système $(K_n)$
vérifie la condition de Mittag-Leffler, puisque chaque $K_n$ est artinien.
D'après Algèbre, lemme \ref{algebra-lemma-ML-exact-sequence}, le passage à la
limite préserve donc l'exactitude. Ainsi,
$\lim R_n \to \lim S_n$ est surjectif, c'est-à-dire que $f$ est surjectif.
\end{proof}

\begin{lemma}
\label{lemma-CLambdahat-pushouts}
La catégorie $\widehat{\mathcal{C}}_\Lambda$ admet des sommes amalgamées.
\end{lemma}

\begin{proof}
Soient $R \to S_1$ et $R \to S_2$ des morphismes de
$\widehat{\mathcal{C}}_\Lambda$. Considérons l'anneau
$C = S_1 \otimes_R S_2$. Cet anneau possède l'idéal maximal de type fini
$\mathfrak m = \mathfrak m_{S_1} \otimes S_2 +
S_1 \otimes \mathfrak m_{S_2}$, de corps résiduel $k$.
Soit $C^\wedge$ le complété de $C$ pour la topologie $\mathfrak m$-adique.
Alors $C^\wedge$ est un anneau noethérien complet pour la topologie définie
par l'idéal maximal $\mathfrak m^\wedge = \mathfrak mC^\wedge$, dont le corps
résiduel est identifié à $k$ ; voir Algèbre, lemme
\ref{algebra-lemma-completion-Noetherian}. Ainsi, $C^\wedge$ est un objet de
$\widehat{\mathcal{C}}_\Lambda$. Les homomorphismes
$S_1 \to C^\wedge$ et $S_2 \to C^\wedge$ font alors de $C^\wedge$ une somme
amalgamée au-dessus de $R$ dans $\widehat{\mathcal{C}}_\Lambda$ (détails omis).
\end{proof}

\noindent
Nous n'aurons pas besoin du lemme suivant.

\begin{lemma}
\label{lemma-CLambdahat-coproducts}
La catégorie $\widehat{\mathcal{C}}_\Lambda$ admet les coproduits de deux
objets.
\end{lemma}

\begin{proof}
Soient $R$ et $S$ des objets de $\widehat{\mathcal{C}}_\Lambda$.
Considérons l'anneau $C = R \otimes_\Lambda S$. Il existe un homomorphisme
surjectif canonique $C \to R \otimes_\Lambda S \to k \otimes_\Lambda k \to k$,
le dernier homomorphisme étant celui de multiplication. Le noyau de
$C \to k$ est un idéal maximal $\mathfrak m$. Remarquons que $\mathfrak m$
est engendré par $\mathfrak m_R C$, $\mathfrak m_S C$ et par un nombre fini
d'éléments de $C$ dont les images engendrent le noyau de
$k \otimes_\Lambda k \to k$. Ainsi, $\mathfrak m$ est un idéal de type fini.
Soit $C^\wedge$ le complété de $C$ pour la topologie $\mathfrak m$-adique.
Alors $C^\wedge$ est un anneau noethérien complet pour la topologie définie
par l'idéal maximal $\mathfrak m^\wedge = \mathfrak mC^\wedge$, de corps
résiduel $k$ ; voir Algèbre, lemme
\ref{algebra-lemma-completion-Noetherian}. Ainsi, $C^\wedge$ est un objet de
$\widehat{\mathcal{C}}_\Lambda$. Les homomorphismes $R \to C^\wedge$ et
$S \to C^\wedge$ font alors de $C^\wedge$ un coproduit dans
$\widehat{\mathcal{C}}_\Lambda$ (détails omis).
\end{proof}

\noindent
Dans une catégorie, un coproduit vide est un objet initial. Dans le cas
classique, $\widehat{\mathcal{C}}_\Lambda$ possède un objet initial, à savoir
$\Lambda$ lui-même. Plus généralement, si $k' = k$, alors le complété
$\Lambda^\wedge$ de $\Lambda$ pour la topologie $\mathfrak m_\Lambda$-adique
est un objet initial. Plus généralement encore, si $k' \subset k$ est
séparable, alors $\widehat{\mathcal{C}}_\Lambda$ possède aussi un objet
initial. En effet, choisissons un polynôme unitaire $P \in \Lambda[T]$ tel que
$k \cong k'[T]/(P')$, où $p' \in k'[T]$ est l'image de $P$. Alors
$R = \Lambda^\wedge[T]/(P)$ est un objet initial ; voir la démonstration du
lemme \ref{lemma-fiber-product-CLambda}.

\medskip\noindent
Si $R$ est un objet initial comme ci-dessus, alors on a
$\mathcal{C}_\Lambda = \mathcal{C}_R$ et
$\widehat{\mathcal{C}}_\Lambda = \widehat{\mathcal{C}}_R$, ce qui ramène en
fait toute la discussion de ce chapitre au cas classique. En revanche, si
$k' \subset k$ est inséparable, il n'existe pas d'objet initial.

\begin{lemma}
\label{lemma-derivations-finite}
Soit $S$ un objet de $\widehat{\mathcal{C}}_\Lambda$.
Alors $\dim_k \text{Der}_\Lambda(S, k) < \infty$.
\end{lemma}

\begin{proof}
Choisissons $x_1, \ldots, x_n \in \mathfrak m_S$ dont les images forment une
base sur $k$ de l'espace cotangent relatif
$\mathfrak m_S/(\mathfrak m_\Lambda S + \mathfrak m_S^2)$.
Choisissons $y_1, \ldots, y_m \in S$ dont les images dans $k$ engendrent $k$
sur $k'$. Nous affirmons que $\dim_k \text{Der}_\Lambda(S, k) \leq n + m$.
Pour le voir, il suffit de démontrer que, si $D(x_i) = 0$ et
$D(y_j) = 0$, alors $D = 0$. Soit $a \in S$. Il existe un polynôme
$P = \sum \lambda_J y^J$, avec $\lambda_J \in \Lambda$, dont l'image dans
$k$ est égale à celle de $a$ dans $k$. On voit alors que
$D(a - P) = D(a) - D(P) = D(a)$ grâce à l'hypothèse
$D(y_j) = 0$ pour tout $j$. On peut donc supposer
$a \in \mathfrak m_S$. Écrivons $a = \sum a_i x_i$, avec $a_i \in S$.
La règle de Leibniz donne
$$
D(a) = \sum x_iD(a_i) + \sum a_iD(x_i) = \sum x_iD(a_i)
$$
car nous avons supposé $D(x_i) = 0$. On a $\sum x_iD(a_i) = 0$, puisque la
multiplication par $x_i$ est nulle sur $k$.
\end{proof}

\begin{lemma}
\label{lemma-derivations-surjective}
Soit $f : R \to S$ un morphisme de $\widehat{\mathcal{C}}_\Lambda$.
Si $\text{Der}_\Lambda(S, k) \to \text{Der}_\Lambda(R, k)$ est injectif,
alors $f$ est surjectif.
\end{lemma}

\begin{proof}
Si $f$ n'est pas surjectif, alors
$\mathfrak m_S/(\mathfrak m_R S + \mathfrak m_S^2)$ est non nul d'après le
lemme \ref{lemma-surjective-cotangent-space}. Par conséquent,
$Q = S/(f(R) + \mathfrak m_R S + \mathfrak m_S^2)$ est lui aussi non nul.
Remarquons que $Q$ est un espace vectoriel sur $k = R/\mathfrak m_R$ par
$f$. Munissons $Q$ d'une structure de $S$-module par
$S \to k$. L'application quotient $D : S \to Q$ est une $R$-dérivation : si
$a_1, a_2 \in S$, on peut écrire $a_1 = f(b_1) + a_1'$ et
$a_2 = f(b_2) + a_2'$, avec $b_1, b_2 \in R$ et
$a_1', a_2' \in \mathfrak m_S$. Alors
Les éléments $b_i$ et $a_i$ ont la même image dans $k$ pour $i = 1, 2$, et
\begin{align*}
a_1a_2 & = (f(b_1) + a_1')(f(b_2) + a_2') \\
& = f(b_1)a_2' + f(b_2)a_1' \\
& = f(b_1)(f(b_2) + a_2') + f(b_2)(f(b_1) + a_1') \\
& = f(b_1)a_2 + f(b_2)a_1
\end{align*}
dans $Q$, ce qui démontre la règle de Leibniz. Ainsi, $D : S \to Q$ est une
$\Lambda$-dérivation dont le composé avec $R \to S$ est nul. Puisque
$Q \not = 0$, il existe aussi des dérivations $D : S \to k$ dont le composé
avec $R \to S$ est nul, c'est-à-dire que
$\text{Der}_\Lambda(S, k) \to \text{Der}_\Lambda(R, k)$ n'est pas injectif.
\end{proof}

\begin{lemma}
\label{lemma-m-adic-topology}
Soit $R$ un objet de $\widehat{\mathcal{C}}_\Lambda$. Soit $(J_n)$ une suite
décroissante d'idéaux telle que $\mathfrak m_R^n \subset J_n$.
Posons $J = \bigcap J_n$. Alors la suite $(J_n/J)$ définit la topologie
$\mathfrak m_{R/J}$-adique sur $R/J$.
\end{lemma}

\begin{proof}
Il est clair que $\mathfrak m_{R/J}^n \subset J_n/J$. Il suffit donc de
montrer que, pour tout $n$, il existe un $N$ tel que
$J_N/J \subset \mathfrak m_{R/J}^n$. Cela équivaut à
$J_N \subset \mathfrak m_R^n + J$. Pour tout $n$, l'anneau
$R/\mathfrak m_R^n$ est artinien ; il existe donc un $N_n$ tel que
$$
J_{N_n} + \mathfrak m_R^n = J_{N_n + 1} + \mathfrak m_R^n = \ldots
$$
Posons $E_n = (J_{N_n} + \mathfrak m_R^n)/\mathfrak m_R^n$.
Posons $E = \lim E_n \subset \lim R/\mathfrak m_R^n = R$.
Remarquons que $E \subset J$ : en effet, pour tout $f \in E$ et tout $m$, on
a $f \in J_m + \mathfrak m_R^n$ pour tout $n \gg 0$, donc $f \in J_m$ par le
théorème d'intersection de Krull ; voir Algèbre, lemme
\ref{algebra-lemma-intersect-powers-ideal-module-zero}. Comme les
homomorphismes de transition $E_n \to E_{n - 1}$ sont tous surjectifs, on voit
que $J$ se surjecte sur $E_n$. Ainsi, $N = N_n$ convient.
\end{proof}

\begin{lemma}
\label{lemma-limit-artinian}
Soit $\ldots \to A_3 \to A_2 \to A_1$ une suite d'homomorphismes d'anneaux
surjectifs dans $\mathcal{C}_\Lambda$. Si
$\dim_k (\mathfrak m_{A_n}/\mathfrak m_{A_n}^2)$ est bornée, alors
$S = \lim A_n$ est un objet de $\widehat{\mathcal{C}}_\Lambda$, et les idéaux
$I_n = \Ker(S \to A_n)$ définissent la topologie $\mathfrak m_S$-adique sur
$S$.
\end{lemma}

\begin{proof}
Nous utiliserons librement le fait que les homomorphismes $S \to A_n$ sont
surjectifs pour tout $n$. Remarquons que les homomorphismes
$\mathfrak m_{A_{n + 1}}/\mathfrak m_{A_{n + 1}}^2 \to
\mathfrak m_{A_n}/\mathfrak m_{A_n}^2$ sont surjectifs ; voir le lemme
\ref{lemma-surjective-cotangent-space}. Par conséquent, pour $n$ assez grand,
la dimension $\dim_k (\mathfrak m_{A_n}/\mathfrak m_{A_n}^2)$ se stabilise à
un entier, disons $r$. Il existe donc
$x_1, \ldots, x_r \in \mathfrak m_S$ dont les images dans $A_n$ engendrent
$\mathfrak m_{A_n}$. Choisissons en outre $y_1, \ldots, y_t \in S$ dont les
images dans $k$ engendrent $k$ sur $\Lambda$. On obtient alors un homomorphisme
d'anneaux
$P = \Lambda[z_1, \ldots, z_{r + t}] \to S$, $z_i \mapsto x_i$ et
$z_{r + j} \mapsto y_j$ tel que le composé $P \to S \to A_n$ soit surjectif
pour tout $n$. Soit $\mathfrak m \subset P$ le noyau de $P \to k$.
Soit $R = P^\wedge$ le complété $\mathfrak m$-adique de $P$ ; c'est un objet
de $\widehat{\mathcal{C}}_\Lambda$. Comme on dispose encore du système
compatible d'homomorphismes surjectifs $R \to A_n$, on obtient un
homomorphisme $R \to S$. Posons $J_n = \Ker(R \to A_n)$.
Posons $J = \bigcap J_n$. Le lemme \ref{lemma-m-adic-topology} montre que
$R/J = \lim R/J_n = \lim A_n = S$ et que les idéaux $J_n/J = I_n$ définissent
la topologie $\mathfrak m$-adique. (Remarquons que, pour tout $n$, on a
$\mathfrak m_R^{N_n} \subset J_n$ pour un certain $N_n$, sans avoir
nécessairement $N_n = n$ ; il peut donc être nécessaire de renuméroter les
idéaux $J_n$ avant d'appliquer le lemme.)
\end{proof}

\begin{lemma}
\label{lemma-power-series}
Soient $R', R \in \Ob(\widehat{\mathcal{C}}_\Lambda)$. Supposons que
$R = R' \oplus I$ pour un idéal $I$ de $R$. Soient
$x_1, \ldots, x_r \in I$ des éléments dont les images forment une base de
$I/\mathfrak m_R I$. Posons $S = R'[[X_1, \ldots, X_r]]$ et considérons
l'homomorphisme de $R'$-algèbres $S \to R$ qui envoie $X_i$ sur $x_i$.
Supposons que, pour tout $n \gg 0$, l'homomorphisme
$S/\mathfrak m_S^n \to R/\mathfrak m_R^n$ admette un inverse à gauche dans
$\mathcal{C}_\Lambda$. Alors $S \to R$ est un isomorphisme.
\end{lemma}

\begin{proof}
Puisque $R = R' \oplus I$, on a
$$
\mathfrak m_R/\mathfrak m_R^2 =
\mathfrak m_{R'}/\mathfrak m_{R'}^2 \oplus I/\mathfrak m_RI
$$
et, de même,
$$
\mathfrak m_S/\mathfrak m_S^2 =
\mathfrak m_{R'}/\mathfrak m_{R'}^2 \oplus \bigoplus kX_i
$$
Par conséquent, pour $n > 1$, l'homomorphisme
$S/\mathfrak m_S^n \to R/\mathfrak m_R^n$ induit un isomorphisme sur les
espaces cotangents. Ainsi, un inverse à gauche
$h_n : R/\mathfrak m_R^n \to S/\mathfrak m_S^n$ est surjectif d'après le
lemme \ref{lemma-surjective-cotangent-space}. Comme $h_n$ est injectif en tant
qu'inverse à gauche, c'est un isomorphisme. Les surjections canoniques
$S/\mathfrak m_S^n \to R/\mathfrak m_R^n$ sont donc toutes des isomorphismes,
ce qui conclut.
\end{proof}




\section{Catégories cofibrées en groupoïdes}
\label{section-preliminary}

\noindent
Pour développer la théorie, nous travaillons avec des catégories {\it cofibrées}
en groupoïdes. Nous supposons connues la définition et les propriétés
fondamentales des catégories {\it fibrées} en groupoïdes, voir
Catégories, section \ref{categories-section-fibred-groupoids}.

\begin{definition}
\label{definition-category-cofibred-groupoids}
Soit $\mathcal{C}$ une catégorie. Une {\it catégorie cofibrée en groupoïdes
au-dessus de $\mathcal{C}$} est une catégorie $\mathcal{F}$ munie d'un foncteur
$p: \mathcal{F} \to \mathcal{C}$ tel que $\mathcal{F}^{opp}$ soit une catégorie
fibrée en groupoïdes au-dessus de $\mathcal{C}^{opp}$ au moyen de
$p^{opp}: \mathcal{F}^{opp} \to \mathcal{C}^{opp}$.
\end{definition}

\noindent
Plus explicitement, le foncteur $p: \mathcal{F} \to \mathcal{C}$ est cofibré
en groupoïdes si les deux conditions suivantes sont satisfaites :
\begin{enumerate}
\item Pour tout morphisme $f: U \to V$ de $\mathcal{C}$ et tout objet
$x$ situé au-dessus de $U$, il existe un morphisme $x \to y$ de $\mathcal{F}$
situé au-dessus de $f$.
\item Pour toute paire de morphismes $a: x \to y$ et $b: x \to z$
de $\mathcal{F}$ et tout morphisme $f: p(y) \to p(z)$ tel que $p(b) = f
\circ p(a)$, il existe un unique morphisme $c: y \to z$ de $\mathcal
F$ situé au-dessus de $f$ tel que $b = c \circ a$.
\end{enumerate}

\begin{remarks}
\label{remarks-cofibered-groupoids}
Tout ce qui concerne les catégories fibrées en groupoïdes se transpose
directement au cadre cofibré. Les remarques suivantes ont pour objet de fixer
les notations. Soit $\mathcal{C}$ une catégorie.
\begin{enumerate}
\item Nous omettons souvent le foncteur $p: \mathcal{F} \to \mathcal{C}$ des
notations.
\item La catégorie fibre au-dessus d'un objet $U$ de $\mathcal{C}$ est notée
$\mathcal{F}(U)$. Ses objets sont ceux de $\mathcal{F}$ situés au-dessus de $U$
et ses morphismes sont ceux de $\mathcal{F}$ situés au-dessus de $\text{id}_U$.
Si $x, y$ sont des objets de $\mathcal{F}(U)$, nous écrivons parfois
$\Mor_U(x, y)$ pour $\Mor_{\mathcal{F}(U)}(x, y)$.
\item Les catégories fibres $\mathcal{F}(U)$ sont des groupoïdes, voir
Catégories, lemme \ref{categories-lemma-fibred-groupoids}.
Ainsi, tous les morphismes de $\mathcal{F}(U)$ sont des isomorphismes.
Nous écrivons parfois $\text{Aut}_U(x)$ pour $\Mor_{\mathcal{F}(U)}(x, x)$.
\item
\label{item-pushforward}
Soit $\mathcal{F}$ une catégorie cofibrée en groupoïdes au-dessus de
$\mathcal{C}$, soit $f: U \to V$ un morphisme de $\mathcal{C}$ et
soit $x \in \Ob(\mathcal{F}(U))$.
Une {\it image directe} de $x$ par $f$ est un morphisme
$x \to y$ de $\mathcal{F}$ situé au-dessus de $f$. Une image directe
est déterminée à un isomorphisme unique près (voir la discussion qui suit
Catégories, définition \ref{categories-definition-cartesian-over-C}).
Nous écrivons parfois $x \to f_*x$ pour désigner « l'image directe » de $x$
par $f$.
\item Un {\it choix d'images directes pour $\mathcal{F}$} consiste à choisir
une image directe de $x$ par $f$ pour toute paire $(x, f)$ comme ci-dessus. Nous
pouvons effectuer un tel choix pour $\mathcal{F}$ à l'aide de l'axiome du choix.
\item Soit $\mathcal{F}$ une catégorie cofibrée en groupoïdes au-dessus de
$\mathcal{C}$. Étant donné un choix d'images directes pour $\mathcal{F}$, il
existe un pseudo-foncteur associé $\mathcal{C} \to \textit{Groupoïdes}$.
Nous n'utiliserons jamais cette construction et n'en donnons donc pas les détails.
\item
\label{item-cofibered-morphism}
Un morphisme de catégories cofibrées en groupoïdes au-dessus de $\mathcal{C}$
est un foncteur commutant aux projections sur $\mathcal{C}$. Si $\mathcal{F}$
et $\mathcal{F}'$ sont des catégories cofibrées en groupoïdes au-dessus de
$\mathcal{C}$, nous notons les morphismes de $\mathcal{F}$ vers $\mathcal{F}'$
par $\Mor_\mathcal{C}(\mathcal{F}, \mathcal{F}')$.
\item
\label{item-definition-cofibered-groupoids-2-category}
Les catégories cofibrées en groupoïdes forment une $(2, 1)$-catégorie
$\text{Cof}(\mathcal{C})$. Ses 1-morphismes sont ceux décrits en
(\ref{item-cofibered-morphism}). Si $p : \mathcal{F} \to C$ et
$p': \mathcal{F}' \to \mathcal{C}$ sont les projections de catégories
cofibrées en groupoïdes et si $\varphi, \psi : \mathcal{F} \to \mathcal{F}'$
sont des $1$-morphismes, un 2-morphisme $t : \varphi \to \psi$ est un morphisme
de foncteurs tel que $p'(t_x) = \text{id}_{p(x)}$ pour tout $x \in \Ob(\mathcal{F})$.
\item
\label{item-construction-associated-cofibered-groupoid}
Soit $F : \mathcal{C} \to \textit{Groupoïdes}$ un foncteur. Il
existe une catégorie cofibrée en groupoïdes $\mathcal{F} \to \mathcal{C}$
associée à $F$ de la façon suivante. Un objet de $\mathcal{F}$ est une paire $(U, x)$
où $U \in \Ob(\mathcal{C})$ et $x \in \Ob(F(U))$. Un
morphisme $(U, x) \to (V, y)$ est une paire $(f, a)$ où
$f \in \Mor_\mathcal{C}(U, V)$ et
$a \in \Mor_{F(V)}(F(f)(x), y)$.
Le foncteur $\mathcal{F} \to \mathcal{C}$ envoie $(U, x)$ sur $U$. Voir
Catégories, section \ref{categories-section-presheaves-groupoids}.
\item
\label{item-associated-functor-isomorphism-classes}
Soit $\mathcal{F}$ cofibrée en groupoïdes au-dessus de $\mathcal{C}$.
Pour $U \in \Ob(\mathcal{C})$, posons $\overline{\mathcal{F}}(U)$ égal à
l'ensemble des classes d'isomorphie des objets de la catégorie $\mathcal{F}(U)$.
Si $f : U \to V$ est un morphisme de $\mathcal{C}$, nous obtenons alors une
application d'ensembles $\overline{\mathcal{F}}(U) \to \overline{\mathcal{F}}(V)$ en
envoyant la classe d'isomorphie de $x$ sur celle d'une image directe
$f_*x$ de $x$ ; voir (\ref{item-pushforward}). Alors
$\overline{\mathcal{F}} : \mathcal{C} \to \textit{Ensembles}$ est un
foncteur. De même, si $\varphi : \mathcal{F} \to \mathcal{G}$ est un
morphisme de catégories cofibrées, nous notons
$\overline{\varphi}: \overline{\mathcal{F}} \to  \overline{\mathcal{G}}$
le morphisme de foncteurs associé.
\item
\label{item-convention-cofibered-sets}
Soit $F: \mathcal{C} \to \textit{Ensembles}$ un foncteur. Nous pouvons regarder un
ensemble comme une catégorie discrète, c'est-à-dire comme un groupoïde dont les
seuls morphismes sont les identités. La construction
(\ref{item-construction-associated-cofibered-groupoid}) associe alors à $F$ une
catégorie cofibrée en ensembles. On obtient ainsi un foncteur pleinement fidèle
de la catégorie des foncteurs $\mathcal{C} \to \textit{Ensembles}$ vers la catégorie
des catégories cofibrées en groupoïdes au-dessus de $\mathcal{C}$.
Nous identifions la catégorie des foncteurs à son image par ce foncteur.
Ainsi, si $F : \mathcal{C} \to \textit{Ensembles}$ est un foncteur, nous notons aussi
par $F$ la catégorie cofibrée en ensembles associée ; et si
$\varphi : F \to G$ est un morphisme de foncteurs, nous notons encore $\varphi$
le morphisme correspondant de catégories cofibrées en ensembles, et réciproquement.
Voir Catégories, section \ref{categories-section-fibred-in-sets}.
\item
\label{item-definition-yoneda}
Soit $U$ un objet de $\mathcal{C}$. Nous notons $\underline{U}$ le
foncteur
$\Mor_\mathcal{C}(U, -): \mathcal{C} \to
\textit{Ensembles}$. Cela définit un foncteur pleinement fidèle de $\mathcal
C^{opp}$ vers la catégorie des foncteurs $\mathcal{C} \to
\textit{Ensembles}$. Ainsi, si $f : U \to V$ est un morphisme, nous pouvons
encore noter $f$ le morphisme induit $\underline{V}
\to \underline{U}$, et réciproquement.
\item
\label{item-fibre-product}
Produits fibrés de catégories cofibrées en groupoïdes : si $\mathcal{F}
\to \mathcal{H}$ et $\mathcal{G} \to \mathcal{H}$ sont des morphismes
de catégories cofibrées en groupoïdes au-dessus de $\mathcal{C}_\Lambda$, une
construction de leur 2-produit fibré est donnée par celle de leur
2-produit fibré comme catégories au-dessus de $\mathcal{C}_\Lambda$, décrite dans
Catégories, lemme \ref{categories-lemma-2-product-categories-over-C}.
\item
\label{item-product}
Produits de catégories cofibrées en groupoïdes : si $\mathcal{F}$ et
$\mathcal{G}$ sont des catégories cofibrées en groupoïdes au-dessus de
$\mathcal{C}_\Lambda$, leur produit est défini comme le $2$-produit fibré
$\mathcal{F} \times_{\mathcal{C}_\Lambda} \mathcal{G}$ décrit dans
Catégories, lemme \ref{categories-lemma-2-product-categories-over-C}.
\item
\label{item-definition-restricting-base-category}
Restriction de la catégorie de base : soit $p : \mathcal{F} \to \mathcal{C}$
une catégorie cofibrée en groupoïdes, et soit $\mathcal{C}'$ une sous-catégorie
pleine de $\mathcal{C}$. La restriction $\mathcal{F}|_{\mathcal{C}'}$
est la sous-catégorie pleine de $\mathcal{F}$ dont les objets sont situés
au-dessus des objets de $\mathcal{C}'$. C'est une catégorie cofibrée en
groupoïdes au moyen du foncteur
$p|_{\mathcal{C}'}: \mathcal{F}|_{\mathcal{C}'} \to \mathcal{C}'$.
\end{enumerate}
\end{remarks}









\section{Foncteurs pro-représentables et catégories de prédéformation}
\sectionmark{Foncteurs pro-représentables et prédéformations}
\label{section-cofibered-groupoids}

\noindent
Notre objectif fondamental est de comprendre les catégories cofibrées en
groupoïdes au-dessus de $\mathcal{C}_\Lambda$ et de
$\widehat{\mathcal{C}}_\Lambda$. Comme $\mathcal{C}_\Lambda$ est une
sous-catégorie pleine de $\widehat{\mathcal{C}}_\Lambda$, nous pouvons
restreindre les catégories cofibrées en groupoïdes au-dessus de
$\widehat{\mathcal{C}}_\Lambda$ à $\mathcal{C}_\Lambda$ ; voir les remarques
\ref{remarks-cofibered-groupoids}
(\ref{item-definition-restricting-base-category}).
Nous pouvons en particulier le faire pour les foncteurs, et notamment pour
les foncteurs représentables. Les foncteurs sur $\mathcal{C}_\Lambda$
ainsi obtenus sont appelés foncteurs pro-représentables.

\begin{definition}
\label{definition-prorepresentable}
Soit $F : \mathcal{C}_\Lambda \to \textit{Ensembles}$ un foncteur.
On dit que $F$ est {\it pro-représentable} s'il existe un isomorphisme
de foncteurs $F \cong \underline{R}|_{\mathcal{C}_\Lambda}$ pour un certain
$R \in \Ob(\widehat{\mathcal{C}}_\Lambda)$.
\end{definition}

\noindent
Remarquons que si $F : \mathcal{C}_\Lambda \to \textit{Ensembles}$ est
pro-représentable par $R \in \Ob(\widehat{\mathcal{C}}_\Lambda)$, alors
$$
F(k) = \Mor_{\widehat{\mathcal{C}}_\Lambda}(R, k) = \{*\}
$$
est réduit à un élément. Les catégories cofibrées en groupoïdes au-dessus de
$\mathcal{C}_\Lambda$ qui interviennent en théorie des déformations
satisferont souvent une condition analogue.

\begin{definition}
\label{definition-predeformation-category}
Une {\it catégorie de prédéformation} $\mathcal{F}$ est une catégorie
cofibrée en groupoïdes au-dessus de $\mathcal{C}_\Lambda$ telle que
$\mathcal{F}(k)$ soit équivalente à une catégorie possédant un seul objet et
un seul morphisme, c'est-à-dire que $\mathcal{F}(k)$ contient au moins un
objet et qu'il existe un unique morphisme entre deux objets quelconques. Un
{\it morphisme de catégories de prédéformation} est un morphisme de catégories
cofibrées en groupoïdes au-dessus de
$\mathcal{C}_\Lambda$.
\end{definition}

\noindent
Une catégorie de prédéformation possède la propriété suivante.
Soit $x_0 \in \Ob(\mathcal{F}(k))$. Tout objet de $\mathcal{F}$ est alors
muni d'un unique morphisme vers $x_0$. En effet, si $x$ est un objet de
$\mathcal{F}$ au-dessus de $A$, nous pouvons choisir une image directe
$x \to q_*x$, où $q : A \to k$ est le morphisme quotient. Il existe un unique
isomorphisme $q_*x \to x_0$, et la composée $x \to q_*x \to x_0$ est le
morphisme recherché.

\begin{remark}
\label{remark-predeformation-functor}
On dit qu'un foncteur $F: \mathcal{C}_\Lambda \to \textit{Ensembles}$ est un
{\it foncteur de prédéformation} si l'ensemble cofibré associé est une
catégorie de prédéformation, c'est-à-dire\ si $F(k)$ est un ensemble à un
seul élément. Ainsi,
si $\mathcal{F}$ est une catégorie de prédéformation, alors
$\overline{\mathcal{F}}$ est un foncteur de prédéformation.
\end{remark}

\begin{remark}
\label{remark-localize-cofibered-groupoid}
Soit $p : \mathcal{F} \to \mathcal{C}_\Lambda$ une catégorie cofibrée en
groupoïdes, et soit $x \in \Ob(\mathcal{F}(k))$. On note
$\mathcal{F}_x$ la catégorie des objets au-dessus de $x$.
Un objet de $\mathcal{F}_x$ est une flèche $y \to x$.
Un morphisme $(y \to x) \to (z \to x)$ dans $\mathcal{F}_x$ est un diagramme
commutatif
$$
\xymatrix{
y \ar[rr] \ar[dr] & & z \ar[dl] \\
& x &
}
$$
Il existe un foncteur d'oubli $\mathcal{F}_x \to \mathcal{F}$. On définit le
foncteur $p_x : \mathcal{F}_x \to \mathcal{C}_\Lambda$ comme la composée
$\mathcal{F}_x \to \mathcal{F} \xrightarrow{p} \mathcal{C}_\Lambda$.
Alors $p_x : \mathcal{F}_x \to \mathcal{C}_\Lambda$ est une catégorie de
prédéformation (démonstration omise). On peut ainsi passer d'une catégorie
cofibrée en groupoïdes arbitraire au-dessus de $\mathcal{C}_\Lambda$ à une
catégorie de prédéformation en tout $x \in \Ob(\mathcal{F}(k))$.
\end{remark}






\section{Objets formels et catégories complétées}
\label{section-formal-objects}

\noindent
Dans cette section, nous expliquons comment passer des catégories cofibrées
en groupoïdes au-dessus de $\mathcal{C}_\Lambda$ aux catégories cofibrées en
groupoïdes au-dessus de $\widehat{\mathcal{C}}_\Lambda$, et réciproquement.

\begin{definition}
\label{definition-formal-objects}
Soit $\mathcal{F}$ une catégorie cofibrée en groupoïdes au-dessus de
$\mathcal{C}_\Lambda$. La {\it catégorie $\widehat{\mathcal{F}}$ des objets
formels de $\mathcal{F}$} est la catégorie dont les objets et les morphismes
sont les suivants.
\begin{enumerate}
\item Un {\it objet formel $\xi = (R, \xi_n, f_n)$ de $\mathcal{F}$} est
constitué d'un objet $R$ de $\widehat{\mathcal{C}}_\Lambda$, d'une famille,
indexée par $n \in \mathbf{N}$, d'objets $\xi_n$ de
$\mathcal{F}(R/\mathfrak m_R^n)$, et de morphismes
$f_n : \xi_{n + 1} \to \xi_n$ situés au-dessus de la projection
$R/\mathfrak m_R^{n + 1} \to R/\mathfrak m_R^n$.
\item Soient $\xi = (R, \xi_n, f_n)$ et $\eta = (S, \eta_n, g_n)$ des
objets formels de $\mathcal{F}$. Un {\it morphisme d'objets formels
$a : \xi \to \eta$} est constitué d'un morphisme $a_0 : R \to S$ dans
$\widehat{\mathcal{C}}_\Lambda$ et d'une famille de morphismes
$a_n : \xi_n \to \eta_n$ de $\mathcal{F}$ situés au-dessus de
$R/\mathfrak m_R^n \to S/\mathfrak m_S^n$,
de telle sorte que, pour tout $n$, le diagramme
$$
\xymatrix{
\xi_{n + 1} \ar[r]_{f_n} \ar[d]_{a_{n + 1}} & \xi_n \ar[d]^{a_n} \\
\eta_{n + 1} \ar[r]^{g_n} & \eta_n
}
$$
soit commutatif.
\end{enumerate}
\end{definition}

\noindent
La catégorie des objets formels est munie d'un foncteur $\widehat{p}:
\widehat{\mathcal{F}} \to \widehat{\mathcal{C}}_\Lambda$ qui envoie un objet
$(R, \xi_n, f_n)$ sur $R$ et un morphisme
$(R, \xi_n, f_n) \to (S, \eta_n, g_n)$ sur le morphisme $R \to S$.

\begin{lemma}
\label{lemma-completion-cofibred}
Soit $p : \mathcal{F} \to \mathcal{C}_\Lambda$ une catégorie cofibrée en
groupoïdes. Alors
$\widehat{p} : \widehat{\mathcal{F}} \to \widehat{\mathcal{C}}_\Lambda$
est une catégorie cofibrée en groupoïdes.
\end{lemma}

\begin{proof}
Soit $R \to S$ un morphisme d'anneaux dans
$\widehat{\mathcal{C}}_\Lambda$. Soit $(R, \xi_n, f_n)$ un objet de
$\widehat{\mathcal{F}}$. Pour tout $n$, choisissons une image directe
$\xi_n \to \eta_n$ de $\xi_n$ le long de
$R/\mathfrak m_R^n \to S/\mathfrak m_S^n$. Pour tout $n$, il existe un unique
morphisme $g_n : \eta_{n + 1} \to \eta_n$ dans $\mathcal{F}$, situé au-dessus
de $S/\mathfrak m_S^{n + 1} \to S/\mathfrak m_S^n$, tel que
$$
\xymatrix{
\xi_{n + 1} \ar[d] \ar[r]_{f_n} & \xi_n \ar[d] \\
\eta_{n + 1} \ar[r]^{g_n} & \eta_n
}
$$
soit commutatif (par le premier axiome d'une catégorie cofibrée en
groupoïdes). Nous obtenons donc un morphisme
$(R, \xi_n, f_n) \to (S, \eta_n, g_n)$ situé au-dessus de $R \to S$ ; le
premier axiome d'une catégorie cofibrée en groupoïdes est ainsi vérifié pour
$\widehat{\mathcal{F}}$. Pour vérifier le second axiome, supposons donnés des
morphismes
$a : (R, \xi_n, f_n) \to (S, \eta_n, g_n)$ et
$b : (R, \xi_n, f_n) \to (T, \theta_n, h_n)$ dans $\widehat{\mathcal{F}}$
et un morphisme $c_0 : S \to T$ dans $\widehat{\mathcal{C}}_\Lambda$ tel que
$c_0 \circ a_0 = b_0$. Le second axiome d'une catégorie cofibrée en groupoïdes
appliqué à $\mathcal{F}$ fournit d'uniques morphismes
$c_n : \eta_n \to \theta_n$, situées au-dessus de
$S/\mathfrak m_S^n \to T/\mathfrak m_T^n$, telles que
$c_n \circ a_n = b_n$. En posant $c = (c_n)_{n \geq 0}$, on obtient le
morphisme recherché $c : (S, \eta_n, g_n) \to (T, \theta_n, h_n)$ dans
$\widehat{\mathcal{F}}$ (nous omettons la vérification de l'égalité
$h_n \circ c_{n + 1} = c_n \circ g_n$).
\end{proof}

\begin{definition}
\label{definition-completion}
Soit $p : \mathcal{F} \to \mathcal{C}_\Lambda$ une catégorie cofibrée en
groupoïdes. La catégorie cofibrée en groupoïdes
$\widehat{p} : \widehat{\mathcal  F} \to \widehat{\mathcal{C}}_\Lambda$
est appelée le {\it complété de $\mathcal{F}$}.
\end{definition}

\noindent
Si $\mathcal{F}$ est une catégorie cofibrée en groupoïdes au-dessus de $\mathcal
C_\Lambda$, nous avons défini $\widehat{\mathcal{F}}(R)$, pour $R \in
\Ob(\widehat{\mathcal{C}}_\Lambda)$, à l'aide de la filtration de $R$
par les puissances de son idéal maximal. Supposons maintenant que
$\mathcal{I} = (I_n)$ soit une filtration de $R$ par des idéaux qui induit la
topologie $\mathfrak{m}_R$-adique. On définit
$\widehat{\mathcal{F}}_\mathcal{I}(R)$ comme la catégorie ayant les objets et
les morphismes suivants :
\begin{enumerate}
\item Un objet est une famille $(\xi_n, f_n)_{n \in \mathbf{N}}$ d'objets
$\xi_n$ de $\mathcal{F}(R/I_n)$ et de morphismes
$f_n : \xi_{n + 1} \to \xi_n$ situés au-dessus des projections
$R/I_{n + 1} \to R/I_n$.
\item Un morphisme $a : (\xi_n, f_n) \to (\eta_n, g_n)$ est une famille de
morphismes $a_n : \xi_n \to \eta_n$ dans $\mathcal{F}(R/I_n)$ telle que,
pour tout $n$, le diagramme
$$
\xymatrix{
\xi_{n + 1} \ar[r]^{f_n} \ar[d]_{a_{n + 1}} & \xi_n \ar[d]^{a_n} \\
\eta_{n + 1} \ar[r]^{g_n} & \eta_n
}
$$
soit commutatif.
\end{enumerate}

\begin{lemma}
\label{lemma-formal-objects-different-filtration}
Dans la situation ci-dessus, $\widehat{\mathcal{F}}_\mathcal{I}(R)$ est
équivalente à la catégorie $\widehat{\mathcal{F}}(R)$.
\end{lemma}

\begin{proof}
On peut définir une équivalence
$\widehat{\mathcal{F}}_\mathcal{I}(R) \to \widehat{\mathcal{F}}(R)$ de la
manière suivante. Pour tout $n$, soit $m(n)$ le plus petit $m$ tel que
$I_m \subset \mathfrak m_R^n$. Étant donné un objet $(\xi_n, f_n)$ de
$\widehat{\mathcal{F}}_\mathcal{I}(R)$, soit $\eta_n$ l'image directe de
$\xi_{m(n)}$ le long de $R/I_{m(n)} \to R/\mathfrak m_R^n$. Soit
$g_n : \eta_{n + 1} \to \eta_n$ l'unique morphisme de $\mathcal{F}$ situé
au-dessus de $R/\mathfrak m_R^{n + 1} \to R/\mathfrak m_R^n$ tel que
$$
\xymatrix{
\xi_{m(n + 1)} \ar[rrr]_{f_{m(n)} \circ \ldots \circ f_{m(n + 1) - 1}} \ar[d]
& & & \xi_{m(n)} \ar[d] \\
\eta_{n + 1} \ar[rrr]^{g_n} & & & \eta_n
}
$$
soit commutatif (l'existence et l'unicité sont assurées par les axiomes d'une
catégorie cofibrée). Le foncteur
$\widehat{\mathcal{F}}_\mathcal{I}(R) \to \widehat{\mathcal{F}}(R)$
envoie $(\xi_n, f_n)$ sur $(R, \eta_n, g_n)$. Nous omettons la vérification
qu'il s'agit bien d'une équivalence de catégories.
\end{proof}

\begin{remark}
\label{remark-different-sequence-ideals}
Soit $p : \mathcal{F} \to \mathcal{C}_\Lambda$ une catégorie cofibrée en
groupoïdes. Supposons que, pour tout
$R \in \Ob(\widehat{\mathcal{C}}_\Lambda)$, on se soit donné une filtration
$\mathcal{I}_R$ de $R$ par des idéaux. Si $\mathcal{I}_R$ induit la topologie
$\mathfrak m_R$-adique sur $R$ pour tout $R$, on peut définir une catégorie
$\widehat{\mathcal{F}}_\mathcal{I}$ en calquant la définition de
$\widehat{\mathcal{F}}$. Cette catégorie est munie d'un morphisme
$\widehat{p}_\mathcal{I} : \widehat{\mathcal{F}}_\mathcal{I} \to
\widehat{\mathcal{C}}_\Lambda$ qui en fait une catégorie cofibrée en
groupoïdes et tel que $\widehat{\mathcal{F}}_\mathcal{I}(R)$ soit isomorphe à
$\widehat{\mathcal{F}}_{\mathcal{I}_R}(R)$, telle qu'elle a été définie
ci-dessus. Les catégories cofibrées en groupoïdes
$\widehat{\mathcal{F}}_\mathcal{I}$ et $\widehat{\mathcal{F}}$ sont
équivalentes : au-dessus d'un objet
$R \in \Ob(\widehat{\mathcal{C}}_\Lambda)$, on utilise l'équivalence du
lemme \ref{lemma-formal-objects-different-filtration}.
\end{remark}

\begin{remark}
\label{remark-completion-functor}
Soit $F: \mathcal{C}_\Lambda \to \textit{Ensembles}$ un foncteur. En identifiant
les foncteurs aux ensembles cofibrés, le complété de $F$ est le foncteur
$\widehat{F} : \widehat{\mathcal{C}}_\Lambda \to \textit{Ensembles}$
donné par $\widehat{F}(S) = \lim F(S/\mathfrak{m}_S^{n})$. Cela concorde avec
la définition de l'article de Schlessinger \cite{Sch}.
\end{remark}

\begin{remark}
\label{remark-restrict-completion}
Soit $\mathcal{F}$ une catégorie cofibrée en groupoïdes au-dessus de
$\mathcal{C}_\Lambda$. Nous affirmons qu'il existe une équivalence canonique
$$
can :
\widehat{\mathcal{F}}|_{\mathcal{C}_\Lambda}
\longrightarrow
\mathcal{F}.
$$
En effet, soient $A \in \Ob(\mathcal{C}_\Lambda)$ et
$(A, \xi_n, f_n)$ un objet de
$\widehat{\mathcal{F}}|_{\mathcal{C}_\Lambda}(A)$. Comme $A$ est artinien,
il existe un plus petit $m \in \mathbf{N}$ tel que
$\mathfrak m_A^m = 0$. Alors $can$ envoie $(A, \xi_n, f_n)$ sur $\xi_m$.
Ce foncteur est une équivalence de catégories cofibrées en groupoïdes en
vertu de Catégories, lemme
\ref{categories-lemma-equivalence-fibred-categories}, car il induit une
équivalence sur toutes les catégories fibres d'après le lemme
\ref{lemma-formal-objects-different-filtration} et parce que la topologie
$\mathfrak m_A$-adique d'un anneau local artinien $A$ provient de l'idéal
nul. Nous identifierons fréquemment $\mathcal{F}$ à une sous-catégorie pleine
de $\widehat{\mathcal{F}}$ au moyen d'un quasi-inverse du foncteur $can$.
\end{remark}

\begin{remark}
\label{remark-completion-morphism}
Soit $\varphi : \mathcal{F} \to \mathcal{G}$ un morphisme de catégories
cofibrées en groupoïdes au-dessus de $\mathcal{C}_\Lambda$. Il en résulte un
morphisme
$\widehat{\varphi}: \widehat{\mathcal{F}} \to \widehat{\mathcal{G}}$
de catégories cofibrées en groupoïdes au-dessus de
$\widehat{\mathcal{C}}_\Lambda$. Il envoie un objet
$\xi = (R, \xi_n, f_n)$ de $\widehat{\mathcal{F}}$ sur
$(R, \varphi(\xi_n), \varphi(f_n))$, et un morphisme
$(a_0 : R \to S, a_n : \xi_n \to \eta_n)$ entre des objets $\xi$ et $\eta$
de $\widehat{\mathcal{F}}$ sur
$(a_0 : R \to S, \varphi(a_n) : \varphi(\xi_n) \to \varphi(\eta_n))$.
Enfin, si $t : \varphi \to \varphi'$ est un $2$-morphisme entre des
$1$-morphismes $\varphi, \varphi': \mathcal{F} \to \mathcal{G}$ de
catégories cofibrées en groupoïdes, on obtient un $2$-morphisme
$\widehat{t} : \widehat{\varphi} \to \widehat{\varphi}'$. Plus précisément,
pour $\xi = (R, \xi_n, f_n)$ comme ci-dessus, on pose
$\widehat{t}_\xi = (t_{\varphi(\xi_n)})$. La complétion définit donc un
foncteur entre $2$-catégories
$$
\widehat{~} :
\text{Cof}(\mathcal{C}_\Lambda)
\longrightarrow
\text{Cof}(\widehat{\mathcal{C}}_\Lambda)
$$
de la $2$-catégorie des catégories cofibrées en groupoïdes au-dessus de
$\mathcal{C}_\Lambda$ vers la $2$-catégorie des catégories cofibrées en
groupoïdes au-dessus de $\widehat{\mathcal{C}}_\Lambda$.
\end{remark}

\begin{remark}
\label{remark-completion-restriction-adjoint}
Nous affirmons que le foncteur de complétion de la remarque
\ref{remark-completion-morphism} et le foncteur de restriction
$|_{\mathcal{C}_\Lambda} : \text{Cof}(\widehat{\mathcal{C}}_\Lambda)
\to \text{Cof}(\mathcal{C}_\Lambda)$ des remarques
\ref{remarks-cofibered-groupoids}
(\ref{item-definition-restricting-base-category})
ont une relation de « 2-adjonction » au sens précis suivant. Soient
$\mathcal{F} \in \Ob(\text{Cof}(\mathcal{C}_\Lambda))$ et
$\mathcal{G} \in \Ob(\text{Cof}(\widehat{\mathcal{C}}_\Lambda))$.
Il existe alors une équivalence de catégories
$$
\Phi :
\Mor_{\mathcal{C}_\Lambda}(
\mathcal{G}|_{\mathcal{C}_\Lambda}, \mathcal{F})
\longrightarrow
\Mor_{\widehat{\mathcal{C}}_\Lambda}(\mathcal{G}, \widehat{\mathcal{F}})
$$
Pour décrire cette équivalence, définissons des morphismes canoniques
$\mathcal{G} \to \widehat{\mathcal{G}|_{\mathcal{C}_\Lambda}}$ et
$\widehat{\mathcal{F}}|_{\mathcal{C}_\Lambda} \to \mathcal{F}$ comme suit :
\begin{enumerate}
\item Soient $R \in \Ob(\widehat{\mathcal{C}}_\Lambda))$ et $\xi$ un objet
de la catégorie fibre $\mathcal{G}(R)$. Choisissons une image directe
$\xi \to \xi_n$ de $\xi$ vers $R/\mathfrak m_R^n$ pour chaque
$n \in \mathbf{N}$, et soit $f_n : \xi_{n + 1} \to \xi_n$ le morphisme
induit. Alors $\mathcal{G} \to \widehat{\mathcal{G}|_{\mathcal{C}_\Lambda}}$
envoie $\xi$ sur
$(R, \xi_n, f_n)$.
\item Il s'agit de l'équivalence
$can : \widehat{\mathcal{F}}|_{\mathcal{C}_\Lambda} \to \mathcal{F}$
de la remarque \ref{remark-restrict-completion}.
\end{enumerate}
Cela étant, l'équivalence
$\Phi : \Mor_{\mathcal{C}_\Lambda}(
\mathcal{G}|_{\mathcal{C}_\Lambda}, \mathcal{F})  \to
\Mor_{\widehat{\mathcal{C}}_\Lambda}(\mathcal{G},
\widehat{\mathcal{F}})$
envoie un morphisme
$\varphi : \mathcal{G}|_{\mathcal{C}_\Lambda} \to \mathcal{F}$
sur
$$
\mathcal{G} \to \widehat{\mathcal{G}|_{\mathcal{C}_\Lambda}}
\xrightarrow{\widehat{\varphi}} \widehat{\mathcal{F}}
$$
Il existe un quasi-inverse
$\Psi :
\Mor_{\widehat{\mathcal{C}}_\Lambda}(
\mathcal{G}, \widehat{\mathcal{F}}) \to
\Mor_{\mathcal{C}_\Lambda}(
\mathcal{G}|_{\mathcal{C}_\Lambda}, \mathcal{F})$
de $\Phi$ qui envoie $\psi : \mathcal{G} \to \widehat{\mathcal{F}}$ sur
$$
\mathcal{G}|_{\mathcal{C}_\Lambda} \xrightarrow{\psi|_{\mathcal{C}_\Lambda}}
\widehat{\mathcal{F}}|_{\mathcal{C}_\Lambda} \to \mathcal{F}.
$$
Nous omettons de vérifier que $\Phi$ et $\Psi$ sont quasi-inverses. Nous
n'abordons pas non plus la fonctorialité de $\Phi$ (car elle nous conduirait
sur le terrain des 3-catégories, que nous voulons éviter à tout prix).
\end{remark}

\begin{remark}
\label{remark-completion-restriction-cofset-adjoint}
Pour une catégorie $\mathcal{C}$, on note $\text{CofSet}(\mathcal{C})$ la
catégorie des ensembles cofibrés au-dessus de $\mathcal{C}$. C'est une
$1$-catégorie isomorphe à la catégorie des foncteurs
$\mathcal{C} \to \textit{Ensembles}$. Voir les remarques
\ref{remarks-cofibered-groupoids}
(\ref{item-convention-cofibered-sets}).
Les foncteurs de complétion et de restriction induisent des foncteurs
$\widehat{~} : \text{CofSet}(\mathcal{C}_\Lambda) \to
\text{CofSet}(\widehat{\mathcal{C}}_\Lambda)$ et
$|_{\mathcal{C}_\Lambda} : \text{CofSet}(\widehat{\mathcal{C}}_\Lambda) \to
\text{CofSet}(\mathcal{C}_\Lambda)$, que l'on désigne par les mêmes symboles.
Considérées comme foncteurs entre catégories d'ensembles cofibrés, la
complétion et la restriction sont adjointes au sens 1-catégorique usuel :
la même construction que dans la remarque
\ref{remark-completion-restriction-adjoint} définit une bijection fonctorielle
$$
\Mor_{\mathcal{C}_\Lambda}(G|_{\mathcal{C}_\Lambda}, F)
\longrightarrow
\Mor_{\widehat{\mathcal{C}}_\Lambda}(G, \widehat{F})
$$
pour $F \in \Ob(\text{CofSet}(\mathcal{C}_\Lambda))$ et
$G \in \Ob(\text{CofSet}(\widehat{\mathcal{C}}_\Lambda))$.
Le morphisme $\widehat{F}|_{\mathcal{C}_\Lambda} \to F$ est, là encore, un
isomorphisme.
\end{remark}

\begin{remark}
\label{remark-restrict-complete-continuous-functor}
Soit $G : \widehat{\mathcal{C}}_\Lambda \to \textit{Ensembles}$ un foncteur qui
commute aux limites. Alors le morphisme
$G \to \widehat{G|_{\mathcal{C}_\Lambda}}$ décrite dans la remarque
\ref{remark-completion-restriction-adjoint} est un isomorphisme. En effet, si
$S$ est un objet de $\widehat{\mathcal{C}}_\Lambda$, on dispose de bijections
canoniques
$$
\widehat{G|_{\mathcal{C}_\Lambda}}(S) =
\lim_n G(S/\mathfrak{m}_S^n) =
G(\lim_n S/\mathfrak{m}_S^n) = G(S).
$$
En particulier, si $R$ est un objet de $\widehat{\mathcal{C}}_\Lambda$, alors
$\underline{R} = \widehat{\underline{R}|_{\mathcal{C}_\Lambda}}$, car le
foncteur représentable $\underline{R}$ commute aux limites par définition des
limites.
\end{remark}

\begin{remark}
\label{remark-formal-objects-yoneda}
Soit $R$ un objet de $\widehat{\mathcal{C}}_\Lambda$. Il définit un foncteur
$\underline{R}: \widehat{\mathcal{C}}_\Lambda \to \textit{Ensembles}$
comme décrit dans les remarques \ref{remarks-cofibered-groupoids}
(\ref{item-definition-yoneda}). Comme d'habitude, on identifie ce foncteur à
l'ensemble cofibré associé. Si $\mathcal{F}$ est une catégorie cofibrée
au-dessus de $\mathcal{C}_\Lambda$, il existe une équivalence de catégories
\begin{equation}
\label{equation-formal-objects-maps}
\Mor_{\mathcal{C}_\Lambda}(
\underline{R}|_{\mathcal{C}_\Lambda}, \mathcal{F})
\longrightarrow
\widehat{\mathcal{F}}(R).
\end{equation}
Elle est donnée par la composée
$$
\Mor_{\mathcal{C}_\Lambda}(
\underline{R}|_{\mathcal{C}_\Lambda}, \mathcal{F})
\xrightarrow{\Phi}
\Mor_{\widehat{\mathcal{C}}_\Lambda}(
\underline{R}, \widehat{\mathcal{F}})
\xrightarrow{\sim}
\widehat{\mathcal{F}}(R)
$$
où $\Phi$ est celle de la remarque
\ref{remark-completion-restriction-adjoint}, et où la seconde équivalence
provient du lemme de 2-Yoneda (l'analogue cofibré de Catégories, lemme
\ref{categories-lemma-yoneda-2category}). Explicitement, l'équivalence envoie
un morphisme
$\varphi : \underline{R}|_{\mathcal{C}_\Lambda} \to \mathcal{F}$
sur l'objet formel
$(R, \varphi(R \to R/\mathfrak{m}_R^n), \varphi(f_n))$ dans
$\widehat{\mathcal{F}}(R)$, où
$f_n : R/\mathfrak m_R^{n + 1} \to R/\mathfrak m_R^n$ est la projection.

\medskip\noindent
Supposons qu'un choix d'images directes ait été fait pour $\mathcal{F}$.
Pour tout $\xi \in \Ob(\widehat{\mathcal{F}}(R))$, nous construisons
explicitement un morphisme
$\underline{\xi} : \underline{R}|_{\mathcal{C}_\Lambda} \to \mathcal{F}$
qui s'envoie sur $\xi$ par (\ref{equation-formal-objects-maps}). Écrivons
$\xi = (R, \xi_n, f_n)$. Se donner un objet $\alpha$ de
$\underline{R}|_{\mathcal{C}_\Lambda}$ revient à se donner un morphisme
$\alpha : R \to A$ de $\widehat{\mathcal{C}}_\Lambda$, avec $A$ artinien.
Soit $m \in \mathbf{N}$ le plus petit entier tel que
$\mathfrak m_A^m = 0$. Alors $\alpha$ se factorise par un unique
$\alpha_m : R/\mathfrak m_R^m \to A$, et l'on peut poser
$\underline{\xi}(\alpha) = \alpha_{m, *}\xi_m$. Nous omettons la description
de $\underline{\xi}$ sur les morphismes, ainsi que la démonstration du fait
que $\underline{\xi}$ s'envoie sur $\xi$ par
(\ref{equation-formal-objects-maps}).

\medskip\noindent
Supposons qu'un choix d'images directes ait été fait pour
$\widehat{\mathcal{F}}$. Dans ce cas, la démonstration de Catégories, lemme
\ref{categories-lemma-yoneda-2category}, fournit un quasi-inverse explicite
$$
\iota :
\widehat{\mathcal{F}}(R) \longrightarrow
\Mor_{\widehat{\mathcal{C}}_\Lambda}(
\underline{R}, \widehat{\mathcal{F}})
$$
de l'équivalence de 2-Yoneda, qui envoie $\xi$ sur le morphisme
$\iota(\xi) : \underline{R} \to \widehat{\mathcal{F}}$ envoyant
$f \in \underline{R}(S) = \Mor_{\mathcal{C}_\Lambda}(R, S)$
sur $f_*\xi$. Un quasi-inverse de (\ref{equation-formal-objects-maps}) est
donc
$$
\widehat{\mathcal{F}}(R)
\xrightarrow{\iota}
\Mor_{\widehat{\mathcal{C}}_\Lambda}(
\underline{R}, \widehat{\mathcal{F}})
\xrightarrow{\Psi}
\Mor_{\mathcal{C}_\Lambda}(
\underline{R}|_{\mathcal{C}_\Lambda}, \mathcal{F})
$$
où $\Psi$ est celle de la remarque
\ref{remark-completion-restriction-adjoint}. Pour
$\xi \in \Ob(\widehat{\mathcal{F}}(R))$, on a
$\Psi(\iota(\xi)) \cong \underline{\xi}$, où $\underline{\xi}$ est comme au
paragraphe précédent, puisque tous deux s'envoient sur $\xi$ par
l'équivalence de catégories (\ref{equation-formal-objects-maps}). En utilisant
$\underline{R} = \widehat{\underline{R}|_{\mathcal{C}_\Lambda}}$ (voir la
remarque \ref{remark-restrict-complete-continuous-functor}) et en développant
les définitions de $\Phi$ et $\Psi$, on conclut que $\iota(\xi)$ est isomorphe
au complété de $\underline{\xi}$.
\end{remark}

\begin{remark}
\label{remark-formal-objects-yoneda-map}
Soit $\mathcal{F}$ une catégorie cofibrée en groupoïdes au-dessus de
$\mathcal{C}_\Lambda$. Soient $\xi = (R, \xi_n, f_n)$ et
$\eta = (S, \eta_n, g_n)$ des objets formels de $\mathcal{F}$. Soit
$a = (a_n) : \xi \to \eta$ un morphisme d'objets formels, c'est-à-dire un
morphisme de $\widehat{\mathcal{F}}$. Soit
$f = \widehat{p}(a) = a_0 : R \to S$ la projection de $a$ dans
$\widehat{\mathcal{C}}_\Lambda$. On obtient alors un diagramme
$2$-commutatif
$$
\xymatrix{
\underline{R}|_{\mathcal{C}_\Lambda} \ar[rd]_{\underline{\xi}} & &
\underline{S}|_{\mathcal{C}_\Lambda} \ar[ll]^f \ar[ld]^{\underline{\eta}} \\
& \mathcal{F}
}
$$
où $\underline{\xi}$ et $\underline{\eta}$ sont les morphismes construits
dans la remarque \ref{remark-formal-objects-yoneda}. Pour le voir, soit
$\alpha : S \to A$ un objet de
$\underline{S}|_{\mathcal{C}_\Lambda}$ (voir loc.\ cit.). Soit
$m \in \mathbf{N}$ le plus petit entier tel que $\mathfrak m_A^m = 0$. On
obtient un diagramme commutatif
$$
\xymatrix{
R \ar[d]^f \ar[r] & R/\mathfrak m_R^m \ar[d]_{f_m} \ar[rd]^{\beta_m} \\
S \ar[r] & S/\mathfrak m_S^m \ar[r]^{\alpha_m} & A
}
$$
tel que la composée des flèches du bas soit $\alpha$. Alors
$\underline{\eta}(\alpha) = \alpha_{m, *}\eta_m$ et
$\underline{\xi}(\alpha \circ f) = \beta_{m, *}\xi_m$. Comme le morphisme
$a_m : \xi_m \to \eta_m$ est situé au-dessus de $f_m$, on obtient un
morphisme canonique
$$
\underline{\xi}(\alpha \circ f) = \beta_{m, *}\xi_m
\longrightarrow
\underline{\eta}(\alpha) = \alpha_{m, *}\eta_m
$$
situé au-dessus de $\text{id}_A$ et tel que
$$
\xymatrix{
\xi_m \ar[r] \ar[d]^{a_m} &  \beta_{m, *}\xi_m \ar[d] \\
\eta_m \ar[r] &  \alpha_{m, *}\eta_m
}
$$
soit commutatif en vertu des axiomes d'une catégorie cofibrée en groupoïdes.
Ceci définit une transformation de foncteurs
$\underline{\xi} \circ f \to \underline{\eta}$ qui témoigne de la
2-commutativité du premier diagramme de cette remarque.
\end{remark}

\begin{remark}
\label{remark-spell-out-formal-object}
Selon la remarque \ref{remark-formal-objects-yoneda}, se donner un objet
formel $\xi$ de $\mathcal{F}$ équivaut à se donner un foncteur
pro-représentable $U : \mathcal{C}_\Lambda \to \textit{Ensembles}$ et un morphisme
$U \to \mathcal{F}$.
\end{remark}




\section{Morphismes lisses}
\label{section-smooth-morphisms}

\noindent
Dans cette section, nous étudions les morphismes lisses de catégories
cofibrées en groupoïdes au-dessus de $\mathcal{C}_\Lambda$.

\begin{definition}
\label{definition-smooth-morphism}
Soit $\varphi : \mathcal{F} \to \mathcal{G}$ un morphisme de catégories
cofibrées en groupoïdes au-dessus de $\mathcal{C}_\Lambda$. On dit que
$\varphi$ est {\it lisse} s'il satisfait la condition suivante. Soient
$B \to A$ un morphisme surjectif d'anneaux dans $\mathcal{C}_\Lambda$,
$y \in
\Ob(\mathcal{G}(B)), x \in \Ob(\mathcal{F}(A))$, et $y
\to \varphi(x)$ un morphisme situé au-dessus de $B \to A$. Alors il
existe $x' \in \Ob(\mathcal{F}(B))$, un morphisme $x' \to x$ situé au-dessus
de $B \to A$, et un morphisme $\varphi(x') \to y$ situé au-dessus de
$\text{id}: B \to B$, tels que le diagramme
$$
\xymatrix{
\varphi(x') \ar[r] \ar[dr] & y \ar[d] \\
& \varphi(x)
}
$$
soit commutatif.
\end{definition}

\begin{lemma}
\label{lemma-smoothness-small-extensions}
Soit $\varphi : \mathcal{F} \to \mathcal{G}$ un morphisme de catégories
cofibrées en groupoïdes au-dessus de $\mathcal{C}_\Lambda$. Alors $\varphi$
est lisse dès que la condition de la définition
\ref{definition-smooth-morphism} est supposée satisfaite seulement pour les
petites extensions $B \to A$.
\end{lemma}

\begin{proof}
Soient $B \to A$ un morphisme surjectif d'anneaux dans
$\mathcal{C}_\Lambda$, $y \in \Ob(\mathcal{G}(B))$,
$x \in \Ob(\mathcal{F}(A))$, et $y \to \varphi(x)$ un morphisme situé
au-dessus de $B \to A$. D'après le lemme
\ref{lemma-factor-small-extension}, on peut factoriser $B \to A$ en petites
extensions $B = B_n \to B_{n-1} \to \ldots \to B_0 = A$. Raisonnons par
récurrence sur $n$. Si $n = 1$, le résultat découle de l'hypothèse. Si
$n > 1$, notons $f : B = B_n \to B_{n - 1}$ et
$g : B_{n - 1} \to B_0 = A$. Choisissons une image directe $y \to f_* y$
de $y$ le long de $f$, de sorte que le morphisme $y \to \varphi(x)$ se
factorise en $y \to f_* y \to \varphi(x)$. Par l'hypothèse de récurrence, on
peut trouver $x_{n - 1} \to x$ situé au-dessus de
$g : B_{n - 1} \to A$ et $a : \varphi(x_{n - 1}) \to f_*y$ situé au-dessus
de $\text{id} : B_{n - 1} \to B_{n - 1}$ tels que
$$
\xymatrix{
\varphi(x_{n - 1}) \ar[r]_-a \ar[dr] & f_*y \ar[d] \\
& \varphi(x)
}
$$
soit commutatif. On peut appliquer l'hypothèse à la composée
$y \to \varphi(x_{n - 1})$ de $y \to f_*y$ et de
$a^{-1} : f_*y \to \varphi(x_{n - 1})$. On obtient
$x_n \to x_{n - 1}$ situé au-dessus de $B_n \to B_{n - 1}$ et
$\varphi(x_n) \to y$ situé au-dessus de $\text{id} : B_n \to B_n$, de sorte
que le diagramme
$$
\xymatrix{
\varphi(x_n) \ar[r] \ar[d] & y \ar[d] \\
\varphi(x_{n - 1}) \ar[r]^-a \ar[dr] & f_*y \ar[d] \\
& \varphi(x)
}
$$
soit commutatif. La composée $x_n \to x_{n - 1} \to x$ et le morphisme
$\varphi(x_n) \to y$ sont alors les morphismes exigés par la définition de la
lissité.
\end{proof}

\begin{remark}
\label{remark-smoothness-2-categorical}
Soit $\varphi : \mathcal{F} \to \mathcal{G}$ un morphisme de catégories
cofibrées en groupoïdes au-dessus de $\mathcal{C}_\Lambda$. Soit $B \to A$
un morphisme d'anneaux dans $\mathcal{C}_\Lambda$. Des choix d'images directes
le long de $B
\to A$ pour les objets des catégories fibres $\mathcal{F}(B)$
et $\mathcal{G}(B)$ déterminent des foncteurs
$\mathcal{F}(B) \to \mathcal{F}(A)$ et
$\mathcal{G}(B) \to \mathcal{G}(A)$ qui s'insèrent dans un diagramme
$2$-commutatif
$$
\xymatrix{
\mathcal{F}(B) \ar[r]^{\varphi} \ar[d] & \mathcal{G}(B) \ar[d] \\
\mathcal{F}(A) \ar[r]^{\varphi}        & \mathcal{G}(A) .
}
$$
Il en résulte un foncteur induit $\mathcal{F}(B) \to \mathcal{F}(A)
\times_{\mathcal{G}(A)} \mathcal{G}(B)$. En développant les définitions, on
voit que $\varphi : \mathcal{F} \to \mathcal{G}$ est lisse si et seulement
si ce foncteur induit est essentiellement surjectif chaque fois que
$B \to A$ est surjectif (ou, de manière équivalente d'après le lemme
\ref{lemma-smoothness-small-extensions}, chaque fois que $B \to A$ est une
petite extension).
\end{remark}

\begin{remark}
\label{remark-compare-smooth-schlessinger}
La caractérisation des morphismes lisses donnée dans la remarque
\ref{remark-smoothness-2-categorical} est analogue à la notion de morphisme
lisse de foncteurs due à Schlessinger ; cf.\ \cite[définition 2.2.]{Sch}. En
fait, lorsque $\mathcal{F}$ et $\mathcal{G}$ sont cofibrées en ensembles,
notre notion équivaut à celle de Schlessinger. Soient en effet, dans ce cas,
$F, G : \mathcal{C}_\Lambda \to \textit{Ensembles}$ les foncteurs correspondants ;
voir les remarques \ref{remarks-cofibered-groupoids}
(\ref{item-convention-cofibered-sets}).
Alors $F \to G$ est lisse si et seulement si, pour toute surjection d'anneaux
$B \to A$ dans $\mathcal{C}_\Lambda$, l'application
$F(B) \to F(A) \times_{G(A)} G(B)$ est surjective.
\end{remark}

\begin{remark}
\label{remark-smooth-to-iso-classes}
Soit $\mathcal{F}$ une catégorie cofibrée en groupoïdes au-dessus de
$\mathcal{C}_\Lambda$. Alors le morphisme
$\mathcal{F} \to \overline{\mathcal{F}}$ est lisse. Supposons en effet que
$f : B \to A$ soit un morphisme d'anneaux dans $\mathcal{C}_\Lambda$. Soient
$x \in \Ob(\mathcal{F}(A))$ et
$\overline{y} \in \overline{\mathcal{F}}(B)$
la classe d'isomorphie de $y \in \Ob(\mathcal{F}(B))$, telle que
$\overline{f_*y} = \overline{x}$. Il suffit alors de prendre $x' = y$, le
morphisme induit $x' = y \to x$ au-dessus de $B \to A$, et l'égalité
$\overline{x'} = \overline{y}$ pour résoudre le problème posé dans la
définition \ref{definition-smooth-morphism}.
\end{remark}

\noindent
Si $R \to S$ est un morphisme d'anneaux de
$\widehat{\mathcal{C}}_\Lambda$, il induit un morphisme
$\underline{S} \to \underline{R}$ entre les foncteurs
$\underline{S}, \underline{R}: \widehat{\mathcal{C}}_\Lambda
\to \textit{Ensembles}$. Dans cette situation, la lissité de la restriction
$\underline{S}|_{\mathcal{C}_\Lambda} \to
\underline{R}|_{\mathcal{C}_\Lambda}$ est une notion familière :

\begin{lemma}
\label{lemma-smooth-morphism-power-series}
Soit $R \to S$ un morphisme d'anneaux dans
$\widehat{\mathcal{C}}_\Lambda$. Alors le morphisme induit
$\underline{S}|_{\mathcal{C}_\Lambda} \to \underline{R}|_{\mathcal{C}_\Lambda}$
est lisse si et seulement si $S$ est un anneau de séries formelles sur $R$.
\end{lemma}

\begin{proof}
Supposons que $S$ soit un anneau de séries formelles sur $R$, disons
$S = R[[x_1, \ldots, x_n]]$. La lissité de
$\underline{S}|_{\mathcal{C}_\Lambda} \to \underline{R}|_{\mathcal{C}_\Lambda}$
signifie ce qui suit (voir la remarque
\ref{remark-compare-smooth-schlessinger}). Étant donnés un morphisme
surjectif d'anneaux $B \to A$ dans $\mathcal{C}_\Lambda$, un morphisme
d'anneaux $R \to B$ et un morphisme d'anneaux $S \to A$ tels que le diagramme
en traits pleins
$$
\xymatrix{
S \ar[r] \ar@{..>}[rd] & A \\
R \ar[u] \ar[r] & B \ar[u]
}
$$
soit commutatif, il existe une flèche en pointillés qui rend le diagramme
commutatif. (Remarquer la similitude avec Algèbre, définition
\ref{algebra-definition-formally-smooth}.) Pour construire la flèche en
pointillés, choisissons des éléments $b_i \in B$ dont les images dans $A$
sont égales aux images de $x_i$ dans $A$. Remarquons que
$b_i \in \mathfrak m_B$, puisque $x_i$ s'envoie sur un élément de
$\mathfrak m_A$. Il existe donc un unique morphisme de $R$-algèbres
$R[[x_1, \ldots, x_n]] \to B$ qui envoie $x_i$ sur $b_i$ et qui peut servir
de flèche en pointillés.

\medskip\noindent
Réciproquement, supposons que
$\underline{S}|_{\mathcal{C}_\Lambda} \to \underline{R}|_{\mathcal{C}_\Lambda}$
soit lisse. Soient $x_1, \ldots, x_n \in S$ des éléments dont les images
forment une base de l'espace cotangent relatif
$\mathfrak m_S/(\mathfrak m_R S + \mathfrak m_S^2)$ de $S$ sur $R$.
Posons $T = R[[X_1, \ldots, X_n]]$. Remarquons que l'on a à la fois
$$
S/(\mathfrak m_R S + \mathfrak m_S^2) \cong
R/\mathfrak m_R[x_1, \ldots, x_n]/(x_ix_j)
$$
et
$$
T/(\mathfrak m_R T + \mathfrak m_T^2) \cong
R/\mathfrak m_R[X_1, \ldots, X_n]/(X_iX_j).
$$
Soit
$S/(\mathfrak m_R S + \mathfrak m_S^2) \to
T/(\mathfrak m_R T + \mathfrak m_T^2)$
l'isomorphisme de $R$-algèbres locales qui envoie la classe de $x_i$ sur la
classe de $X_i$. Soit
$f_1 : S \to T/(\mathfrak m_R T + \mathfrak m_T^2)$ la composée
$S \to S/(\mathfrak m_R S + \mathfrak m_S^2)
\to T/(\mathfrak m_R T + \mathfrak m_T^2)$.
L'hypothèse selon laquelle
$\underline{S}|_{\mathcal{C}_\Lambda} \to \underline{R}|_{\mathcal{C}_\Lambda}$
est lisse signifie que l'on peut relever $f_1$ en un morphisme
$f_2 : S \to T/\mathfrak{m}_T^2$, puis en un morphisme
$f_3 : S \to T/\mathfrak{m}_T^3$, et ainsi de suite pour tout $n \geq 1$.
On obtient donc un morphisme induit $f : S \to T = \lim T/\mathfrak m_T^n$
de $R$-algèbres locales. Par notre choix de $f_1$, le morphisme $f$ induit
un isomorphisme
$\mathfrak m_S/(\mathfrak m_R S + \mathfrak m_S^2) \to
\mathfrak m_T/(\mathfrak m_R T + \mathfrak m_T^2)$
d'espaces cotangents relatifs. Le lemme
\ref{lemma-surjective-cotangent-space} montre donc que $f$ est surjective
(où l'on considère $f$ comme un morphisme dans
$\widehat{\mathcal{C}}_R$). Choisissons des antécédents $y_i \in S$ des
$X_i \in T$ par $f$. Comme $T$ est un anneau de séries formelles sur $R$, il
existe un homomorphisme de $R$-algèbres locales $s : T \to S$ envoyant $X_i$
sur $y_i$. Par construction, $f \circ s = \text{id}$ ; donc $s$ est
injectif. Mais $s$ induit un isomorphisme sur les espaces cotangents relatifs,
puisque $f$ en induit un ; il est donc également surjectif, de nouveau par le
lemme \ref{lemma-surjective-cotangent-space}. Ainsi $s$ et $f$ sont des
isomorphismes.
\end{proof}

\noindent
Les morphismes lisses possèdent les propriétés fonctorielles suivantes.

\begin{lemma}
\label{lemma-smooth-properties}
Soient $\varphi : \mathcal{F} \to \mathcal{G}$ et $\psi : \mathcal{G}
\to \mathcal{H}$ des morphismes de catégories cofibrées en groupoïdes
au-dessus de $\mathcal{C}_\Lambda$.
\begin{enumerate}
\item Si $\varphi$ et $\psi$ sont lisses, alors $\psi \circ \varphi$ est
lisse.
\item Si $\varphi$ est essentiellement surjectif et si
$\psi \circ \varphi$ est lisse, alors $\psi$ est lisse.
\item Si $\mathcal{G}' \to \mathcal{G}$ est un morphisme de catégories
cofibrées en groupoïdes et si $\varphi$ est lisse, alors
$\mathcal{F} \times_\mathcal{G} \mathcal{G}' \to \mathcal{G}'$ est lisse.
\end{enumerate}
\end{lemma}

\begin{proof}
Les assertions (1) et (2) résultent immédiatement des définitions. Nous
omettons la démonstration de (3). Indications : utiliser la formulation de la
lissité donnée dans la remarque \ref{remark-smoothness-2-categorical} ainsi
que le fait que $\mathcal{F} \times_\mathcal{G} \mathcal{G}'$ est le
$2$-produit fibré ; voir les remarques \ref{remarks-cofibered-groupoids}
(\ref{item-fibre-product}).
\end{proof}

\begin{lemma}
\label{lemma-smooth-morphism-essentially-surjective}
Soit $\varphi : \mathcal{F} \to \mathcal{G}$ un morphisme lisse de
catégories cofibrées en groupoïdes au-dessus de $\mathcal{C}_\Lambda$.
Supposons que $\varphi : \mathcal{F}(k) \to \mathcal{G}(k)$ soit
essentiellement surjectif. Alors $\varphi : \mathcal{F} \to \mathcal{G}$ et
$\widehat{\varphi} : \widehat{\mathcal{F}} \to \widehat{\mathcal{G}}$
sont essentiellement surjectifs.
\end{lemma}

\begin{proof}
Soit $y$ un objet de $\mathcal{G}$ situé au-dessus de
$A \in \Ob(\mathcal{C}_\Lambda)$. Soit $y \to y_0$ une image directe de $y$
le long de $A \to k$. L'hypothèse de surjectivité essentielle de
$\varphi : \mathcal{F}(k) \to \mathcal{G}(k)$ donne un objet $x_0$ de
$\mathcal{F}$ situé au-dessus de $k$ et un isomorphisme
$y_0 \to \varphi(x_0)$. La lissité de $\varphi$ implique qu'il existe un
objet $x$ de $\mathcal{F}$ au-dessus de $A$ dont l'image $\varphi(x)$ est
isomorphe à $y$. Ainsi, $\varphi : \mathcal{F} \to \mathcal{G}$ est
essentiellement surjectif.

\medskip\noindent
Soit $\eta = (R, \eta_n, g_n)$ un objet de $\widehat{\mathcal{G}}$.
Construisons un objet $\xi$ de $\widehat{\mathcal{F}}$ muni d'un isomorphisme
$\eta \to \varphi(\xi)$. Par l'hypothèse de surjectivité essentielle de
$\varphi : \mathcal{F}(k) \to \mathcal{G}(k)$, il existe un morphisme
$\eta_1 \to \varphi(\xi_1)$ dans $\mathcal{G}(k)$ pour un certain
$\xi_1 \in \Ob(\mathcal{F}(k))$. Le morphisme
$\eta_2 \xrightarrow{g_1} \eta_1 \to \varphi(\xi_1)$
est situé au-dessus du morphisme surjectif d'anneaux
$R/\mathfrak m_R^2 \to k$ ; la lissité de $\varphi$ donne donc un objet
$\xi_2 \in \Ob(\mathcal{F}(R/\mathfrak m_R^2))$, un morphisme
$f_1: \xi_2 \to \xi_1$ situé au-dessus de $R/\mathfrak m_R^2 \to k$, et un
morphisme $\eta_2 \to \varphi(\xi_2)$ tels que
$$
\xymatrix{
\varphi(\xi_2)  \ar[r]^{\varphi(f_1)} &  \varphi(\xi_{1})   \\
\eta_2   \ar[u] \ar[r]^{g_1}  & \eta_1  \ar[u] \\
}
$$
soit commutatif. En poursuivant ainsi, on construit un objet
$\xi = (R, \xi_n, f_n)$ de $\widehat{\mathcal{F}}$ et un morphisme
$\eta \to \varphi(\xi) = (R, \varphi(\xi_n), \varphi(f_n))$ dans
$\widehat{\mathcal{G}}(R)$.
\end{proof}

\noindent
Nous chercherons plus loin à construire des morphismes lisses de foncteurs
pro-représentables vers des catégories de prédéformation $\mathcal{F}$.
D'après la discussion de la remarque \ref{remark-formal-objects-yoneda}, ces
morphismes correspondent à certains objets formels de $\mathcal{F}$ : ce
sont, plus précisément, les objets formels dits versels de $\mathcal{F}$.

\begin{definition}
\label{definition-versal}
Soit $\mathcal{F}$ une catégorie cofibrée en groupoïdes. Soit $\xi$ un objet
formel de $\mathcal{F}$ situé au-dessus de
$R \in \Ob(\widehat{\mathcal{C}}_\Lambda)$. On dit que $\xi$ est {\it versel}
si le morphisme correspondant
$\underline{\xi}: \underline{R}|_{\mathcal{C}_\Lambda} \to \mathcal{F}$
de la remarque \ref{remark-formal-objects-yoneda} est lisse.
\end{definition}

\begin{remark}
\label{remark-versal-object}
Soit $\mathcal{F}$ une catégorie cofibrée en groupoïdes au-dessus de $\mathcal
C_\Lambda$, et soit $\xi$ un objet formel de $\mathcal{F}$. Il
résulte de la définition de la lissité que la versalité de $\xi$ équivaut à la
condition suivante. Si
$$
\xymatrix{
& y \ar[d] \\
\xi \ar[r] & x
}
$$
est un diagramme dans $\widehat{\mathcal{F}}$ tel que $y \to x$ soit situé
au-dessus d'un morphisme surjectif $B \to A$ d'anneaux artiniens (on peut
supposer qu'il s'agit d'une petite extension), alors il existe un morphisme
$\xi \to y$ tel que
$$
\xymatrix{
& y \ar[d] \\
\xi \ar[r] \ar[ur] & x
}
$$
soit commutatif. En particulier, la condition pour $\xi$ d'être versel ne
dépend pas des choix d'images directes effectués lors de la construction de
$\underline{\xi} : \underline{R}|_{\mathcal{C}_\Lambda} \to \mathcal{F}$
dans la remarque \ref{remark-formal-objects-yoneda}.
\end{remark}

\begin{lemma}
\label{lemma-versal-object-quasi-initial}
Soit $\mathcal{F}$ une catégorie de prédéformation. Soit $\xi$ un objet
formel versel de $\mathcal{F}$. Pour tout objet formel $\eta$ de
$\widehat{\mathcal{F}}$, il existe un morphisme $\xi \to \eta$.
\end{lemma}

\begin{proof}
Par hypothèse, le morphisme
$\underline{\xi} : \underline{R}|_{\mathcal{C}_\Lambda} \to \mathcal{F}$
est lisse. Alors
$\iota(\xi) : \underline{R} \to \widehat{\mathcal{F}}$
est le complété de $\underline{\xi}$ ; voir la remarque
\ref{remark-formal-objects-yoneda}. D'après le lemme
\ref{lemma-smooth-morphism-essentially-surjective}, il existe un objet $f$ de
$\underline{R}$ tel que $\iota(\xi)(f) = \eta$. Ainsi, $f$ est un morphisme
d'anneaux $f : R \to S$ dans $\widehat{\mathcal{C}}_\Lambda$. L'égalité
$\iota(\xi)(f) = \eta$ signifie que $f_*\xi \cong \eta$, ce qui signifie
exactement qu'il existe un morphisme $\xi \to \eta$ situé au-dessus de $f$.
\end{proof}





\section{Catégories lisses ou non obstruées}
\label{section-smooth}

\noindent
Soit $p : \mathcal{F} \to \mathcal{C}_\Lambda$ une catégorie cofibrée en
groupoïdes. On peut considérer $\mathcal{C}_\Lambda$ comme une catégorie
cofibrée en groupoïdes au-dessus de $\mathcal{C}_\Lambda$ au moyen du foncteur
identité. De cette manière,
$p : \mathcal{F} \longrightarrow \mathcal{C}_\Lambda$
devient un morphisme de catégories cofibrées en groupoïdes au-dessus de
$\mathcal{C}_\Lambda$.

\begin{definition}
\label{definition-cofibered-groupoid-projection-smooth}
Soit $p : \mathcal{F} \to \mathcal{C}_\Lambda$ une catégorie cofibrée en
groupoïdes. On dit que $\mathcal{F}$ est {\it lisse} ou {\it non obstruée} si
son morphisme structural $p$ est lisse au sens de la définition
\ref{definition-smooth-morphism}.
\end{definition}

\noindent
C'est la notion « absolue » de lissité pour une catégorie cofibrée en
groupoïdes au-dessus de $\mathcal{C}_\Lambda$, bien qu'il soit plus correct de
dire que $\mathcal{F}$ est lisse sur $\Lambda$. Il faut manier avec prudence
l'expression « {\it $\mathcal{F}$ est non obstruée} » : il peut arriver que
$\mathcal{F}$ possède une théorie de l'obstruction dont les espaces
d'obstructions ne sont pas nuls, alors même que $\mathcal{F}$ est lisse.

\begin{remark}
\label{remark-smooth-on-top}
Supposons que $\mathcal{F}$ soit une catégorie de prédéformation admettant un
morphisme lisse $\varphi : \mathcal U \to \mathcal{F}$ depuis une catégorie
de prédéformation $\mathcal{U}$. D'après le lemme
\ref{lemma-smooth-morphism-essentially-surjective}, $\varphi$ est
essentiellement surjectif ; le lemme \ref{lemma-smooth-properties} montre donc
que $p: \mathcal{F} \to \mathcal{C}_\Lambda$ est lisse si et seulement si la
composée $\mathcal U \xrightarrow{\varphi} \mathcal{F} \xrightarrow{p}
\mathcal{C}_\Lambda$ est lisse, c'est-à-dire\ que $\mathcal{F}$ est lisse si
et seulement si $\mathcal{U}$ est lisse.
\end{remark}

\begin{lemma}
\label{lemma-smooth}
Soit $R \in \Ob(\widehat{\mathcal{C}}_\Lambda)$. Les conditions suivantes
sont équivalentes :
\begin{enumerate}
\item $\underline{R}|_{\mathcal{C}_\Lambda}$ est lisse ;
\item $\Lambda \to R$ est formellement lisse pour la topologie
$\mathfrak m_R$-adique ;
\item $\Lambda \to R$ est plat et $R \otimes_\Lambda k'$ est
géométriquement régulier sur $k'$ ;
\item $\Lambda \to R$ est plat et $k' \to R \otimes_\Lambda k'$ est
formellement lisse pour la topologie $\mathfrak m_R$-adique.
\end{enumerate}
Dans le cas classique, elles équivalent également à :
\begin{enumerate}
\item[(5)] $R$ est isomorphe à $\Lambda[[x_1, \ldots, x_n]]$ pour un certain
$n$.
\end{enumerate}
\end{lemma}

\begin{proof}
La lissité de
$p : \underline{R}|_{\mathcal{C}_\Lambda} \to \mathcal{C}_\Lambda$
signifie qu'étant donnés $B \to A$ surjectif dans $\mathcal{C}_\Lambda$ et
$R \to A$, on peut trouver la flèche en pointillés dans le diagramme
$$
\xymatrix{
R \ar[r] \ar@{-->}[rd] & A \\
\Lambda \ar[r] \ar[u] & B \ar[u]
}
$$
Cela est certainement vrai si $\Lambda \to R$ est formellement lisse pour la
topologie $\mathfrak m_R$-adique ; voir Compléments d'algèbre, définitions
\ref{more-algebra-definition-formally-smooth-adic} et
\ref{more-algebra-definition-formally-smooth}. Réciproquement, si cela est
vrai, Compléments d'algèbre, lemme \ref{more-algebra-lemma-fs-local}, montre
que $\Lambda \to R$ est formellement lisse pour la topologie
$\mathfrak m_R$-adique. Ainsi, (1) et (2) sont équivalentes.

\medskip\noindent
L'équivalence de (2), (3) et (4) est Compléments d'algèbre, proposition
\ref{more-algebra-proposition-fs-flat-fibre-fs}. L'équivalence avec (5)
résulte, par exemple, du lemme \ref{lemma-smooth-morphism-power-series} et du
fait que, dans le cas classique, $\mathcal{C}_\Lambda$ s'identifie à
$\underline{\Lambda}|_{\mathcal{C}_\Lambda}$.
\end{proof}

\begin{lemma}
\label{lemma-smooth-power-series-classical}
Soit $\mathcal{F}$ une catégorie de prédéformation. Soit $\xi$ un objet
formel versel de $\mathcal{F}$ situé au-dessus de
$R \in \Ob(\widehat{\mathcal{C}}_\Lambda)$. Les conditions suivantes sont
équivalentes :
\begin{enumerate}
\item $\mathcal{F}$ est non obstruée ;
\item $\Lambda \to R$ est formellement lisse pour la topologie
$\mathfrak m_R$-adique.
\end{enumerate}
Dans le cas classique, elles équivalent également à :
\begin{enumerate}
\item[(3)] $R \cong \Lambda[[x_1, \ldots, x_n]]$ pour un certain $n$.\
\end{enumerate}
\end{lemma}

\begin{proof}
Si (1) est vérifiée, c'est-à-dire si $\mathcal{F}$ est non obstruée, alors la
composée
$$
\underline{R}|_{\mathcal{C}_\Lambda}
\xrightarrow{\underline{\xi}}
\mathcal{F}
\to
\mathcal{C}_\Lambda
$$
est lisse ; voir le lemme \ref{lemma-smooth-properties}. Le lemme
\ref{lemma-smooth} donne donc (2). Réciproquement, si (2) est vérifiée, la
composée est lisse ; de plus, la première flèche est essentiellement
surjective d'après le lemme \ref{lemma-versal-object-quasi-initial}. Le lemme
\ref{lemma-smooth-properties} montre donc que la seconde flèche est lisse, ce
qui signifie par définition que $\mathcal{F}$ est non obstruée. Dans le cas
classique, l'équivalence avec (3) résulte du lemme \ref{lemma-smooth}.
\end{proof}

\begin{lemma}
\label{lemma-exists-smooth}
Il existe $R \in \Ob(\widehat{\mathcal{C}}_\Lambda)$ tel que les conditions
équivalentes du lemme \ref{lemma-smooth} soient satisfaites et que, de plus,
$H_1(L_{k/\Lambda}) = \mathfrak m_R/\mathfrak m_R^2$ et
$\Omega_{R/\Lambda} \otimes_R k = \Omega_{k/\Lambda}$.
\end{lemma}

\begin{proof}
Dans le cas classique, on choisit $R = \Lambda$. Plus généralement, si
l'extension de corps résiduels $k/k'$ est séparable, il existe une unique
extension finie étale $\Lambda^\wedge \to R$ (Algèbre, lemmes
\ref{algebra-lemma-complete-henselian} et
\ref{algebra-lemma-henselian-cat-finite-etale}) du complété
$\Lambda^\wedge$ de $\Lambda$ qui induit l'extension $k/k'$ sur les corps
résiduels.

\medskip\noindent
Dans le cas général, procédons comme suit. Choisissons une $\Lambda$-algèbre
lisse $P$ et une surjection de $\Lambda$-algèbres $P \to k$. (On peut, par
exemple, prendre pour $P$ une algèbre de polynômes.) Notons $\mathfrak m_P$ le
noyau de $P \to k$. La suite de Jacobi--Zariski, voir
(\ref{equation-sequence-extended}) et Algèbre, lemme
\ref{algebra-lemma-exact-sequence-NL}, est une suite exacte
$$
0 \to H_1(\NL_{k/\Lambda}) \to
\mathfrak m_P/\mathfrak m_P^2 \to
\Omega_{P/\Lambda} \otimes_P k \to
\Omega_{k/\Lambda} \to 0
$$
Le $0$ de gauche provient de ce que $P/k$ est lisse :
$\NL_{P/\Lambda}$ est alors quasi-isomorphe à un module projectif de type fini placé
en degré $0$, et donc $H_1(\NL_{P/\Lambda} \otimes_P k) = 0$. Supposons que
$f \in \mathfrak m_P$ s'envoie sur un élément non nul de
$\Omega_{P/\Lambda} \otimes_P k$. En posant $P' = P/(f)$, on dispose d'une
surjection de $\Lambda$-algèbres $P' \to k$. Remarquons que $P'$ est lisse en
$\mathfrak m_{P'}$ ; cela résulte de Compléments sur les morphismes, lemme
\ref{more-morphisms-lemma-slice-smooth-given-element}. Après avoir remplacé
$P$ par une localisation principale de $P'$, on voit donc que
$\dim(\mathfrak m_P/\mathfrak m_P^2)$ diminue. En répétant l'opération un
nombre fini de fois, on peut supposer que l'application
$\mathfrak m_P/\mathfrak m_P^2 \to
\Omega_{P/\Lambda} \otimes_P k$ est nulle, de sorte que la suite exacte se
décompose en les deux isomorphismes
$H_1(L_{k/\Lambda}) = \mathfrak m_P/\mathfrak m_P^2$ et
$\Omega_{P/\Lambda} \otimes_P k = \Omega_{k/\Lambda}$.

\medskip\noindent
Soit $R$ le complété $\mathfrak m_P$-adique de $P$. Alors $R$ est un objet de
$\widehat{\mathcal{C}}_\Lambda$ : c'est un anneau local noethérien complet
(voir Algèbre, lemme
\ref{algebra-lemma-completion-Noetherian-Noetherian}) dont le corps résiduel
s'identifie à $k$. Nous affirmons que $R$ convient.

\medskip\noindent
Remarquons d'abord que le morphisme $P \to R$ induit les isomorphismes
$\mathfrak m_P/\mathfrak m_P^2 = \mathfrak m_R/\mathfrak m_R^2$ et
$\Omega_{P/\Lambda} \otimes_P k = \Omega_{R/\Lambda} \otimes_R k$. En effet,
$\mathfrak m_P/\mathfrak m_P^2$ et
$\Omega_{P/\Lambda} \otimes_P k$ ne dépendent que de la $\Lambda$-algèbre
$P/\mathfrak m_P^2$ ; voir Algèbre, lemme
\ref{algebra-lemma-differential-mod-power-ideal}. Il en va de même pour $R$,
et l'on a $P/\mathfrak m_P^2 = R/\mathfrak m_R^2$. La fonctorialité de la
suite de Jacobi--Zariski (\ref{equation-sequence-functorial}) permet d'en
déduire $H_1(L_{k/\Lambda}) = \mathfrak m_R/\mathfrak m_R^2$ et
$\Omega_{R/\Lambda} \otimes_R k = \Omega_{k/\Lambda}$, puisque ces égalités
sont vraies pour $P$.

\medskip\noindent
Enfin, puisque $\Lambda \to P$ est lisse, Algèbre, proposition
\ref{algebra-proposition-smooth-formally-smooth}, montre que
$\Lambda \to P$ est formellement lisse. Alors $\Lambda \to P$ est
formellement lisse pour
la topologie $\mathfrak m_P$-adique d'après Compléments d'algèbre, lemme
\ref{more-algebra-lemma-formally-smooth}. Le complété $R$ hérite de cette
propriété par Compléments d'algèbre, lemme
\ref{more-algebra-lemma-formally-smooth-completion}, ce qui achève la
démonstration. En fait, on peut montrer que, chaque fois que
$\underline{R}|_{\mathcal{C}_\Lambda}$ est lisse, $R$ est isomorphe au
complété d'une algèbre lisse sur $\Lambda$, mais nous n'utiliserons pas ce
fait.
\end{proof}

\begin{example}
\label{example-smooth-explicit}
Voici un exemple plus explicite d'un $R$ comme dans le lemme
\ref{lemma-exists-smooth}. Soient $p$ un nombre premier et
$n \in \mathbf{N}$. Posons
$\Lambda = \mathbf{F}_p(t_1, t_2, \ldots, t_n)$ et
$k = \mathbf{F}_p(x_1, \ldots, x_n)$, l'application $\Lambda \to k$ étant
donnée par $t_i \mapsto x_i^p$. On peut alors prendre
$$
R = \Lambda[x_1, \ldots, x_n]^\wedge_{(x_1^p - t_1, \ldots, x_n^p - t_n)}
$$
On ne peut pas faire « mieux » dans cet exemple, c'est-à-dire que l'on ne
peut pas approcher $\mathcal{C}_\Lambda$ par un objet lisse plus petit de
$\widehat{\mathcal{C}}_\Lambda$ (on peut montrer que la dimension de $R$ doit
être au moins $n$, puisque l'application
$\Omega_{R/\Lambda} \otimes_R k \to \Omega_{k/\Lambda}$ est surjective). Nous
reviendrons plus loin sur ce phénomène de manière plus détaillée.
\end{example}







\section{Conditions de Schlessinger}
\label{section-schlessinger-conditions}

\noindent
Dans la suite, nous considérerons souvent des produits fibrés
$A_1 \times_A A_2$ d'anneaux de la catégorie $\mathcal{C}_\Lambda$. Nous
avons vu dans l'exemple \ref{example-fibre-product} qu'un tel produit fibré
n'est pas nécessairement un objet de $\mathcal{C}_\Lambda$. Toutefois, dans
presque tous les cas ci-dessous, l'un des deux morphismes $A_i \to A$ est
surjectif, et $A_1 \times_A A_2$ est alors un objet de
$\mathcal{C}_\Lambda$ d'après le lemme \ref{lemma-fiber-product-CLambda}.
Nous utiliserons ce résultat sans autre mention.

\medskip\noindent
On note $k[\epsilon]$ l'algèbre des nombres duaux sur $k$. Plus généralement,
pour un $k$-espace vectoriel $V$, on note $k[V]$ la $k$-algèbre dont l'espace
vectoriel sous-jacent est $k \oplus V$ et dont la multiplication est donnée
par $(a, v) \cdot (a', v') = (aa', av' + a'v)$. Lorsque $V = k$, $k[V]$ est
l'algèbre des nombres duaux sur $k$. Pour tout $k$-espace vectoriel $V$ de
dimension finie, l'anneau $k[V]$ appartient à $\mathcal{C}_\Lambda$.

\begin{definition}
\label{definition-S1-S2}
Soit $\mathcal{F}$ une catégorie cofibrée en groupoïdes au-dessus de $\mathcal
C_\Lambda$. On définit les {\it conditions (S1) et (S2)} pour
$\mathcal{F}$ comme suit :
\begin{enumerate}
\item[(S1)] Tout diagramme de $\mathcal{F}$
$$
\vcenter{
\xymatrix{
           & x_2 \ar[d] \\
x_1 \ar[r] & x
}
}
\quad\text{au-dessus de}\quad
\vcenter{
\xymatrix{
           & A_2 \ar[d] \\
A_1 \ar[r] & A
}
}
$$
dans $\mathcal{C}_\Lambda$, avec $A_2 \to A$ surjectif, peut être complété en
un diagramme commutatif
$$
\vcenter{
\xymatrix{
y \ar[r] \ar[d] & x_2 \ar[d] \\
x_1 \ar[r]      & x
}
}
\quad\text{au-dessus de}\quad
\vcenter{
\xymatrix{
A_1 \times_A A_2 \ar[r] \ar[d] & A_2 \ar[d] \\
A_1 \ar[r]      & A.
}
}
$$
\item[(S2)]
La condition (S1) est satisfaite pour les diagrammes de $\mathcal{F}$ situés
au-dessus d'un diagramme de $\mathcal{C}_\Lambda$ de la forme
$$
\xymatrix{
          & k[\epsilon] \ar[d] \\
A  \ar[r] & k.
}
$$
De plus, si l'on a deux diagrammes commutatifs dans $\mathcal{F}$
$$
\vcenter{
\xymatrix{
y \ar[r]_c \ar[d]_a & x_\epsilon \ar[d]^e \\
x \ar[r]^d          & x_0
}
}
\quad\text{et}\quad
\vcenter{
\xymatrix{
y' \ar[r]_{c'} \ar[d]_{a'} & x_\epsilon \ar[d]^e \\
x \ar[r]^d                 & x_0
}
}
\quad\text{au-dessus de}\quad
\vcenter{
\xymatrix{
A \times_k k[\epsilon] \ar[r] \ar[d] & k[\epsilon] \ar[d] \\
A  \ar[r] & k
}
}
$$
alors il existe un morphisme $b : y \to y'$ dans
$\mathcal{F}(A \times_k k[\epsilon])$ tel que $a = a' \circ b$.
\end{enumerate}
\end{definition}

\noindent
On peut expliquer en partie le sens des conditions (S1) et (S2) à l'aide des
catégories fibres. Supposons que $f_1 : A_1 \to A$ et $f_2 : A_2 \to A$
soient des morphismes d'anneaux dans $\mathcal{C}_\Lambda$, avec $f_2$
surjectif. Notons $p_i : A_1 \times_A A_2 \to A_i$ les projections.
Supposons fixé un choix d'images directes pour $\mathcal{F}$. Le diagramme
commutatif d'anneaux se traduit alors par un diagramme $2$-commutatif
$$
\xymatrix{
\mathcal{F}(A_1 \times_A A_2) \ar[r]_-{p_{2, *}} \ar[d]_{p_{1, *}} &
\mathcal{F}(A_2) \ar[d]^{f_{2, *}} \\
\mathcal{F}(A_1) \ar[r]^{f_{1, *}} & \mathcal{F}(A)
}
$$
de catégories fibres, d'où un foncteur
\begin{equation}
\label{equation-compare}
\mathcal{F}(A_1 \times_A A_2) \to
\mathcal{F}(A_1) \times_{\mathcal{F}(A)} \mathcal{F}(A_2)
\end{equation}
à valeurs dans le $2$-produit fibré des catégories. La condition (S1) exige
que ce foncteur soit essentiellement surjectif. La première partie de la
condition (S2) exige que ce foncteur soit essentiellement surjectif lorsque
$f_2$ est le morphisme $k[\epsilon] \to k$. En outre, dans ce cas, la seconde
partie de (S2) implique que deux objets qui deviennent isomorphes dans la
cible sont isomorphes dans la source (mais elle n'est {\it pas} équivalente à
cette assertion). L'avantage de formuler les conditions comme dans la
définition est qu'aucun choix n'est nécessaire.

\begin{lemma}
\label{lemma-S1-small-extensions}
Soit $\mathcal{F}$ une catégorie cofibrée en groupoïdes au-dessus de $\mathcal
C_\Lambda$. Alors $\mathcal{F}$ satisfait (S1) dès que la condition
(S1) est supposée satisfaite seulement lorsque $A_2 \to A$ est une petite
extension.
\end{lemma}

\begin{proof}
Démonstration omise. Indications : appliquer le lemme
\ref{lemma-factor-small-extension} et procéder par récurrence comme dans la
démonstration du lemme \ref{lemma-smoothness-small-extensions}.
\end{proof}

\begin{remark}
\label{remark-compare-S1-S2-schlessinger}
Lorsque $\mathcal{F}$ est cofibrée en ensembles, les conditions (S1) et (S2)
sont exactement les conditions (H1) et (H2) de l'article de Schlessinger
\cite{Sch}. En effet, pour un foncteur $F: \mathcal{C}_\Lambda \to
\textit{Ensembles}$, les conditions (S1) et (S2) s'énoncent ainsi :
\begin{enumerate}
\item [(S1)] Si $A_1 \to A$ et $A_2 \to A$ sont des morphismes dans
$\mathcal{C}_\Lambda$, avec $A_2 \to A$ surjectif, alors l'application
induite $F(A_1 \times_A A_2) \to F(A_1) \times_{F(A)} F(A_2)$ est surjective.
\item [(S2)] Si $A \to k$ est un morphisme dans $\mathcal{C}_\Lambda$,
alors l'application induite
$F(A \times_k k[\epsilon]) \to F(A) \times_{F(k)} F(k[\epsilon])$
est bijective.
\end{enumerate}
L'injectivité de l'application
$F(A \times_k k[\epsilon]) \to F(A) \times_{F(k)} F(k[\epsilon])$
résulte de la seconde partie de la condition (S2) et du fait que les
morphismes sont des identités.
\end{remark}

\begin{lemma}
\label{lemma-S2-extensions}
Soit $\mathcal{F}$ une catégorie cofibrée en groupoïdes au-dessus de
$\mathcal{C}_\Lambda$. Si $\mathcal{F}$ satisfait (S2), alors la condition
(S2) reste valable lorsque $k[\epsilon]$ est remplacé par $k[V]$, pour tout
$k$-espace vectoriel $V$ de dimension finie.
\end{lemma}

\begin{proof}
Lorsque $\mathcal{F}$ est cofibrée en ensembles, c'est-à-dire correspond à un
foncteur $F : \mathcal{C}_\Lambda \to \textit{Ensembles}$, cela résulte de la
description de (S2) pour $F$ donnée dans la remarque
\ref{remark-compare-S1-S2-schlessinger} et du fait que
$k[V] \cong k[\epsilon] \times_k \ldots \times_k k[\epsilon]$
avec $\dim_k V$ facteurs. C'est le cas des foncteurs que nous utiliserons dans
la suite du chapitre.

\medskip\noindent
Démontrons le cas général par récurrence sur $\dim(V)$. Si $\dim(V) = 1$,
alors $k[V] \cong k[\epsilon]$ et le résultat est vrai par hypothèse. Si
$\dim(V) > 1$, écrivons $V = V' \oplus k\epsilon$. Choisissons un diagramme
$$
\vcenter{
\xymatrix{
& x_V \ar[d] \\
x \ar[r] & x_0
}
}
\quad\text{au-dessus de}\quad
\vcenter{
\xymatrix{
& k[V] \ar[d] \\
A \ar[r] & k
}
}
$$
Choisissons un morphisme $x_V \to x_{V'}$ situé au-dessus de
$k[V] \to k[V']$ et un morphisme $x_V \to x_\epsilon$ situé au-dessus de
$k[V] \to k[\epsilon]$. Remarquons que le morphisme $x_V \to x_0$ se
factorise à la fois en $x_V \to x_{V'} \to x_0$ et en
$x_V \to x_\epsilon \to x_0$. Par l'hypothèse de récurrence, on peut trouver
un diagramme
$$
\vcenter{
\xymatrix{
y' \ar[d] \ar[r] & x_{V'} \ar[d] \\
x \ar[r] & x_0
}
}
\quad\text{au-dessus de}\quad
\vcenter{
\xymatrix{
A \times_k k[V'] \ar[d] \ar[r] & k[V'] \ar[d] \\
A \ar[r] & k
}
}
$$
Cela nous donne un diagramme commutatif
$$
\vcenter{
\xymatrix{
& x_\epsilon \ar[d] \\
y' \ar[r] & x_0
}
}
\quad\text{au-dessus de}\quad
\vcenter{
\xymatrix{
& k[\epsilon] \ar[d] \\
A \times_k k[V'] \ar[r] & k
}
}
$$
La condition (S2) fournit donc un diagramme commutatif
$$
\vcenter{
\xymatrix{
y \ar[d] \ar[r] & x_\epsilon \ar[d] \\
y' \ar[r] & x_0
}
}
\quad\text{au-dessus de}\quad
\vcenter{
\xymatrix{
(A \times_k k[V']) \times_k k[\epsilon] \ar[d] \ar[r] & k[\epsilon] \ar[d] \\
A \times_k k[V'] \ar[r] & k
}
}
$$
Remarquons que
$(A \times_k k[V']) \times_k k[\epsilon] = A \times_k k[V' \oplus k\epsilon]
= A \times_k k[V]$. Nous affirmons que $y$ s'insère dans le diagramme
commutatif voulu. Pour le voir, soit $y \to y_V$ un morphisme situé au-dessus
de $A \times_k k[V] \to k[V]$. Les morphismes
$y \to y' \to x_{V'}$ et $y \to x_\epsilon$ se factorisent par le morphisme
$y \to y_V$ (d'après les axiomes des catégories cofibrées en groupoïdes).
Ainsi, $y_V$ et $x_V$ s'insèrent tous deux dans des diagrammes commutatifs
$$
\vcenter{
\xymatrix{
y_V \ar[r] \ar[d] & x_\epsilon \ar[d] \\
x_{V'} \ar[r]          & x_0
}
}
\quad\text{et}\quad
\vcenter{
\xymatrix{
x_V \ar[r] \ar[d] & x_\epsilon \ar[d] \\
x_{V'} \ar[r]                 & x_0
}
}
$$
La seconde partie de (S2) donne donc un isomorphisme $y_V \to x_V$
compatible avec $y_V \to x_{V'}$ et $x_V \to x_{V'}$, et en particulier
avec les morphismes vers $x_0$. La composée $y \to y_V \to x_V$ s'insère
alors dans le diagramme commutatif requis
$$
\vcenter{
\xymatrix{
y \ar[r] \ar[d] & x_V \ar[d] \\
x \ar[r] & x_0
}
}
\quad\text{au-dessus de}\quad
\vcenter{
\xymatrix{
A \times_k k[V] \ar[d] \ar[r] & k[V] \ar[d] \\
A \ar[r] & k
}
}
$$
On voit ainsi que la première partie de $(S2)$ reste vraie lorsque
$k[\epsilon]$ est remplacé par $k[V]$.

\medskip\noindent
Pour démontrer la seconde partie, supposons donnés deux diagrammes
commutatifs
$$
\vcenter{
\xymatrix{
y \ar[r] \ar[d] & x_V \ar[d] \\
x \ar[r] & x_0
}
}
\quad\text{et}\quad
\vcenter{
\xymatrix{
y' \ar[r] \ar[d] & x_V \ar[d] \\
x \ar[r] & x_0
}
}
\quad\text{au-dessus de}\quad
\vcenter{
\xymatrix{
A \times_k k[V] \ar[d] \ar[r] & k[V] \ar[d] \\
A \ar[r] & k
}
}
$$
Nous utiliserons les morphismes $x_V \to x_{V'} \to x_0$ et
$x_V \to x_\epsilon \to x_0$ introduits au premier paragraphe de la
démonstration. Choisissons des morphismes $y \to y_{V'}$ et
$y' \to y'_{V'}$ situés au-dessus de
$A \times_k k[V] \to A \times_k k[V']$. Les axiomes d'une catégorie cofibrée
impliquent que l'on peut trouver des diagrammes commutatifs
$$
\vcenter{
\xymatrix{
y_{V'} \ar[r] \ar[d] & x_{V'} \ar[d] \\
x \ar[r] & x_0
}
}
\quad\text{et}\quad
\vcenter{
\xymatrix{
y'_{V'} \ar[r] \ar[d] & x_{V'} \ar[d] \\
x \ar[r] & x_0
}
}
\quad\text{au-dessus de}\quad
\vcenter{
\xymatrix{
A \times_k k[V'] \ar[d] \ar[r] & k[V'] \ar[d] \\
A \ar[r] & k
}
}
$$
Par l'hypothèse de récurrence, on obtient un isomorphisme
$b : y_{V'} \to y'_{V'}$ compatible avec les morphismes $y_{V'} \to x$ et
$y'_{V'} \to x$, donc en particulier avec les morphismes vers $x_0$. On a
alors des diagrammes commutatifs
$$
\vcenter{
\xymatrix{
y \ar[r] \ar[d] & x_\epsilon \ar[d] \\
y'_{V'} \ar[r] & x_0
}
}
\quad\text{et}\quad
\vcenter{
\xymatrix{
y' \ar[r] \ar[d] & x_\epsilon \ar[d] \\
y'_{V'} \ar[r] & x_0
}
}
\quad\text{au-dessus de}\quad
\vcenter{
\xymatrix{
A \times_k k[\epsilon] \ar[d] \ar[r] & k[\epsilon] \ar[d] \\
A \ar[r] & k
}
}
$$
où le morphisme $y \to y'_{V'}$ est la composée
$y \to y_{V'} \xrightarrow{b} y'_{V'}$, et où les morphismes
$y \to x_\epsilon$ et $y' \to x_\epsilon$ sont les composées des morphismes
$y \to x_V$ et $y' \to x_V$ avec le morphisme $x_V \to x_\epsilon$. La
seconde partie de (S2) assure alors l'existence d'un isomorphisme $y \to y'$
compatible avec les morphismes vers $y'_{V'}$, donc en particulier avec
les morphismes vers $x$ (puisque $b$ était compatible avec les morphismes
vers $x$).
\end{proof}

\begin{lemma}
\label{lemma-S1-S2-associated-functor}
Soit $\mathcal{F}$ une catégorie cofibrée en groupoïdes au-dessus de
$\mathcal{C}_\Lambda$.
\begin{enumerate}
\item Si $\mathcal{F}$ satisfait (S1), alors $\overline{\mathcal{F}}$ la
satisfait aussi.
\item Si $\mathcal{F}$ satisfait (S2), alors $\overline{\mathcal{F}}$ la
satisfait aussi, pourvu que l'une au moins des conditions suivantes soit
vérifiée :
\begin{enumerate}
\item $\mathcal{F}$ est une catégorie de prédéformation ;
\item la catégorie $\mathcal{F}(k)$ est un ensemble ou un ensoïde ;
\item pour tout morphisme $x_\epsilon \to x_0$ de $\mathcal{F}$ situé
au-dessus de $k[\epsilon] \to k$, l'application d'image directe
$\text{Aut}_{k[\epsilon]}(x_\epsilon) \to \text{Aut}_k(x_0)$
est surjective.
\end{enumerate}
\end{enumerate}
\end{lemma}

\begin{proof}
Supposons que $\mathcal{F}$ satisfasse (S1). Soient des morphismes d'anneaux
$f_i : A_i \to A$ dans $\mathcal{C}_\Lambda$, avec $f_2$ surjectif, et des
$x_i \in \mathcal{F}(A_i)$ tels que les images directes
$f_{1, *}(x_1)$ et $f_{2, *}(x_2)$ soient isomorphes. On peut alors choisir un
objet $x$ de $\mathcal{F}$ au-dessus de $A$ isomorphe à chacune de ces images,
et l'on obtient un diagramme comme dans (S1). On trouve donc un objet $y$ de
$\mathcal{F}$ au-dessus de $A_1 \times_A A_2$ dont l'image directe vers
$A_1$, respectivement $A_2$, est isomorphe à $x_1$, respectivement $x_2$.
Ainsi, (S1) est vérifiée pour $\overline{\mathcal{F}}$.

\medskip\noindent
Supposons que $\mathcal{F}$ satisfasse (S2). La première partie de (S2) pour
$\overline{\mathcal{F}}$ résulte de l'argument précédent. La seconde partie
de (S2) pour $\overline{\mathcal{F}}$ signifie que l'application
$$
\overline{\mathcal{F}}(A \times_k k[\epsilon]) \to
\overline{\mathcal{F}}(A)
\times_{\overline{\mathcal{F}}(k)} \overline{\mathcal{F}}(k[\epsilon])
$$
est injective pour tout anneau $A$ de $\mathcal{C}_\Lambda$. Supposons que
$y, y' \in \mathcal{F}(A \times_k k[\epsilon])$. En utilisant les axiomes
des catégories cofibrées, on peut choisir des diagrammes commutatifs
$$
\vcenter{
\xymatrix{
y \ar[r]_c \ar[d]_a & x_\epsilon \ar[d]^e \\
x \ar[r]^d          & x_0
}
}
\quad\text{et}\quad
\vcenter{
\xymatrix{
y' \ar[r]_{c'} \ar[d]_{a'} & x'_\epsilon \ar[d]^{e'} \\
x' \ar[r]^{d'}                 & x'_0
}
}
\quad\text{au-dessus de}\quad
\vcenter{
\xymatrix{
A \times_k k[\epsilon] \ar[d] \ar[r] & k[\epsilon] \ar[d] \\
A \ar[r] & k
}
}
$$
Supposons qu'il existe des isomorphismes
$\alpha : x \to x'$ dans $\mathcal{F}(A)$ et
$\beta : x_\epsilon \to x'_\epsilon$ dans $\mathcal{F}(k[\epsilon])$. Cela
signifie également qu'il existe un isomorphisme $\gamma : x_0 \to x'_0$
compatible avec $\alpha$. Pour démontrer (S2) pour
$\overline{\mathcal{F}}$, il faut montrer qu'il existe un isomorphisme
$y \to y'$ dans $\mathcal{F}(A \times_k k[\epsilon])$. D'après (S2) pour
$\mathcal{F}$, un tel morphisme existera si l'on peut choisir les
isomorphismes $\alpha$, $\beta$ et $\gamma$ de telle sorte que
$$
\xymatrix{
x \ar[d]^\alpha \ar[r] & x_0 \ar[d]^\gamma &
x_\epsilon \ar[d]^\beta \ar[l]^e \\
x' \ar[r] & x'_0 & x'_\epsilon \ar[l]_{e'}
}
$$
soit commutatif (on peut alors remplacer $x$ par $x'$ et $x_\epsilon$ par
$x'_\epsilon$ dans le diagramme affiché précédent). Le carré de gauche est
commutatif par notre choix de $\gamma$. On peut factoriser $e' \circ \beta$
en $\gamma' \circ e$ pour un second morphisme
$\gamma' : x_0 \to x'_0$. Il reste à savoir si l'on peut faire en sorte que
$\gamma = \gamma'$. C'est clair si $\mathcal{F}(k)$ est un ensemble ou un
ensoïde. De plus, si
$\text{Aut}_{k[\epsilon]}(x_\epsilon) \to \text{Aut}_k(x_0)$
est surjective, on peut modifier le choix de $\beta$ en le précomposant avec
un automorphisme de $x_\epsilon$ dont l'image est
$\gamma^{-1} \circ \gamma'$, ce qui donne le résultat.
\end{proof}

\begin{lemma}
\label{lemma-S1-S2-localize}
Soit $\mathcal{F}$ une catégorie cofibrée en groupoïdes au-dessus de
$\mathcal{C}_\Lambda$. Soit $x_0 \in \Ob(\mathcal{F}(k))$. Soit
$\mathcal{F}_{x_0}$ la catégorie cofibrée en groupoïdes au-dessus de
$\mathcal{C}_\Lambda$ construite dans la remarque
\ref{remark-localize-cofibered-groupoid}.
\begin{enumerate}
\item Si $\mathcal{F}$ satisfait (S1), alors $\mathcal{F}_{x_0}$ la satisfait
aussi.
\item Si $\mathcal{F}$ satisfait (S2), alors $\mathcal{F}_{x_0}$ la satisfait
aussi.
\end{enumerate}
\end{lemma}

\begin{proof}
Tout diagramme de $\mathcal{F}_{x_0}$ comme dans la définition
\ref{definition-S1-S2} donne un diagramme dans $\mathcal{F}$, et le résultat
fourni par la condition (S1) ou (S2) pour ce diagramme dans $\mathcal{F}$ peut
également être considéré comme un résultat dans $\mathcal{F}_{x_0}$.
\end{proof}

\begin{lemma}
\label{lemma-lifting-section}
Soit $p: \mathcal{F} \to \mathcal{C}_\Lambda$ une catégorie cofibrée en
groupoïdes. Considérons un diagramme de $\mathcal{F}$
$$
\vcenter{
\xymatrix{
y \ar[r] \ar[d]_a & x_\epsilon \ar[d]_e \\
x \ar[r]^d        & x_0
}
}
\quad\text{au-dessus de}\quad
\vcenter{
\xymatrix{
A \times_k k[\epsilon] \ar[r] \ar[d] & k[\epsilon] \ar[d] \\
A \ar[r] & k.
}
}
$$
dans $\mathcal{C}_\Lambda$. Supposons que $\mathcal{F}$ satisfasse (S2).
Alors il existe un morphisme $s : x \to y$ tel que
$a \circ s = \text{id}_x$ si et seulement s'il existe un morphisme
$s_\epsilon : x \to x_\epsilon$ tel que $e \circ s_\epsilon = d$.
\end{lemma}

\begin{proof}
Le sens direct est clair. Réciproquement, supposons qu'il existe un morphisme
$s_\epsilon : x \to x_\epsilon$ tel que $e \circ s_\epsilon = d$.
Remarquons que $p(s_\epsilon) : A \to k[\epsilon]$ est un morphisme d'anneaux
compatible avec le morphisme $A \to k$. On obtient donc
$$
\sigma = (\text{id}_A, p(s_\epsilon)) : A \to A \times_k k[\epsilon].
$$
Choisissons une image directe $x \to \sigma_*x$. Par construction, on peut
factoriser $s_\epsilon$ en $x \to \sigma_*x \to x_\epsilon$. De plus, comme
$\sigma$ est une section de $A \times_k k[\epsilon] \to A$, on obtient un
morphisme $\sigma_*x \to x$ tel que $x \to \sigma_*x \to x$ soit
$\text{id}_x$. L'égalité $e \circ s_\epsilon = d$ montre que le diagramme
$$
\xymatrix{
\sigma_*x \ar[r] \ar[d] & x_\epsilon \ar[d]_e \\
x \ar[r]^d        & x_0
}
$$
est commutatif. La condition (S2) fournit donc un morphisme
$\sigma_*x \to y$ tel que $\sigma_*x \to y \to x$ soit le morphisme donné
$\sigma_*x \to x$. Il suffit maintenant de prendre pour solution la composée
$a : x \to y$ égale à $x \to \sigma_*x \to y$.
\end{proof}

\begin{lemma}
\label{lemma-lifting-along-small-extension}
Considérons un diagramme commutatif dans une catégorie de prédéformation
$\mathcal{F}$
$$
\vcenter{
\xymatrix{
y \ar[r] \ar[d] & x_2 \ar[d]^{a_2} \\
x_1 \ar[r]^{a_1}        & x
}
}
\quad\text{au-dessus de}
\vcenter{
\xymatrix{
A_1 \times_A A_2 \ar[r] \ar[d] & A_2 \ar[d]^{f_2} \\
A_1 \ar[r]^{f_1} & A
}
}
$$
dans $\mathcal{C}_\Lambda$, où $f_2 : A_2 \to A$ est une petite extension.
Supposons qu'il existe un morphisme $h : A_1 \to A_2$ tel que
$f_2 = f_1 \circ h$. Soit $I = \Ker(f_2)$. Considérons le morphisme d'anneaux
$$
g : A_1 \times_A A_2 \longrightarrow k[I] = k \oplus I, \quad
(u, v) \longmapsto \overline{u} \oplus (v - h(u))
$$
Choisissons une image directe $y \to g_*y$. Supposons que $\mathcal{F}$
satisfasse (S2). S'il existe un morphisme $x_1 \to g_*y$, alors il existe un
morphisme $b: x_1 \to x_2$ tel que $a_1 =  a_2 \circ b$.
\end{lemma}

\begin{proof}
Remarquons que
$\text{id}_{A_1} \times g : A_1 \times_A A_2 \to A_1 \times_k k[I]$
est un isomorphisme et que $k[I] \cong k[\epsilon]$. On a donc un diagramme
$$
\vcenter{
\xymatrix{
y \ar[r] \ar[d] & g_*y \ar[d] \\
x_1 \ar[r]        & x_0
}
}
\quad\text{au-dessus de}\quad
\vcenter{
\xymatrix{
A_1 \times_k k[\epsilon] \ar[r] \ar[d] & k[\epsilon] \ar[d] \\
A_1 \ar[r] & k.
}
}
$$
où $x_0$ est un objet de $\mathcal{F}$ situé au-dessus de $k$ (tout objet de
$\mathcal{F}$ possède un unique morphisme vers $x_0$ ; voir la discussion qui
suit la définition \ref{definition-predeformation-category}). Si l'on dispose
d'un morphisme $x_1 \to g_*y$, le lemme \ref{lemma-lifting-section} fournit
une section $s : x_1 \to y$ du morphisme $y \to x_1$. En la composant
avec le morphisme $y \to x_2$, on obtient $b : x_1 \to x_2$, qui vérifie
$a_1 =  a_2 \circ b$ puisque le diagramme du lemme est commutatif et que $s$
est une section.
\end{proof}





\section{Espaces tangents des foncteurs}
\label{section-tangent-spaces-functors}

\noindent
Soit $R$ un anneau. On note $\text{Mod}_R$ la catégorie des $R$-modules et
$\text{Mod}^{fg}_R$ la catégorie des $R$-modules de type fini.

\begin{definition}
\label{definition-linear}
Soit $L: \text{Mod}^{fg}_R \to \text{Mod}_R$, respectivement
$L: \text{Mod}_R \to \text{Mod}_R$, un foncteur. On dit que $L$ est
{\it $R$-linéaire} si, pour toute paire d'objets $M, N$ de
$\text{Mod}^{fg}_R$, respectivement de $\text{Mod}_R$, l'application
$$
L : \Hom_R(M, N) \longrightarrow \Hom_R(L(M), L(N))
$$
est un morphisme de $R$-modules.
\end{definition}

\begin{remark}
\label{remark-linear-enriched-over-modules}
On peut définir la notion de $R$-linéarité pour tout foncteur entre catégories
enrichies sur $\text{Mod}_R$. Nous avons formulé la définition spécifiquement
pour les foncteurs $L: \text{Mod}^{fg}_R \to \text{Mod}_R$ et
$L: \text{Mod}_R \to \text{Mod}_R$, car ce sont les cas dont nous avons eu
besoin jusqu'ici.
\end{remark}

\begin{remark}
\label{remark-linear-functor}
Si $L: \text{Mod}^{fg}_R \to \text{Mod}_R$ est un foncteur $R$-linéaire,
alors $L$ préserve les produits finis et envoie le module nul sur le module
nul ; voir Homologie, lemme \ref{homology-lemma-additive-additive}.
Réciproquement, si un foncteur $\text{Mod}^{fg}_R \to \textit{Ensembles}$ préserve
les produits finis et envoie le module nul sur un ensemble à un seul élément, il admet un
unique relèvement en un foncteur $R$-linéaire ; voir le lemme
\ref{lemma-linear-functor}.
\end{remark}

\begin{lemma}
\label{lemma-linear-functor}
Soit $L: \text{Mod}^{fg}_R \to \textit{Ensembles}$, respectivement
$L: \text{Mod}_R \to \textit{Ensembles}$, un foncteur. Supposons que $L(0)$ soit
un ensemble à un seul élément et que $L$ préserve les produits finis. Il existe alors un
unique foncteur $R$-linéaire
$\widetilde{L} : \text{Mod}^{fg}_R \to \text{Mod}_R$,
respectivement $\widetilde{L} : \text{Mod}_R \to \text{Mod}_R$, tel que
$$
\vcenter{
\xymatrix{
& \text{Mod}_R \ar[dr]^{\text{forget}} &   \\
\text{Mod}^{fg}_R  \ar[ur]^{\widetilde{L}} \ar[rr]^{L} &  &
\textit{Ensembles}
}
}
\quad\text{resp.}\quad
\vcenter{
\xymatrix{
& \text{Mod}_R \ar[dr]^{\text{forget}} &   \\
\text{Mod}_R  \ar[ur]^{\widetilde{L}} \ar[rr]^{L} &  &
\textit{Ensembles}
}
}
$$
soit commutatif.
\end{lemma}

\begin{proof}
Nous ne le démontrons que dans le cas
$L: \text{Mod}^{fg}_R \to \textit{Ensembles}$. Soit $M$ un $R$-module de type
fini. On définit $\widetilde{L}(M)$ comme l'ensemble $L(M)$ muni de la
structure de $R$-module suivante.

\medskip\noindent
Multiplication : si $r \in R$, la multiplication par $r$ sur $L(M)$ est, par
définition, l'application $L(M) \to L(M)$ induite par l'homothétie
$r \cdot : M \to M$.

\medskip\noindent
Addition : l'application somme $M \times M \to M: (m_1, m_2) \mapsto m_1 + m_2$
induit une application
$L(M \times M) \to L(M)$. Par hypothèse, $L(M \times M)$ est canoniquement
isomorphe à $L(M) \times L(M)$. L'addition sur $L(M)$ est définie par
l'application $L(M) \times L(M) \cong L(M \times M) \to L(M)$.

\medskip\noindent
Zéro : il existe un unique morphisme $0 \to M$. L'élément nul de $L(M)$
est l'image de $L(0) \to L(M)$.

\medskip\noindent
Nous omettons de vérifier que cela définit un $R$-module
$\widetilde{L}(M)$, et que cette structure est l'unique pour laquelle la
construction soit $R$-linéaire et fonctorielle en $M$.
\end{proof}

\begin{lemma}
\label{lemma-morphism-linear-functors}
Soient $L_1, L_2: \text{Mod}^{fg}_R \to \textit{Ensembles}$ des foncteurs qui
envoient $0$ sur un ensemble à un seul élément et préservent les produits finis. Soit
$t : L_1 \to L_2$ un morphisme de foncteurs. Alors $t$ induit un morphisme
$\widetilde{t} : \widetilde{L}_1 \to \widetilde{L}_2$ entre les foncteurs
fournis par le lemme \ref{lemma-linear-functor}, donné simplement par
$\widetilde{t}_M = t_M: \widetilde{L}_1(M) \to \widetilde{L}_2(M)$ pour tout
$M \in \Ob(\text{Mod}^{fg}_R)$. Autrement dit,
$t_M: \widetilde{L}_1(M) \to \widetilde{L}_2(M)$ est un morphisme de
$R$-modules.
\end{lemma}

\begin{proof}
Omis.
\end{proof}

\noindent
Dans le cas où $R = K$ est un corps, un foncteur $K$-linéaire
$L : \text{Mod}^{fg}_K \to \text{Mod}_K$ est déterminé par sa valeur $L(K)$.

\begin{lemma}
\label{lemma-linear-functor-over-field}
Soit $K$ un corps. Soit $L: \text{Mod}^{fg}_K \to
\text{Mod}_K$ un foncteur
$K$-linéaire. Alors $L$ est isomorphe au foncteur
$L(K) \otimes_K - : \text{Mod}^{fg}_K \to
\text{Mod}_K$.
\end{lemma}

\begin{proof}
Pour $V \in \Ob(\text{Mod}^{fg}_K)$, l'isomorphisme
$L(K) \otimes_K V \to L(V)$ est donné sur les tenseurs purs par
$x \otimes v \mapsto L(f_v)(x)$, où $f_v: K \to V$ est l'application
$K$-linéaire envoyant $1 \mapsto v$. Lorsque $V = K$, c'est l'isomorphisme
$L(K) \otimes_K K \to L(K)$ donné par la multiplication, c'est-à-dire par
l'action de $K$. Pour $V$
quelconque, c'est un isomorphisme en vertu du cas $V = K$ et du fait que $L$
commute aux produits finis (remarque \ref{remark-linear-functor}).
\end{proof}

\noindent
Pour un anneau $R$ et un $R$-module $M$, soit $R[M]$ la $R$-algèbre dont le
$R$-module sous-jacent est $R \oplus M$ et dont la multiplication est donnée
par $(r, m) \cdot (r', m') = (rr', rm' + r'm)$. Lorsque $M = R$, c'est
l'algèbre des nombres duaux sur $R$, que l'on note $R[\epsilon]$.

\medskip\noindent
Soit maintenant $S$ un anneau, et supposons que $R$ soit une $S$-algèbre.
L'affectation $M \mapsto R[M]$ détermine alors un foncteur
$\text{Mod}_R \to S\text{-Alg}/R$, où $S\text{-Alg}/R$ désigne la catégorie
des $S$-algèbres au-dessus de $R$. Remarquons que $S\text{-Alg}/R$ admet les
produits finis : si $A_1 \to R$ et $A_2 \to R$ sont deux objets, alors
$A_1 \times_R A_2$ est un produit.

\begin{lemma}
\label{lemma-preserves-products}
Soit $R$ une $S$-algèbre. Alors le foncteur
$\text{Mod}_R \to S\text{-Alg}/R$ décrit ci-dessus préserve les produits
finis.
\end{lemma}

\begin{proof}
Il s'agit simplement de l'assertion suivante : si $M$ et $N$ sont des
$R$-modules, le morphisme $R[M \times N] \to R[M] \times_R R[N]$ est un
isomorphisme dans $S\text{-Alg}/R$.
\end{proof}

\begin{lemma}
\label{lemma-tangent-space-functor}
Soit $R$ une $S$-algèbre, et soit $\mathcal{C}$ une sous-catégorie strictement
pleine de $S\text{-Alg}/R$ contenant $R[M]$ pour tout
$M \in \Ob(\text{Mod}^{fg}_R)$. Soit
$F: \mathcal{C} \to \textit{Ensembles}$ un foncteur. Supposons que $F(R)$ soit un
ensemble à un seul élément et que, pour tous $M, N \in
\Ob(\text{Mod}^{fg}_R)$, l'application
induite
$$
F(R[M] \times_R R[N]) \to F(R[M]) \times F(R[N])
$$
soit une bijection. Alors $F(R[M])$ possède une structure naturelle de
$R$-module pour tout $M
\in \Ob(\text{Mod}^{fg}_R)$.
\end{lemma}

\begin{proof}
Remarquons que $R \cong R[0]$ et
$R[M] \times_R R[N] \cong R[M \times N]$ ; ainsi, $R$ et
$R[M] \times_R R[N]$ sont des objets de $\mathcal{C}$ par nos hypothèses sur
$\mathcal{C}$. Les conditions imposées à $F$ ont donc un sens. Le foncteur
$\text{Mod}_R \to S\text{-Alg}/R$ du lemme \ref{lemma-preserves-products} se
restreint à un foncteur $\text{Mod}^{fg}_R \to \mathcal{C}$ par l'hypothèse
sur $\mathcal{C}$. Soit $L$ la composée
$\text{Mod}^{fg}_R \to \mathcal{C} \to \textit{Ensembles}$, c'est-à-dire
$L(M) = F(R[M])$. Alors $L$ préserve les produits finis d'après le lemme
\ref{lemma-preserves-products} et l'hypothèse sur $F$. Le lemme
\ref{lemma-linear-functor} montre donc que $L(M) = F(R[M])$ possède une
structure naturelle de $R$-module pour tout
$M \in \Ob(\text{Mod}^{fg}_R)$.
\end{proof}

\begin{definition}
\label{definition-tangent-space-over-R}
Soit $\mathcal{C}$ une catégorie comme dans le lemme
\ref{lemma-tangent-space-functor}. Soit
$F : \mathcal{C} \to \textit{Ensembles}$ un foncteur tel que $F(R)$ soit un
ensemble à un seul élément. L'{\it espace tangent $TF$ de $F$} est $F(R[\epsilon])$.
\end{definition}

\noindent
Lorsque $F : \mathcal{C} \to \textit{Ensembles}$ satisfait les hypothèses du
lemme \ref{lemma-tangent-space-functor}, l'espace tangent $TF$ possède une
structure naturelle de $R$-module.

\begin{example}
\label{example-tangent-space-functor}
Puisque $\mathcal{C}_\Lambda$ contient tous les $k[V]$ pour les espaces
vectoriels $V$ de dimension finie, la définition
\ref{definition-tangent-space-over-R} s'applique avec $S = \Lambda$, $R = k$,
$\mathcal{C} = \mathcal{C}_\Lambda$ et
$F : \mathcal{C}_\Lambda \to \textit{Ensembles}$ un foncteur de prédéformation.
L'espace tangent est $TF = F(k[\epsilon])$.
\end{example}

\begin{example}
\label{example-tangent-space-prorepresentable-functor}
Calculons l'espace tangent de l'exemple
\ref{example-tangent-space-functor} lorsque
$F : \mathcal{C}_\Lambda \to \textit{Ensembles}$ est un foncteur
pro-représentable, disons $F = \underline{S}|_{\mathcal{C}_\Lambda}$ pour
$S \in \Ob(\widehat{\mathcal{C}}_\Lambda)$. Alors $F$ commute aux limites
arbitraires et satisfait donc les hypothèses du lemme
\ref{lemma-tangent-space-functor}. On calcule
$$
TF = F(k[\epsilon]) = \Mor_{\mathcal{C}_\Lambda}(S, k[\epsilon])
= \text{Der}_\Lambda(S, k)
$$
et, plus généralement, pour un $k$-espace vectoriel $V$ de dimension finie,
on a
$$
F(k[V]) = \Mor_{\mathcal{C}_\Lambda}(S, k[V]) = \text{Der}_\Lambda(S, V).
$$
Explicitement, un morphisme de $\Lambda$-algèbres $f : S \to k[V]$
compatible avec les augmentations $q : S \to k$ et $k[V] \to k$ correspond
à la dérivation $D$ définie par $s \mapsto f(s) - q(s)$. Réciproquement, une
$\Lambda$-dérivation $D : S \to V$ correspond au morphisme
$f : S \to k[V]$ de $\mathcal{C}_\Lambda$ défini par
$f(s) = q(s) + D(s)$. Ces identifications étant fonctorielles, les structures
de $k$-espace vectoriel de $TF$ et de $\text{Der}_\Lambda(S, k)$ se
correspondent (voir le lemme \ref{lemma-morphism-linear-functors}). Il résulte
du lemme \ref{lemma-derivations-finite} que $\dim_k TF$ est finie.
\end{example}

\begin{example}
\label{example-tangent-space-classical-prorepresentable-functor}
Le calcul de l'exemple
\ref{example-tangent-space-prorepresentable-functor} se simplifie dans le cas
classique. Dans ce cas, l'espace tangent du foncteur
$F = \underline{S}|_{\mathcal{C}_\Lambda}$ est simplement l'espace cotangent
relatif de $S$ sur $\Lambda$, soit $TF = T_{S/\Lambda}$. Cela reste en fait
valable, plus généralement, lorsque l'extension de corps $k/k'$ est
séparable. Voir Exercices, exercice
\ref{exercises-exercise-tangent-space-Zariski}.
\end{example}

\begin{lemma}
\label{lemma-morphism-tangent-spaces}
Soient $F, G: \mathcal{C} \to \textit{Ensembles}$ des foncteurs satisfaisant les
hypothèses du lemme \ref{lemma-tangent-space-functor}. Soit $t : F \to G$ un
morphisme de foncteurs. Pour tout $M \in \Ob(\text{Mod}^{fg}_R)$,
l'application $t_{R[M]}: F(R[M]) \to G(R[M])$ est un morphisme de
$R$-modules, où $F(R[M])$ et $G(R[M])$ sont munis des structures de
$R$-modules du lemme \ref{lemma-tangent-space-functor}. En particulier,
$t_{R[\epsilon]} : TF \to TG$ est un morphisme de $R$-modules.
\end{lemma}

\begin{proof}
Cela résulte du lemme \ref{lemma-morphism-linear-functors}.
\end{proof}

\begin{example}
\label{example-tangent-space-map-prorepresentable-functor}
Supposons que $f : R \to S$ soit un morphisme d'anneaux dans
$\widehat{\mathcal{C}}_\Lambda$. Posons
$F = \underline{R}|_{\mathcal{C}_\Lambda}$ et
$G = \underline{S}|_{\mathcal{C}_\Lambda}$. Le morphisme d'anneaux $f$ induit
une transformation de foncteurs $G \to F$. Le lemme
\ref{lemma-morphism-tangent-spaces} donne une application $k$-linéaire
$TG \to TF$. Il s'agit de l'application
$$
TG = \text{Der}_\Lambda(S, k) \longrightarrow \text{Der}_\Lambda(R, k) = TF
$$
comme il résulte des identifications canoniques
$F(k[V]) = \text{Der}_\Lambda(R, V)$ et
$G(k[V]) = \text{Der}_\Lambda(S, V)$ de l'exemple
\ref{example-tangent-space-prorepresentable-functor}, ainsi que de la règle de
calcul de l'application sur les espaces tangents.
\end{example}

\begin{lemma}
\label{lemma-tangent-space-tensor}
Soit $F: \mathcal{C} \to \textit{Ensembles}$ un foncteur satisfaisant les
hypothèses du lemme \ref{lemma-tangent-space-functor}. Supposons que $R = K$
soit un corps. Alors $F(K[V]) \cong TF \otimes_K V$ pour tout $K$-espace
vectoriel $V$ de dimension finie.
\end{lemma}

\begin{proof}
Cela résulte du lemme \ref{lemma-linear-functor-over-field}.
\end{proof}






\section{Espaces tangents des catégories de prédéformation}
\label{section-tangent-spaces}

\noindent
Nous définirons les espaces tangents des foncteurs de prédéformation à l'aide
de la définition générale \ref{definition-tangent-space-over-R}. Nous l'avons
explicitée dans l'exemple \ref{example-tangent-space-functor}. Elle s'applique
aux catégories de prédéformation en considérant le foncteur associé des
classes d'isomorphie.

\begin{definition}
\label{definition-tangent-space}
Soit $\mathcal{F}$ une catégorie de prédéformation. L'{\it espace tangent
$T \mathcal{F}$ de $\mathcal{F}$} est l'ensemble
$\overline{\mathcal{F}}(k[\epsilon])$ des classes d'isomorphie d'objets de la
catégorie fibre $\mathcal
F(k[\epsilon])$.
\end{definition}

\noindent
Ainsi, $T \mathcal{F}$ n'est autre que l'espace tangent du foncteur associé
$\overline{\mathcal{F}}: \mathcal{C}_\Lambda \to \textit{Ensembles}$. Il possède
une structure naturelle d'espace vectoriel lorsque $\mathcal{F}$ satisfait
(S2), ou plus généralement dès que $\overline{\mathcal{F}}$ la satisfait.

\begin{lemma}
\label{lemma-tangent-space-vector-space}
Soit $\mathcal{F}$ une catégorie de prédéformation telle que
$\overline{\mathcal{F}}$ satisfasse (S2)\footnote{Par exemple, si
$\mathcal{F}$ satisfait (S2), voir le lemme
\ref{lemma-S1-S2-associated-functor}.}. Alors $T \mathcal{F}$ possède une
structure naturelle de $k$-espace vectoriel. Pour tout espace vectoriel $V$
de dimension finie, on a
$\overline{\mathcal{F}}(k[V]) = T\mathcal{F} \otimes_k V$
fonctoriellement en $V$.
\end{lemma}

\begin{proof}
Posons
$F = \overline{\mathcal{F}} : \mathcal{C}_\Lambda \to \textit{Ensembles}$.
C'est un foncteur de prédéformation et $F$ satisfait (S2). D'après le lemme
\ref{lemma-S2-extensions} (et la traduction de la remarque
\ref{remark-compare-S1-S2-schlessinger}), on voit que
$$
F(A \times_k k[V]) \longrightarrow F(A) \times F(k[V])
$$
est une bijection pour tout espace vectoriel $V$ de dimension finie et tout
$A \in \Ob(\mathcal{C}_\Lambda)$. En particulier, si $A = k[W]$, on voit que
$F(k[W] \times_k k[V]) = F(k[W]) \times F(k[V])$. Autrement dit, les
hypothèses du lemme \ref{lemma-tangent-space-functor} sont satisfaites, et
$TF = T \mathcal{F}$ possède une structure naturelle de $k$-espace
vectoriel. La dernière assertion résulte du lemme
\ref{lemma-tangent-space-tensor}.
\end{proof}

\noindent
Un morphisme de catégories de prédéformation induit une application sur les
espaces tangents.

\begin{definition}
\label{definition-differential}
Soit $\varphi : \mathcal{F} \to \mathcal{G}$ un morphisme de catégories de
prédéformation. La {\it différentielle
$d \varphi : T \mathcal{F} \to T \mathcal{G}$ de $\varphi$} est l'application
obtenue en évaluant le morphisme de foncteurs
$\overline{\varphi}: \overline{\mathcal{F}} \to \overline{\mathcal{G}}$ en
$A = k[\epsilon]$.
\end{definition}

\begin{lemma}
\label{lemma-k-linear-differential}
Soit $\varphi : \mathcal{F} \to \mathcal{G}$ un morphisme de catégories de
prédéformation. Supposons que $\overline{\mathcal{F}}$ et
$\overline{\mathcal{G}}$ satisfassent toutes deux (S2). Alors
$d \varphi : T \mathcal{F} \to T \mathcal{G}$ est $k$-linéaire.
\end{lemma}

\begin{proof}
Dans la démonstration du lemme \ref{lemma-tangent-space-vector-space}, nous
avons vu que $\overline{\mathcal{F}}$ et $\overline{\mathcal{G}}$ satisfont
les hypothèses du lemme \ref{lemma-tangent-space-functor}. Le résultat découle
donc du lemme \ref{lemma-morphism-tangent-spaces}.
\end{proof}

\begin{remark}
\label{remark-tangent-space-cofibered-groupoid}
On peut étendre les notions d'espace tangent et de différentielle aux
catégories cofibrées en groupoïdes arbitraires comme suit. Soit $\mathcal{F}$
une catégorie cofibrée en groupoïdes au-dessus de $\mathcal{C}_\Lambda$, et
soit $x \in \Ob(\mathcal{F}(k))$. Comme dans la remarque
\ref{remark-localize-cofibered-groupoid}, on obtient une catégorie de
prédéformation $\mathcal{F}_x$. On définit
$$
T_x\mathcal{F} = T\mathcal{F}_x
$$
comme l'{\it espace tangent de $\mathcal{F}$ en $x$}. Si
$\varphi : \mathcal{F} \to \mathcal{G}$ est un morphisme de catégories
cofibrées en groupoïdes au-dessus de $\mathcal{C}_\Lambda$ et si
$x \in \Ob(\mathcal{F}(k))$, il existe un morphisme induit
$\varphi_x: \mathcal{F}_x \to \mathcal{G}_{\varphi(x)}$. On définit la
{\it différentielle
$d_x \varphi : T_x \mathcal{F} \to T_{\varphi(x)} \mathcal{G}$
de $\varphi$ en $x$} comme l'application
$d \varphi_x: T \mathcal{F}_x \to T \mathcal{G}_{\varphi(x)}$.
Si $\mathcal{F}$ et $\mathcal{G}$ satisfont toutes deux (S2), alors tous ces
espaces tangents possèdent une structure naturelle de $k$-espace vectoriel et
toutes les différentielles
$d_x \varphi : T_x \mathcal{F} \to T_{\varphi(x)} \mathcal{G}$
sont $k$-linéaires (utiliser les lemmes \ref{lemma-S1-S2-localize} et
\ref{lemma-k-linear-differential}).
\end{remark}

\noindent
Les observations suivantes sont sans intérêt dans le cas classique ou lorsque
$k/k'$ est une extension de corps séparable, car
$\text{Der}_\Lambda(k, k)$ et $\text{Der}_\Lambda(k, V)$ sont alors nuls. Il
existe une identification canonique
$$
\Mor_{\mathcal{C}_\Lambda}(k, k[\epsilon]) =
\text{Der}_\Lambda(k, k).
$$
En effet, pour $D \in \text{Der}_\Lambda(k, k)$, soit
$f_D : k \to k[\epsilon]$ le morphisme $a \mapsto a + D(a)\epsilon$. Plus
généralement, pour un espace vectoriel $V$ de dimension finie sur $k$, on a
$$
\Mor_{\mathcal{C}_\Lambda}(k, k[V]) =
\text{Der}_\Lambda(k, V)
$$
et nous utiliserons la même notation $f_D$ pour le morphisme associé à la
dérivation $D$. On a également
$$
\Mor_{\mathcal{C}_\Lambda}(k[W], k[V]) =
\Hom_k(V, W) \oplus \text{Der}_\Lambda(k, V)
$$
où $(\varphi, D)$ correspond au morphisme
$f_{\varphi, D} : a + w \mapsto a + \varphi(w) + D(a)$. Nous noterons parfois
$f_{1, D} : a + v \to a + v + D(a)$ l'automorphisme de $k[V]$ déterminé par
la dérivation $D : k \to V$. Remarquons que
$f_{1, D} \circ f_{1, D'} = f_{1, D + D'}$.

\medskip\noindent
Soit $\mathcal{F}$ une catégorie de prédéformation au-dessus de
$\mathcal{C}_\Lambda$. Soit $x_0 \in \Ob(\mathcal{F}(k))$. D'après ce qui
précède, il existe une application canonique
$$
\gamma_V :
\text{Der}_\Lambda(k, V)
\longrightarrow
\overline{\mathcal{F}}(k[V])
$$
définie par $D \mapsto f_{D, *}(x_0)$. Il existe en outre une action
$$
a_V : \text{Der}_\Lambda(k, V) \times \overline{\mathcal{F}}(k[V])
\longrightarrow
\overline{\mathcal{F}}(k[V])
$$
définie par $(D, x) \mapsto f_{1, D, *}(x)$. Ces deux applications sont
compatibles, c'est-à-dire que
$f_{1, D, *}f_{D', *}x_0 = f_{D + D', *}x_0$, comme le montre le calcul de
leurs composées. Remarquons que les applications $\gamma_V$ et $a_V$ ne
dépendent pas du choix de $x_0$, puisque $x_0$ est unique à isomorphisme près.

\begin{lemma}
\label{lemma-action-linear}
Soit $\mathcal{F}$ une catégorie de prédéformation au-dessus de
$\mathcal{C}_\Lambda$. Si $\overline{\mathcal{F}}$ satisfait (S2), alors les
applications $\gamma_V$ sont $k$-linéaires et l'on a
$a_V(D, x) = x + \gamma_V(D)$.
\end{lemma}

\begin{proof}
Dans la démonstration du lemme \ref{lemma-tangent-space-vector-space}, nous
avons vu que le foncteur $V \mapsto \overline{\mathcal{F}}(k[V])$ envoie $0$
sur un ensemble à un seul élément et les produits sur des produits. Il en va de même du
foncteur $V \mapsto \text{Der}_\Lambda(k, V)$. Le lemme
\ref{lemma-morphism-linear-functors} montre donc que $\gamma_V$ est linéaire.
Soit $D : k \to V$ une $\Lambda$-dérivation. Posons
$D_1 : k \to V^{\oplus 2}$ égal à $a \mapsto (D(a), 0)$. Alors
$$
\xymatrix{
k[V \times V] \ar[r]_{+} \ar[d]^{f_{1, D_1}} & k[V] \ar[d]^{f_{1, D}} \\
k[V \times V] \ar[r]^{+} & k[V]
}
$$
est commutatif. En développant les définitions et en utilisant
$\overline{F}(V \times V) = \overline{F}(V) \times \overline{F}(V)$, cela
signifie que $a_D(x_1) + x_2 = a_D(x_1 + x_2)$ pour tous
$x_1, x_2 \in \overline{F}(V)$. Il suffit donc de montrer que
$a_V(D, 0) = 0 + \gamma_V(D)$, où $0 \in \overline{F}(V)$ est le vecteur
nul. Par définition, celui-ci est l'élément $f_{0, *}(x_0)$. Comme
$f_D = f_{1, D} \circ f_0$, le résultat voulu s'ensuit.
\end{proof}

\noindent
Deux cas particuliers des constructions précédentes sont l'application
\begin{equation}
\label{equation-map}
\gamma : \text{Der}_\Lambda(k, k) \longrightarrow T\mathcal{F}
\end{equation}
ainsi que l'action
\begin{equation}
\label{equation-action}
a : \text{Der}_\Lambda(k, k) \times T\mathcal{F} \longrightarrow T\mathcal{F}
\end{equation}
définies pour toute catégorie de prédéformation $\mathcal{F}$. Remarquons que,
si $\varphi : \mathcal{F} \to \mathcal{G}$ est un morphisme de catégories de
prédéformation, on obtient les diagrammes commutatifs
$$
\vcenter{
\xymatrix{
\text{Der}_\Lambda(k, k) \ar[r]_-\gamma \ar[rd]_\gamma &
T\mathcal{F} \ar[d]_{d\varphi} \\
& T\mathcal{G}
}
}
\quad\text{et}\quad
\vcenter{
\xymatrix{
\text{Der}_\Lambda(k, k) \times T\mathcal{F} \ar[r]_-a
\ar[d]_{1 \times d\varphi} &
T\mathcal{F} \ar[d]^{d\varphi} \\
\text{Der}_\Lambda(k, k) \times T\mathcal{G} \ar[r]^-a &
T\mathcal{G}
}
}
$$










\section{Objets formels versels}
\label{section-versal-objects}

\noindent
L'existence d'un objet formel versel impose à $\mathcal{F}$ de posséder la
propriété (S1).

\begin{lemma}
\label{lemma-versal-object-S1}
Soit $\mathcal{F}$ une catégorie de prédéformation. Supposons que
$\mathcal{F}$ possède un objet formel versel. Alors $\mathcal{F}$ satisfait
(S1).
\end{lemma}

\begin{proof}
Soit $\xi$ un objet formel versel de $\mathcal{F}$. Soit
$$
\xymatrix{
           & x_2 \ar[d] \\
x_1 \ar[r] & x
}
$$
un diagramme dans $\mathcal{F}$ tel que $x_2 \to x$ soit situé au-dessus d'un
morphisme surjectif d'anneaux. Comme le morphisme naturel
$\widehat{\mathcal{F}}|_{\mathcal{C}_\Lambda} \xrightarrow{\sim} \mathcal{F}$
est une équivalence (voir la remarque \ref{remark-restrict-completion}), on
peut aussi considérer ce diagramme comme un diagramme dans
$\widehat{\mathcal{F}}$. D'après le lemme
\ref{lemma-versal-object-quasi-initial}, il existe un morphisme
$\xi \to x_1$ ; la remarque \ref{remark-versal-object} fournit donc également
un morphisme $\xi \to x_2$ qui rend le diagramme
$$
\xymatrix{
\xi \ar[r] \ar[d]          & x_2 \ar[d] \\
x_1 \ar[r] & x
}
$$
commutatif. Si $x_1 \to x$ et $x_2 \to x$ sont situés au-dessus des
morphismes d'anneaux $A_1 \to A$ et $A_2 \to A$, l'image directe de $\xi$
vers $A_1 \times_A A_2$ fournit un objet $y$ comme l'exige (S1).
\end{proof}

\noindent
Lorsque notre catégorie cofibrée satisfait (S1) et (S2), on peut caractériser
les objets formels versels de la manière suivante.

\begin{lemma}
\label{lemma-versal-criterion}
Soit $\mathcal{F}$ une catégorie de prédéformation satisfaisant (S1) et (S2).
Soit $\xi$ un objet formel de $\mathcal{F}$ correspondant à
$\underline{\xi} : \underline{R}|_{\mathcal{C}_\Lambda} \to \mathcal{F}$ ;
voir la remarque \ref{remark-formal-objects-yoneda}. Alors $\xi$ est versel
si et seulement si les deux conditions suivantes sont satisfaites :
\begin{enumerate}
\item l'application
$d\underline{\xi} : T\underline{R}|_{\mathcal{C}_\Lambda} \to T\mathcal{F}$
sur les espaces tangents est surjective ;
\item étant donné un diagramme dans $\widehat{\mathcal{F}}$
$$
\vcenter{
\xymatrix{
            &  y \ar[d] \\
\xi \ar[r]  &  x
}
}
\quad\text{au-dessus de}\quad
\vcenter{
\xymatrix{
         &   B  \ar[d]^{f} \\
R \ar[r] &   A
}
}
$$
au-dessus d'un diagramme de $\widehat{\mathcal{C}}_\Lambda$ où $B \to A$ est
une petite extension d'anneaux artiniens, il existe un morphisme d'anneaux
$R \to B$ tel que
$$
\xymatrix{
         &   B  \ar[d]^{f} \\
R \ar[ur] \ar[r] &   A
}
$$
soit commutatif.
\end{enumerate}
\end{lemma}

\begin{proof}
Si $\xi$ est versel, (1) découle du lemme
\ref{lemma-smooth-morphism-essentially-surjective} et (2) de la remarque
\ref{remark-versal-object}. Supposons (1) et (2) satisfaites. D'après la
remarque \ref{remark-versal-object}, il faut montrer que, pour un diagramme de
$\widehat{\mathcal{F}}$ comme dans (2), il existe $\xi \to y$ tel que
$$
\xymatrix{
            &  y \ar[d] \\
\xi \ar[ur] \ar[r]  &  x
}
$$
soit commutatif. Soit $b : R \to B$ le morphisme fourni par (2). Posons
$y' = b_*\xi$ et choisissons une factorisation $\xi \to y' \to x$, située
au-dessus de $R \to B \to A$, du morphisme donné $\xi \to x$. La condition
(S1) fournit un diagramme commutatif
$$
\vcenter{
\xymatrix{
z  \ar[r] \ar[d]          &  y \ar[d] \\
y' \ar[r]  &  x
}
}
\quad\text{au-dessus de}\quad
\vcenter{
\xymatrix{
B \times_A B \ar[d] \ar[r] &   B  \ar[d]^{f} \\
B \ar[r]^{f} &   A .
}
}
$$
Posons $I = \Ker(f)$. Soit $\overline{g} : B \times_A B \to k[I]$ le
morphisme d'anneaux $(u, v) \mapsto \overline{u} \oplus (v - u)$ ; cf.\ le
lemme \ref{lemma-lifting-along-small-extension}. Par (1), il existe un
morphisme $\xi \to \overline{g}_*z$ situé au-dessus d'un morphisme d'anneaux
$i : R \to k[\epsilon]$. Choisissons un quotient artinien
$b_1 : R \to B_1$ tel que $b : R \to B$ et $i : R \to k[\epsilon]$ se
factorisent tous deux par $R \to B_1$, donnant ainsi
$h : B_1 \to B$ et $i' : B_1 \to k[\epsilon]$. Choisissons une image directe
$y_1 = b_{1, *}\xi$, une factorisation $\xi \to y_1 \to y'$ située au-dessus
de $R \to B_1 \to B$ de $\xi \to y'$, et une factorisation
$\xi \to y_1 \to \overline{g}_*z$ située au-dessus de
$R \to B_1 \to k[\epsilon]$ de $\xi \to \overline{g}_*z$. Une nouvelle
application de (S1) donne
$$
\vcenter{
\xymatrix{
z_1  \ar[r] \ar[d] & z \ar[r] \ar[d] & y \ar[d] \\
y_1 \ar[r] & y' \ar[r] &  x
}
}
\quad\text{au-dessus de}\quad
\vcenter{
\xymatrix{
B_1 \times_A B \ar[d] \ar[r] & B \times_A B \ar[r] \ar[d] & B \ar[d]^{f} \\
B_1 \ar[r]  & B \ar[r] &  A .
}
}
$$
Remarquons que le morphisme $g : B_1 \times_A B \to k[I]$ du lemme
\ref{lemma-lifting-along-small-extension} (définie à l'aide de $h$) est la
composée de $B_1 \times_A B \to B \times_A B$ et du morphisme
$\overline{g}$ ci-dessus. Par construction, il existe un morphisme
$y_1 \to g_*z_1 \cong \overline{g}_*z$ ! Le lemme
\ref{lemma-lifting-along-small-extension} s'applique donc (aux rectangles
extérieurs des diagrammes précédents) et fournit un morphisme $y_1 \to y$ ;
en le précomposant avec $\xi \to y_1$, on obtient le morphisme recherché
$\xi \to y$.
\end{proof}

\noindent
Si $\mathcal{F}$ possède la propriété (S1), alors le « plus grand quotient sur
lequel un relèvement existe » existe. Voici un énoncé précis.

\begin{lemma}
\label{lemma-largest-closed-where-lift}
Soit $\mathcal{F}$ une catégorie cofibrée en groupoïdes au-dessus de
$\mathcal{C}_\Lambda$ qui satisfait (S1). Soit $B \to A$ une surjection dans
$\mathcal{C}_\Lambda$, de noyau $I$ annulé par $\mathfrak m_B$. Soit
$x \in \mathcal{F}(A)$. L'ensemble des idéaux
$$
\mathcal{J} = \{ J \subset I \mid
\text{il existe un }y \to x\text{ au-dessus de }B/J \to A\}
$$
possède un plus petit élément.
\end{lemma}

\begin{proof}
Remarquons que $\mathcal{J}$ n'est pas vide, puisque
$I \in \mathcal{J}$. De plus, si $J \in \mathcal{J}$ et
$J \subset J' \subset I$, alors $J' \in \mathcal{J}$, car on peut prendre
l'image directe de l'objet $y$ en un objet $y'$ au-dessus de $B/J'$. Soient
$J$ et $K$ des éléments de l'ensemble affiché. Nous affirmons que
$J \cap K \in \mathcal{J}$, ce qui démontrera le lemme. Comme $I$ est un
$k$-espace vectoriel, on peut trouver un idéal $J \subset J' \subset I$ tel
que $J \cap K = J' \cap K$ et $J' + K = I$. D'après ce qui précède, on peut
remplacer $J$ par $J'$ et supposer $J + K = I$. Dans ce cas,
$$
A/(J \cap K) = A/J \times_{A/I} A/K.
$$
L'existence d'un élément $z \in \mathcal{F}(A/(J \cap K))$ qui s'envoie sur
$x$ résulte donc, par (S1), de l'existence supposée des éléments au-dessus de
$A/J$ et de $A/K$.
\end{proof}

\noindent
Nous améliorerons plus loin le résultat suivant.

\begin{lemma}
\label{lemma-versal-object-existence}
Soit $\mathcal{F}$ une catégorie cofibrée en groupoïdes au-dessus de
$\mathcal{C}_\Lambda$. Supposons les conditions suivantes satisfaites :
\begin{enumerate}
\item $\mathcal{F}$ est une catégorie de prédéformation.
\item $\mathcal{F}$ satisfait (S1).
\item $\mathcal{F}$ satisfait (S2).
\item $\dim_k T\mathcal{F}$ est finie.
\end{enumerate}
Alors $\mathcal{F}$ possède un objet formel versel.
\end{lemma}

\begin{proof}
Supposons (1), (2), (3) et (4) satisfaites. Choisissons un objet
$R \in \Ob(\widehat{\mathcal{C}}_\Lambda)$ tel que
$\underline{R}|_{\mathcal{C}_\Lambda}$ soit lisse. Voir le lemme
\ref{lemma-exists-smooth}. Soit $r = \dim_k T\mathcal{F}$ et posons
$S = R[[X_1, \ldots, X_r]]$.

\medskip\noindent
Nous allons construire par récurrence, pour $n \geq 2$, des paires
$(J_n, f_{n - 1} : \xi_n \to \xi_{n - 1})$, où les $J_n \subset S$ forment
une suite décroissante d'idéaux et où
$f_{n - 1} : \xi_n \to \xi_{n - 1}$ est un morphisme de $\mathcal{F}$ situé
au-dessus de la projection $S/J_n \to S/J_{n - 1}$.

\medskip\noindent
Étape 1. Soit $J_1 = \mathfrak m_S$. Soit $\xi_1$ l'unique objet de
$\mathcal{F}$ au-dessus de $k = S/J_1 = S/\mathfrak m_S$, à isomorphisme
unique près.

\medskip\noindent
Étape 2. Soit $J_2 = \mathfrak m_S^2 + \mathfrak{m}_R S$. Alors
$S/J_2 = k[V]$, où $V = kX_1 \oplus \ldots \oplus kX_r$. La condition (S2)
pour $\overline{\mathcal{F}}$ donne une bijection
$$
\overline{\mathcal{F}}(S/J_2)
\longrightarrow
T\mathcal{F} \otimes_k V,
$$
voir les lemmes \ref{lemma-S1-S2-associated-functor} et
\ref{lemma-tangent-space-vector-space}. Choisissons une base
$\theta_1, \ldots, \theta_r$ de $T\mathcal{F}$ et posons
$\xi_2 = \sum \theta_i \otimes X_i \in \Ob(\mathcal{F}(S/J_2))$. Ce choix a
pour effet que
$$
d\xi_2 :
\Mor_{\mathcal{C}_\Lambda}(S/J_2, k[\epsilon])
\longrightarrow
T\mathcal{F}
$$
est surjective. Soit $f_1 : \xi_2 \to \xi_1$ l'unique morphisme.

\medskip\noindent
Étape de récurrence. Supposons construite
$(J_n, f_{n - 1} : \xi_n \to \xi_{n - 1})$ pour un certain $n \geq 2$. Il
existe un plus petit élément $J_{n + 1}$ de l'ensemble des idéaux
$J \subset S$ satisfaisant : (a) $\mathfrak m_S J_n \subset J \subset J_n$
et (b) il existe un morphisme $\xi_{n + 1} \to \xi_n$ situé au-dessus de
$S/J \to S/J_n$ ; voir le lemme \ref{lemma-largest-closed-where-lift}. Soit
$f_n : \xi_{n + 1} \to \xi_n$ un morphisme quelconque de $\mathcal{F}$ situé
au-dessus de $S/J_{n + 1} \to S/J_n$.

\medskip\noindent
Posons $J = \bigcap J_n$, $\overline{S} = S/J$ et
$\overline{J}_n = J_n/J$. D'après le lemme \ref{lemma-m-adic-topology}, la
suite d'idéaux $(\overline{J}_n)$ induit la topologie
$\mathfrak m_{\overline{S}}$-adique sur $\overline{S}$. Comme
$(\xi_n, f_n)$ est un objet de
$\widehat{\mathcal{F}}_\mathcal{I}(\overline{S})$, où $\mathcal{I}$ est la
filtration $(\overline{J}_n)$ de $\overline{S}$, on voit que
$(\xi_n, f_n)$ induit un objet $\xi$ de
$\widehat{\mathcal{F}}(\overline{S})$ ; voir le lemme
\ref{lemma-formal-objects-different-filtration}.

\medskip\noindent
Montrons que $\xi$ est versel. Pour la versalité, il suffit de vérifier les
conditions (1) et (2) du lemme \ref{lemma-versal-criterion}. La condition (1)
résulte de notre choix de $\xi_2$ à l'étape 2. Considérons un diagramme dans
$\widehat{\mathcal{F}}$
$$
\vcenter{
\xymatrix{
            &  y \ar[d] \\
\xi \ar[r]  &  x
}
}
\quad\text{au-dessus de}\quad
\vcenter{
\xymatrix{
         &   B  \ar[d]^{f} \\
\overline{S} \ar[r] &   A
}
}
$$
au-dessus d'un diagramme de $\widehat{\mathcal{C}}_\Lambda$, où
$f: B \to A$ est une petite extension d'anneaux artiniens. Il faut montrer
qu'il existe un morphisme $\overline{S} \to B$ s'insérant dans le
diagramme de droite. Choisissons $n$ tel que $\overline{S} \to A$ se
factorise par $\overline{S} \to S/J_n$. C'est possible puisque la suite
$(\overline{J}_n)$ induit la topologie $\mathfrak m_{\overline{S}}$-adique,
comme nous l'avons vu. L'image directe de $\xi$ le long de
$\overline{S} \to S/J_n$ est $\xi_n$. On peut factoriser $\xi \to x$ en
$\xi \to \xi_n \to x$ ; on obtient donc un diagramme dans $\mathcal{F}$
$$
\vcenter{
\xymatrix{
            &  y \ar[d] \\
\xi_n \ar[r]  &  x
}
}
\quad\text{au-dessus de}\quad
\vcenter{
\xymatrix{
         &   B  \ar[d]^{f} \\
S/J_n \ar[r] &   A .
}
}
$$
Pour vérifier la condition (2) du lemme \ref{lemma-versal-criterion}, il
suffit de compléter le diagramme
$$
\xymatrix{
S/J_{n + 1} \ar[d] \ar@{-->}[r] & B \ar[d]^{f} \\
S/J_n   \ar[r] & A
}
$$
ou, de manière équivalente, le diagramme
$$
\xymatrix{
  &  S/J_n \times_A B \ar[d]^{p_1} \\
S/J_{n + 1} \ar@{-->}[ur] \ar[r] &  S/J_n.
}
$$
Si $p_1$ admet une section, la conclusion est acquise. Sinon, d'après le lemme
\ref{lemma-fiber-product-CLambda} (2), $p_1$ est une petite extension ; le
lemme \ref{lemma-essential-surjection} (4) montre donc que $p_1$ est une
surjection essentielle. Rappelons que $S = R[[X_1, \ldots, X_r]]$ et que
nous avons choisi $R$ de sorte que
$\underline{R}|_{\mathcal{C}_\Lambda}$ soit lisse. Il existe donc un
morphisme $h : R \to B$ relevant le morphisme
$R \to S \to S/J_n \to A$. Par la propriété universelle d'un anneau de
séries formelles, il existe un morphisme de $R$-algèbres
$h : S = R[[X_1, \ldots, X_r]] \to B$ relevant le morphisme donné
$S \to S/J_n \to A$. Celui-ci induit un morphisme
$g: S \to S/J_n \times_A B$ qui rend commutatif le carré en traits pleins du
diagramme
$$
\xymatrix{
S \ar[d] \ar[r]_-g  &  S/J_n \times_A B \ar[d]^{p_1} \\
S/J_{n + 1} \ar@{-->}[ur] \ar[r] &  S/J_n
}
$$
Alors $g$ est une surjection, puisque $p_1$ est une surjection essentielle.
Nous affirmons que l'idéal $K = \Ker(g)$ de $S$ satisfait les conditions (a)
et (b) de la construction de $J_{n + 1}$ à l'étape de récurrence précédente.
En effet, $K \subset J_n$ est clair et
$\mathfrak m_SJ_n \subset K$, puisque $p_1$ est une petite extension ; cela
démontre (a). En appliquant (S1) à
$$
\xymatrix{
            &  y \ar[d] \\
\xi_n \ar[r]  &  x,
}
$$
on obtient un relèvement de $\xi_n$ à
$S/K \cong S/J_n \times_A B$, donc (b) est satisfaite. Puisque $J_{n + 1}$
était le plus petit idéal possédant les propriétés (a) et (b), on a
$J_{n + 1} \subset K$. Le morphisme recherché
$S/J_{n+1} \to S/K \cong S/J_n \times_A B$ existe donc.
\end{proof}

\begin{remark}
\label{remark-compare-schlessinger-H3-pre}
Soit $F : \mathcal{C}_\Lambda \to \textit{Ensembles}$ un foncteur de
prédéformation satisfaisant (S1) et (S2). La condition
$\dim_k TF < \infty$ est précisément la condition (H3) de l'article de
Schlessinger. Rappelons que (S1) et (S2) correspondent aux conditions (H1) et
(H2) ; voir la remarque \ref{remark-compare-S1-S2-schlessinger}. Le lemme
\ref{lemma-versal-object-existence} nous dit donc que
$$
(H1) + (H2) + (H3)
\Rightarrow
\text{ there exists a versal formal object}
$$
pour les foncteurs de prédéformation. Nous établirons le lien avec les
enveloppes dans la remarque \ref{remark-compare-schlessinger-H3}.
\end{remark}





\section{Objets formels versels minimaux}
\label{section-minimal-versal}

\noindent
Nous effectuons quelques préparatifs afin de comprendre la (non-)unicité des
objets formels versels. Il s'avère que, si une catégorie de prédéformation
possède un objet formel versel, elle possède un objet formel versel minimal,
et que deux tels objets sont isomorphes. En outre, tous les objets formels
versels sont « plus ou moins » les mêmes, quitte à remplacer l'anneau de base
par une extension en séries formelles.

\medskip\noindent
Soit $\mathcal{F}$ une catégorie cofibrée en groupoïdes au-dessus de
$\mathcal{C}_\Lambda$. Pour tout objet $x$ de $\mathcal{F}$ situé au-dessus
de $A \in \Ob(\mathcal{C}_\Lambda)$, considérons la catégorie
$\mathcal{S}_x$ dont les objets sont
$$
\Ob(\mathcal{S}_x) =
\{x' \to x \mid x' \to x\text{ lies over }A' \subset A\}
$$
et dont les morphismes sont les morphismes au-dessus de $x$. Pour tout
$y \to x$ dans $\mathcal{F}$ situé au-dessus de $f : B \to A$ dans
$\mathcal{C}_\Lambda$, il existe un foncteur
$f_* : \mathcal{S}_y \to \mathcal{S}_x$ défini comme suit. Étant donné
$y' \to y$ situé au-dessus de $B' \subset B$, posons $A' = f(B')$ et soit
$y' \to x'$ au-dessus de $B' \to f(B')$ l'image directe de $y'$. Les axiomes
d'une catégorie cofibrée en groupoïdes donnent un unique morphisme
$x' \to x$ situé au-dessus de $f(B') \to A$ tel que
$$
\xymatrix{
y' \ar[d] \ar[r] & x' \ar[d] \\
y \ar[r] & x
}
$$
soit commutatif. Alors $x' \to x$ est un objet de $\mathcal{S}_x$. On dit
qu'un objet $x' \to x$ de $\mathcal{S}_x$ est {\it minimal} si tout morphisme
$(x'_1 \to x) \to (x' \to x)$ dans $\mathcal{S}_x$ est un isomorphisme,
c'est-à-dire si $x'$ et $x'_1$ sont définis sur le même sous-anneau de $A$.
Comme $A$ est de longueur finie comme $\Lambda$-module, les objets minimaux
existent toujours.

\begin{lemma}
\label{lemma-smallest-where-descends}
Soit $\mathcal{F}$ une catégorie cofibrée en groupoïdes au-dessus de
$\mathcal{C}_\Lambda$ qui satisfait (S1).
\begin{enumerate}
\item Pour $y \to x$ dans $\mathcal{F}$, l'image d'un objet minimal de
$\mathcal{S}_y$ est un objet minimal de $\mathcal{S}_x$.
\item Pour $y \to x$ dans $\mathcal{F}$ situé au-dessus d'une surjection
$f : B \to A$ dans $\mathcal{C}_\Lambda$, tout objet minimal de
$\mathcal{S}_x$ est l'image d'un objet minimal de $\mathcal{S}_y$.
\end{enumerate}
\end{lemma}

\begin{proof}
Démonstration de (1). Supposons $y \to x$ situé au-dessus de $f : B \to A$.
Soit $y' \to y$, situé au-dessus de $B' \subset B$, un objet minimal de
$\mathcal{S}_y$. Soit
$$
\vcenter{
\xymatrix{
y' \ar[d] \ar[r] & x' \ar[d] \\
y \ar[r] & x
}
}
\quad\text{au-dessus de}\quad
\vcenter{
\xymatrix{
B' \ar[d] \ar[r] & f(B') \ar[d] \\
B \ar[r] & A
}
}
$$
le diagramme de la construction de $f_*$ ci-dessus. Supposons que
$(x'' \to x) \to (x' \to x)$ soit un morphisme de $\mathcal{S}_x$, avec
$x'' \to x'$ situé au-dessus de $A'' \subset f(B')$. Par (S1), il existe
$y'' \to y'$ situé au-dessus de $B' \times_{f(B')} A'' \to B'$. La
minimalité de $y' \to y$ implique que
$B' \times_{f(B')} A'' \to B'$ est un isomorphisme, donc que
$A'' = f(B')$ ; ainsi, $x' \to x$ est minimal.

\medskip\noindent
Démonstration de (2). Supposons $f : B \to A$ surjective et $y \to x$ situé
au-dessus de $f$. Soit $x' \to x$ un objet minimal de $\mathcal{S}_x$ situé
au-dessus de $A' \subset A$. Par (S1), il existe $y' \to y$ situé au-dessus
de $B' = f^{-1}(A') = B \times_A A' \to B$ dont l'image dans
$\mathcal{S}_x$ est $x' \to x$. Ainsi, $f_*(y' \to y) = x' \to x$.
Choisissons un morphisme $(y'' \to y) \to (y' \to y)$ dans
$\mathcal{S}_y$ tel que $y'' \to y$ soit minimal (c'est possible d'après la
remarque sur les longueurs qui précède le lemme). Alors $f_*(y'' \to y)$ est
un objet de $\mathcal{S}_x$ qui s'envoie sur $x' \to x$ (par fonctorialité de
$f_*$), donc est isomorphe à $x' \to x$ par minimalité de $x' \to x$.
\end{proof}

\begin{lemma}
\label{lemma-smallest-where-descends-versal}
Soit $\mathcal{F}$ une catégorie cofibrée en groupoïdes au-dessus de
$\mathcal{C}_\Lambda$ qui satisfait (S1). Soit $\xi$ un objet formel versel de
$\mathcal{F}$ situé au-dessus de $R$. Il existe un morphisme $\xi' \to \xi$
situé au-dessus de $R' \subset R$ ayant les propriétés de minimalité
suivantes :
\begin{enumerate}
\item pour tout $f : R \to A$, avec $A \in \Ob(\mathcal{C}_\Lambda)$, les
images directes
$$
\vcenter{
\xymatrix{
\xi' \ar[d] \ar[r] & x' \ar[d] \\
\xi \ar[r] & x
}
}
\quad\text{au-dessus de}\quad
\vcenter{
\xymatrix{
R' \ar[d] \ar[r] & f(R') \ar[d] \\
R \ar[r] & A
}
}
$$
produisent un objet minimal $x' \to x$ de $\mathcal{S}_x$ ;
\item pour tout morphisme d'objets formels $\xi'' \to \xi'$, le morphisme
correspondant $R'' \to R'$ est surjectif.
\end{enumerate}
\end{lemma}

\begin{proof}
Écrivons $\xi = (R, \xi_n, f_n)$. Posons $R'_1 = k$ et
$\xi'_1 = \xi_1$. Supposons construits des objets minimaux
$\xi'_m \to \xi_m$ de $\mathcal{S}_{\xi_m}$ situés au-dessus de
$R'_m \subset R/\mathfrak m_R^m$ pour $m \leq n$, ainsi que des morphismes
$f'_m : \xi'_{m + 1} \to \xi'_m$ compatibles avec $f_m$
pour $m \leq n - 1$. D'après
le lemme \ref{lemma-smallest-where-descends} (2), il existe un objet minimal
$\xi'_{n + 1} \to \xi_{n + 1}$ situé au-dessus de
$R'_{n + 1} \subset R/\mathfrak m_R^{n + 1}$ dont l'image est
$\xi'_n \to \xi_n$ au-dessus de $R'_n \subset R/\mathfrak m_R^n$. Cela donne
le diagramme commutatif
$$
\xymatrix{
\xi'_{n + 1} \ar[r]_{f'_n} \ar[d] & \xi'_n \ar[d] \\
\xi_{n + 1} \ar[r]^{f_n} & \xi_n
}
$$
par construction. De plus, le morphisme d'anneaux
$R'_{n + 1} \to R'_n$ est surjectif. Posons $R' = \lim_n R'_n$. Alors
$R' \to R$ est injectif.

\medskip\noindent
Toutefois, il n'est pas clair a priori que $R'$ soit noethérien. Pour le
montrer, on utilise la versalité de $\xi$. Celle-ci implique l'existence d'un
morphisme $\xi \to \xi'_n$ dans $\widehat{\mathcal{F}}$ ; voir le lemme
\ref{lemma-versal-object-quasi-initial}. Le morphisme correspondant
$R \to R'_n$ est nécessairement surjective (puisque
$\xi'_n \to \xi_n$ est minimal dans $\mathcal{S}_{\xi_n}$). Les dimensions
des espaces cotangents sont donc bornées, et le lemme
\ref{lemma-limit-artinian} implique que $R'$ est noethérien, c'est-à-dire un
objet de $\widehat{\mathcal{C}}_\Lambda$. D'après le lemme
\ref{lemma-formal-objects-different-filtration} (et le résultat sur les
filtrations du lemme \ref{lemma-limit-artinian}), la suite d'éléments
$\xi'_n$ définit un objet formel $\xi'$ au-dessus de $R'$, et l'on dispose
d'un morphisme $\xi' \to \xi$.

\medskip\noindent
Par construction, (1) est vérifiée pour $R \to R/\mathfrak m_R^n$, pour tout
$n$. Comme tout morphisme $R \to A$ de (1) se factorise par
$R \to R/\mathfrak m_R^n \to A$, la propriété (1) pour $x' \to x$ au-dessus
de $f(R) \subset A$ résulte, par le lemme
\ref{lemma-smallest-where-descends} (1), de la minimalité de
$\xi'_n \to \xi_n$ au-dessus de $R'_n \to R/\mathfrak m_R^n$.

\medskip\noindent
Si $R'' \to R'$ comme dans (2) n'était pas surjectif, alors
$R'' \to R' \to R'_n$ ne serait pas surjectif pour un certain $n$, et
$\xi'_n \to \xi_n$ ne serait pas minimal, contradiction. Cela démontre (2).
\end{proof}

\begin{lemma}
\label{lemma-descends-versal}
Soit $\mathcal{F}$ une catégorie cofibrée en groupoïdes au-dessus de
$\mathcal{C}_\Lambda$ qui satisfait (S1). Soit $\xi$ un objet formel versel de
$\mathcal{F}$ situé au-dessus de $R$. Soit $\xi' \to \xi$ un morphisme
d'objets formels situé au-dessus de $R' \subset R$, construit comme dans le
lemme \ref{lemma-smallest-where-descends-versal}. Alors
$$
R \cong R'[[x_1, \ldots, x_r]]
$$
est un anneau de séries formelles sur $R'$. De plus, $\xi'$ est lui aussi un
objet formel versel.
\end{lemma}

\begin{proof}
D'après le lemme \ref{lemma-versal-object-quasi-initial}, il existe un
morphisme $\xi \to \xi'$. Selon le lemme
\ref{lemma-smallest-where-descends-versal}, le morphisme correspondant
$f : R \to R'$ induit une surjection $f|_{R'} : R' \to R'$. C'est un
isomorphisme d'après Algèbre, lemme
\ref{algebra-lemma-surjective-endo-noetherian-ring-is-iso}. Ainsi,
$I = \Ker(f)$ est un idéal de $R$ tel que $R = R' \oplus I$. Soient
$x_1, \ldots, x_n \in I$ des éléments qui forment une base de
$I/\mathfrak m_RI$. Considérons le morphisme
$S = R'[[X_1, \ldots, X_r]] \to R$ qui envoie $X_i$ sur $x_i$. Pour tout
$n \geq 1$, on obtient une surjection de $R'$-algèbres artiniennes
$B = S/\mathfrak m_S^n \to R/\mathfrak m_R^n = A$. Notons
$y \in \Ob(\mathcal{F}(B)$, respectivement $x \in \Ob(\mathcal{F}(A))$,
l'image directe de $\xi'$ le long de $R' \to S \to B$, respectivement de
$R' \to S \to A$. Remarquons que $x$ est aussi l'image directe de $\xi$ le
long de $R \to A$, puisque $\xi$ est l'image directe de $\xi'$ le long de
$R' \to R$. On dispose donc d'un diagramme en traits pleins
$$
\vcenter{
\xymatrix{
& y \ar[d] \\
\xi \ar[r] \ar@{..>}[ru] & x
}
}
\quad\text{au-dessus de}\quad
\vcenter{
\xymatrix{
& S/\mathfrak m_S^n \ar[d] \\
R \ar[r] \ar@{..>}[ru] & R/\mathfrak m_R^n
}
}
$$
Comme $\xi$ est versel, la remarque \ref{remark-versal-object} fournit les
flèches en pointillés qui complètent ces diagrammes. En particulier, les
morphismes $S/\mathfrak m_S^n \to R/\mathfrak m_R^n$ admettent des sections
$h_n : R/\mathfrak m_R^n \to S/\mathfrak m_S^n$. Il résulte du lemme
\ref{lemma-power-series} que $S \to R$ est un isomorphisme.

\medskip\noindent
Puisque $\xi$ est une image directe de $\xi'$ le long de $R' \to R$, la
remarque \ref{remark-formal-objects-yoneda-map} fournit un diagramme
commutatif
$$
\xymatrix{
\underline{R}|_{\mathcal{C}_\Lambda} \ar[rr] \ar[rd]_{\underline{\xi}} & &
\underline{R'}|_{\mathcal{C}_\Lambda} \ar[ld]^{\underline{\xi'}} \\
& \mathcal{F}
}
$$
Comme $R' \to R$ admet un inverse à gauche (à savoir $R \to R/I = R'$),
$\underline{R}|_{\mathcal{C}_\Lambda} \to
\underline{R'}|_{\mathcal{C}_\Lambda}$ est essentiellement surjectif. Le
lemme \ref{lemma-smooth-properties} montre donc que $\underline{\xi'}$ est
lisse, c'est-à-dire que $\xi'$ est un objet formel versel.
\end{proof}

\noindent
Les lemmes précédents motivent la définition suivante.

\begin{definition}
\label{definition-minimal-versal}
Soit $\mathcal{F}$ une catégorie de prédéformation. On dit qu'un objet formel
versel $\xi$ de $\mathcal{F}$ est {\it minimal}\footnote{Cette terminologie
n'est peut-être pas standard. Beaucoup d'auteurs relient cette notion à des
propriétés des espaces tangents. Nous établirons ce lien dans la section
\ref{section-miniversal-objects-existence}.} si, pour tout morphisme d'objets
formels $\xi' \to \xi$, le morphisme sous-jacent sur les anneaux est
surjective. Un objet formel versel minimal est parfois appelé {\it miniversel}.
\end{definition}

\noindent
Les résultats de cette section montrent que cette définition est raisonnable.
Tout d'abord, l'existence d'un objet formel versel implique que
$\mathcal{F}$ satisfait (S1). Les lemmes précédents montrent ensuite qu'il
existe un objet formel versel minimal. Enfin, deux objets formels versels
minimaux quelconques sont isomorphes. Voici un résumé de nos résultats (avec
des démonstrations détaillées).

\begin{lemma}
\label{lemma-minimal-versal}
Soit $\mathcal{F}$ une catégorie de prédéformation qui possède un objet
formel versel. Alors :
\begin{enumerate}
\item $\mathcal{F}$ possède un objet formel versel minimal ;
\item les objets versels minimaux sont uniques à isomorphisme près ;
\item tout objet versel est l'image directe d'un objet versel minimal le long
d'une extension en séries formelles.
\end{enumerate}
\end{lemma}

\begin{proof}
Supposons que $\mathcal{F}$ possède un objet formel versel $\xi$ au-dessus de
$R$. Elle satisfait alors (S1) ; voir le lemme
\ref{lemma-versal-object-S1}. Soit $\xi' \to \xi$ au-dessus de
$R' \subset R$ l'un des morphismes construits dans le lemme
\ref{lemma-smallest-where-descends-versal}. D'après le lemme
\ref{lemma-descends-versal}, $\xi'$ est versel ; il est donc, par
construction, un objet formel versel minimal. Cela démontre (1). De plus,
$R \cong R'[[x_1, \ldots, x_n]]$, ce qui démontre (3).

\medskip\noindent
Supposons que $\xi_i/R_i$ soient deux objets formels versels minimaux. Le
lemme \ref{lemma-versal-object-quasi-initial} fournit des morphismes
$\xi_1 \to \xi_2$ et $\xi_2 \to \xi_1$. Les morphismes d'anneaux
correspondants $f : R_1 \to R_2$ et $g : R_2 \to R_1$ sont surjectifs par
minimalité. Les composées $g \circ f : R_1 \to R_1$ et
$f \circ g : R_2 \to R_2$ sont donc des isomorphismes d'après Algèbre, lemme
\ref{algebra-lemma-surjective-endo-noetherian-ring-is-iso}. Ainsi, $f$ et $g$
sont des isomorphismes, et par conséquent les morphismes
$\xi_1 \to \xi_2$ et $\xi_2 \to \xi_1$ sont des isomorphismes (car
$\widehat{\mathcal{F}}$ est cofibrée en groupoïdes d'après le lemme
\ref{lemma-completion-cofibred}). Cela démontre (2) et achève la démonstration
du lemme.
\end{proof}








\section{Objets formels miniversels et espaces tangents}
\label{section-miniversal-objects-existence}

\noindent
La notion générale de minimalité introduite dans la définition
\ref{definition-minimal-versal} peut parfois se déduire du comportement sur
les espaces tangents. Soit $\xi$ un objet formel de la catégorie de
prédéformation $\mathcal{F}$, et soit
$\underline{\xi} : \underline{R}|_{\mathcal{C}_\Lambda} \to \mathcal{F}$
le morphisme correspondant. On peut alors considérer la condition
\begin{equation}
\label{equation-bijective}
d\underline{\xi} : \text{Der}_\Lambda(R, k) \to T\mathcal{F} 
\text{ is bijective}
\end{equation}
et la condition
\begin{equation}
\label{equation-bijective-orbits}
d\underline{\xi} : \text{Der}_\Lambda(R, k) \to T\mathcal{F} 
\text{ is bijective on }\text{Der}_\Lambda(k, k)\text{-orbits.}
\end{equation}
Nous utilisons ici l'identification
$T\underline{R}|_{\mathcal{C}_\Lambda} = \text{Der}_\Lambda(R, k)$ de
l'exemple \ref{example-tangent-space-prorepresentable-functor}, ainsi que
l'action (\ref{equation-action}) des dérivations sur les espaces tangents. Si
$k' \subset k$ est séparable, alors $\text{Der}_\Lambda(k, k) = 0$ et les deux
conditions sont équivalentes. Il s'avère qu'en présence de la condition (S2),
un objet formel versel est minimal si et seulement si $\underline{\xi}$
satisfait (\ref{equation-bijective-orbits}). De plus, si $\underline{\xi}$
satisfait (\ref{equation-bijective}), alors $\mathcal{F}$ satisfait (S2).

\begin{lemma}
\label{lemma-miniversal-object-unique}
Soit $\mathcal{F}$ une catégorie de prédéformation. Soit $\xi$ un objet
formel versel de $\mathcal{F}$ satisfaisant
(\ref{equation-bijective-orbits}). Alors $\xi$ est un objet formel versel
minimal. En particulier, de tels $\xi$ sont uniques à isomorphisme près.
\end{lemma}

\begin{proof}
Si $\xi$ n'est pas minimal, il existe un morphisme $\xi' \to \xi$ situé
au-dessus de $R' \to R$ tel que $R = R'[[x_1, \ldots, x_n]]$ avec $n > 0$ ;
voir le lemme \ref{lemma-minimal-versal}. Ainsi, $d\underline{\xi}$ se
factorise en
$$
\text{Der}_\Lambda(R, k) \to
\text{Der}_\Lambda(R', k) \to T\mathcal{F}
$$
et l'on voit que (\ref{equation-bijective-orbits}) ne peut être satisfaite,
car $D : f \mapsto \partial/\partial x_1(f) \bmod \mathfrak m_R$ est un
élément du noyau de la première flèche qui n'appartient pas à l'image de
$\text{Der}_\Lambda(k, k) \to \text{Der}_\Lambda(R, k)$.
\end{proof}

\begin{lemma}
\label{lemma-miniversal-object-existence-1}
Soit $\mathcal{F}$ une catégorie de prédéformation. Soit $\xi$ un objet
formel versel de $\mathcal{F}$ satisfaisant (\ref{equation-bijective}). Alors :
\begin{enumerate}
\item $\mathcal{F}$ satisfait (S1).
\item $\mathcal{F}$ satisfait (S2).
\item $\dim_k T\mathcal{F}$ est finie.
\end{enumerate}
\end{lemma}

\begin{proof}
La condition (S1) résulte du lemme \ref{lemma-versal-object-S1}. La première
partie de (S2) est satisfaite puisque (S1) l'est. Soient
$$
\vcenter{
\xymatrix{
y \ar[r]_c \ar[d]_a & x_\epsilon \ar[d]^e \\
x \ar[r]^d          & x_0
}
}
\quad\text{et}\quad
\vcenter{
\xymatrix{
y' \ar[r]_{c'} \ar[d]_{a'} & x_\epsilon \ar[d]^e \\
x \ar[r]^d                 & x_0
}
}
\quad\text{au-dessus de}\quad
\vcenter{
\xymatrix{
A \times_k k[\epsilon] \ar[r] \ar[d] & k[\epsilon] \ar[d] \\
A  \ar[r] & k
}
}
$$
des diagrammes comme dans la seconde partie de (S2). Comme ci-dessus, on peut
trouver des morphismes $b : \xi \to y$ et $b' : \xi \to y'$ tels que
$$
\xymatrix{
\xi \ar[r]^{b'} \ar[d]_b          & y' \ar[d]^{a'} \\
y \ar[r]^{a} & x
}
$$
soit commutatif. Notons $p : \mathcal{F} \to \mathcal{C}_\Lambda$ le
morphisme structural. Écrivons $\widehat{p}(\xi) = R$, c'est-à-dire que $\xi$
est situé au-dessus de $R \in \Ob(\widehat{\mathcal{C}}_\Lambda)$. L'image
directe de $\xi$ par $p(c) \circ p(b)$ est $x_\epsilon$, et l'image directe de
$\xi$ par $p(c') \circ p(b')$ est $x_\epsilon$. Puisque $\xi$ satisfait
(\ref{equation-bijective}), on a
$p(c) \circ p(b) = p(c') \circ p(b')$ comme morphismes
$R \to k[\epsilon]$. Par conséquent, $p(b) = p(b')$ comme morphismes
$R \to A \times_k k[\epsilon]$. Ainsi, $y$ et $y'$ sont isomorphes à l'image
directe de $\xi$ le long de ce morphisme, et l'on obtient l'unique
morphisme recherché $y \to y'$ au-dessus de $A \times_k k[\epsilon]$,
compatible avec $b$ et $b'$.

\medskip\noindent
Enfin, l'exemple \ref{example-tangent-space-prorepresentable-functor} montre
que $\dim_k T\mathcal{F} = \dim_k T\underline{R}|_{\mathcal{C}_\Lambda}$
est finie.
\end{proof}

\begin{example}
\label{example-do-not-get-S2}
Il existe des catégories de prédéformation qui possèdent un objet formel
versel satisfaisant (\ref{equation-bijective-orbits}) mais ne satisfont pas
(S2). Un exemple immédiat consiste à prendre
$F = \underline{k[\epsilon]}/G$, où
$G \subset \text{Aut}_{\mathcal{C}_\Lambda}(k[\epsilon])$ est un sous-groupe
fini non trivial. En effet, le morphisme $\underline{k[\epsilon]} \to F$ est
lisse, mais l'espace tangent de $F$ ne possède pas de structure naturelle de
$k$-espace vectoriel (puisqu'il est le quotient d'un $k$-espace vectoriel par
un groupe fini).
\end{example}

\begin{lemma}
\label{lemma-construct-bijective-orbits}
Soit $\mathcal{F}$ une catégorie de prédéformation satisfaisant (S2) et
possédant un objet formel versel. Alors son objet formel versel minimal
satisfait (\ref{equation-bijective-orbits}).
\end{lemma}

\begin{proof}
Soit $\xi$ un objet formel versel minimal pour $\mathcal{F}$ ; voir le lemme
\ref{lemma-minimal-versal}. Supposons $\xi$ situé au-dessus de
$R \in \Ob(\widehat{\mathcal{C}}_\Lambda)$. Pour interpréter
(\ref{equation-bijective-orbits}), remarquons que $T\mathcal{F}$ possède une
structure naturelle de $k$-espace vectoriel (voir le lemme
\ref{lemma-tangent-space-vector-space}), que
$d\underline{\xi} : \text{Der}_\Lambda(R, k) \to T\mathcal{F}$ est linéaire
(voir le lemme \ref{lemma-k-linear-differential}), et que l'action de
$\text{Der}_\Lambda(k, k)$ est donnée par l'addition (voir le lemme
\ref{lemma-action-linear}). Considérons le diagramme
$$
\xymatrix{
& \Hom_k(\mathfrak m_R/\mathfrak m_R^2, k) \\
K \ar[r] & \text{Der}_\Lambda(R, k) \ar[r]^{d\underline{\xi}} \ar[u] &
T\mathcal{F} \\
& \text{Der}_\Lambda(k, k) \ar[u] \ar[ru]
}
$$
L'espace vectoriel $K$ est le noyau de $d\underline{\xi}$. Remarquons que la
colonne médiane est exacte en son milieu, puisqu'elle est duale de la suite
(\ref{equation-sequence}). Si (\ref{equation-bijective-orbits}) n'est pas
satisfaite, on peut trouver un élément non nul $D \in K$ dont l'image dans
$\Hom_k(\mathfrak m_R/\mathfrak m_R^2, k)$ n'est pas nulle. Il existe donc
$t \in \mathfrak m_R$ tel que $D(t) = 1$. Posons
$R' = \{a \in R \mid D(a) = 0\}$. Comme $D$ est une dérivation, c'est un
sous-anneau de $R$. Puisque $D(t) = 1$, le morphisme $R' \to k$ est
surjective (comparer avec la démonstration du lemme
\ref{lemma-essential-surjection}). Remarquons que
$\mathfrak m_{R'} = \Ker(D : \mathfrak m_R \to k)$ est un idéal de $R$ et
que $\mathfrak m_R^2 \subset \mathfrak m_{R'}$. Ainsi,
$$
\mathfrak m_R/\mathfrak m_R^2 =
\mathfrak m_{R'}/\mathfrak m_R^2 + k\overline{t}
$$
ce qui implique que le morphisme
$$
R'/\mathfrak m_R^2 \times_k k[\epsilon] \to R/\mathfrak m_R^2
$$
qui envoie $\epsilon$ sur $\overline{t}$ est un isomorphisme. En particulier,
il existe un morphisme $R/\mathfrak m_R^2 \to R'/\mathfrak m_R^2$.

\medskip\noindent
Soit $\xi \to y$ un morphisme situé au-dessus de
$R \to R/\mathfrak m_R^2$. Soit $y \to x$ un morphisme situé au-dessus de
$R/\mathfrak m_R^2 \to R'/\mathfrak m_R^2$. Soit $y \to x_\epsilon$ un
morphisme situé au-dessus de $R/\mathfrak m_R^2 \to k[\epsilon]$. Soit $x_0$
l'unique objet de $\mathcal{F}$ au-dessus de $k$, à isomorphisme unique près.
Les axiomes d'une catégorie cofibrée en groupoïdes donnent un diagramme
commutatif
$$
\vcenter{
\xymatrix{
y \ar[r] \ar[d] & x_\epsilon \ar[d] \\
x \ar[r]   & x_0
}
}
\quad\text{au-dessus de}\quad
\vcenter{
\xymatrix{
R'/\mathfrak m_R^2 \times_k k[\epsilon] \ar[r] \ar[d] & k[\epsilon] \ar[d] \\
R'/\mathfrak m_R^2 \ar[r] & k.
}
}
$$
Comme $D \in K$, on voit que $x_\epsilon$ est isomorphe à
$0 \in \mathcal{F}(k[\epsilon])$, c'est-à-dire que $x_\epsilon$ est l'image
directe de $x_0$ par $k \to k[\epsilon], a \mapsto a$. Le lemme
\ref{lemma-lifting-section} montre donc qu'il existe un morphisme $x \to y$.
Comme
$\text{length}_\Lambda(R'/\mathfrak m_R^2) <
\text{length}_\Lambda(R/\mathfrak m_R^2)$
le morphisme d'anneaux correspondant
$R'/\mathfrak m_R^2 \to R/\mathfrak m_R^2$ n'est pas surjectif. Cela
contredit la minimalité de $\xi/R$ ; voir la partie (1) du lemme
\ref{lemma-smallest-where-descends-versal}. Cette contradiction montre qu'un
tel $D$ ne peut exister, d'où le résultat.
\end{proof}

\begin{theorem}
\label{theorem-miniversal-object-existence}
Soit $\mathcal{F}$ une catégorie de prédéformation. Considérons les conditions
suivantes :
\begin{enumerate}
\item $\mathcal{F}$ possède un objet formel versel minimal satisfaisant
(\ref{equation-bijective}) ;
\item $\mathcal{F}$ possède un objet formel versel minimal satisfaisant
(\ref{equation-bijective-orbits}) ;
\item les conditions suivantes sont satisfaites :
\begin{enumerate}
\item $\mathcal{F}$ satisfait (S1).
\item $\mathcal{F}$ satisfait (S2).
\item $\dim_k T\mathcal{F}$ est finie.
\end{enumerate}
\end{enumerate}
On a toujours
$$
(1) \Rightarrow (3) \Rightarrow (2).
$$
Si $k' \subset k$ est séparable, les trois conditions sont équivalentes.
\end{theorem}

\begin{proof}
Le lemme \ref{lemma-miniversal-object-existence-1} montre que
(1) $\Rightarrow$ (3). Les lemmes \ref{lemma-versal-object-existence} et
\ref{lemma-construct-bijective-orbits} montrent que
(3) $\Rightarrow$ (2). Si $k' \subset k$ est séparable, alors
$\text{Der}_\Lambda(k, k) = 0$, et l'on voit que
(\ref{equation-bijective}) $=$ (\ref{equation-bijective-orbits}),
c'est-à-dire que (1) et (2) sont la même condition.

\medskip\noindent
Une autre démonstration de (3) $\Rightarrow$ (1) dans le cas classique
consiste à compléter légèrement la démonstration du lemme
\ref{lemma-versal-object-existence} pour voir que l'on peut construire
directement, dans ce cas, un objet versel qui satisfait
(\ref{equation-bijective}). Cela évite d'utiliser le lemme
\ref{lemma-versal-object-existence} dans le cas classique. Nous omettons les
détails.
\end{proof}

\begin{remark}
\label{remark-compare-schlessinger-H3}
Soit $F : \mathcal{C}_\Lambda \to \textit{Ensembles}$ un foncteur de
prédéformation satisfaisant (S1), (S2) et $\dim_k TF < \infty$. Rappelons que
ces conditions correspondent aux conditions (H1), (H2) et (H3) de l'article
de Schlessinger ; voir la remarque
\ref{remark-compare-schlessinger-H3-pre}. Dans le cas classique (ou si
$k' \subset k$ est séparable), nous introduisons à la suite de Schlessinger
la notion d'enveloppe : une {\it enveloppe} est un objet formel versel
$\xi \in \widehat{F}(R)$ tel que
$d\underline{\xi} : T\underline{R}|_{\mathcal{C}_\Lambda} \to TF$
soit un isomorphisme, c'est-à-dire que (\ref{equation-bijective}) soit
satisfaite. Ainsi, le théorème \ref{theorem-miniversal-object-existence} nous
dit que
$$
(H1) + (H2) + (H3)
\Rightarrow
\text{ there exists a hull}
$$
dans le cas classique. Autrement dit, notre théorème retrouve le théorème de
Schlessinger sur l'existence des enveloppes.
\end{remark}

\begin{remark}
\label{remark-compose-minimal-into-iso-classes}
Soit $\mathcal{F}$ une catégorie de prédéformation. Rappelons que
$\mathcal{F} \to \overline{\mathcal{F}}$ est lisse ; voir la remarque
\ref{remark-smooth-to-iso-classes}. Par conséquent, si
$\xi \in \widehat{\mathcal{F}}(R)$ est un objet formel versel, la composée
$$
\underline{R}|_{\mathcal{C}_\Lambda} \longrightarrow
\mathcal{F} \longrightarrow \overline{\mathcal{F}}
$$
est lisse (lemme \ref{lemma-smooth-properties}), et l'on conclut que l'image
$\overline{\xi}$ de $\xi$ dans $\overline{\mathcal{F}}$ est un objet formel
versel. Si (\ref{equation-bijective}) est satisfaite, alors
$\overline{\xi}$ induit un isomorphisme
$T\underline{R}|_{\mathcal{C}_\Lambda} \to T\overline{\mathcal{F}}$
car $\mathcal{F} \to \overline{\mathcal{F}}$ identifie les espaces tangents.
Dans ce cas, $\overline{\xi}$ est donc une enveloppe de
$\overline{\mathcal{F}}$ ; voir la remarque
\ref{remark-compare-schlessinger-H3}. D'après le théorème
\ref{theorem-miniversal-object-existence}, on peut toujours trouver un tel
$\xi$ si $k' \subset k$ est séparable et si $\mathcal{F}$ est une catégorie
de prédéformation satisfaisant (S1), (S2) et
$\dim_k T\mathcal{F} < \infty$.
\end{remark}

\begin{example}
\label{example-smooth-continued}
Dans le lemme \ref{lemma-exists-smooth}, nous avons construit des objets
$R \in \widehat{\mathcal{C}}_\Lambda$ tels que
$\underline{R}|_{\mathcal{C}_\Lambda}$ soit lisse et que
$$
H_1(L_{k/\Lambda}) = \mathfrak m_R/\mathfrak m_R^2
\quad\text{et}\quad
\Omega_{R/\Lambda} \otimes_R k = \Omega_{k/\Lambda}
$$
Réinterprétons ceci à l'aide du théorème précédent. Considérons
$\mathcal{F} = \mathcal{C}_\Lambda$ comme une catégorie cofibrée en
groupoïdes au-dessus d'elle-même (au moyen du foncteur identité). Alors
$\mathcal{F}$ est une catégorie de prédéformation, satisfait (S1) et (S2), et
l'on a $T\mathcal{F} = 0$. Ainsi, $\mathcal{F}$ satisfait la condition (3) du
théorème \ref{theorem-miniversal-object-existence}. Le théorème implique que
(2) est satisfaite, c'est-à-dire que l'on peut trouver un objet formel versel
minimal $\xi \in \widehat{\mathcal{F}}(S)$, situé au-dessus d'un certain
$S \in \widehat{\mathcal{C}}_\Lambda$, et satisfaisant
(\ref{equation-bijective-orbits}). Le lemme \ref{lemma-smooth} montre que
$\Lambda \to S$ est formellement lisse pour la topologie
$\mathfrak m_S$-adique (car
$\underline{\xi} : \underline{R}|_{\mathcal{C}_\Lambda} \to
\mathcal{F} = \mathcal{C}_\Lambda$ est lisse). La condition
(\ref{equation-bijective-orbits}) dit maintenant que
$\text{Der}_\Lambda(S, k) \to 0$ est bijective sur les orbites de
$\text{Der}_\Lambda(k, k)$. Cela signifie que l'injection
$\text{Der}_\Lambda(k, k) \to \text{Der}_\Lambda(S, k)$ est également
surjective. Autrement dit,
$\Omega_{S/\Lambda} \otimes_S k = \Omega_{k/\Lambda}$. Puisque
$\Lambda \to S$ est formellement lisse pour la topologie
$\mathfrak m_S$-adique, on peut appliquer le lemme
\ref{more-algebra-lemma-formally-smooth-JZ} de Compléments d'algèbre ; la suite
exacte (\ref{equation-sequence-extended}) se traduit alors par les deux identifications
$$
H_1(L_{k/\Lambda}) = \mathfrak m_S/\mathfrak m_S^2
\quad\text{et}\quad
\Omega_{S/\Lambda} \otimes_S k = \Omega_{k/\Lambda}
$$
En remontant l'argument, on voit que le $R$ construit dans le lemme
\ref{lemma-exists-smooth} porte un objet versel minimal. Par l'unicité des
objets versels minimaux (lemme \ref{lemma-minimal-versal}), on conclut aussi
que $R \cong S$, c'est-à-dire que les deux constructions donnent le même
résultat.
\end{example}







\section{Conditions de Rim--Schlessinger et catégories de déformation}
\sectionmark{Conditions de Rim--Schlessinger}
\label{section-RS-condition}

\noindent
Il existe une propriété très naturelle des catégories cofibrées en groupoïdes
au-dessus de $\mathcal{C}_\Lambda$, qui est facile à vérifier en pratique et
qui entraîne les propriétés (S1) et (S2) de Schlessinger introduites plus haut.

\begin{definition}
\label{definition-RS}
Soit $\mathcal{F}$ une catégorie cofibrée en groupoïdes au-dessus de $\mathcal
C_\Lambda$. On dit que $\mathcal{F}$ satisfait à la {\it condition (RS)} si,
pour tout diagramme de $\mathcal{F}$
$$
\vcenter{
\xymatrix{
           & x_2 \ar[d] \\
x_1 \ar[r] & x
}
}
\quad\text{au-dessus de}\quad
\vcenter{
\xymatrix{
           & A_2 \ar[d] \\
A_1 \ar[r] & A
}
}
$$
dans $\mathcal{C}_\Lambda$, où $A_2 \to A$ est surjectif, il existe un produit
fibré $x_1 \times_x x_2$ dans $\mathcal{F}$ tel que le diagramme
$$
\vcenter{
\xymatrix{
x_1 \times_x x_2 \ar[r] \ar[d] & x_2 \ar[d] \\
x_1 \ar[r]      & x
}
}
\quad\text{lies over}\quad
\vcenter{
\xymatrix{
A_1 \times_A A_2 \ar[r] \ar[d] & A_2 \ar[d] \\
A_1 \ar[r]      & A.
}
}
$$
\end{definition}

\begin{lemma}
\label{lemma-RS-fiber-square}
Soit $\mathcal{F}$ une catégorie cofibrée en groupoïdes au-dessus de
$\mathcal{C}_\Lambda$ satisfaisant à (RS). Étant donné un diagramme commutatif
de $\mathcal{F}$
$$
\vcenter{
\xymatrix{
y \ar[r] \ar[d] & x_2 \ar[d]   \\
x_1 \ar[r]      & x
}
}
\quad\text{au-dessus de}\quad
\vcenter{
\xymatrix{
A_1 \times_A A_2 \ar[r] \ar[d] & A_2 \ar[d] \\
A_1 \ar[r]      & A.
}
}
$$
où $A_2 \to A$ est surjectif, ce diagramme est cartésien.
\end{lemma}

\begin{proof}
Puisque $\mathcal{F}$ satisfait à (RS), il existe un diagramme de produit fibré
$$
\vcenter{
\xymatrix{
x_1 \times_x x_2 \ar[r] \ar[d] & x_2 \ar[d] \\
x_1 \ar[r]      & x
}
}
\quad\text{au-dessus de}\quad
\vcenter{
\xymatrix{
A_1 \times_A A_2 \ar[r] \ar[d] & A_2 \ar[d] \\
A_1 \ar[r]      & A.
}
}
$$
Le morphisme induit $y \to x_1 \times_x x_2$ est situé au-dessus de
$\text{id} : A_1 \times_A A_1 \to A_1 \times_A A_1$ ; c'est donc un
isomorphisme.
\end{proof}

\begin{lemma}
\label{lemma-RS-small-extension}
Soit $\mathcal{F}$ une catégorie cofibrée en groupoïdes au-dessus de $\mathcal
C_\Lambda$. Alors $\mathcal{F}$ satisfait à (RS) s'il suffit de supposer que
la condition de la définition \ref{definition-RS} est vérifiée lorsque
$A_2 \to A$ est une petite extension.
\end{lemma}

\begin{proof}
Appliquer le lemme \ref{lemma-factor-small-extension}. La démonstration est
analogue à celle du lemme \ref{lemma-smoothness-small-extensions}.
\end{proof}

\begin{lemma}
\label{lemma-RS-2-categorical}
Soit $\mathcal{F}$ une catégorie cofibrée en groupoïdes au-dessus de
$\mathcal{C}_\Lambda$. Les assertions suivantes sont équivalentes :
\begin{enumerate}
\item $\mathcal{F}$ satisfait à (RS) ;
\item le foncteur
$\mathcal{F}(A_1 \times_A A_2) \to
\mathcal{F}(A_1) \times_{\mathcal{F}(A)} \mathcal{F}(A_2)$
de (\ref{equation-compare}) est une équivalence de catégories dès que
$A_2 \to A$ est surjectif ;
\item il en est de même qu'en (2) dès que $A_2 \to A$ est une petite extension.
\end{enumerate}
\end{lemma}

\begin{proof}
Supposons (1). Le lemme \ref{lemma-RS-fiber-square} montre que tout objet de
$\mathcal{F}(A_1 \times_A A_2)$ est de la forme $x_1 \times_x x_2$. De plus,
$$
\Mor_{A_1 \times_A A_2}(x_1 \times_x x_2, y_1 \times_y y_2) =
\Mor_{A_1}(x_1, y_1) \times_{\Mor_A(x, y)}
\Mor_{A_2}(x_2, y_2).
$$
Ainsi, $\mathcal{F}(A_1 \times_A A_2)$ est un $2$-produit fibré de
$\mathcal{F}(A_1)$ et $\mathcal{F}(A_2)$ au-dessus de $\mathcal{F}(A)$ d'après
Catégories, remarque \ref{categories-remark-other-description-2-fibre-product}.
Autrement dit, (2) est vérifiée.

\medskip\noindent
L'implication (2) $\Rightarrow$ (3) est immédiate.

\medskip\noindent
Supposons (3). Soient $q_1 : A_1 \to A$ et $q_2 : A_2 \to A$, avec $q_2$
une petite extension. Nous utiliserons la description du $2$-produit fibré
$\mathcal{F}(A_1) \times_{\mathcal{F}(A)} \mathcal{F}(A_2)$ donnée dans
Catégories, remarque \ref{categories-remark-other-description-2-fibre-product}.
Soit donc $y \in \mathcal{F}(A_1 \times_A A_2)$ correspondant à
$(x_1, x_2, x, a_1 : x_1 \to x, a_2 : x_2 \to x)$.
Soit $z$ un objet de $\mathcal{F}$ au-dessus de $C$. Alors
\begin{align*}
\Mor_\mathcal{F}(z, y) & =
\{(f, \alpha) \mid f : C \to A_1 \times_A A_2,
\alpha : f_*z \to y\} \\
& = \{(f_1, f_2, \alpha_1, \alpha_2) \mid
f_i : C \to A_i, \ \alpha_i : f_{i, *}z \to x_i, \\
& \quad\quad\quad\quad
q_1 \circ f_1 = q_2 \circ f_2, \ q_{1, *} \alpha_1 = q_{2, *}\alpha_2\} \\
& =
\Mor_\mathcal{F}(z, x_1) \times_{\Mor_\mathcal{F}(z, x)}
\Mor_\mathcal{F}(z, x_2)
\end{align*}
d'où $y$ est un produit fibré de $x_1$ et $x_2$ au-dessus de $x$. Ainsi,
$\mathcal{F}$ satisfait à (RS) lorsque $A_2 \to A$ est une petite extension.
Le lemme \ref{lemma-RS-small-extension} montre donc que (RS) est vérifiée.
\end{proof}

\begin{remark}
\label{remark-compare-schlessinger-H4}
Lorsque $\mathcal{F}$ est cofibrée en ensembles, la condition (RS) est
exactement la condition (H4) de l'article de Schlessinger
\cite[théorème 2.11]{Sch}. Plus précisément, pour un foncteur
$F: \mathcal{C}_\Lambda \to \textit{Ensembles}$, la condition (RS) affirme ceci :
si $A_1 \to A$ et $A_2 \to A$ sont des morphismes de
$\mathcal{C}_\Lambda$, avec $A_2 \to A$ surjectif, alors l'application induite
$F(A_1 \times_A A_2) \to F(A_1) \times_{F(A)} F(A_2)$ est bijectif.
\end{remark}

\begin{lemma}
\label{lemma-RS-implies-S1-S2}
Soit $\mathcal{F}$ une catégorie cofibrée en groupoïdes au-dessus de
$\mathcal{C}_\Lambda$. La condition (RS) pour $\mathcal{F}$ entraîne à la
fois (S1) et (S2) pour $\mathcal{F}$.
\end{lemma}

\begin{proof}
La reformulation du lemme \ref{lemma-RS-2-categorical} et l'explication de
(S1) qui suit la définition \ref{definition-S1-S2} montrent immédiatement que
(RS) entraîne (S1). Cela démontre la première partie de (S2). La seconde partie
de (S2) résulte du lemme \ref{lemma-RS-fiber-square}, qui donne
$y = x_1 \times_{d, x_0, e} x_2 = y'$ lorsque $y, y'$ sont comme dans la
seconde partie de la définition de (S2), dans la définition
\ref{definition-S1-S2}. (En fait, le morphisme $y \to y'$ est compatible à la
fois avec $a, a'$ et avec $c, c'$ !)
\end{proof}

\noindent
Le lemme suivant est l'analogue du lemme
\ref{lemma-S1-S2-associated-functor}. Rappelons que, si $\mathcal{F}$ est une
catégorie cofibrée en groupoïdes au-dessus de $\mathcal{C}_\Lambda$ et si $x$
est un objet de $\mathcal{F}$ au-dessus de $A$, on note
$\text{Aut}_A(x) = \Mor_A(x, x) = \Mor_{\mathcal{F}(A)}(x, x)$.
Si $x' \to x$ est un morphisme de $\mathcal{F}$ au-dessus de $A' \to A$,
il existe un homomorphisme de groupes bien défini
$\text{Aut}_{A'}(x') \to \text{Aut}_A(x)$.

\begin{lemma}
\label{lemma-RS-associated-functor}
Soit $\mathcal{F}$ une catégorie cofibrée en groupoïdes au-dessus de
$\mathcal{C}_\Lambda$ satisfaisant à (RS). Les conditions suivantes sont
équivalentes :
\begin{enumerate}
\item $\overline{\mathcal{F}}$ satisfait à (RS).
\item Soient $f_1: A_1 \to A$ et $f_2: A_2 \to A$ des morphismes
d'anneaux de $\mathcal{C}_\Lambda$, avec $f_2$ surjectif. L'application induite
entre ensembles de classes d'isomorphie
$$
\overline{\mathcal{F}(A_1) \times_{\mathcal{F}(A)} \mathcal{F}(A_2)}
\to \overline{\mathcal{F}}(A_1) \times_{\overline{\mathcal{F}}(A)}
\overline{\mathcal{F}}(A_2)
$$
est injective.
\item Pour tout morphisme $x' \to x$ de $\mathcal{F}$ au-dessus d'un
morphisme d'anneaux surjectif $A' \to A$, l'homomorphisme
$\text{Aut}_{A'}(x') \to \text{Aut}_A(x)$ est surjective.
\item Pour tout morphisme $x' \to x$ de $\mathcal{F}$ au-dessus d'une petite
extension $A' \to A$, l'homomorphisme
$\text{Aut}_{A'}(x') \to \text{Aut}_A(x)$ est surjective.
\end{enumerate}
\end{lemma}

\begin{proof}
Nous montrons que (1) est équivalente à (2), puis que (2) est équivalente à
(3). L'équivalence de (3) et (4) résulte du lemme
\ref{lemma-factor-small-extension}.

\medskip\noindent
Soient $f_1: A_1 \to A$ et $f_2: A_2 \to A$ des morphismes d'anneaux de
$\mathcal{C}_\Lambda$, avec $f_2$ surjectif. D'après la remarque
\ref{remark-compare-schlessinger-H4}, $\overline{\mathcal{F}}$ satisfait à
(RS) si et seulement si l'application
$$
\overline{\mathcal{F}}(A_1 \times_A A_2) \to \overline{\mathcal
F}(A_1) \times_{\overline{\mathcal{F}}(A)} \overline{\mathcal{F}}(A_2)
$$
est bijective pour tous $f_1, f_2$ de cette sorte. Cette application est au
moins surjective, car telle est la condition (S1), et
$\overline{\mathcal{F}}$ satisfait à (S1) d'après les lemmes
\ref{lemma-RS-implies-S1-S2} et \ref{lemma-S1-S2-associated-functor}. De plus,
elle se factorise sous la forme
$$
\overline{\mathcal{F}}(A_1 \times_A A_2)
\longrightarrow
\overline{\mathcal{F}(A_1) \times_{\mathcal{F}(A)} \mathcal{F}(A_2)}
\longrightarrow
\overline{\mathcal{F}}(A_1) \times_{\overline{\mathcal{F}}(A)}
\overline{\mathcal{F}}(A_2),
$$
où la première flèche est une bijection puisque
$$
\mathcal{F}(A_1 \times_A A_2)
\longrightarrow
\mathcal{F}(A_1) \times_{\mathcal{F}(A)} \mathcal{F}(A_2)
$$
est une équivalence par (RS) pour $\mathcal{F}$. Ainsi, (1) est équivalente à
(2).

\medskip\noindent
Supposons (2). Soit $x' \to x$ un morphisme de $\mathcal{F}$ au-dessus d'un
morphisme d'anneaux surjectif $f: A' \to A$. Soit
$a \in \text{Aut}_A(x)$. Les objets
$$
(x', x', a : x \to x), \ (x', x', \text{id} : x \to x)
$$
de
$\mathcal{F}(A') \times_{\mathcal{F}(A)} \mathcal{F}(A')$
ont la même image dans
$\overline{\mathcal{F}}(A') \times_{\overline{\mathcal{F}}(A)}
\overline{\mathcal{F}}(A')$. D'après (2), il existe des morphismes
$b_1, b_2 : x' \to x'$ tels que
$$
\xymatrix{
x \ar[r]_a \ar[d]_{f_*b_1} & x \ar[d]^{f_*b_2} \\
x \ar[r]^{\text{id}} & x
}
$$
soit commutatif. Ainsi, l'image de
$b_2^{-1} \circ b_1 \in \text{Aut}_{A'}(x')$ est
$a \in \text{Aut}_A(x)$. L'assertion (3) est donc vérifiée.

\medskip\noindent
Supposons (3). Supposons que
$$
(x_1, x_2, a : (f_1)_*x_1 \to (f_2)_*x_2),
\ (x'_1, x'_2, a' : (f_1)_*x'_1 \to (f_2)_*x'_2)
$$
soient des objets de
$\mathcal{F}(A_1) \times_{\mathcal{F}(A)} \mathcal{F}(A_2)$
ayant la même image dans
$\overline{\mathcal{F}}(A_1) \times_{\overline{\mathcal{F}}(A)}
\overline{\mathcal{F}}(A_2)$. Il existe alors des morphismes $b_1: x_1 \to
x'_1$ dans $\mathcal{F}(A_1)$ et $b_2: x_2 \to x'_2$ dans $\mathcal
F(A_2)$. D'après (3), on peut modifier $b_2$ par un automorphisme de $x_2$
au-dessus de $A_2$ de telle sorte que le diagramme
$$
\xymatrix{
(f_1)_*x_1 \ar[r]_a \ar[d]_{(f_1)_*b_1} & (f_2)_*x_2 \ar[d]^{(f_2)_*b_2} \\
(f_1)_*x'_1 \ar[r]^{a'} & (f_2)_*x'_2.
}
$$
soit commutatif. Cela démontre $(x_1, x_2, a) \cong (x'_1, x'_2, a')$ dans
$\overline{\mathcal{F}(A_1) \times_{\mathcal{F}(A)} \mathcal{F}(A_2)}$.
Ainsi, (2) est vérifiée.
\end{proof}

\noindent
Définissons enfin la notion de catégorie de déformation.

\begin{definition}
\label{definition-deformation-category}
Une {\it catégorie de déformation} est une catégorie de prédéformation
$\mathcal{F}$ satisfaisant à (RS). Un morphisme de catégories de déformation
est un morphisme de catégories au-dessus de $\mathcal{C}_\Lambda$.
\end{definition}

\begin{remark}
\label{remark-deformation-functor}
On dit qu'un foncteur $F: \mathcal{C}_\Lambda \to \textit{Ensembles}$ est un
{\it foncteur de déformation} si l'ensemble cofibré associé est une catégorie
de déformation, c'est-à-dire si $F(k)$ est un ensemble à un seul élément et si $F$
satisfait à (RS). Si $\mathcal{F}$ est une catégorie de déformation, alors
$\overline{\mathcal{F}}$ est un foncteur de prédéformation, mais pas
nécessairement un foncteur de déformation, comme le montre le lemme
\ref{lemma-RS-associated-functor}.
\end{remark}

\begin{example}
\label{example-prorepresentable-deformation-functor}
Un foncteur pro-représentable $F$ est un foncteur de déformation. En effet,
supposons $R \in \Ob(\widehat{\mathcal{C}}_\Lambda)$ et
$F(A) = \Mor_{\widehat{\mathcal{C}}_\Lambda}(R, A)$. Il existe un unique
morphisme $R \to k$, de sorte que $F(k)$ est un ensemble à un seul élément. Comme
$$
\Hom_\Lambda(R, A_1 \times_A A_2) =
\Hom_\Lambda(R, A_1) \times_{\Hom_\Lambda(R, A)}
\Hom_\Lambda(R, A_2)
$$
il en est de même pour les morphismes de $\widehat{\mathcal{C}}_\Lambda$, et
l'on voit que $F$ satisfait à (RS).
\end{example}

\noindent
Voici l'une de nos remarques typiques sur le passage d'une catégorie cofibrée
en groupoïdes à la catégorie de prédéformation en un point au-dessus de $k$ :
ce passage préserve (RS).

\begin{lemma}
\label{lemma-localize-RS}
Soit $\mathcal{F}$ une catégorie cofibrée en groupoïdes au-dessus de
$\mathcal{C}_\Lambda$. Soit $x_0 \in \Ob(\mathcal{F}(k))$. Soit
$\mathcal{F}_{x_0}$ la catégorie cofibrée en groupoïdes au-dessus de
$\mathcal{C}_\Lambda$ construite dans la remarque
\ref{remark-localize-cofibered-groupoid}. Si $\mathcal{F}$ satisfait à (RS),
alors $\mathcal{F}_{x_0}$ y satisfait aussi. En particulier,
$\mathcal{F}_{x_0}$ est une catégorie de déformation.
\end{lemma}

\begin{proof}
Tout diagramme comme dans la définition \ref{definition-RS} au sein de
$\mathcal{F}_{x_0}$ donne lieu à un diagramme dans $\mathcal{F}$, et le
résultat de (RS) pour ce diagramme dans $\mathcal{F}$ peut également être
considéré comme un résultat pour $\mathcal{F}_{x_0}$.
\end{proof}

\noindent
Le lemme suivant est l'analogue du fait que les $2$-produits fibrés de champs
algébriques sont des champs algébriques.

\begin{lemma}
\label{lemma-RS-fiber-product-morphisms}
Supposons que
$$
\xymatrix{
\mathcal{H} \times_\mathcal{F} \mathcal{G} \ar[r] \ar[d] &
\mathcal{G} \ar[d]^g \\
\mathcal{H} \ar[r]^f & \mathcal{F}
}
$$
soit un $2$-produit fibré de catégories cofibrées en groupoïdes au-dessus de
$\mathcal{C}_\Lambda$. Si $\mathcal{F}, \mathcal{G}, \mathcal{H}$ satisfont
toutes à (RS), alors $\mathcal{H} \times_\mathcal{F} \mathcal{G}$ satisfait à
(RS).
\end{lemma}

\begin{proof}
Si $A$ est un objet de $\mathcal{C}_\Lambda$, un objet de la catégorie fibre
de $\mathcal{H} \times_\mathcal{F} \mathcal{G}$ au-dessus de $A$ est un
triplet $(u, v, a)$, où $u \in \mathcal{H}(A)$, $v \in \mathcal{G}(A)$ et
$a : f(u) \to g(v)$ est un morphisme de $\mathcal{F}(A)$. Considérons un
diagramme de $\mathcal{H} \times_\mathcal{F} \mathcal{G}$
$$
\vcenter{
\xymatrix{
           & (u_2, v_2, a_2) \ar[d] \\
(u_1, v_1, a_1) \ar[r] & (u, v, a)
}
}
\quad\text{au-dessus de}\quad
\vcenter{
\xymatrix{
           & A_2 \ar[d] \\
A_1 \ar[r] & A
}
}
$$
dans $\mathcal{C}_\Lambda$, avec $A_2 \to A$ surjectif. Puisque
$\mathcal{H}$ et $\mathcal{G}$ satisfont à (RS), il existe des produits fibrés
$u_1 \times_u u_2$ et $v_1 \times_v v_2$ au-dessus de
$A_1 \times_A A_2$. Puisque $\mathcal{F}$ satisfait à (RS), le lemme
\ref{lemma-RS-fiber-square} montre que
$$
\vcenter{
\xymatrix{
f(u_1 \times_u u_2) \ar[r] \ar[d] & f(u_2) \ar[d] \\
f(u_1) \ar[r] & f(u)
}
}
\quad\text{et}\quad
\vcenter{
\xymatrix{
g(v_1 \times_v v_2) \ar[r] \ar[d] & g(v_2) \ar[d] \\
g(v_1) \ar[r] & g(v)
}
}
$$
sont tous deux des carrés cartésiens dans $\mathcal{F}$. On peut donc voir
$a_1 \times_a a_2$ comme un morphisme de $f(u_1 \times_u u_2)$ vers
$g(v_1 \times_v v_2)$ au-dessus de $A_1 \times_A A_2$. Il s'ensuit que
$$
\xymatrix{
 (u_1 \times_u u_2, v_1 \times_v v_2, a_{1} \times_a a_2) \ar[d] \ar[r] &
(u_2, v_2, a_2) \ar[d] \\
(u_1, v_1, a_1) \ar[r] & (u, v, a)
}
$$
est un carré cartésien dans $\mathcal{H} \times_\mathcal{F} \mathcal{G}$,
comme voulu.
\end{proof}


















\section{Relèvements d'objets}
\label{section-lifts}

\noindent
Le contenu de cette section est que, pour une catégorie de déformation,
l'ensemble des relèvements le long d'une petite extension est un espace
principal homogène sous l'espace tangent.

\begin{definition}
\label{definition-lifts}
Soit $\mathcal{F}$ une catégorie cofibrée en groupoïdes au-dessus de
$\mathcal{C}_\Lambda$. Soit $f: A' \to A$ un morphisme de
$\mathcal{C}_\Lambda$. Soit $x \in \mathcal{F}(A)$. La catégorie
$\textit{Lift}(x, f)$ des relèvements de $x$ le long de $f$ est la catégorie
dont les objets et les morphismes sont les suivants.
\begin{enumerate}
\item Objets : un {\it relèvement de $x$ le long de $f$} est un morphisme
$x' \to x$ au-dessus de $f$.
\item Morphismes : un {\it morphisme de relèvements} de
$a_1 : x'_1 \to x$ vers $a_2 : x'_2 \to x$ est un morphisme
$b : x'_1 \to x'_2$ dans $\mathcal{F}(A')$ tel que $a_2 = a_1 \circ b$.
\end{enumerate}
L'ensemble $\text{Lift}(x, f)$ des relèvements de $x$ le long de $f$ est
l'ensemble des classes d'isomorphie de $\textit{Lift}(x, f)$.
\end{definition}

\begin{remark}
\label{remark-omit-arrow}
Lorsque le morphisme $f: A' \to A$ ressort clairement du contexte, on peut
écrire $\textit{Lift}(x, A')$ et $\text{Lift}(x, A')$ au lieu de
$\textit{Lift}(x, f)$ et $\text{Lift}(x, f)$.
\end{remark}

\begin{remark}
\label{remark-tangent-space-lifting}
Soit $\mathcal{F}$ une catégorie cofibrée en groupoïdes au-dessus de
$\mathcal{C}_\Lambda$. Soit $x_0 \in \Ob(\mathcal{F}(k))$. Soit $V$ un espace
vectoriel de dimension finie. Alors $\text{Lift}(x_0, k[V])$ est l'ensemble
des classes d'isomorphie de $\mathcal{F}_{x_0}(k[V])$, où
$\mathcal{F}_{x_0}$ est la catégorie de prédéformation des objets de
$\mathcal{F}$ au-dessus de $x_0$ ; voir la remarque
\ref{remark-localize-cofibered-groupoid}. Par conséquent, si $\mathcal{F}$
satisfait à (S2), il en est de même de $\mathcal{F}_{x_0}$ (voir le lemme
\ref{lemma-S1-S2-localize}), et le lemme
\ref{lemma-tangent-space-vector-space} donne
$$
\text{Lift}(x_0, k[V]) = T\mathcal{F}_{x_0} \otimes_k V
$$
comme espaces vectoriels sur $k$.
\end{remark}

\begin{remark}
\label{remark-lift-bijections}
Soit $\mathcal{F}$ une catégorie cofibrée en groupoïdes au-dessus de $\mathcal
C_\Lambda$ satisfaisant à (RS). Soit
$$
\xymatrix{
A_1 \times_A A_2 \ar[r] \ar[d] & A_2 \ar[d] \\
A_1 \ar[r] & A
}
$$
un carré cartésien dans $\mathcal{C}_\Lambda$, tel que $A_1 \to A$ ou
$A_2 \to A$ soit surjectif. Soit $x \in \Ob(\mathcal{F}(A))$. Étant donnés
des relèvements $x_1 \to x$ et $x_2 \to x$ de $x$ à $A_1$ et $A_2$, (RS)
fournit un relèvement $x_1 \times_x x_2 \to x$ de $x$ à
$A_1 \times_A A_2$. Réciproquement, d'après le lemme
\ref{lemma-RS-fiber-square}, tout relèvement de $x$ à
$A_1 \times_A A_2$ est de cette forme. On obtient donc une bijection
$$
\text{Lift}(x, A_1) \times \text{Lift}(x, A_2)
\longrightarrow
\text{Lift}(x, A_1 \times_A A_2).
$$
De même, si $x_1 \to x$ est un relèvement fixé de $x$ à $A_1$, il existe une
bijection
$$
\text{Lift}(x_1, A_1 \times_A A_2)
\longrightarrow
\text{Lift}(x, A_2).
$$
Soit maintenant
$$
\xymatrix{
A_1' \times_A A_2 \ar[r] \ar[d] & A_1 \times_A A_2 \ar[r] \ar[d] & A_2
\ar[d] \\
A_1' \ar[r] & A_1 \ar[r] & A
}
$$
une composée de carrés cartésiens dans $\mathcal{C}_\Lambda$, où
$A'_1 \to A_1$ et $A_1 \to A$ sont tous deux surjectifs. Soit $x_1 \to x$ un
morphisme au-dessus de $A_1 \to A$. D'après ce qui précède, on a alors les
bijections
\begin{align*}
\text{Lift}(x_1, A_1' \times_A A_2)
& = \text{Lift}(x_1, A_1') \times \text{Lift}(x_1, A_1 \times_A A_2) \\
& = \text{Lift}(x_1, A_1') \times \text{Lift}(x, A_2).
\end{align*}
\end{remark}

\begin{lemma}
\label{lemma-free-transitive-action}
Soit $\mathcal{F}$ une catégorie de déformation. Soit $A' \to A$ un
morphisme d'anneaux surjectif de $\mathcal{C}_\Lambda$, dont le noyau $I$
est annulé par $\mathfrak m_{A'}$. Soit $x \in \Ob(\mathcal{F}(A))$. Si
$\text{Lift}(x, A')$ est non vide, alors $T\mathcal{F} \otimes_k I$ agit sur
$\text{Lift}(x, A')$ et en fait un espace principal homogène.
\end{lemma}

\begin{proof}
Considérons le morphisme d'anneaux $g : A' \times_A A' \to k[I]$ défini
par la règle $g(a_1, a_2) = \overline{a_1} \oplus a_2 - a_1$ (comparer au
lemme \ref{lemma-lifting-along-small-extension}). Il existe un isomorphisme
$$
A' \times_A A' \xrightarrow{\sim} A' \times_k k[I]
$$
donné par $(a_1, a_2) \mapsto (a_1, g(a_1, a_2))$. Cet isomorphisme est
compatible avec les projections sur le premier facteur $A'$, et donc avec celles de
$A' \times_A A'$ et $A' \times_k k[I]$ vers $A$. Il existe ainsi une bijection
\begin{equation}
\label{equation-one}
\text{Lift}(x, A' \times_A A')
\longrightarrow
\text{Lift}(x, A' \times_k k[I])
\end{equation}
D'après la remarque \ref{remark-lift-bijections}, il existe une bijection
\begin{equation}
\label{equation-two}
\text{Lift}(x, A') \times \text{Lift}(x, A')
\longrightarrow
\text{Lift}(x, A' \times_A A')
\end{equation}
On a un diagramme commutatif
$$
\xymatrix{
A' \times_k k[I] \ar[r] \ar[d] & A \times_k k[I] \ar[r] \ar[d] & k[I] \ar[d] \\
A' \ar[r] & A \ar[r] & k.
}
$$
Ainsi, si l'on choisit une image directe $x \to x_0$ de $x$ le long de
$A \to k$, la fin de la remarque \ref{remark-lift-bijections} fournit une
bijection
\begin{equation}
\label{equation-three}
\text{Lift}(x, A' \times_k k[I])
\longrightarrow
\text{Lift}(x, A') \times \text{Lift}(x_0, k[I])
\end{equation}
En composant (\ref{equation-two}), (\ref{equation-one}) et
(\ref{equation-three}), on obtient une bijection
$$
\Phi :
\text{Lift}(x, A') \times \text{Lift}(x, A')
\longrightarrow
\text{Lift}(x, A') \times \text{Lift}(x_0, k[I]).
$$
Cette bijection est compatible avec les projections sur les premiers facteurs. La remarque
\ref{remark-tangent-space-lifting} donne
$\text{Lift}(x_0, k[I]) = T\mathcal{F} \otimes_k I$. Si $\text{pr}_2$ est la
seconde projection de $\text{Lift}(x, A') \times \text{Lift}(x, A')$, on
obtient une application
$$
a = \text{pr}_2 \circ \Phi^{-1} :
\text{Lift}(x, A') \times (T\mathcal{F} \otimes_k I)
\longrightarrow
\text{Lift}(x, A').
$$
En explicitant ce qui précède, on voit que $a(x' \to x, \theta)$ est l'unique
relèvement $x'' \to x$ tel que $g_*(x', x'') = \theta$ dans
$\text{Lift}(x_0, k[I]) = T\mathcal{F} \otimes_k I$. Pour voir qu'il s'agit
d'une action de $T\mathcal{F} \otimes_k I$ sur $\text{Lift}(x, A')$, il faut
montrer ceci : si $x', x'', x'''$ sont des relèvements de $x$ et si
$g_*(x', x'') = \theta$, $g_*(x'', x''') = \theta'$, alors
$g_*(x', x''') = \theta + \theta'$. Cela résulte du diagramme commutatif
$$
\xymatrix{
A' \times_A A' \times_A A'
\ar[rrrrr]_-{(a_1, a_2, a_3) \mapsto (g(a_1, a_2), g(a_2, a_3))}
\ar[rrrrrd]_{(a_1, a_2, a_3) \mapsto g(a_1, a_3)} & & & & &
k[I] \times_k k[I] = k[I \times I] \ar[d]^{+} \\
& & & & & k[I]
}
$$
L'action est libre et transitive puisque $\Phi$ est bijective.
\end{proof}

\begin{remark}
\label{remark-free-transitive-action-functorial}
L'action du lemme \ref{lemma-free-transitive-action} est fonctorielle. Soit
$\varphi : \mathcal{F} \to \mathcal{G}$ un morphisme de catégories de
déformation. Soit $A' \to A$ un morphisme d'anneaux surjectif dont le
noyau $I$ est annulé par $\mathfrak m_{A'}$. Soit
$x \in \Ob(\mathcal{F}(A))$. Dans cette situation, $\varphi$ induit les
flèches verticales du diagramme commutatif suivant
$$
\xymatrix{
\text{Lift}(x, A') \times (T\mathcal{F} \otimes_k I)
\ar[d]_{(\varphi, d\varphi \otimes \text{id}_I)} \ar[r] &
\text{Lift}(x, A') \ar[d]^\varphi \\
\text{Lift}(\varphi(x), A') \times (T\mathcal{G} \otimes_k I) \ar[r] &
\text{Lift}(\varphi(x), A')
}
$$
La commutativité résulte du fait que chacune des applications
(\ref{equation-two}), (\ref{equation-one}) et (\ref{equation-three}) de la
démonstration du lemme \ref{lemma-free-transitive-action} donne lieu à un
diagramme commutatif analogue.
\end{remark}






\section{Théorème de Schlessinger sur les foncteurs pro-représentables}
\label{section-schlessingers-theorem}

\noindent
Nous en déduisons le théorème de Schlessinger qui caractérise les foncteurs
pro-représentables sur $\mathcal{C}_\Lambda$.

\begin{lemma}
\label{lemma-minimal-smooth-morphism-functors}
Soient $F, G: \mathcal{C}_\Lambda \to \textit{Ensembles}$ des foncteurs de
déformation. Soit $\varphi : F \to G$ un morphisme lisse qui induit un
isomorphisme $d\varphi : TF \to TG$ d'espaces tangents. Alors $\varphi$ est un
isomorphisme.
\end{lemma}

\begin{proof}
Nous montrons que $F(A) \to G(A)$ est une bijection pour tout $A \in
\Ob(\mathcal{C}_\Lambda)$, par récurrence sur
$\text{length}_A(A)$. Pour $A = k$, l'assertion résulte de
l'hypothèse que $F$ et $G$ sont des foncteurs de déformation. Supposons
l'assertion vraie pour les anneaux de longueur strictement inférieure à $n$ et
soit $A'$ un anneau de longueur $n$. Choisissons une petite extension
$f : A' \to A$. On a un diagramme commutatif
$$
\xymatrix{
F(A') \ar[r] \ar[d]_{F(f)} & G(A') \ar[d]^{G(f)} \\
F(A) \ar[r]^{\sim} & G(A)
}
$$
où l'application $F(A) \to G(A)$ est une bijection. Par la lissité de $F
\to G$, $F(A') \to G(A')$ est surjective (lemme
\ref{lemma-smooth-morphism-essentially-surjective}). Il suffit donc de vérifier
la bijectivité sur les fibres $F(f)^{-1}(x) \to
G(f)^{-1}(\varphi(x))$, pour les $x \in F(A)$ tels que $F(f)^{-1}(x)$ soit
non vide. Ces fibres sont précisément $\text{Lift}(x, A')$ et
$\text{Lift}(\varphi(x), A')$, et l'hypothèse fournit un isomorphisme
$d\varphi \otimes \text{id} :
TF \otimes_k \Ker(f)  \to TG \otimes_k \Ker(f)$.
Ainsi, d'après le lemme \ref{lemma-free-transitive-action} et la remarque
\ref{remark-free-transitive-action-functorial}, pour tout $x \in F(A)$ tel que
$F(f)^{-1}(x)$ soit non vide, l'application
$F(f)^{-1}(x) \to G(f)^{-1}(\varphi(x))$ est équivariante pour les actions
libres et transitives de
$TF \otimes_k \Ker(f)$. Elle est donc bijective.
\end{proof}

\noindent
Notons que, lorsque l'extension $k' \subset k$ est séparable, la condition (c)
du théorème ci-dessous est vide.

\begin{theorem}
\label{theorem-Schlessinger-prorepresentability}
Soit $F: \mathcal{C}_\Lambda \to \textit{Ensembles}$ un foncteur. Alors $F$ est
pro-représentable si et seulement si (a) $F$ est un foncteur de déformation,
(b) $\dim_k TF$ est fini, et (c) $\gamma : \text{Der}_\Lambda(k, k) \to TF$
est injective.
\end{theorem}

\begin{proof}
Supposons $F$ pro-représenté par
$R \in \widehat{\mathcal{C}}_\Lambda$. L'exemple
\ref{example-prorepresentable-deformation-functor} montre que $F$ est un
foncteur de déformation. L'exemple
\ref{example-tangent-space-prorepresentable-functor} montre que
$\dim_k TF$ est fini. Enfin, l'exemple
\ref{example-tangent-space-map-prorepresentable-functor} identifie
$\text{Der}_\Lambda(k, k) \to TF$ à
$\text{Der}_\Lambda(k, k) \to \text{Der}_\Lambda(R, k)$, qui est injectif
puisque $R \to k$ est surjectif.

\medskip\noindent
Réciproquement, supposons (a), (b) et (c) vérifiées. Le lemme
\ref{lemma-RS-implies-S1-S2} montre que (S1) et (S2) sont vérifiées. Le théorème
\ref{theorem-miniversal-object-existence} fournit donc un objet formel versel
minimal $\xi$ de $F$ tel que (\ref{equation-bijective-orbits}) soit vérifiée.
Disons que $\xi$ est situé au-dessus de $R$. L'application
$$
d\underline{\xi} : \text{Der}_\Lambda(R, k) \to T\mathcal{F}
$$
est bijective sur les orbites de $\text{Der}_\Lambda(k, k)$. Puisque l'action
de $\text{Der}_\Lambda(k, k)$ sur le membre de gauche est libre d'après (c) et
le lemme \ref{lemma-action-linear}, cette application est bijective. Le lemme
\ref{lemma-minimal-smooth-morphism-functors} montre donc que
$\underline{\xi}$ est un isomorphisme.
\end{proof}




\section{Automorphismes infinitésimaux}
\label{section-infinitesimal-automorphisms}

\noindent
Soit $\mathcal{F}$ une catégorie cofibrée en groupoïdes au-dessus de
$\mathcal{C}_\Lambda$. Étant donné un morphisme $x' \to x$ de $\mathcal{F}$
au-dessus de $A' \to A$, on a un homomorphisme de groupes induit
$$
\text{Aut}_{A'}(x') \to \text{Aut}_A(x).
$$
Le lemme \ref{lemma-RS-associated-functor} affirme que le conoyau de cet
homomorphisme détermine si la condition (RS) sur $\mathcal{F}$ passe à
$\overline{\mathcal{F}}$. Dans cette section, nous étudions le noyau de cet
homomorphisme. Nous verrons qu'il mesure lui aussi l'écart entre $\mathcal{F}$
et $\overline{\mathcal{F}}$.

\begin{definition}
\label{definition-relative-infinitesimal-auts}
Soit $\mathcal{F}$ une catégorie cofibrée en groupoïdes au-dessus de $\mathcal
C_\Lambda$. Soit $x' \to x$ un morphisme de $\mathcal{F}$ au-dessus de
$A' \to A$. Le noyau
$$
\text{Inf}(x'/x) = \Ker(\text{Aut}_{A'}(x') \to \text{Aut}_A(x))
$$
est le {\it groupe des automorphismes infinitésimaux de $x'$ au-dessus de $x$}.
\end{definition}

\begin{definition}
\label{definition-infinitesimal-auts}
Soit $\mathcal{F}$ une catégorie cofibrée en groupoïdes au-dessus de $\mathcal
C_\Lambda$. Soit $x_0 \in \Ob(\mathcal{F}(k))$. Supposons choisie une image
directe $x_0 \to x_0'$ de $x_0$ le long du morphisme
$k \to k[\epsilon], a \mapsto a$. Il existe alors un unique morphisme
$x'_0 \to x_0$ tel que $x_0 \to x_0' \to x_0$ soit l'identité de $x_0$.
Alors
$$
\text{Inf}_{x_0}(\mathcal F) = \text{Inf}(x'_0/x_0)
$$
est le {\it groupe des automorphismes infinitésimaux de $x_0$}.
\end{definition}

\begin{remark}
\label{remark-choice-pushforward-immaterial-infinitesimal-aut}
À isomorphisme canonique près, $\text{Inf}_{x_0}(\mathcal{F})$ ne dépend pas
du choix de l'image directe $x_0 \to x_0'$, car deux images directes
quelconques sont canoniquement isomorphes. De plus, si
$y_0 \in \mathcal{F}(k)$ et $x_0 \cong y_0$ dans $\mathcal{F}(k)$, alors
$\text{Inf}_{x_0}(\mathcal{F}) \cong \text{Inf}_{y_0}(\mathcal{F})$,
où l'isomorphisme ne dépend que du choix d'un isomorphisme $x_0 \to y_0$. En
particulier, $\text{Aut}_k(x_0)$ agit sur $\text{Inf}_{x_0}(\mathcal{F})$.
\end{remark}

\begin{remark}
\label{remark-trivial-aut-point}
Supposons que $\mathcal{F}$ soit une catégorie de prédéformation. Alors
\begin{enumerate}
\item pour $x_0 \in \Ob(\mathcal{F}(k))$, le groupe d'automorphismes
$\text{Aut}_k(x_0)$ est trivial et, par conséquent,
$\text{Inf}_{x_0}(\mathcal{F}) = \text{Aut}_{k[\epsilon]}(x'_0)$ ;
\item pour $x_0, y_0 \in \Ob(\mathcal{F}(k))$, il existe un unique
isomorphisme $x_0 \to y_0$, et donc une identification canonique
$\text{Inf}_{x_0}(\mathcal{F}) = \text{Inf}_{y_0}(\mathcal{F})$.
\end{enumerate}
Puisque $\mathcal{F}(k)$ est non vide, en choisissant
$x_0 \in \Ob(\mathcal{F}(k))$ et en posant
$$
\text{Inf}(\mathcal{F}) = \text{Inf}_{x_0}(\mathcal{F})
$$
on obtient un {\it groupe des automorphismes infinitésimaux de $\mathcal{F}$}
bien défini. Avec cette notation, on a
$\text{Inf}(\mathcal{F}_{x_0}) = \text{Inf}_{x_0}(\mathcal{F})$.
Comparer à l'égalité
$T\mathcal{F}_{x_0} = T_{x_0}\mathcal{F}$
de la remarque \ref{remark-tangent-space-cofibered-groupoid}.
\end{remark}

\noindent
Nous verrons que $\text{Inf}_{x_0}(\mathcal{F})$ est muni d'une structure
naturelle d'espace vectoriel sur $k$ lorsque $\mathcal{F}$ satisfait à (RS).
Nous verrons en même temps que, si $\mathcal{F}$ satisfait à (RS), les
automorphismes infinitésimaux
$\text{Inf}(x'/x)$ d'un morphisme $x' \to x$ au-dessus d'une petite extension
sont gouvernés par $\text{Inf}_{x_0}(\mathcal{F})$, où $x_0$ est une image
directe de $x$ dans $\mathcal{F}(k)$. À cette fin, nous introduisons comme suit
le foncteur des automorphismes de tout objet $x \in \Ob(\mathcal{F})$.

\begin{definition}
\label{definition-automorphism-functor}
Soit $p : \mathcal{F} \to \mathcal{C}$ une catégorie cofibrée en groupoïdes
au-dessus d'une catégorie de base arbitraire $\mathcal{C}$. Supposons effectué
un choix d'images directes. Soient $x \in \Ob(\mathcal{F})$ et $U = p(x)$.
Notons $U/\mathcal{C}$ la catégorie des objets sous $U$. Le {\it foncteur des
automorphismes de $x$} est le foncteur
$\mathit{Aut}(x) : U/\mathcal{C} \to \textit{Ensembles}$ qui envoie un objet
$f : U \to V$ sur $\text{Aut}_V(f_*x)$ et un morphisme
$$
\xymatrix{
V' \ar[rr] &                    & V\\
          & U \ar[ul]^{f'}  \ar[ur]_f &
}
$$
sur l'homomorphisme de groupes
$\text{Aut}_{V'}(f'_*x) \to \text{Aut}_V(f_*x)$
provenant de l'unique morphisme $f'_*x \to f_*x$ situé au-dessus de
$V' \to V$ et compatible avec $x \to f'_*x$ et $x \to f_*x$.
\end{definition}

\noindent
Nous nous intéresserons aux foncteurs des automorphismes d'objets d'une
catégorie cofibrée en groupoïdes $\mathcal{F}$ au-dessus de
$\mathcal{C}_\Lambda$. Si $A \in \Ob(\mathcal{C}_\Lambda)$, la catégorie
$A/\mathcal{C}_\Lambda$ n'est autre que la catégorie $\mathcal{C}_A$,
c'est-à-dire la catégorie définie dans la section \ref{section-CLambda}, où
l'on prend $\Lambda = A$ et $k = A/\mathfrak m_A$. Par conséquent, le foncteur
des automorphismes d'un objet $x \in \Ob(\mathcal{F}(A))$ est un foncteur
$\mathit{Aut}(x) : \mathcal{C}_A \to \textit{Ensembles}$.

\medskip\noindent
Le lemme suivant pourrait se déduire du lemme
\ref{lemma-RS-fiber-product-morphisms} en considérant l'« inertie » d'une
catégorie cofibrée en groupoïdes ; voir par exemple Champs, section
\ref{stacks-section-the-inertia-stack}, et Catégories, section
\ref{categories-section-inertia}. Il est toutefois plus facile de le vérifier
directement.

\begin{lemma}
\label{lemma-Aut-functor-RS}
Soit $\mathcal{F}$ une catégorie cofibrée en groupoïdes au-dessus de
$\mathcal{C}_\Lambda$ satisfaisant à (RS). Soit
$x \in \Ob(\mathcal{F}(A))$. Alors
$\mathit{Aut}(x): \mathcal{C}_A \to \textit{Ensembles}$ satisfait à (RS).
\end{lemma}

\begin{proof}
Le fait que $\mathit{Aut}(x)$ satisfait à (RS) résulte de la pleine fidélité du
foncteur
$\mathcal{F}(A_1 \times_A A_2) \to
\mathcal{F}(A_1) \times_{\mathcal{F}(A)} \mathcal{F}(A_2)$ du lemme
\ref{lemma-RS-2-categorical}.
\end{proof}

\begin{lemma}
\label{lemma-Aut-functor-tangent-space}
Soit $\mathcal{F}$ une catégorie cofibrée en groupoïdes au-dessus de
$\mathcal{C}_\Lambda$ satisfaisant à (RS). Soit
$x \in \Ob(\mathcal{F}(A))$. Soit $x_0$ une image directe de $x$ dans
$\mathcal{F}(k)$.
\begin{enumerate}
\item $T_{\text{id}_{x_0}} \mathit{Aut}(x)$ possède une structure naturelle
d'espace vectoriel sur $k$ telle que l'addition coïncide avec la composition
dans $T_{\text{id}_{x_0}} \mathit{Aut}(x)$. En particulier, la composition
dans $T_{\text{id}_{x_0}} \mathit{Aut}(x)$ est commutative.
\item Il existe un isomorphisme canonique
$T_{\text{id}_{x_0}} \mathit{Aut}(x) \to
T_{\text{id}_{x_0}} \mathit{Aut}(x_0)$
d'espaces vectoriels sur $k$.
\end{enumerate}
\end{lemma}

\begin{proof}
Appliquons la remarque \ref{remark-localize-cofibered-groupoid} au foncteur
$\mathit{Aut}(x) : \mathcal{C}_A \to \textit{Ensembles}$ et à l'élément
$\text{id}_{x_0} \in \mathit{Aut}(x)(k)$. On obtient un foncteur de
prédéformation $F = \mathit{Aut}(x)_{\text{id}_{x_0}}$. D'après les lemmes
\ref{lemma-Aut-functor-RS} et \ref{lemma-localize-RS}, $F$ est un foncteur de
déformation. Par définition,
$T_{\text{id}_{x_0}} \mathit{Aut}(x) = TF = F(k[\epsilon])$
possède la structure naturelle d'espace vectoriel sur $k$ décrite dans le
lemme \ref{lemma-tangent-space-functor}.

\medskip\noindent
L'addition est définie comme la composée
$$
F(k[\epsilon]) \times F(k[\epsilon]) \longrightarrow
F(k[\epsilon] \times_k k[\epsilon]) \longrightarrow
F(k[\epsilon])
$$
où la première application est l'inverse de la bijection garantie par (RS) et
la seconde est induite par le morphisme de $k$-algèbres
$k[\epsilon] \times_k k[\epsilon] \to k[\epsilon]$
qui envoie $(\epsilon, 0)$ et $(0, \epsilon)$ sur $\epsilon$. Si $A \to B$
est un morphisme d'anneaux de $\mathcal{C}_\Lambda$, alors
$F(A) \to F(B)$ est un homomorphisme de groupes, où
$F(A) = \mathit{Aut}(x)_{\text{id}_{x_0}}(A)$ et
$F(B) = \mathit{Aut}(x)_{\text{id}_{x_0}}(B)$ sont des groupes pour la
composition. On en déduit que
$+ : F(k[\epsilon]) \times F(k[\epsilon])\to F(k[\epsilon])$
est un homomorphisme de groupes, $F(k[\epsilon])$ étant considéré comme groupe pour la
composition. En prenant $\text{id} \in F(k[\epsilon])$ pour élément neutre, on
voit que $+(v, \text{id}) =
+(\text{id}, v) = v$ pour tout
$v \in F(k[\epsilon])$, car $(\text{id}, v)$ est l'image directe de $v$ le
long du morphisme d'anneaux
$k[\epsilon] \to k[\epsilon] \times_k k[\epsilon]$ défini par
$\epsilon \mapsto (\epsilon, 0)$. En général, si $G$ est un groupe de loi
$\circ$ et si $+ : G \times G \to G$ est un homomorphisme de groupes tel que
$+(g, 1) = +(1, g) = g$, où $1$ est l'élément neutre de $G$, alors
$+ = \circ$. Cela montre que l'addition de la structure d'espace vectoriel sur
$k$ de $F(k[\epsilon])$ coïncide avec la composition.

\medskip\noindent
Enfin, (2) s'obtient en explicitant les définitions. En effet,
$T_{\text{id}_{x_0}} \mathit{Aut}(x)$ est l'ensemble des automorphismes
$\alpha$ de l'image directe de $x$ le long de $A \to k \to k[\epsilon]$ qui
sont triviaux modulo $\epsilon$. D'autre part,
$T_{\text{id}_{x_0}} \mathit{Aut}(x_0)$ est l'ensemble des automorphismes de
l'image directe de $x_0$ le long de $k \to k[\epsilon]$ qui sont triviaux
modulo $\epsilon$. Puisque $x_0$ est l'image directe de $x$ le long de
$A \to k$, le résultat est clair.
\end{proof}

\begin{remark}
\label{remark-infaut-lifting-equalities}
Signalons quelques relations élémentaires entre les groupes d'automorphismes
infinitésimaux, les relèvements et les espaces tangents aux foncteurs des
automorphismes. Soit $\mathcal{F}$ une catégorie cofibrée en groupoïdes
au-dessus de $\mathcal{C}_\Lambda$. Soit $x' \to x$ un morphisme situé au-dessus
d'un morphisme d'anneaux $A' \to A$. Les définitions donnent alors l'égalité
$$
\text{Inf}(x'/x) = \text{Lift}(\text{id}_x, A')
$$
où il s'agit des relèvements de $\text{id}_x$ comme objet de
$\mathit{Aut}(x')$. Si $x_0 \in \Ob(\mathcal{F}(k))$ et si $x'_0$ est son
image directe dans $\mathcal{F}(k[\epsilon])$, alors, en appliquant ce qui précède
à $x'_0 \to x_0$, on obtient
$$
\text{Inf}_{x_0}(\mathcal{F}) =
\text{Lift}(\text{id}_{x_0}, k[\epsilon]) =
T_{\text{id}_{x_0}} \mathit{Aut}(x_0),
$$
la dernière égalité résultant directement des définitions.
\end{remark}

\begin{lemma}
\label{lemma-infaut-vector-space}
Soit $\mathcal{F}$ une catégorie cofibrée en groupoïdes au-dessus de
$\mathcal{C}_\Lambda$ satisfaisant à (RS). Soit
$x_0 \in \Ob(\mathcal{F}(k))$. Alors, comme ensemble,
$\text{Inf}_{x_0}(\mathcal{F})$ est égal à
$T_{\text{id}_{x_0}} \mathit{Aut}(x_0)$ ; il possède donc une structure
naturelle d'espace vectoriel sur $k$ telle que l'addition coïncide avec la
composition des automorphismes.
\end{lemma}

\begin{proof}
L'égalité d'ensembles est celle obtenue à la fin de la remarque
\ref{remark-infaut-lifting-equalities}, et l'assertion sur la structure
d'espace vectoriel résulte du lemme \ref{lemma-Aut-functor-tangent-space}.
\end{proof}

\begin{lemma}
\label{lemma-k-linear-infaut}
Soit $\varphi : \mathcal{F} \to \mathcal{G}$ un morphisme de catégories
cofibrées en groupoïdes au-dessus de $\mathcal{C}_\Lambda$ satisfaisant à
(RS). Soit $x_0 \in \Ob(\mathcal{F}(k))$. Alors $\varphi$ induit une
application $k$-linéaire
$\text{Inf}_{x_0}(\mathcal{F}) \to \text{Inf}_{\varphi(x_0)}(\mathcal{G})$.
\end{lemma}

\begin{proof}
Il est clair que $\varphi$ induit un morphisme de foncteurs
$\mathit{Aut}(x_0) \to \mathit{Aut}(\varphi(x_0))$
qui envoie l'identité sur l'identité. Le résultat découle donc de celui sur les
espaces tangents ; voir le lemme \ref{lemma-k-linear-differential}.
\end{proof}

\begin{lemma}
\label{lemma-lifted-automorphisms-torsor}
Soit $\mathcal{F}$ une catégorie cofibrée en groupoïdes au-dessus de
$\mathcal{C}_\Lambda$ satisfaisant à (RS). Soit $x' \to x$ un morphisme
situé au-dessus d'un morphisme d'anneaux surjectif $A' \to A$, de noyau $I$
annulé par $\mathfrak m_{A'}$. Soit $x_0$ une image directe de $x$ dans
$\mathcal{F}(k)$. Alors $\text{Inf}(x'/x)$ est un espace principal homogène
sous $T_{\text{id}_{x_0}} \mathit{Aut}(x') \otimes_k I
= \text{Inf}_{x_0}(\mathcal{F}) \otimes_k I$.
\end{lemma}

\begin{proof}
C'est simplement l'analogue du lemme \ref{lemma-free-transitive-action} dans
le cadre des faisceaux d'automorphismes. Plus précisément, appliquons la
remarque \ref{remark-localize-cofibered-groupoid} au foncteur
$\mathit{Aut}(x') : \mathcal{C}_{A'} \to \textit{Ensembles}$ et à l'élément
$\text{id}_{x_0} \in \mathit{Aut}(x)(k)$. On obtient un foncteur de
prédéformation $F = \mathit{Aut}(x')_{\text{id}_{x_0}}$. D'après les lemmes
\ref{lemma-Aut-functor-RS} et \ref{lemma-localize-RS}, $F$ est un foncteur de
déformation. Le lemme \ref{lemma-free-transitive-action} donne donc une action
libre et transitive de $TF \otimes_k I$ sur
$\text{Lift}(\text{id}_x, A')$. Comme
$\text{Lift}(\text{id}_x, A')$ est un groupe, cet ensemble est toujours non
vide. Notons que l'on a les égalités d'espaces vectoriels
$$
TF = T_{\text{id}_{x_0}} \mathit{Aut}(x') \otimes_k I =
\text{Inf}_{x_0}(\mathcal{F}) \otimes_k I
$$
d'après le lemme \ref{lemma-Aut-functor-tangent-space}. L'égalité
$\text{Inf}(x'/x) = \text{Lift}(\text{id}_x, A')$ de la remarque
\ref{remark-infaut-lifting-equalities} achève la démonstration.
\end{proof}

\begin{lemma}
\label{lemma-infaut-trivial}
Soit $\mathcal{F}$ une catégorie cofibrée en groupoïdes au-dessus de
$\mathcal{C}_\Lambda$ satisfaisant à (RS). Soit $x' \to x$ un morphisme de
$\mathcal{F}$ situé au-dessus d'un morphisme d'anneaux surjectif. Soit $x_0$ une
image directe de $x$ dans $\mathcal{F}(k)$. Si
$\text{Inf}_{x_0}(\mathcal{F}) = 0$, alors $\text{Inf}(x'/x) = 0$.
\end{lemma}

\begin{proof}
Cela résulte des lemmes \ref{lemma-factor-small-extension} et
\ref{lemma-lifted-automorphisms-torsor}.
\end{proof}

\begin{lemma}
\label{lemma-infdef-trivial}
Soit $\mathcal{F}$ une catégorie cofibrée en groupoïdes au-dessus de
$\mathcal{C}_\Lambda$ satisfaisant à (RS). Soit
$x_0 \in \Ob(\mathcal{F}(k))$. Alors $\text{Inf}_{x_0}(\mathcal{F}) = 0$ si
et seulement si le morphisme naturel
$\mathcal{F}_{x_0} \to \overline{\mathcal{F}_{x_0}}$ de catégories cofibrées
en groupoïdes est une équivalence.
\end{lemma}

\begin{proof}
Le morphisme $\mathcal{F}_{x_0} \to \overline{\mathcal{F}_{x_0}}$ est une
équivalence si et seulement si $\mathcal{F}_{x_0}$ est fibrée en ensoïdes ;
cf.\ Catégories, section \ref{categories-section-fibred-in-setoids} (un
ensoïde est, par définition, un groupoïde dans lequel l'identité est le seul
automorphisme de chaque objet). Montrons que
$\text{Inf}_{x_0}(\mathcal{F}) = 0$ si et seulement si cette condition est
vérifiée pour $\mathcal{F}_{x_0}$. Il est évident que, si
$\mathcal{F}_{x_0}$ est fibrée en ensoïdes, alors
$\text{Inf}_{x_0}(\mathcal{F}) = 0$. Réciproquement, supposons
$\text{Inf}_{x_0}(\mathcal{F}) = 0$. Soit $A$ un objet de
$\mathcal{C}_\Lambda$. D'après le lemme \ref{lemma-infaut-trivial},
$\text{Inf}(x/x_0) = 0$ pour tout objet $x \to x_0$ de
$\mathcal{F}_{x_0}(A)$. Puisque, par définition, $\text{Inf}(x/x_0)$ est le
groupe des automorphismes de $x \to x_0$ dans $\mathcal{F}_{x_0}(A)$, cela
montre que $\mathcal{F}_{x_0}(A)$ est un ensoïde.
\end{proof}








\section{Applications}
\label{section-applications}

\noindent
Rassemblons quelques résultats sur les catégories de déformation dont nous
nous servirons plus loin.

\begin{lemma}
\label{lemma-deformation-categories-fiber-product-morphisms}
Soient $f : \mathcal{H} \to \mathcal{F}$ et
$g : \mathcal{G} \to \mathcal{F}$ des $1$-morphismes entre catégories de
déformation. Alors
\begin{enumerate}
\item $\mathcal{W} = \mathcal{H} \times_\mathcal{F} \mathcal{G}$ est une
catégorie de déformation ;
\item on a une suite exacte à $6$ termes d'espaces vectoriels
$$
0 \to \text{Inf}(\mathcal{W})
\to \text{Inf}(\mathcal{H}) \oplus \text{Inf}(\mathcal{G})
\to \text{Inf}(\mathcal{F}) \to
T\mathcal{W} \to T\mathcal{H} \oplus T\mathcal{G} \to T\mathcal{F}
$$
\end{enumerate}
\end{lemma}

\begin{proof}
La partie (1) résulte du lemme \ref{lemma-RS-fiber-product-morphisms} et du fait
que $\mathcal{W}(k)$ est le produit fibré de deux ensoïdes ayant une unique
classe d'isomorphie au-dessus d'un ensoïde ayant une unique classe
d'isomorphie.

\medskip\noindent
Partie (2). Soit $w_0 \in \Ob(\mathcal{W}(k))$, et soient $x_0, y_0, z_0$
l'image de $w_0$ dans $\mathcal{F}, \mathcal{H}, \mathcal{G}$. Alors
$\text{Inf}(\mathcal{W}) = \text{Inf}_{w_0}(\mathcal{W})$, et de même pour
$\mathcal{H}$, $\mathcal{G}$ et $\mathcal{F}$ ; voir la remarque
\ref{remark-trivial-aut-point}. Les lemmes \ref{lemma-k-linear-differential}
et \ref{lemma-k-linear-infaut} fournissent toutes les applications linéaires,
sauf le « morphisme bord »
$\delta : \text{Inf}_{x_0}(\mathcal{F}) \to T\mathcal{W}$. Nous insérerons
plus tard les signes convenables.

\medskip\noindent
Construction de $\delta$. Choisissons une image directe $w_0 \to w'_0$ le
long de $k \to k[\epsilon]$. Notons $x'_0, y'_0, z'_0$ les images de $w'_0$
dans $\mathcal{F}, \mathcal{H}, \mathcal{G}$. On obtient notamment des
isomorphismes $b' : f(y'_0) \to x'_0$ et $c' : x'_0 \to g(z'_0)$. Notons
$b : f(y_0) \to x_0$ et $c : x_0 \to g(z_0)$ les images directes le long de
$k[\epsilon] \to k$. Cela signifie que
$w'_0 = (k[\epsilon], y'_0, z'_0, c' \circ b')$ et
$w_0 = (k, y_0, z_0, c \circ b)$, eu égard à la forme explicite du produit
fibré de catégories ; voir les remarques \ref{remarks-cofibered-groupoids}
(\ref{item-fibre-product}). Étant donné $\alpha : x'_0 \to x'_0$, posons
$\delta(\alpha) = (k[\epsilon], y'_0, z'_0, c' \circ \alpha \circ b')$
qui est bien un objet de $\mathcal{W}$ au-dessus de $k[\epsilon]$ et qui est
muni d'un morphisme
$(k[\epsilon], y'_0, z'_0, c' \circ \alpha \circ b') \to w_0$ au-dessus de
$k[\epsilon] \to k$, puisque l'image directe de $\alpha$ est l'identité
au-dessus de $k$. Plus généralement, pour tout espace vectoriel sur $k$, disons
$V$, on peut définir une application
$$
\text{Lift}(\text{id}_{x_0}, k[V])
\longrightarrow
\text{Lift}(w_0, k[V])
$$
au moyen d'exactement les mêmes formules. Cette construction est fonctorielle
en l'espace vectoriel $V$ (nous omettons les détails). Le lemme
\ref{lemma-morphism-linear-functors} montre donc que $\delta$ est $k$-linéaire.

\medskip\noindent
Ces applications étant construites, il est immédiat de montrer que la suite
est exacte. L'injectivité de la première application résulte du fait que
$f \times g : \mathcal{W} \to \mathcal{H} \times \mathcal{G}$ est fidèle. Si
$(\beta, \gamma) \in
\text{Inf}_{y_0}(\mathcal{H}) \oplus \text{Inf}_{z_0}(\mathcal{G})$
ont la même image dans $\text{Inf}_{x_0}(\mathcal{F})$, alors
$(\beta, \gamma)$ définit un automorphisme de
$w'_0 = (k[\epsilon], y'_0, z'_0, c' \circ b')$, d'où l'exactitude au
deuxième terme. Si $\alpha$ ci-dessus donne la déformation triviale
$(k[\epsilon], y'_0, z'_0, c' \circ \alpha \circ b')$
de $w_0$, alors l'isomorphisme
$w'_0 = (k[\epsilon], y'_0, z'_0, c' \circ b') \to
(k[\epsilon], y'_0, z'_0, c' \circ \alpha \circ b')$
produit une paire $(\beta, \gamma)$ qui est un antécédent de $\alpha$. Si
$w = (k[\epsilon], y, z, \phi)$ est une déformation de $w_0$ telle que
$y'_0 \cong y$ et $z \cong z'_0$, alors le morphisme
$$
f(y'_0) \to f(y) \xrightarrow{\phi} g(z) \to g(z'_0)
$$
définit un automorphisme $\alpha$ qui a pour image $w$ par $\delta$.
Enfin, si $y$ et $z$ sont des déformations de $y_0$ et $z_0$ et s'il existe un isomorphisme
$\phi : f(y) \to g(z)$ de déformations de
$f(y_0) = x_0 = g(z_0)$, on obtient un antécédent
$w = (k[\epsilon], y, z, \phi)$ de $(x, y)$ dans $T\mathcal{W}$. Cela achève
la démonstration.
\end{proof}

\begin{lemma}
\label{lemma-map-fibre-products-smooth}
Soient $\mathcal{H}_1 \to \mathcal{G}$, $\mathcal{H}_2 \to \mathcal{G}$ et
$\mathcal{G} \to \mathcal{F}$ des morphismes de catégories cofibrées en
groupoïdes au-dessus de $\mathcal{C}_\Lambda$. Supposons que
\begin{enumerate}
\item $\mathcal{F}$ et $\mathcal{G}$ soient des catégories de déformation ;
\item $T\mathcal{G} \to T\mathcal{F}$ soit injectif ;
\item $\text{Inf}(\mathcal{G}) \to \text{Inf}(\mathcal{F})$ soit surjectif.
\end{enumerate}
Alors $\mathcal{H}_1 \times_\mathcal{G} \mathcal{H}_2 \to
\mathcal{H}_1 \times_\mathcal{F} \mathcal{H}_2$ est lisse.
\end{lemma}

\begin{proof}
Notons $p_i : \mathcal{H}_i \to \mathcal{G}$ et
$q : \mathcal{G} \to \mathcal{F}$ les morphismes donnés. Soit $A' \to A$ une
petite extension de $\mathcal{C}_\Lambda$. Un objet de
$\mathcal{H}_1 \times_\mathcal{F} \mathcal{H}_2$ au-dessus de $A'$ est un
triplet $(x'_1, x'_2, a')$, où $x'_i$ est un objet de $\mathcal{H}_i$
au-dessus de $A'$ et où $a' : q(p_1(x'_1)) \to q(p_2(x'_2))$ est un morphisme
de la catégorie fibre de $\mathcal{F}$ au-dessus de $A'$. L'image directe le
long de $A' \to A$ donne $(x_1, x_2, a)$. Un relèvement fixé de cet objet en
$\mathcal{H}_1 \times_\mathcal{G} \mathcal{H}_2$ au-dessus de $A$
consiste en la donnée d'un morphisme $b : p_1(x_1) \to p_2(x_2)$ situé au-dessus de $A$ tel que
$q(b) = a$. Il faut donc montrer que l'on peut relever $b$ en un morphisme
$b' : p_1(x'_1) \to p_2(x'_2)$ dont l'image par $q$ est $a'$.

\medskip\noindent
Observons que l'on peut considérer
$$
p_1(x'_1) \to p_1(x_1) \xrightarrow{b} p_2(x_2)
\quad\text{et}\quad
p_2(x'_2) \to p_2(x_2)
$$
comme deux objets de $\textit{Lift}(p_2(x_2), A' \to A)$. Le foncteur $q$
envoie ces objets sur les deux objets
$$
q(p_1(x'_1)) \to q(p_1(x_1)) \xrightarrow{b} q(p_2(x_2))
\quad\text{et}\quad
q(p_2(x'_2)) \to q(p_2(x_2))
$$
de $\textit{Lift}(q(p_2(x_2)), A' \to A)$, lesquels sont isomorphes au moyen
du morphisme $a' : q(p_1(x'_1)) \to q(p_2(x'_2))$. D'autre part, le foncteur
$$
q : \textit{Lift}(p_2(x_2), A' \to A) \to
\textit{Lift}(q(p_2(x_2)), A' \to A)
$$
induit une injection sur les classes d'isomorphie d'après le lemme
\ref{lemma-free-transitive-action} et notre hypothèse sur les espaces tangents.
Il existe donc un morphisme $b' : p_1(x_1') \to p_2(x'_2)$ dont l'image
directe sur $A$ est $b$. Il peut toutefois être nécessaire d'ajuster le choix
de $b'$ afin d'obtenir $q(b') = a'$. Il suffit pour cela de voir que
$q : \text{Inf}(p_2(x'_2)/p_2(x_2)) \to \text{Inf}(q(p_2(x'_2))/q(p_2(x_2)))$
est surjectif. Cela résulte de notre hypothèse sur les automorphismes
infinitésimaux et du lemme \ref{lemma-lifted-automorphisms-torsor}.
\end{proof}

\begin{lemma}
\label{lemma-easy-check-smooth}
Soit $f : \mathcal{F} \to \mathcal{G}$ un morphisme de catégories de
déformation. Soit $x_0 \in \Ob(\mathcal{F}(k))$, d'image
$y_0 \in \Ob(\mathcal{G}(k))$. Si
\begin{enumerate}
\item l'application $T\mathcal{F} \to T\mathcal{G}$ est surjective ;
\item pour toute petite extension $A' \to A$ de $\mathcal{C}_\Lambda$ et tout
$x \in \mathcal{F}(A)$ d'image $y \in \mathcal{G}(A)$, s'il existe un
relèvement de $y$ à $A'$, alors il existe un relèvement de $x$ à $A'$,
\end{enumerate}
alors $\mathcal{F} \to \mathcal{G}$ est lisse (et réciproquement).
\end{lemma}

\begin{proof}
Soit $A' \to A$ une petite extension. Soit $x \in \mathcal{F}(A)$. Soit
$y' \to f(x)$ un morphisme de $\mathcal{G}$ au-dessus de $A' \to A$.
Considérons le foncteur $\text{Lift}(A', x) \to \text{Lift}(A', f(x))$ induit
par $f$. Il faut montrer qu'il existe un objet $x' \to x$ de
$\text{Lift}(A', x)$ dont l'image est $y' \to f(x)$ ; voir le lemme
\ref{lemma-smoothness-small-extensions}. La condition (2) montre que
$\text{Lift}(A', x)$ n'est pas la catégorie vide. La condition (1) et le lemme
\ref{lemma-free-transitive-action} montrent alors que l'application sur les
classes d'isomorphie est surjective, comme voulu.
\end{proof}

\begin{lemma}
\label{lemma-map-between-smooth}
Soient $\mathcal{F} \to \mathcal{G} \to \mathcal{H}$ des morphismes de
catégories cofibrées en groupoïdes au-dessus de $\mathcal{C}_\Lambda$. Si
\begin{enumerate}
\item $\mathcal{F}$ et $\mathcal{G}$ sont des catégories de déformation ;
\item l'application $T\mathcal{F} \to T\mathcal{G}$ est surjective ;
\item $\mathcal{F} \to \mathcal{H}$ est lisse.
\end{enumerate}
Alors $\mathcal{F} \to \mathcal{G}$ est lisse.
\end{lemma}

\begin{proof}
Soit $A' \to A$ une petite extension de $\mathcal{C}_\Lambda$, et soit
$x \in \mathcal{F}(A)$ d'image $y \in \mathcal{G}(A)$. Supposons qu'il existe
un relèvement $y' \in \mathcal{G}(A')$. D'après le lemme
\ref{lemma-easy-check-smooth}, il suffit de vérifier que $x$ possède aussi un
relèvement. Prenons l'image $z' \in \mathcal{H}(A')$ de $y'$. Puisque
$\mathcal{F} \to \mathcal{H}$ est lisse, il existe
$x' \in \mathcal{F}(A')$ ayant pour images à la fois
$x \in \mathcal{F}(A)$ et $z' \in \mathcal{H}(A')$ ; voir la définition
\ref{definition-smooth-morphism}. Cela achève la démonstration.
\end{proof}





\section{Groupoïdes en foncteurs sur une catégorie arbitraire}
\label{section-groupoids-arbitrary}

\noindent
Commençons par des généralités sur les groupoïdes en foncteurs sur une
catégorie arbitraire. Dans la section suivante, nous passerons à la catégorie
$\mathcal{C}_\Lambda$. Par souci de clarté, nous appellerons parfois
« catégorie groupoïde » un groupoïde ordinaire, c'est-à-dire une catégorie
dont tous les morphismes sont des isomorphismes.

\begin{definition}
\label{definition-groupoid-in-functors}
Soit $\mathcal{C}$ une catégorie. La {\it catégorie des groupoïdes en
foncteurs sur $\mathcal{C}$} est la catégorie dont les objets et les morphismes
sont les suivants.
\begin{enumerate}
\item Objets : un {\it groupoïde en foncteurs sur $\mathcal{C}$} est un
quintuplet
$(U, R, s, t, c)$, où $U, R : \mathcal{C} \to \textit{Ensembles}$ sont des
foncteurs, et $s, t : R \to U$ et $c : R \times_{s, U, t} R \to R$ sont
des morphismes ayant la propriété suivante : pour tout objet $T$ de
$\mathcal{C}$, le quintuplet
$$
(U(T), R(T), s, t, c)
$$
est une catégorie groupoïde.
\item Morphismes : un {\it morphisme
$(U, R, s, t, c) \to (U', R', s', t', c')$ de groupoïdes en foncteurs sur
$\mathcal{C}$} consiste en des morphismes $U \to U'$ et $R \to R'$ ayant la
propriété suivante : pour tout objet $T$ de $\mathcal{C}$, les applications
induites $U(T) \to U'(T)$ et $R(T) \to R'(T)$ définissent un foncteur entre
catégories groupoïdes
$$
(U(T), R(T), s, t, c) \to (U'(T), R'(T), s', t', c').
$$
\end{enumerate}
\end{definition}

\begin{remark}
\label{remark-confusion-groupoids-in-functors}
Un groupoïde en foncteurs sur $\mathcal{C}$ équivaut à la donnée d'un foncteur
$\mathcal{C} \to \textit{Groupoïdes}$, et un morphisme de groupoïdes en
foncteurs sur $\mathcal{C}$ équivaut à un morphisme des foncteurs
correspondants $\mathcal{C} \to \textit{Groupoïdes}$ (où l'on considère
$\textit{Groupoïdes}$ comme une 1-catégorie). Pour nos
besoins, il est toutefois plus commode d'employer la terminologie des
groupoïdes en foncteurs. En effet, considérer un groupoïde en foncteurs comme
le foncteur correspondant $\mathcal{C} \to \textit{Groupoïdes}$, ou, de manière
équivalente, comme la catégorie cofibrée en groupoïdes associée à ce foncteur,
peut prêter à confusion (remarque
\ref{remark-smooth-groupoid-in-functors-warning}).
\end{remark}

\begin{remark}
\label{remark-identity-inverse}
Soit $(U, R, s, t, c)$ un groupoïde en foncteurs sur une catégorie
$\mathcal{C}$. Il existe d'uniques morphismes $e : U \to R$ et $i : R \to R$
tels que, pour tout objet $T$ de $\mathcal{C}$, $e: U(T) \to R(T)$ envoie
$x \in U(T)$ sur la flèche identique de $x$, et $i: R(T) \to R(T)$ envoie
$a \in U(T)$ sur l'inverse de $a$ dans la catégorie groupoïde
$(U(T), R(T), s, t, c)$. Nous appellerons parfois $s$, $t$, $c$, $e$ et $i$
la « source », le « but », la « composition », l'« identité » et l'« inverse ».
\end{remark}

\begin{definition}
\label{definition-representable}
Soit $\mathcal{C}$ une catégorie. Un groupoïde en foncteurs sur $\mathcal{C}$
est {\it représentable} s'il est isomorphe à un groupoïde de la forme
$(\underline{U}, \underline{R}, s, t, c)$, où $U$ et $R$ sont des objets de
$\mathcal{C}$ et où la somme amalgamée $R \amalg_{s, U, t} R$ existe.
\end{definition}

\begin{remark}
\label{remark-reason-existence-coproduct}
Ainsi, un groupoïde représentable en foncteurs sur $\mathcal{C}$ est donné par
des objets $U$ et $R$ de $\mathcal{C}$, ainsi que par des morphismes
$s, t : U \to R$ et $c : R \to R \amalg_{s, U, t} R$, tels que
$(\underline{U}, \underline{R}, s, t, c)$ satisfasse à la condition de la
définition \ref{definition-groupoid-in-functors}. On exige l'existence de la
somme amalgamée $R \amalg_{s, U, t} R$ afin que le morphisme de composition
$c$ soit défini au niveau des morphismes de $\mathcal{C}$. Cette exigence sera
toujours satisfaite plus bas lorsque nous considérerons des groupoïdes
représentables en foncteurs sur $\widehat{\mathcal{C}}_\Lambda$, car, d'après
le lemme \ref{lemma-CLambdahat-pushouts}, la catégorie
$\widehat{\mathcal{C}}_\Lambda$ admet des sommes amalgamées.
\end{remark}

\begin{remark}
\label{remark-simplify-terminology}
Nous dirons « {\it soit $(\underline{U}, \underline{R}, s, t, c)$ un
groupoïde en foncteurs sur $\mathcal{C}$} » pour signifier que nous disposons
d'un groupoïde représentable en foncteurs. Cela signifie donc que $U$ et $R$
sont des objets de $\mathcal{C}$, qu'il existe des morphismes
$s, t : U \to R$, que la somme amalgamée $R \amalg_{s, U, t} R$ existe, qu'il
existe un morphisme $c : R \to R \amalg_{s, U, t} R$, et que
$(\underline{U}, \underline{R}, s, t, c)$ est un groupoïde en foncteurs sur
$\mathcal{C}$.
\end{remark}

\noindent
Introduisons une notation pour la restriction des groupoïdes en foncteurs. Elle
sera utile plus bas lorsque nous restreindrons de $\widehat{\mathcal
C}_\Lambda$ à $\mathcal{C}_\Lambda$.

\begin{definition}
\label{definition-restricting-groupoids-in-functors}
Soit $(U, R, s, t, c)$ un groupoïde en foncteurs sur une catégorie
$\mathcal{C}$. Soit $\mathcal{C}'$ une sous-catégorie de $\mathcal{C}$. La
{\it restriction $(U, R, s, t, c)|_{\mathcal{C}'}$ de $(U, R, s, t, c)$ à
$\mathcal{C}'$} est le groupoïde en foncteurs sur $\mathcal{C}'$ donné par
$(U|_{\mathcal{C}'}, R|_{\mathcal
C'}, s|_{\mathcal{C}'}, t|_{\mathcal{C}'}, c|_{\mathcal{C}'})$.
\end{definition}

\begin{remark}
\label{remark-notation-restriction}
Dans la situation de la définition
\ref{definition-restricting-groupoids-in-functors}, nous notons souvent
$s|_{\mathcal{C}'}, t|_{\mathcal{C}'}, c|_{\mathcal{C}'}$ simplement $s, t, c$.
\end{remark}

\begin{definition}
\label{definition-quotient}
Soit $(U, R, s, t, c)$ un groupoïde en foncteurs sur une catégorie
$\mathcal{C}$.
\begin{enumerate}
\item La correspondance $T \mapsto  (U(T), R(T), s, t, c)$ définit un foncteur
$\mathcal{C} \to \textit{Groupoïdes}$. La {\it catégorie quotient cofibrée en
groupoïdes $[U/R] \to \mathcal{C}$} est la catégorie cofibrée en groupoïdes
au-dessus de $\mathcal{C}$ associée à ce foncteur (comme dans les remarques
\ref{remarks-cofibered-groupoids}
(\ref{item-construction-associated-cofibered-groupoid})).
\item Le {\it morphisme quotient $U \to [U/R]$} est le morphisme de catégories
cofibrées en groupoïdes au-dessus de $\mathcal{C}$ défini par les règles
suivantes :
\begin{enumerate}
\item $x \in U(T)$ est envoyé sur l'objet $(T, x) \in \Ob([U/R](T))$ ;
\item $x \in U(T)$ et $f : T \to T'$ donnent lieu au morphisme
$(f, \text{id}_{U(f)(x)}): (T, x) \to (T, U(f)(x))$ au-dessus de
$f : T \to T'$.
\end{enumerate}
\end{enumerate}
\end{definition}





\section{Groupoïdes en foncteurs sur la catégorie de base}
\label{section-prorepresentable-groupoids-in-functors}

\noindent
Dans cette section, nous étudions les groupoïdes en foncteurs sur
$\mathcal{C}_\Lambda$. Notre but ultime est de montrer que les groupoïdes en
foncteurs sur $\mathcal{C}_\Lambda$ qui sont pro-représentables fournissent des
« présentations » de catégories de déformation qui se comportent bien, de la
même manière que les groupoïdes lisses en espaces algébriques fournissent des
présentations de champs algébriques ; cf.\ Champs algébriques, section
\ref{algebraic-section-stack-to-presentation}.

\begin{definition}
\label{definition-prorepresentable-groupoid-in-functors}
Un groupoïde en foncteurs sur $\mathcal{C}_\Lambda$ est
{\it pro-représentable} s'il est isomorphe à
$(\underline{R_0}, \underline{R_1}, s, t, c)|_{\mathcal{C}_\Lambda}$
pour un certain groupoïde en foncteurs représentable
$(\underline{R_0}, \underline{R_1}, s, t, c)$ sur la catégorie
$\widehat{\mathcal{C}}_\Lambda$.
\end{definition}

\noindent
Soit $(U, R, s, t, c)$ un groupoïde en foncteurs sur
$\mathcal{C}_\Lambda$. En passant aux complétés, on obtient un quintuplet
$(\widehat{U}, \widehat{R}, \widehat{s}, \widehat{t}, \widehat{c})$. D'après
la remarque \ref{remark-completion-restriction-cofset-adjoint}, la complétion,
vue comme foncteur sur $\text{CofSet}(\mathcal{C}_\Lambda)$, est un adjoint à
droite ; elle commute donc aux limites. En particulier, il existe un
isomorphisme canonique
$$
\widehat{R \times_{s, U, t} R}
\longrightarrow
\widehat{R} \times_{\widehat{s}, \widehat{U}, \widehat{t}} \widehat{R},
$$
de sorte que $\widehat{c}$ peut être considéré comme un foncteur
$\widehat{R} \times_{\widehat{s}, \widehat{U}, \widehat{t}} \widehat{R} \to
\widehat{R}$. Alors
$(\widehat{U}, \widehat{R}, \widehat{s}, \widehat{t}, \widehat{c})$
est un groupoïde en foncteurs sur $\widehat{\mathcal{C}}_\Lambda$, dont les
morphismes identité et inverse sont les complétés de ceux de
$(U, R, s, t, c)$.

\begin{definition}
\label{definition-completion-groupoid-in-functors}
Soit $(U, R, s, t, c)$ un groupoïde en foncteurs sur
$\mathcal{C}_\Lambda$. Le {\it complété $(U, R, s, t, c)^{\wedge}$ de
$(U, R, s, t, c)$} est le groupoïde en foncteurs
$(\widehat{U}, \widehat{R}, \widehat{s}, \widehat{t}, \widehat{c})$ sur
$\widehat{\mathcal{C}}_\Lambda$ décrit ci-dessus.
\end{definition}

\begin{remark}
\label{remark-groupoid-in-functors-complete-restrict}
Soit $(U, R, s, t, c)$ un groupoïde en foncteurs sur
$\mathcal{C}_\Lambda$. Il existe alors un isomorphisme canonique
$(U, R, s, t, c)^{\wedge}|_{\mathcal{C}_\Lambda} \cong (U, R, s, t, c)$ ; voir
la remarque \ref{remark-restrict-completion}. D'autre part, soit
$(U, R, s, t, c)$ un groupoïde en foncteurs sur
$\widehat{\mathcal{C}}_\Lambda$ tel que
$U, R : \widehat{\mathcal{C}}_\Lambda \to \textit{Ensembles}$
commutent tous deux aux limites, par exemple si $U, R$ sont représentables.
Il existe alors un isomorphisme canonique
$((U, R, s, t, c)|_{\mathcal{C}_\Lambda})^{\wedge} \cong (U, R, s, t, c)$.
Cela résulte de la remarque \ref{remark-restrict-complete-continuous-functor}.
\end{remark}

\begin{lemma}
\label{lemma-groupoid-in-functors-prorep-equivalences}
Soit $(U, R, s, t, c)$ un groupoïde en foncteurs sur
$\mathcal{C}_\Lambda$.
\begin{enumerate}
\item $(U, R, s, t, c)$ est pro-représentable si et seulement si son complété
est représentable comme groupoïde en foncteurs sur
$\widehat{\mathcal{C}}_\Lambda$.
\item $(U, R, s, t, c)$ est pro-représentable si et seulement si $U$ et $R$
sont pro-représentables.
\end{enumerate}
\end{lemma}

\begin{proof}
La partie (1) résulte de la remarque
\ref{remark-groupoid-in-functors-complete-restrict}. Pour (2), l'implication
directe découle immédiatement de la définition d'un groupoïde en foncteurs
pro-représentable. Réciproquement, supposons $U$ et $R$
pro-représentables, et écrivons $U \cong \underline{R_0}|_{\mathcal{C}_\Lambda}$ et
$R \cong \underline{R_1}|_{\mathcal{C}_\Lambda}$ pour des objets $R_0$ et
$R_1$ de $\widehat{\mathcal{C}}_\Lambda$. Puisque
$\underline{R_0} \cong \widehat{\underline{R_0}|_{\mathcal{C}_\Lambda}}$ et
$\underline{R_1} \cong \widehat{\underline{R_1}|_{\mathcal{C}_\Lambda}}$
d'après la remarque \ref{remark-restrict-complete-continuous-functor}, on voit
que le complété $(U, R, s, t, c)^\wedge$ est un groupoïde en foncteurs de la
forme
$(\underline{R_0}, \underline{R_1}, \widehat{s}, \widehat{t}, \widehat{c})$.
Par le lemme \ref{lemma-CLambdahat-pushouts}, la somme amalgamée
$\underline{R_1} \times_{\widehat{s}, \underline{R_1}, \widehat{t}}
\underline{R_1}$ existe. Ainsi,
$(\underline{R_0}, \underline{R_1}, \widehat{s}, \widehat{t}, \widehat{c})$
est un groupoïde représentable en foncteurs sur
$\widehat{\mathcal{C}}_\Lambda$. Enfin, la restriction
$(\underline{R_0}, \underline{R_1}, s, t, c)|_{\mathcal{C}_\Lambda}$
redonne $(U, R, s, t, c)$ d'après la remarque
\ref{remark-groupoid-in-functors-complete-restrict} ; ainsi,
$(U, R, s, t, c)$ est pro-représentable par définition.
\end{proof}






\section{Groupoïdes en foncteurs lisses sur la catégorie de base}
\label{section-smooth-minimal-groupoids-in-functors}

\noindent
La notion de lissité des groupoïdes en foncteurs sur $\mathcal{C}_\Lambda$ est
définie comme suit.

\begin{definition}
\label{definition-smooth-groupoid-in-functors}
Soit $(U, R, s, t, c)$ un groupoïde en foncteurs sur
$\mathcal{C}_\Lambda$. On dit que $(U, R, s, t, c)$ est {\it lisse} si
$s, t: R \to U$ sont lisses.
\end{definition}

\begin{remark}
\label{remark-smooth-groupoid-in-functors-warning}
Notons que cette terminologie peut prêter à confusion : si
$(U, R, s, t, c)$ est un groupoïde en foncteurs lisse, alors le quotient
$[U/R]$ n'est pas nécessairement une catégorie cofibrée en groupoïdes lisse au
sens de la définition \ref{definition-cofibered-groupoid-projection-smooth}.
Cependant, la lissité de $(U, R, s, t, c)$ entraîne la lissité du morphisme
quotient $U \to [U/R]$ (et lui est en fait équivalente), comme nous le verrons
dans le lemme \ref{lemma-smooth-quotient-morphism}.
\end{remark}

\begin{remark}
\label{remark-smooth-power-series-prorepresentable-smooth-groupoid-in-functors}
Soit $(\underline{R_0}, \underline{R_1}, s, t, c)|_{\mathcal{C}_\Lambda}$ un
groupoïde en foncteurs pro-représentable sur $\mathcal{C}_\Lambda$. Alors
$(\underline{R_0}, \underline{R_1}, s, t, c)|_{\mathcal{C}_\Lambda}$ est
lisse si et seulement si $R_1$ est un anneau de séries formelles sur $R_0$
pour chacune des structures définies par $s$ et par $t$. Cela résulte du lemme
\ref{lemma-smooth-morphism-power-series}.
\end{remark}

\begin{lemma}
\label{lemma-smooth-quotient-morphism}
Soit $(U, R, s, t, c)$ un groupoïde en foncteurs sur
$\mathcal{C}_\Lambda$. Les assertions suivantes sont équivalentes :
\begin{enumerate}
\item Le groupoïde en foncteurs $(U, R, s, t, c)$ est lisse.
\item Le morphisme $s : R \to U$ est lisse.
\item Le morphisme $t : R \to U$ est lisse.
\item Le morphisme quotient $U \to [U/R]$ est lisse.
\end{enumerate}
\end{lemma}

\begin{proof}
L'assertion (2) est équivalente à (3), puisque l'inverse $i: R \to R$ de
$(U, R, s, t, c)$ est un isomorphisme et que $t = s \circ i$. Par définition,
(1) équivaut à la conjonction de (2) et (3), donc à chacune d'elles séparément.

\medskip\noindent
Montrons enfin que (2) est équivalente à (4). En explicitant les définitions :
\begin{enumerate}
\item[(2)] La lissité de $s: R \to U$ équivaut à la condition suivante : si
$f: B \to A$ est un morphisme d'anneaux surjectif dans
$\mathcal{C}_\Lambda$, si $a \in R(A)$ et si $y \in U(B)$ vérifie
$s(a) = U(f)(y)$, alors il existe $a' \in R(B)$ tel que
$R(f)(a') = a$ et $s(a') = y$.
\item[(4)] La lissité de $U \to [U/R]$ équivaut à la condition suivante : si
$f: B \to A$ est un morphisme d'anneaux surjectif dans
$\mathcal{C}_\Lambda$ et si $(f, a) : (B, y) \to (A, x)$ est un morphisme de
$[U/R]$, alors il existe $x' \in U(B)$ et $b \in R(B)$, avec
$s(b) = x'$, $t(b) = y$, tels que $c(a, R(f)(b)) = e(x)$. Ici,
$e : U \to R$ désigne l'identité, et la notation $(f, a)$ est celle des
remarques \ref{remarks-cofibered-groupoids}
(\ref{item-construction-associated-cofibered-groupoid}) ; en particulier,
$a \in R(A)$, avec $s(a) = U(f)(y)$ et $t(a) = x$.
\end{enumerate}
Supposons (4) vérifiée et soient $f, a, y$ comme en (2). Posons $x = t(a)$,
de sorte que l'on ait un morphisme $(f, a): (B, y) \to (A, x)$. Alors (4)
fournit $x'$ et $b$, et $a' = i(b)$ satisfait aux exigences de (2).
Réciproquement, supposons (2) vérifiée et soit
$(f, a): (B, y) \to (A, x)$ comme en (4). Alors (2) fournit
$a' \in R(B)$, et $x' = t(a')$ et $b = i(a')$ satisfont aux exigences de (4).
\end{proof}





\section{Catégories de déformation comme quotients de groupoïdes en foncteurs}
\label{section-deformation-categories-as-quotients}

\noindent
Nous étudions des conditions sur un groupoïde en foncteurs sur
$\mathcal{C}_\Lambda$ qui garantissent que son quotient est une catégorie de
déformation, et nous calculons l'espace tangent et l'espace des automorphismes
infinitésimaux d'un tel quotient.

\begin{lemma}
\label{lemma-smooth-RS-groupoid-in-functors-quotient}
Soit $(U, R, s, t, c)$ un groupoïde en foncteurs lisse sur
$\mathcal{C}_\Lambda$. Supposons que $U$ et $R$ satisfassent à (RS). Alors
$[U/R]$ satisfait à (RS).
\end{lemma}

\begin{proof}
Soit
$$
\xymatrix{
                           &     (A_2, x_2) \ar[d]^{(f_2, a_2)} \\
(A_1, x_1) \ar[r]^{(f_1, a_1)} &     (A, x)
}
$$
un diagramme de $[U/R]$ tel que $f_2: A_2 \to A$ soit surjectif. La notation
est celle des remarques \ref{remarks-cofibered-groupoids}
(\ref{item-construction-associated-cofibered-groupoid}). Ainsi,
$f_1: A_1 \to A, f_2: A_2 \to A$ sont des morphismes de
$\mathcal{C}_\Lambda$, $x \in U(A)$, $x_1 \in U(A_1)$, $x_2 \in U(A_2)$, et
$a_1, a_2 \in R(A)$, avec $s(a_1) = U(f_1)(x_1)$, $t(a_1) = x$ et
$s(a_2) = U(f_2)(x_2)$, $t(a_2) = x$. Construisons comme suit un produit fibré
au-dessus de $A_1 \times_A A_2$ pour ce diagramme de $[U/R]$.

\medskip\noindent
Soit $a = c(i(a_1), a_2)$, où $i: R \to R$ est le morphisme inverse. Alors
$a \in R(A)$, $x_2 \in U(A_2)$ et $s(a) = U(f_2)(x_2)$. On a donc un élément
$(a, x_2) \in R(A) \times_{s, U(A), U(f_2)} U(A_2)$. Par la lissité de
$s : R \to U$, il existe un élément $\widetilde{a} \in R(A_2)$ tel que
$R(f_2)(\widetilde{a}) = a$ et $s(\widetilde{a}) = x_2$. En particulier,
$U(f_2)(t(\widetilde{a})) = t(a) = U(f_1)(x_1)$. Ainsi, $x_1$ et
$t(\widetilde{a})$ définissent un élément
$$
(x_1, t(\widetilde{a})) \in U(A_1) \times_{U(A)} U(A_2).
$$
Puisque $U$ satisfait à (RS), on a une identification
$U(A_1) \times_{U(A)} U(A_2) = U(A_1 \times_A A_2)$. Notons
$x_1 \times t(\widetilde{a}) \in U(A_1 \times_A A_2)$ l'élément correspondant
à $(x_1, t(\widetilde{a})) \in U(A_1) \times_{U(A)} U(A_2)$. Soient
$p_1, p_2$ les projections de $A_1 \times_A A_2$. Nous affirmons que
$$
\xymatrix{
(A_1 \times_A A_2, x_1 \times t(\widetilde{a})) \ar[d]_{(p_1, e(x_1))}
\ar[rr]_-{(p_2, i(\widetilde{a}))} & & (A_2, x_2) \ar[d]^{(f_2, a_2)} \\
(A_1, x_1) \ar[rr]^{(f_1, a_1)} & & (A, x)
}
$$
est un carré cartésien dans $[U/R]$. (Notons que $e: U \to R$ désigne
l'identité.)

\medskip\noindent
Le diagramme est commutatif parce que
$c(a_2, R(f_2)(i(\widetilde{a}))) = c(a_2, i(a)) = a_1$.
Pour vérifier qu'il est cartésien, soit
$$
\xymatrix{
(B, z) \ar[d]_{(g_1, b_1)} \ar[rr]_{(g_2, b_2)} & & (A_2, x_2)
\ar[d]^{(f_2, a_2)} \\
(A_1, x_1) \ar[rr]^{(f_1, a_1)} & & (A, x)
}
$$
un diagramme commutatif de $[U/R]$. Nous allons montrer qu'il existe un unique
morphisme
$(g, b) : (B, z) \to (A_1 \times_A A_2, x_1 \times t(\widetilde{a}))$
compatible aux morphismes vers $(A_1, x_1)$ et $(A_2, x_2)$. Nécessairement,
$g = (g_1, g_2) : B \to A_1 \times_A A_2$. Puisque, par hypothèse, $R$
satisfait à (RS), on a une identification
$R(A_1 \times_A A_2) = R(A_1) \times_{R(A)} R(A_2)$. On peut donc écrire
$b = (b'_1, b'_2)$ pour certains $b'_1 \in R(A_1)$, $b'_2 \in R(A_2)$ ayant
la même image dans $R(A)$. Alors
$((g_1, g_2), (b'_1, b'_2)) : (B, z) \to
(A_1 \times_A A_2, x_1 \times t(\widetilde{a}))$
est compatible avec les projections si et seulement si $b'_1 = b_1$ et
$b'_2 = c(\widetilde{a}, b_2)$, ce qui démontre l'existence et l'unicité.
\end{proof}

\begin{lemma}
\label{lemma-deformation-groupoid-quotient}
Soit $(U, R, s, t, c)$ un groupoïde en foncteurs lisse sur
$\mathcal{C}_\Lambda$. Supposons que $U$ et $R$ soient des foncteurs de
déformation. Alors :
\begin{enumerate}
\item Le quotient $[U/R]$ est une catégorie de déformation.
\item L'espace tangent de $[U/R]$ est
$$
T[U/R] = \Coker(ds-dt: TR \to TU).
$$
\item L'espace des automorphismes infinitésimaux de $[U/R]$ est
$$
\text{Inf}([U/R]) =
\Ker(ds \oplus dt : TR \to TU \oplus TU).
$$
\end{enumerate}
\end{lemma}

\begin{proof}
Puisque $U$ et $R$ sont des foncteurs de déformation, $[U/R]$ est une
catégorie de prédéformation. Comme (RS) est vérifiée par définition pour les
foncteurs de déformation, le lemme
\ref{lemma-smooth-RS-groupoid-in-functors-quotient} montre que (RS) est
vérifiée pour [U/R]. Ainsi, $[U/R]$ est une catégorie de déformation. Les
assertions (2) et (3) résultent directement des définitions.
\end{proof}







\section{Présentations des catégories cofibrées en groupoïdes}
\label{section-presentation-categories-cofibred-in-groupoids}

\noindent
Une présentation se définit comme suit.

\begin{definition}
\label{definition-presentation}
Soit $\mathcal{F}$ une catégorie cofibrée en groupoïdes au-dessus d'une
catégorie $\mathcal{C}$. Soit $(U, R, s, t, c)$ un groupoïde en foncteurs sur
$\mathcal{C}$. Une {\it présentation de $\mathcal{F}$ par
$(U, R, s, t, c)$} est une équivalence
$\varphi : [U/R] \to \mathcal{F}$ de catégories cofibrées en groupoïdes
au-dessus de $\mathcal{C}$.
\end{definition}

\noindent
Les deux lemmes généraux suivants serviront à construire des présentations.

\begin{lemma}
\label{lemma-presentation-construction}
Soit $\mathcal{F}$ une catégorie cofibrée en groupoïdes au-dessus d'une
catégorie $\mathcal{C}$. Soit $U : \mathcal{C} \to \textit{Ensembles}$ un
foncteur. Soit $f : U \to \mathcal{F}$ un morphisme de catégories cofibrées
en groupoïdes au-dessus de $\mathcal{C}$. Définissons $R, s, t, c$ comme suit :
\begin{enumerate}
\item $R : \mathcal{C} \to \textit{Ensembles}$ est le foncteur
$U \times_{f, \mathcal{F}, f} U$.
\item $t, s : R \to U$ sont respectivement la première et la seconde
projection.
\item $c : R \times_{s, U, t} R \to R$ est le morphisme donné par la
projection sur le premier et le dernier facteur de
$U \times_{f, \mathcal{F}, f} U \times_{f, \mathcal{F}, f} U$
sous l'isomorphisme canonique
$R \times_{s, U, t} R \to
U \times_{f, \mathcal{F}, f} U \times_{f, \mathcal{F}, f} U$.
\end{enumerate}
Alors $(U, R, s, t, c)$ est un groupoïde en foncteurs sur $\mathcal{C}$.
\end{lemma}

\begin{proof}
La démonstration est omise.
\end{proof}

\begin{lemma}
\label{lemma-presentation-morphism}
Soit $\mathcal{F}$ une catégorie cofibrée en groupoïdes au-dessus d'une
catégorie $\mathcal{C}$. Soit $U : \mathcal{C} \to \textit{Ensembles}$ un
foncteur. Soit $f : U \to \mathcal{F}$ un morphisme de catégories cofibrées
en groupoïdes au-dessus de $\mathcal{C}$. Soit $(U, R, s, t, c)$ le groupoïde
en foncteurs sur $\mathcal{C}$ construit à partir de
$f : U \to \mathcal{F}$ dans le lemme
\ref{lemma-presentation-construction}. Il existe alors un morphisme naturel
$[f] : [U/R] \to \mathcal{F}$ tel que :
\begin{enumerate}
\item $[f]: [U/R] \to \mathcal{F}$ est pleinement fidèle.
\item $[f]: [U/R] \to \mathcal{F}$ est une équivalence si et seulement si
$f : U \to \mathcal{F}$ est essentiellement surjectif.
\end{enumerate}
\end{lemma}

\begin{proof}
La démonstration est omise.
\end{proof}







\section{Présentations des catégories de déformation}
\label{section-presentation-deformation-categories}

\noindent
D'après le lemme suivant, un morphisme lisse d'un foncteur de prédéformation
vers une catégorie de prédéformation $\mathcal{F}$ donne lieu à une
présentation de $\mathcal{F}$ par un groupoïde en foncteurs lisse.

\begin{lemma}
\label{lemma-smooth-groupoid-in-functors-construction}
Soit $\mathcal{F}$ une catégorie cofibrée en groupoïdes au-dessus de
$\mathcal{C}_\Lambda$. Soit $U : \mathcal{C}_\Lambda \to \textit{Ensembles}$ un
foncteur. Soit $f : U \to \mathcal{F}$ un morphisme lisse de catégories
cofibrées en groupoïdes. Alors :
\begin{enumerate}
\item Si $(U, R, s, t, c)$ est le groupoïde en foncteurs sur
$\mathcal{C}_\Lambda$ construit à partir de $f : U \to \mathcal{F}$ dans le
lemme \ref{lemma-presentation-construction}, alors $(U, R, s, t, c)$ est
lisse.
\item Si $f : U(k) \to \mathcal{F}(k)$ est essentiellement surjectif, alors le
morphisme $[f] : [U/R] \to \mathcal{F}$ du lemme
\ref{lemma-presentation-morphism} est une équivalence.
\end{enumerate}
\end{lemma}

\begin{proof}
La construction du lemme \ref{lemma-presentation-construction} donne un
diagramme commutatif
$$
\xymatrix{
R = U \times_{f, \mathcal{F}, f} U \ar[r]_-s \ar[d]_t & U
\ar[d]^f \\
U \ar[r]^f & \mathcal{F}
}
$$
où $t, s$ sont respectivement la première et la seconde projection. Les morphismes $t, s$
sont donc lisses d'après le lemme \ref{lemma-smooth-properties}. Ainsi, (1)
est vérifiée.

\medskip\noindent
Si l'hypothèse de (2) est vérifiée, le lemme
\ref{lemma-smooth-morphism-essentially-surjective} montre que le morphisme
$f : U \to \mathcal{F}$ est essentiellement surjectif. Le lemme
\ref{lemma-presentation-morphism} montre donc que le morphisme
$[f] : [U/R] \to \mathcal{F}$ est une équivalence.
\end{proof}

\begin{lemma}
\label{lemma-deformation-functor-diagonal}
Soit $\mathcal{F}$ une catégorie de déformation. Soit
$U : \mathcal{C}_\Lambda \to \textit{Ensembles}$ un foncteur de déformation. Soit
$f: U \to \mathcal{F}$ un morphisme de catégories cofibrées en groupoïdes.
Alors $U \times_{f, \mathcal{F}, f} U$ est un foncteur de déformation dont
l'espace tangent s'insère dans une suite exacte d'espaces vectoriels sur $k$
$$
0 \to \text{Inf}(\mathcal{F}) \to
T(U \times_{f, \mathcal{F}, f} U) \to TU \oplus TU
$$
\end{lemma}

\begin{proof}
Cela résulte du lemme
\ref{lemma-deformation-categories-fiber-product-morphisms} et du fait que
$\text{Inf}(U) = (0)$.
\end{proof}

\begin{lemma}
\label{lemma-prorepresentable-groupoid-in-functors-construction}
Soit $\mathcal{F}$ une catégorie de déformation. Soit
$U : \mathcal{C}_\Lambda \to \textit{Ensembles}$ un foncteur pro-représentable.
Soit $f : U \to \mathcal{F}$ un morphisme de catégories cofibrées en
groupoïdes. Soit $(U, R, s, t, c)$ le groupoïde en foncteurs sur
$\mathcal{C}_\Lambda$ construit à partir de $f : U \to \mathcal{F}$ dans le
lemme \ref{lemma-presentation-construction}. Si
$\dim_k \text{Inf}(\mathcal{F}) < \infty$, alors $(U, R, s, t, c)$ est
pro-représentable.
\end{lemma}

\begin{proof}
Notons que $U$ est un foncteur de déformation d'après l'exemple
\ref{example-prorepresentable-deformation-functor}. Le lemme
\ref{lemma-deformation-functor-diagonal} montre que
$R = U \times_{f, \mathcal{F}, f} U$ est un foncteur de déformation dont
l'espace tangent $TR = T(U \times_{f, \mathcal{F}, f} U)$ s'insère dans une
suite exacte $0 \to \text{Inf}(\mathcal{F}) \to TR \to TU \oplus TU$. Le
premier espace est de dimension finie par hypothèse, et $TU$ est de dimension
finie d'après l'exemple
\ref{example-tangent-space-prorepresentable-functor} ; ainsi,
$\dim TR < \infty$. L'application
$\gamma : \text{Der}_\Lambda(k, k) \to TR$ (voir (\ref{equation-map})) est
injective, car son composé avec $TR \to TU$ est injectif, d'après le théorème
\ref{theorem-Schlessinger-prorepresentability} appliqué au foncteur
pro-représentable $U$. Le théorème
\ref{theorem-Schlessinger-prorepresentability} montre donc que $R$ est
pro-représentable. Il résulte alors du lemme
\ref{lemma-groupoid-in-functors-prorep-equivalences} que
$(U, R, s, t, c)$ est pro-représentable.
\end{proof}

\begin{theorem}
\label{theorem-presentation-deformation-groupoid}
Soit $\mathcal{F}$ une catégorie cofibrée en groupoïdes au-dessus de
$\mathcal{C}_\Lambda$. Alors $\mathcal{F}$ admet une présentation par un
groupoïde en foncteurs lisse et pro-représentable sur $\mathcal{C}_\Lambda$ si et
seulement si les conditions suivantes sont vérifiées :
\begin{enumerate}
\item $\mathcal{F}$ est une catégorie de déformation.
\item $\dim_k T\mathcal{F}$ est finie.
\item $\dim_k \text{Inf}(\mathcal{F})$ est finie.
\end{enumerate}
\end{theorem}

\begin{proof}
Rappelons qu'un foncteur pro-représentable est un foncteur de déformation ;
voir l'exemple \ref{example-prorepresentable-deformation-functor}. Ainsi, si
$\mathcal{F}$ est équivalente à un groupoïde en foncteurs lisse et
pro-représentable, les conditions (1), (2) et (3) résultent des assertions
(1), (2) et (3) du lemme \ref{lemma-deformation-groupoid-quotient}.

\medskip\noindent
Réciproquement, supposons les conditions (1), (2) et (3) vérifiées. La
condition (1) entraîne que (S1) et (S2) sont vérifiées ; voir le lemme
\ref{lemma-RS-implies-S1-S2}. D'après le lemme
\ref{lemma-versal-object-existence}, il existe un objet formel versel $\xi$.
En posant $U = \underline{R}|_{\mathcal{C}_\Lambda}$, le morphisme associé
$\underline{\xi} : U \to \mathcal{F}$ est lisse (c'est la définition d'un
objet formel versel). Soit $(U, R, s, t, c)$ le groupoïde en foncteurs construit
dans le lemme \ref{lemma-presentation-construction} à partir du morphisme
$\underline{\xi}$. Le lemme
\ref{lemma-smooth-groupoid-in-functors-construction} montre que
$(U, R, s, t, c)$ est un groupoïde en foncteurs lisse et que
$[U/R] \to \mathcal{F}$ est une équivalence. Le lemme
\ref{lemma-prorepresentable-groupoid-in-functors-construction} montre que
$(U, R, s, t, c)$ est pro-représentable. Ainsi,
$[U/R] \to \mathcal{F}$ est la présentation cherchée de $\mathcal{F}$.
\end{proof}








\section{Remarques sur la minimalité}
\label{section-minimality}

\noindent
Le théorème principal de ce chapitre est le théorème
\ref{theorem-presentation-deformation-groupoid} ci-dessus. Il décrit
complètement les catégories cofibrées en groupoïdes au-dessus de
$\mathcal{C}_\Lambda$ qui admettent une présentation par un groupoïde en foncteurs
lisse et pro-représentable. Dans cette section, nous expliquons brièvement
comment la minimalité étudiée dans les sections \ref{section-minimal-versal} et
\ref{section-miniversal-objects-existence} permet d'obtenir une présentation
lisse et pro-représentable « minimale ».

\begin{definition}
\label{definition-minimal-groupoid-in-functors}
Soit $(U, R, s, t, c)$ un groupoïde en foncteurs lisse et pro-représentable sur
$\mathcal{C}_\Lambda$.
\begin{enumerate}
\item On dit que $(U, R, s, t, c)$ est {\it normalisé} si le groupoïde
$(U(k[\epsilon]), R(k[\epsilon]), s, t, c)$ est totalement discontinu,
c'est-à-dire s'il n'existe aucun morphisme entre deux objets distincts.
\item On dit que $(U, R, s, t, c)$ est {\it minimal} si $U \to [U/R]$ est
donné par un objet formel versel minimal de $[U/R]$.
\end{enumerate}
\end{definition}

\noindent
La différence entre les deux notions est liée à celle entre les conditions
(\ref{equation-bijective}) et (\ref{equation-bijective-orbits}), et elle
disparaît lorsque $k' \subset k$ est séparable. En outre, un groupoïde en foncteurs
lisse, pro-représentable et normalisé est minimal, comme le montre le lemme
suivant. Voici un énoncé précis.

\begin{lemma}
\label{lemma-characterize-minimal-groupoid-in-functors}
Soit $(U, R, s, t, c)$ un groupoïde en foncteurs lisse et pro-représentable sur
$\mathcal{C}_\Lambda$.
\begin{enumerate}
\item $(U, R, s, t, c)$ est normalisé si et seulement si le morphisme
$U \to [U/R]$ induit un isomorphisme sur les espaces tangents ;
\item $(U, R, s, t, c)$ est minimal si et seulement si le noyau de
$TU \to T[U/R]$ est contenu dans l'image de
$\text{Der}_\Lambda(k, k) \to TU$.
\end{enumerate}
\end{lemma}

\begin{proof}
La partie (1) résulte immédiatement des définitions. Pour voir (2), posons
$\mathcal{F} = [U/R]$. Puisque $\mathcal{F}$ possède une présentation, c'est
une catégorie de déformation ; voir le théorème
\ref{theorem-presentation-deformation-groupoid}. Elle satisfait en particulier
à (RS), (S1) et (S2) ; voir le lemme \ref{lemma-RS-implies-S1-S2}. Rappelons
que les objets formels versels minimaux sont uniques à isomorphisme près ; voir
le lemme \ref{lemma-minimal-versal}. D'après le théorème
\ref{theorem-miniversal-object-existence}, un objet versel minimal induit un
morphisme
$\underline{\xi} : \underline{R}|_{\mathcal{C}_\Lambda} \to \mathcal{F}$
qui satisfait à (\ref{equation-bijective-orbits}). Puisque
$U \cong \underline{R}|_{\mathcal{C}_\Lambda}$ au-dessus de $\mathcal{F}$, on
voit que $TU \to T\mathcal{F} = T[U/R]$ satisfait à la propriété énoncée dans
le lemme.
\end{proof}

\noindent
Le quotient d'un groupoïde en foncteurs minimal et pro-représentable sur $\mathcal
C_\Lambda$ n'admet pas d'autoéquivalence qui ne soit un automorphisme. Pour le
démontrer, commençons par le lemme suivant.

\begin{lemma}
\label{lemma-surjective-morphism-prorepresentable-functor}
Soit $U: \mathcal{C}_\Lambda \to \textit{Ensembles}$ un foncteur
pro-représentable. Soit $\varphi : U \to U$ un morphisme tel que
$d\varphi : TU \to TU$ soit un isomorphisme. Alors $\varphi$ est un
isomorphisme.
\end{lemma}

\begin{proof}
Si $U \cong \underline{R}|_{\mathcal{C}_\Lambda}$ pour un certain
$R \in \Ob(\widehat{\mathcal{C}}_\Lambda)$, la complétion de $\varphi$ donne un
morphisme $\underline{R} \to \underline{R}$. Si $f: R \to R$ est le morphisme
correspondant dans $\widehat{\mathcal{C}}_\Lambda$, alors $f$ induit un
isomorphisme $\text{Der}_\Lambda(R, k) \to \text{Der}_\Lambda(R, k)$ ; voir
l'exemple \ref{example-tangent-space-map-prorepresentable-functor}. En
particulier, $f$ est surjectif d'après le lemme
\ref{lemma-derivations-surjective}. Or un endomorphisme surjectif d'un anneau
noethérien est un isomorphisme (voir Algèbre, lemme
\ref{algebra-lemma-surjective-endo-noetherian-ring-is-iso}) ; on conclut que
$f$, puis $\underline{R}
\to \underline{R}$, puis $\varphi : U \to U$ sont des isomorphismes.
\end{proof}

\begin{lemma}
\label{lemma-minimal-prorepresentable-groupoid-autoequivalence}
Soit $(U, R, s, t, c)$ un groupoïde en foncteurs minimal, lisse et
pro-représentable sur $\mathcal{C}_\Lambda$. Si $\varphi : [U/R] \to [U/R]$ est une
équivalence de catégories cofibrées en groupoïdes, alors $\varphi$ est un
isomorphisme.
\end{lemma}

\begin{proof}
Un morphisme $\varphi : [U/R] \to [U/R]$ équivaut à un morphisme
$\varphi : (U, R, s, t, c) \to (U, R, s, t, c)$ de groupoïdes en foncteurs
au-dessus de $\mathcal{C}_\Lambda$, au sens de la définition
\ref{definition-groupoid-in-functors}. Notons $\phi : U \to U$ et
$\psi : R \to R$ les morphismes correspondants. Puisque le diagramme
$$
\xymatrix{
& \text{Der}_\Lambda(k, k) \ar[dr]_\gamma \ar[dl]^\gamma \\
TU \ar[rr]_{d\phi} \ar[d] & & TU \ar[d]  \\
T[U/R] \ar[rr]^{d\varphi} & & T[U/R]
}
$$
est commutatif, que $d\varphi$ est bijectif, et que nous disposons de la
caractérisation de la minimalité du lemme
\ref{lemma-characterize-minimal-groupoid-in-functors}, on conclut que $d\phi$
est injectif (et donc bijectif pour des raisons de dimension). Par le lemme
\ref{lemma-surjective-morphism-prorepresentable-functor},
$\phi : U \to U$ est donc un isomorphisme. Un argument analogue, faisant
intervenir la suite exacte
$$
0 \to \text{Inf}([U/R]) \to TR \to TU \oplus TU
$$
du lemme \ref{lemma-deformation-functor-diagonal}, montre que
$\psi : R \to R$ est un isomorphisme. Cela résulte aussi du fait que
$R = U \times_{[U/R]} U$ et que $\varphi$ et $\phi$ sont des isomorphismes.
\end{proof}

\begin{lemma}
\label{lemma-minimal-prorepresentable-groupoid-equivalence}
Soient $(U, R, s, t, c)$ et $(U', R', s', t', c')$ des groupoïdes en foncteurs
minimaux, lisses et pro-représentables sur $\mathcal{C}_\Lambda$. Si
$\varphi : [U/R] \to [U'/R']$ est une équivalence de catégories cofibrées en
groupoïdes, alors $\varphi$ est un isomorphisme.
\end{lemma}

\begin{proof}
Soit $\psi : [U'/R'] \to [U/R]$ un quasi-inverse de $\varphi$. Alors
$\psi \circ \varphi$ et $\varphi \circ \psi$ sont des isomorphismes d'après
le lemme \ref{lemma-minimal-prorepresentable-groupoid-autoequivalence} ; ainsi,
$\varphi$ et $\psi$ sont des isomorphismes.
\end{proof}

\noindent
Le lemme suivant résume une partie de ce que nous avons vu plus haut dans ce
chapitre.

\begin{lemma}
\label{lemma-minimal-groupoid-in-functors-construction}
Soit $\mathcal{F}$ une catégorie de déformation telle que
$\dim_k T\mathcal{F} <\infty$ et
$\dim_k \text{Inf}(\mathcal{F}) < \infty$. Il existe alors un objet formel
versel minimal $\xi$ de $\mathcal{F}$. Supposons $\xi$ au-dessus de
$R \in \Ob(\widehat{\mathcal{C}}_\Lambda)$. Soit
$U = \underline{R}|_{\mathcal{C}_\Lambda}$. Soit
$f = \underline{\xi} : U \to \mathcal{F}$ le morphisme associé. Soit
$(U, R, s, t, c)$ le groupoïde en foncteurs sur $\mathcal{C}_\Lambda$
construit à partir de $f : U \to \mathcal{F}$ dans le lemme
\ref{lemma-presentation-construction}. Alors $(U, R, s, t, c)$ est un
groupoïde en foncteurs minimal, lisse et pro-représentable sur
$\mathcal{C}_\Lambda$, et il existe une équivalence
$[U/R] \to \mathcal{F}$.
\end{lemma}

\begin{proof}
Comme $\mathcal{F}$ est une catégorie de déformation, elle satisfait à (S1)
et (S2) ; voir le lemme \ref{lemma-RS-implies-S1-S2}. D'après le lemme
\ref{lemma-versal-object-existence}, il existe un objet formel versel. D'après
le lemme \ref{lemma-minimal-versal}, il existe un objet formel versel minimal
$\xi/R$ comme dans l'énoncé. En posant
$U = \underline{R}|_{\mathcal{C}_\Lambda}$, le morphisme associé
$\underline{\xi} : U \to \mathcal{F}$ est lisse (c'est la définition d'un
objet formel versel). Soit $(U, R, s, t, c)$ le groupoïde en foncteurs construit
dans le lemme \ref{lemma-presentation-construction} à partir du morphisme
$\underline{\xi}$. D'après le lemme
\ref{lemma-smooth-groupoid-in-functors-construction},
$(U, R, s, t, c)$ est un groupoïde en foncteurs lisse et
$[U/R] \to \mathcal{F}$ est une équivalence. D'après le lemme
\ref{lemma-prorepresentable-groupoid-in-functors-construction},
$(U, R, s, t, c)$ est pro-représentable. Enfin, $(U, R, s, t, c)$ est minimal,
car $U \to [U/R] = \mathcal{F}$ correspond à l'objet formel versel minimal
$\xi$.
\end{proof}

\noindent
Les présentations par des groupoïdes en foncteurs minimaux et pro-représentables
satisfont à la propriété d'unicité suivante.

\begin{lemma}
\label{lemma-minimal-presentations-equivalent}
Soit $\mathcal{F}$ une catégorie cofibrée en groupoïdes au-dessus de
$\mathcal{C}_\Lambda$. Supposons qu'il existe des présentations de
$\mathcal{F}$ par des groupoïdes en foncteurs minimaux, lisses et
pro-représentables $(U, R, s, t, c)$ et $(U', R', s', t', c')$. Alors
$(U, R, s, t, c)$ et $(U', R', s', t', c')$ sont isomorphes.
\end{lemma}

\begin{proof}
Cela résulte du lemme
\ref{lemma-minimal-prorepresentable-groupoid-equivalence} et de l'observation
qu'un morphisme $[U/R] \to [U'/R']$ équivaut à un morphisme de groupoïdes en
foncteurs (par notre construction explicite de $[U/R]$ dans la définition
\ref{definition-quotient}).
\end{proof}

\noindent
En résumé, nous avons démontré le théorème suivant.

\begin{theorem}
\label{theorem-minimal-smooth-prorepresentable-presentations}
Soit $\mathcal{F}$ une catégorie cofibrée en groupoïdes au-dessus de
$\mathcal{C}_\Lambda$. Considérons les conditions suivantes :
\begin{enumerate}
\item $\mathcal{F}$ admet une présentation par un groupoïde en foncteurs normalisé,
lisse et pro-représentable sur $\mathcal{C}_\Lambda$ ;
\item $\mathcal{F}$ admet une présentation par un groupoïde en foncteurs lisse
et pro-représentable sur $\mathcal{C}_\Lambda$ ;
\item $\mathcal{F}$ admet une présentation par un groupoïde en foncteurs minimal,
lisse et pro-représentable sur $\mathcal{C}_\Lambda$ ;
\item $\mathcal{F}$ satisfait aux conditions suivantes :
\begin{enumerate}
\item $\mathcal{F}$ est une catégorie de déformation.
\item $\dim_k T\mathcal{F}$ est finie.
\item $\dim_k \text{Inf}(\mathcal{F})$ est finie.
\end{enumerate}
\end{enumerate}
Alors (2), (3) et (4) sont équivalentes et sont entraînées par (1). Si
$k' \subset k$ est séparable, (1), (2), (3) et (4) sont toutes équivalentes.
De plus, les groupoïdes en foncteurs minimaux, lisses et pro-représentables qui
fournissent une présentation de $\mathcal{F}$ sont uniques à isomorphisme près.
\end{theorem}

\begin{proof}
Le lemme \ref{lemma-characterize-minimal-groupoid-in-functors} montre que (1)
entraîne (3), et qu'elle lui est équivalente si $k' \subset k$ est séparable.
Il est clair que (3) entraîne (2). Le théorème
\ref{theorem-presentation-deformation-groupoid} montre que (2) entraîne (4).
Le lemme \ref{lemma-minimal-groupoid-in-functors-construction} montre que (4)
entraîne (3). Cela démontre toutes les implications. La dernière assertion
d'unicité résulte du lemme \ref{lemma-minimal-presentations-equivalent}.
\end{proof}








\section{Unicité des anneaux versels}
\label{section-uniqueness}

\noindent
Étant donnés $R, S$ dans $\widehat{\mathcal{C}}_\Lambda$, on dit que des
morphismes $f, g : R \to S$ sont {\it formellement homotopes} s'il existe un
$r \geq 0$ et des morphismes $h : R \to R[[t_1, \ldots, t_r]]$ et
$k : R[[t_1, \ldots, t_r]] \to S$ dans $\widehat{\mathcal{C}}_\Lambda$ tels
que, pour tout $a \in R$, on ait
\begin{enumerate}
\item $h(a) \bmod (t_1, \ldots, t_r) = a$,
\item $f(a) = k(a)$,
\item $g(a) = k(h(a))$.
\end{enumerate}
On dit que $(r, h, k)$ est une {\it homotopie formelle} entre $f$ et $g$.

\begin{lemma}
\label{lemma-homotopy}
La relation « être formellement homotopes » est une relation d'équivalence sur
les ensembles de morphismes de $\widehat{\mathcal{C}}_\Lambda$.
\end{lemma}

\begin{proof}
Supposons donnés $r \geq 1$ et deux morphismes
$h_1, h_2 : R \to R[[t_1, \ldots, t_r]]$ tels que
$h_1(a) \bmod (t_1, \ldots, t_r) =  h_2(a) \bmod (t_1, \ldots, t_r) = a$
pour tout $a \in R$, ainsi qu'un morphisme
$k : R[[t_1, \ldots, t_r]] \to S$. Nous affirmons que $k \circ h_1$ est
formellement homotope à $k \circ h_2$. La symétrie inhérente à cette assertion
montrera que notre notion est symétrique. En effet, le morphisme
$$
\Psi :
R[[t_1, \ldots, t_r]]
\longrightarrow
R[[t_1, \ldots, t_r]],\quad
\sum a_I t^I \longmapsto \sum h_1(a_I)t^I
$$
est un isomorphisme. Posons $h(a) = \Psi^{-1}(h_2(a))$ pour $a \in R$ et
$k' = k \circ \Psi$ ; alors $(r, h, k')$ est une homotopie formelle entre
$k \circ h_1$ et $k \circ h_2$, ce qui démontre l'assertion.

\medskip\noindent
Supposons donnés trois morphismes $f_1, f_2, f_3 : R \to S$ comme ci-dessus,
une homotopie formelle $(r_1, h_1, k_1)$ entre $f_1$ et $f_2$, et une
homotopie formelle $(r_2, h_2, k_2)$ entre $f_3$ et $f_2$ (!). Après
renumérotation des coordonnées, on peut supposer
$h_2 : R \to R[[t_{r_1 + 1}, \ldots, t_{r_1 + r_2}]]$
et
$k_2 : R[[t_{r_1 + 1}, \ldots, t_{r_1 + r_2}]] \to S$.
En choisissant un isomorphisme convenable
$$
R[[t_1, \ldots, t_{r_1 + r_2}]]
\longrightarrow
R[[t_{r_1 + 1}, \ldots, t_{r_1 + r_2}]]
\widehat{\otimes}_{h_2, R, h_1}
R[[t_1, \ldots, t_{r_1}]]
$$
on peut insérer ces morphismes dans un diagramme commutatif
$$
\xymatrix{
R \ar[r]_{h_1} \ar[d]^{h_2} &
R[[t_1, \ldots, t_{r_1}]] \ar[d]^{h_2'} \\
R[[t_{r_1 + 1}, \ldots, t_{r_1 + r_2}]] \ar[r]^{h_1'} &
R[[t_1, \ldots, t_{r_1 + r_2}]]
}
$$
avec $h_2'(t_i) = t_i$ pour $1 \leq i \leq r_1$ et
$h_1'(t_i) = t_i$ pour $r_1 + 1 \leq i \leq r_2$. Nous omettons certains
détails. Puisque ce diagramme présente une somme amalgamée dans la catégorie
$\widehat{\mathcal{C}}_\Lambda$ (voir la démonstration du lemme
\ref{lemma-CLambdahat-pushouts}) et que
$k_1 \circ h_1 = f_2 = k_2 \circ h_2$, on conclut qu'il existe un morphisme
$$
k : R[[t_1, \ldots, t_{r_1 + r_2}]] \to S
$$
avec $k_1 = k \circ h_2'$ et $k_2 = k \circ h_1'$. Notons
$h = h_1' \circ h_2 = h_2' \circ h_1$. Alors
\begin{enumerate}
\item $k(h_1'(a)) = k_2(a) = f_3(a)$ ;
\item $k(h_2'(a)) = k_1(a) = f_1(a)$.
\end{enumerate}
L'assertion du premier paragraphe de la démonstration montre donc que $f_1$ et
$f_3$ sont formellement homotopes.
\end{proof}

\begin{lemma}
\label{lemma-composition-homotopic}
Dans la catégorie $\widehat{\mathcal{C}}_\Lambda$, si
$f_1, f_2 : R \to S$ sont formellement homotopes et si $g : S \to S'$ est un
morphisme, alors $g \circ f_1$ et $g \circ f_2$ sont formellement homotopes.
\end{lemma}

\begin{proof}
En effet, si $(r, h, k)$ est une homotopie formelle entre $f_1$ et $f_2$,
alors $(r, h, g \circ k)$ est une homotopie formelle entre $g \circ f_1$ et
$g \circ f_2$.
\end{proof}

\begin{lemma}
\label{lemma-versal-unique-up-to-homotopy}
Soit $\mathcal{F}$ une catégorie de déformation au-dessus de
$\mathcal{C}_\Lambda$ telle que $\dim_k T\mathcal{F} < \infty$ et
$\dim_k \text{Inf}(\mathcal{F}) < \infty$. Soit $\xi$ un objet formel versel
au-dessus de $R$. Soit $\eta$ un objet formel au-dessus de $S$. Alors deux
morphismes quelconques
$$
f, g : R \to S
$$
tels que $f_*\xi \cong \eta \cong g_*\xi$ sont formellement homotopes.
\end{lemma}

\begin{proof}
D'après le théorème \ref{theorem-presentation-deformation-groupoid} et sa
démonstration, $\mathcal{F}$ admet une présentation par le groupoïde en foncteurs
suivant :
$$
(\underline{R}, \underline{R_1}, s, t, c, e, i)|_{\mathcal{C}_\Lambda}
$$
Il est lisse et pro-représentable sur $\mathcal{C}_\lambda$ et présente $\mathcal{F}$. Alors les
morphismes $s : R \to R_1$ et $t : R \to R_1$ sont des homomorphismes
d'anneaux formellement lisses, et $e : R_1 \to R$ est une section. En
particulier, on peut choisir un isomorphisme
$R_1 = R[[t_1, \ldots, t_r]]$ pour un certain $r \geq 0$, tel que $s$ soit
l'inclusion $R \subset R[[t_1, \ldots, t_r]]$ et que $t$ corresponde à un
morphisme $h : R \to R[[t_1, \ldots, t_r]]$ vérifiant
$h(a) \bmod (t_1, \ldots, t_r) = a$ pour tout $a \in R$. L'existence de
l'isomorphisme $\alpha : f_*\xi \to g_*\xi$ signifie exactement qu'il existe
un morphisme $k : R_1 \to S$ tel que $f = k \circ s$ et $g = k \circ t$.
Cela signifie exactement que $(r, h, k)$ est une homotopie formelle entre $f$
et $g$.
\end{proof}

\begin{lemma}
\label{lemma-homotopic-minimal-prime}
Dans la catégorie $\widehat{\mathcal{C}}_\Lambda$, si
$f_1, f_2 : R \to S$ sont formellement homotopes et si
$\mathfrak p \subset R$ est un idéal premier minimal, alors
$f_1(\mathfrak p)S = f_2(\mathfrak p)S$ comme idéaux.
\end{lemma}

\begin{proof}
Supposons que $(r, h, k)$ soit une homotopie formelle entre $f_1$ et $f_2$.
Nous affirmons que
$\mathfrak pR[[t_1, \ldots, t_r]] = h(\mathfrak p)R[[t_1, \ldots, t_r]]$.
L'assertion implique le lemme en composant encore avec $k$. Pour démontrer
l'assertion, observons que l'application
$\mathfrak p \mapsto \mathfrak pR[[t_1, \ldots, t_r]]$
est une bijection de l'ensemble des idéaux premiers minimaux de $R$ sur celui
des idéaux premiers minimaux de $R[[t_1, \ldots, t_r]]$. Enfin,
$h(\mathfrak p)R[[t_1, \ldots, t_r]]$ est un idéal premier minimal puisque $h$ est
plat ; il est donc de la forme $\mathfrak q R[[t_1, \ldots, t_r]]$ pour un
certain idéal premier minimal $\mathfrak q \subset R$, d'après ce qui précède.
Mais, puisque $h \bmod (t_1, \ldots, t_r) = \text{id}_R$ par définition d'une
homotopie formelle, on conclut que $\mathfrak q = \mathfrak p$, comme voulu.
\end{proof}











\section{Changement du corps résiduel}
\label{section-change-of-field}

\noindent
Dans cette section, nous examinons rapidement ce qui se passe lorsque l'on
remplace le corps résiduel $k$ par une extension finie. Soit $\Lambda$ un
anneau noethérien et soit $\Lambda \to k$ un homomorphisme fini, où $k$ est un
corps. Dans tout ce chapitre, nous avons noté $\mathcal{C}_\Lambda$ la
catégorie des $\Lambda$-algèbres locales artiniennes dont le corps résiduel est
identifié à $k$ ; voir la définition \ref{definition-CLambda}. Cependant,
puisque nous étudierons dans cette section ce qui arrive lorsque l'on change
$k$, nous utiliserons plutôt la notation $\mathcal{C}_{\Lambda, k}$ afin de
faire apparaître la dépendance en $k$.

\begin{situation}
\label{situation-change-of-fields}
Soit $\Lambda$ un anneau noethérien, et soient $\Lambda \to k \to l$ des
homomorphismes d'anneaux finis, où $k$ et $l$ sont des corps. Ainsi, $l/k$ est
une extension finie de corps. Un objet typique de
$\mathcal{C}_{\Lambda, l}$ sera noté $B$, et un objet typique de
$\mathcal{C}_{\Lambda, k}$ sera noté $A$. Définissons
\begin{equation}
\label{equation-comparison}
\mathcal{C}_{\Lambda, l} \longrightarrow \mathcal{C}_{\Lambda, k},
\quad
B \longmapsto B \times_l k
\end{equation}
Étant donnée une catégorie cofibrée en groupoïdes
$p : \mathcal{F} \to \mathcal{C}_{\Lambda, k}$, on obtient une catégorie
cofibrée en groupoïdes associée
$$
p_{l/k} : \mathcal{F}_{l/k} \longrightarrow \mathcal{C}_{\Lambda, l}
$$
en posant $\mathcal{F}_{l/k}(B) = \mathcal{F}(B \times_l k)$.
\end{situation}

\noindent
Le foncteur (\ref{equation-comparison}) est bien défini : puisque
$B \times_l k \subset B$, on a
\begin{align*}
[k : k']\ \text{length}_{B \times_l k}(B \times_l k) & =
\text{length}_\Lambda(B \times_l k) \\
& \leq \text{length}_\Lambda(B) \\
& = [l : k']\ \text{length}_B(B) < \infty
\end{align*}
(voir le lemme \ref{lemma-length}) ; par conséquent, $B \times_l k$ est
artinien (voir Algèbre, lemme
\ref{algebra-lemma-artinian-finite-length}). Ainsi, $B \times_l k$ est un
anneau local artinien de corps résiduel $k$. Notons que
(\ref{equation-comparison}) commute aux produits fibrés
$$
(B_1 \times_B B_2) \times_l k =
(B_1 \times_l k) \times_{(B \times_l k)} (B_2 \times_l k)
$$
et transforme les homomorphismes d'anneaux surjectifs en homomorphismes
d'anneaux surjectifs. Nous employons la notation « lourde »
$\mathcal{F}_{l/k}$ afin d'éviter toute confusion avec la construction de la
remarque \ref{remark-localize-cofibered-groupoid}. Voici quelques observations
élémentaires.

\begin{lemma}
\label{lemma-elementary-properties-change-of-field}
Avec les notations et les hypothèses de la situation
\ref{situation-change-of-fields}.
\begin{enumerate}
\item On a $\overline{\mathcal{F}_{l/k}} = (\overline{\mathcal{F}})_{l/k}$.
\item Si $\mathcal{F}$ est une catégorie de prédéformation, alors
$\mathcal{F}_{l/k}$ est une catégorie de prédéformation.
\item Si $\mathcal{F}$ satisfait à (S1), alors $\mathcal{F}_{l/k}$ satisfait à
(S1).
\item Si $\mathcal{F}$ satisfait à (S2), alors $\mathcal{F}_{l/k}$ satisfait à
(S2).
\item Si $\mathcal{F}$ satisfait à (RS), alors $\mathcal{F}_{l/k}$ satisfait à
(RS).
\end{enumerate}
\end{lemma}

\begin{proof}
La partie (1) résulte immédiatement des définitions.

\medskip\noindent
Puisque $\mathcal{F}_{l/k}(l) = \mathcal{F}(k)$, la partie (2) résulte de la
définition ; voir la définition \ref{definition-predeformation-category}.

\medskip\noindent
La partie (3) résulte du fait que le foncteur (\ref{equation-comparison})
commute aux produits fibrés et transforme les morphismes surjectifs en
morphismes surjectifs ; voir la définition \ref{definition-S1-S2}.

\medskip\noindent
Partie (4). Pour le voir, considérons un diagramme
$$
\xymatrix{
          & l[\epsilon] \ar[d] \\
B  \ar[r] & l
}
$$
de $\mathcal{C}_{\Lambda, l}$ comme dans la définition
\ref{definition-S1-S2}. En appliquant le foncteur
(\ref{equation-comparison}), on obtient
$$
\xymatrix{
          & k[l\epsilon] \ar[d] \\
B \times_l k  \ar[r] & k
}
$$
où $l\epsilon$ désigne l'espace vectoriel de dimension finie sur $k$
$l\epsilon \subset l[\epsilon]$. D'après le lemme
\ref{lemma-S2-extensions}, la condition (S2) est aussi satisfaite par $\mathcal{F}$
pour ce diagramme. Ainsi, (S2) est satisfaite par $\mathcal{F}_{l/k}$.

\medskip\noindent
La partie (5) résulte de la caractérisation de (RS) donnée dans le lemme
\ref{lemma-RS-2-categorical}, partie (2), et du fait que
(\ref{equation-comparison}) commute aux produits fibrés.
\end{proof}

\noindent
Le lemme suivant s'applique notamment lorsque $\mathcal{F}$ satisfait à (S2)
et est une catégorie de prédéformation ; voir le lemme
\ref{lemma-S1-S2-associated-functor}.

\begin{lemma}
\label{lemma-tangent-space-change-of-field}
Avec les notations et les hypothèses de la situation
\ref{situation-change-of-fields}. Supposons que $\mathcal{F}$ soit une
catégorie de prédéformation et que $\overline{\mathcal{F}}$ satisfasse à (S2).
Il existe alors un isomorphisme canonique $l$-linéaire
$$
T\mathcal{F} \otimes_k l \longrightarrow T\mathcal{F}_{l/k}
$$
entre les espaces tangents.
\end{lemma}

\begin{proof}
D'après le lemme \ref{lemma-elementary-properties-change-of-field}, on peut
remplacer $\mathcal{F}$ par $\overline{\mathcal{F}}$. De plus,
$T\mathcal{F}$ et $T\mathcal{F}_{l/k}$ possèdent respectivement une structure
canonique d'espace vectoriel sur $k$ et sur $l$ ; voir le lemme
\ref{lemma-tangent-space-vector-space}. Alors
$$
T\mathcal{F}_{l/k} = \mathcal{F}_{l/k}(l[\epsilon])
= \mathcal{F}(k[l\epsilon]) = T\mathcal{F} \otimes_k l
$$
où la dernière égalité résulte du lemme
\ref{lemma-tangent-space-vector-space}. Plus généralement, étant donné un
espace vectoriel de dimension finie sur $l$, noté $V$, on a
$$
\mathcal{F}_{l/k}(l[V]) = \mathcal{F}(k[V_k]) = T\mathcal{F} \otimes_k V_k
$$
où $V_k$ désigne $V$ considéré comme espace vectoriel sur $k$. On en déduit
que les foncteurs $V \mapsto \mathcal{F}_{l/k}(l[V])$ et
$V \mapsto T\mathcal{F} \otimes_k V_k$ sont canoniquement identifiés comme
foncteurs à valeurs dans la catégorie des ensembles. Le lemme
\ref{lemma-linear-functor} montre qu'il existe au plus une façon de rendre l'un
ou l'autre foncteur $l$-linéaire. Les isomorphismes sont donc compatibles avec les
structures d'espaces vectoriels sur $l$, ce qui conclut.
\end{proof}

\begin{lemma}
\label{lemma-inf-aut-change-of-field}
Avec les notations et les hypothèses de la situation
\ref{situation-change-of-fields}. Supposons que $\mathcal{F}$ soit une
catégorie de déformation. Il existe alors un isomorphisme canonique
$l$-linéaire
$$
\text{Inf}(\mathcal{F}) \otimes_k l
\longrightarrow
\text{Inf}(\mathcal{F}_{l/k})
$$
entre les espaces d'automorphismes infinitésimaux.
\end{lemma}

\begin{proof}
Soit $x_0 \in \Ob(\mathcal{F}(k))$, et notons $x_{l, 0}$ l'objet correspondant
de $\mathcal{F}_{l/k}$ au-dessus de $l$. Rappelons que
$\text{Inf}(\mathcal{F}) = \text{Inf}_{x_0}(\mathcal{F})$ et
$\text{Inf}(\mathcal{F}_{l/k}) = \text{Inf}_{x_{l, 0}}(\mathcal{F}_{l/k})$ ;
voir la remarque \ref{remark-trivial-aut-point}. Rappelons que la structure
d'espace vectoriel sur $\text{Inf}_{x_0}(\mathcal{F})$ provient de son
identification à l'espace tangent du foncteur $\mathit{Aut}(x_0)$ défini sur la
catégorie $\mathcal{C}_{k, k}$ des $k$-algèbres locales artiniennes de corps
résiduel $k$. De même, $\text{Inf}_{x_{l, 0}}(\mathcal{F}_{l/k})$ est l'espace
tangent de $\mathit{Aut}(x_{l, 0})$, défini sur la catégorie
$\mathcal{C}_{l, l}$ des $l$-algèbres locales artiniennes de corps résiduel
$l$. En explicitant les définitions, on voit que $\mathit{Aut}(x_{l, 0})$ est
la restriction de $\mathit{Aut}(x_0)_{l/k}$ (qui est défini sur
$\mathcal{C}_{k, l}$) à $\mathcal{C}_{l, l}$. Puisqu'il n'y a aucune
différence entre l'espace tangent de $\mathit{Aut}(x_0)_{l/k}$ considéré comme
foncteur sur $\mathcal{C}_{k, l}$ ou sur $\mathcal{C}_{l, l}$, le lemme résulte
du lemme \ref{lemma-tangent-space-change-of-field} et du fait que
$\mathit{Aut}(x_0)$ satisfait à (RS) d'après le lemme
\ref{lemma-Aut-functor-RS} (d'où (S2) par le lemme
\ref{lemma-RS-implies-S1-S2}).
\end{proof}

\begin{lemma}
\label{lemma-change-of-fields-smooth}
Avec les notations et les hypothèses de la situation
\ref{situation-change-of-fields}. Si $\mathcal{F} \to \mathcal{G}$ est un
morphisme lisse de catégories cofibrées en groupoïdes au-dessus de
$\mathcal{C}_{\Lambda, k}$, alors
$\mathcal{F}_{l/k} \to \mathcal{G}_{l/k}$ est un morphisme lisse de catégories
cofibrées en groupoïdes au-dessus de $\mathcal{C}_{\Lambda, l}$.
\end{lemma}

\begin{proof}
Cela résulte immédiatement des définitions et du fait que
(\ref{equation-comparison}) préserve les surjections.
\end{proof}

\noindent
On pourrait dire bien davantage sur la relation entre $\mathcal{F}$ et
$\mathcal{F}_{l/k}$ (notamment sur la relation entre les déformations
verselles), et nous ajouterons ici ces compléments au besoin.

\begin{lemma}
\label{lemma-change-of-field-versal-ring}
Avec les notations et les hypothèses de la situation
\ref{situation-change-of-fields}. Soit $\xi$ un objet formel versel de
$\mathcal{F}$ au-dessus de
$R \in \Ob(\widehat{\mathcal{C}}_{\Lambda, k})$. Il existe alors
\begin{enumerate}
\item un $S \in \Ob(\widehat{\mathcal{C}}_{\Lambda, l})$ et un homomorphisme
local de $\Lambda$-algèbres $R \to S$, formellement lisse pour la topologie
$\mathfrak m_S$-adique et induisant sur les corps résiduels l'extension donnée
$l/k$ ;
\item un objet formel versel de $\mathcal{F}_{l/k}$ au-dessus de $S$.
\end{enumerate}
\end{lemma}

\begin{proof}
Construction de $S$. Choisissons une surjection
$R[x_1, \ldots, x_n] \to l$ de $R$-algèbres. Son noyau est un idéal maximal
$\mathfrak m$.
Prenons pour $S$ le complété $\mathfrak m$-adique de l'anneau noethérien
$R[x_1, \ldots, x_n]$. Alors $S$ appartient à
$\widehat{\mathcal{C}}_{\Lambda, l}$ d'après Algèbre, lemme
\ref{algebra-lemma-completion-Noetherian-Noetherian}. Le morphisme $R \to S$
est formellement lisse pour la topologie $\mathfrak m_S$-adique d'après
Compléments d'algèbre, lemmes \ref{more-algebra-lemma-formally-smooth} et
\ref{more-algebra-lemma-formally-smooth-completion}, et parce que
$R \to R[x_1, \ldots, x_n]$ est formellement lisse. (Comparer avec la
démonstration du lemme \ref{lemma-exists-smooth}.)

\medskip\noindent
Puisque $\xi$ est versel, la transformation
$\underline{\xi} : \underline{R}|_{\mathcal{C}_{\Lambda, k}} \to \mathcal{F}$
est lisse. D'après le lemme \ref{lemma-change-of-fields-smooth}, le morphisme
induit
$$
(\underline{R}|_{\mathcal{C}_{\Lambda, k}})_{l/k}
\longrightarrow
\mathcal{F}_{l/k}
$$
est lisse. Il suffit donc de construire un morphisme lisse
$\underline{S}|_{\mathcal{C}_{\Lambda, l}} \to
(\underline{R}|_{\mathcal{C}_{\Lambda, k}})_{l/k}$. Donner un tel morphisme
revient à donner, pour tout objet $B$ de $\mathcal{C}_{\Lambda, l}$, une
application d'ensembles
$$
\Mor_{\widehat{\mathcal{C}}_{\Lambda, l}}(S, B)
\longrightarrow
\Mor_{\widehat{\mathcal{C}}_{\Lambda, k}}(R, B \times_l k)
$$
fonctorielle en $B$. Un élément $\varphi : S \to B$ du membre de gauche est
envoyé sur le composé $R \to S \to B$, dont l'image est contenue dans la
sous-$\Lambda$-algèbre $B \times_l k$. La lissité du morphisme signifie
qu'étant donnés une surjection $B' \to B$ et un diagramme commutatif
$$
\xymatrix{
S \ar[r] & B \ar@{=}[r] & B \\
R \ar[u] \ar[r] & B' \times_l k \ar[u] \ar[r] & B' \ar[u]
}
$$
il faut trouver un homomorphisme d'anneaux $S \to B'$ s'insérant dans le
rectangle extérieur. L'existence de ce morphisme est garantie par le fait que
nous avons choisi $R \to S$ formellement lisse pour la topologie $\mathfrak m_S$-adique ; voir
Compléments d'algèbre, lemme \ref{more-algebra-lemma-lift-continuous}.
\end{proof}





















\begin{multicols}{2}[\section{Autres chapitres}]
\noindent
Préliminaires
\begin{enumerate}
\item \hyperref[introduction-section-phantom]{Introduction}
\item \hyperref[conventions-section-phantom]{Conventions}
\item \hyperref[sets-section-phantom]{Théorie des ensembles}
\item \hyperref[categories-section-phantom]{Catégories}
\item \hyperref[topology-section-phantom]{Topologie}
\item \hyperref[sheaves-section-phantom]{Faisceaux sur les espaces}
\item \hyperref[sites-section-phantom]{Sites et faisceaux}
\item \hyperref[stacks-section-phantom]{Champs}
\item \hyperref[fields-section-phantom]{Corps}
\item \hyperref[algebra-section-phantom]{Algèbre commutative}
\item \hyperref[brauer-section-phantom]{Groupes de Brauer}
\item \hyperref[homology-section-phantom]{Algèbre homologique}
\item \hyperref[derived-section-phantom]{Catégories dérivées}
\item \hyperref[simplicial-section-phantom]{Méthodes simpliciales}
\item \hyperref[more-algebra-section-phantom]{Compléments d'algèbre}
\item \hyperref[smoothing-section-phantom]{Lissage des morphismes d'anneaux}
\item \hyperref[modules-section-phantom]{Faisceaux de modules}
\item \hyperref[sites-modules-section-phantom]{Modules sur les sites}
\item \hyperref[injectives-section-phantom]{Objets injectifs}
\item \hyperref[cohomology-section-phantom]{Cohomologie des faisceaux}
\item \hyperref[sites-cohomology-section-phantom]{Cohomologie sur les sites}
\item \hyperref[dga-section-phantom]{Algèbre différentielle graduée}
\item \hyperref[dpa-section-phantom]{Algèbre à puissances divisées}
\item \hyperref[sdga-section-phantom]{Faisceaux différentiels gradués}
\item \hyperref[hypercovering-section-phantom]{Hyperrecouvrements}
\end{enumerate}
Schémas
\begin{enumerate}
\setcounter{enumi}{25}
\item \hyperref[schemes-section-phantom]{Schémas}
\item \hyperref[constructions-section-phantom]{Constructions de schémas}
\item \hyperref[properties-section-phantom]{Propriétés des schémas}
\item \hyperref[morphisms-section-phantom]{Morphismes de schémas}
\item \hyperref[coherent-section-phantom]{Cohomologie des schémas}
\item \hyperref[divisors-section-phantom]{Diviseurs}
\item \hyperref[limits-section-phantom]{Limites de schémas}
\item \hyperref[varieties-section-phantom]{Variétés}
\item \hyperref[topologies-section-phantom]{Topologies sur les schémas}
\item \hyperref[descent-section-phantom]{Descente}
\item \hyperref[perfect-section-phantom]{Catégories dérivées de schémas}
\item \hyperref[more-morphisms-section-phantom]{Compléments sur les morphismes}
\item \hyperref[flat-section-phantom]{Compléments sur la platitude}
\item \hyperref[groupoids-section-phantom]{Schémas en groupoïdes}
\item \hyperref[more-groupoids-section-phantom]{Compléments sur les schémas en groupoïdes}
\item \hyperref[etale-section-phantom]{Morphismes étales de schémas}
\end{enumerate}
Thèmes de théorie des schémas
\begin{enumerate}
\setcounter{enumi}{41}
\item \hyperref[chow-section-phantom]{Homologie de Chow}
\item \hyperref[intersection-section-phantom]{Théorie de l'intersection}
\item \hyperref[pic-section-phantom]{Schémas de Picard des courbes}
\item \hyperref[weil-section-phantom]{Théories cohomologiques de Weil}
\item \hyperref[adequate-section-phantom]{Modules adéquats}
\item \hyperref[dualizing-section-phantom]{Complexes dualisants}
\item \hyperref[duality-section-phantom]{Dualité pour les schémas}
\item \hyperref[discriminant-section-phantom]{Discriminants et différentes}
\item \hyperref[derham-section-phantom]{Cohomologie de de Rham}
\item \hyperref[local-cohomology-section-phantom]{Cohomologie locale}
\item \hyperref[algebraization-section-phantom]{Géométrie algébrique et formelle}
\item \hyperref[curves-section-phantom]{Courbes algébriques}
\item \hyperref[resolve-section-phantom]{Résolution des surfaces}
\item \hyperref[models-section-phantom]{Réduction semi-stable}
\item \hyperref[functors-section-phantom]{Foncteurs et morphismes}
\item \hyperref[equiv-section-phantom]{Catégories dérivées de variétés}
\item \hyperref[pione-section-phantom]{Groupes fondamentaux de schémas}
\item \hyperref[etale-cohomology-section-phantom]{Cohomologie étale}
\item \hyperref[crystalline-section-phantom]{Cohomologie cristalline}
\item \hyperref[proetale-section-phantom]{Cohomologie pro-étale}
\item \hyperref[relative-cycles-section-phantom]{Cycles relatifs}
\item \hyperref[more-etale-section-phantom]{Compléments de cohomologie étale}
\item \hyperref[trace-section-phantom]{La formule des traces}
\end{enumerate}
Espaces algébriques
\begin{enumerate}
\setcounter{enumi}{64}
\item \hyperref[spaces-section-phantom]{Espaces algébriques}
\item \hyperref[spaces-properties-section-phantom]{Propriétés des espaces algébriques}
\item \hyperref[spaces-morphisms-section-phantom]{Morphismes d'espaces algébriques}
\item \hyperref[decent-spaces-section-phantom]{Espaces algébriques décents}
\item \hyperref[spaces-cohomology-section-phantom]{Cohomologie des espaces algébriques}
\item \hyperref[spaces-limits-section-phantom]{Limites d'espaces algébriques}
\item \hyperref[spaces-divisors-section-phantom]{Diviseurs sur les espaces algébriques}
\item \hyperref[spaces-over-fields-section-phantom]{Espaces algébriques sur des corps}
\item \hyperref[spaces-topologies-section-phantom]{Topologies sur les espaces algébriques}
\item \hyperref[spaces-descent-section-phantom]{Descente et espaces algébriques}
\item \hyperref[spaces-perfect-section-phantom]{Catégories dérivées d'espaces}
\item \hyperref[spaces-more-morphisms-section-phantom]{Compléments sur les morphismes d'espaces}
\item \hyperref[spaces-flat-section-phantom]{Platitude sur les espaces algébriques}
\item \hyperref[spaces-groupoids-section-phantom]{Groupoïdes dans les espaces algébriques}
\item \hyperref[spaces-more-groupoids-section-phantom]{Compléments sur les groupoïdes dans les espaces}
\item \hyperref[bootstrap-section-phantom]{Amorçage}
\item \hyperref[spaces-pushouts-section-phantom]{Sommes amalgamées d'espaces algébriques}
\end{enumerate}
Thèmes de géométrie
\begin{enumerate}
\setcounter{enumi}{81}
\item \hyperref[spaces-chow-section-phantom]{Groupes de Chow des espaces}
\item \hyperref[groupoids-quotients-section-phantom]{Quotients de groupoïdes}
\item \hyperref[spaces-more-cohomology-section-phantom]{Compléments sur la cohomologie des espaces}
\item \hyperref[spaces-simplicial-section-phantom]{Espaces simpliciaux}
\item \hyperref[spaces-duality-section-phantom]{Dualité pour les espaces}
\item \hyperref[formal-spaces-section-phantom]{Espaces algébriques formels}
\item \hyperref[restricted-section-phantom]{Algébrisation des espaces formels}
\item \hyperref[spaces-resolve-section-phantom]{Résolution des surfaces revisitée}
\end{enumerate}
Théorie des déformations
\begin{enumerate}
\setcounter{enumi}{89}
\item \hyperref[formal-defos-section-phantom]{Théorie formelle des déformations}
\item \hyperref[defos-section-phantom]{Théorie des déformations}
\item \hyperref[cotangent-section-phantom]{Le complexe cotangent}
\item \hyperref[examples-defos-section-phantom]{Problèmes de déformation}
\end{enumerate}
Champs algébriques
\begin{enumerate}
\setcounter{enumi}{93}
\item \hyperref[algebraic-section-phantom]{Champs algébriques}
\item \hyperref[examples-stacks-section-phantom]{Exemples de champs}
\item \hyperref[stacks-sheaves-section-phantom]{Faisceaux sur les champs algébriques}
\item \hyperref[criteria-section-phantom]{Critères de représentabilité}
\item \hyperref[artin-section-phantom]{Axiomes d'Artin}
\item \hyperref[quot-section-phantom]{Espaces de Quot et de Hilbert}
\item \hyperref[stacks-properties-section-phantom]{Propriétés des champs algébriques}
\item \hyperref[stacks-morphisms-section-phantom]{Morphismes de champs algébriques}
\item \hyperref[stacks-limits-section-phantom]{Limites de champs algébriques}
\item \hyperref[stacks-cohomology-section-phantom]{Cohomologie des champs algébriques}
\item \hyperref[stacks-perfect-section-phantom]{Catégories dérivées de champs}
\item \hyperref[stacks-introduction-section-phantom]{Introduction aux champs algébriques}
\item \hyperref[stacks-more-morphisms-section-phantom]{Compléments sur les morphismes de champs}
\item \hyperref[stacks-geometry-section-phantom]{La géométrie des champs}
\end{enumerate}
Thèmes de théorie des espaces de modules
\begin{enumerate}
\setcounter{enumi}{107}
\item \hyperref[moduli-section-phantom]{Champs de modules}
\item \hyperref[moduli-curves-section-phantom]{Espaces de modules de courbes}
\end{enumerate}
Divers
\begin{enumerate}
\setcounter{enumi}{109}
\item \hyperref[examples-section-phantom]{Exemples}
\item \hyperref[exercises-section-phantom]{Exercices}
\item \hyperref[guide-section-phantom]{Guide bibliographique}
\item \hyperref[desirables-section-phantom]{Desiderata}
\item \hyperref[coding-section-phantom]{Style de programmation}
\item \hyperref[obsolete-section-phantom]{Obsolète}
\item \hyperref[fdl-section-phantom]{Licence de documentation libre GNU}
\item \hyperref[index-section-phantom]{Index généré automatiquement}
\end{enumerate}
\end{multicols}


\bibliography{my}
\bibliographystyle{amsalpha}

\end{document}
